%!TEX root = ../larxxia.tex

\chapter{Eigenvalues and eigenvectors in general}
\label{ch:gee}

\minitoc





\paragraph{Population modelling} \index{population modelling}
Suppose two \idx{species} of animals interact: how do their populations evolve in time?  
Let \(y(t)\) and~\(z(t)\) be the number of \idx{female} animals in each of the species at time~\(t\) in years (biologists usually just count females in population models as females usually determine reproduction). 
Modelling might deduce that the populations interact according to the rule that one year later the population is \(y(t+1)=2y(t)-4z(t)\) and \(z(t+1)=-y(t)+2z(t)\): that is, if it was not for the other species, then for each species the number of females would double every year (since then \(y(t+1)=2y(t)\) and \(z(t+1)=2z(t)\)); but the presence of the other species causes competition for resources that decreases each of these growths via the~\(-4z(t)\) and~\(-y(t)\) terms.  

Question: can we find special solutions in the form \((y,z)=\xv \lambda^{t}\) for some constant~\(\lambda\) and nonzero \(\xv=(x_1,x_2)\)?  
Let's try by substituting \(y=x_1\lambda^{t}\) and \(z=x_2\lambda^{t}\) into the  equations:
\begin{eqnarray*}
&&y(t+1)=2y(t)-4z(t)\,,\quad z(t+1)=-y(t)+2z(t)
\\&\iff& 
x_1 \lambda^{t+1}=2x_1\lambda^{t}-4x_2\lambda^{t}
\,,\quad
x_2 \lambda^{t+1}=-x_1\lambda^{t}+2x_2\lambda^{t}
\\&\iff& 2x_1-4x_2=\lambda x_1
\,,\quad
-x_1+2x_2=\lambda x_2
\end{eqnarray*}
after dividing by the factor~\(\lambda^{t}\) (assuming constant~\(\lambda\) is nonzero).
Then form these last two equations as the matrix-vector equation
\begin{equation*}
\begin{bmatrix} 2&-4\\-1&2 \end{bmatrix}\xv=\lambda\xv\,.
\end{equation*}
That is, this substitution \((y,z)=\xv \lambda^{t}\) shows the question about finding solutions of the population equations reduces to solving \(A\xv=\lambda \xv\)\,, called an \idx{eigen-problem}.

This chapter develops linear algebra for such \idx{eigen-problem}s that empowers us to construct the general solution for the population \(y(t)\) and~\(z(t)\). 
Here, a general solution is, in terms of two constants~\(c_1\) and~\(c_2\), that one \idx{species} has female population \(y(t)=2c_14^t+2c_2\) whereas the second species has female population \(z(t)=-c_14^t+c_2\)\,.

%Modelling might deduce the populations interact according to the pair  of \idx{differential equation}s \(\dot y=y-4z\) and \(\dot z=-y+z\) (where the overdot denotes~\(d/dt\)); that is, if it was not for the other species the number of females would both grow according to \(\dot y=y\) and \(\dot z=z\)\,, but the other species decreases their growth via the~\(-2z\) and~\(-y\) terms.  
%
%Question: can we find special solutions in the form \((y,z)=\xv e^{\lambda t}\)?  
%Substitute \(y=x_1e^{\lambda t}\) and \(z=x_2e^{\lambda t}\) into the differential equations:
%\begin{eqnarray*}
%&&\dot y=y-4z\,,\quad \dot z=-y+z
%\\&\iff& 
%x_1\lambda e^{\lambda t}=x_1e^{\lambda t}-4x_2e^{\lambda t}
%\,,\quad
%x_2\lambda e^{\lambda t}=-x_1e^{\lambda t}+x_2e^{\lambda t}
%\\&\iff& x_1-4x_2=\lambda x_1
%\,,\quad
%-x_1+x_2=\lambda x_2
%\end{eqnarray*}
%after dividing by the exponential~\(e^{\lambda t}\) (it is always nonzero).
%Forming as a matrix-vector equation these last two equations are
%\begin{equation*}
%\begin{bmatrix} 1&-4\\-1&1 \end{bmatrix}\xv=\lambda\xv\,.
%\end{equation*}
%That is, this substitution \((y,z)=\xv e^{\lambda t}\) shows the question about solutions of differential equations reduces to solving \(A\xv=\lambda \xv\)\,, called an eigen-problem.


\index{eigenvalue|(}
\index{eigenvector|(}
\index{eigenspace|(}


\paragraph{The basic eigen-problem}
Recall from \cref{sec:iee} that the \idx{eigen-problem} equation \(A\xv=\lambda\xv\) is just asking the following: 
can we find directions~\xv\ such that matrix~\(A\) acting on~\xv\ is in the same direction as~\xv? 
That is, when is~\(A\xv\) the same as~\(\lambda\xv\) for some proportionality constant~\(\lambda\)?
Now \(\xv=\ov\) is always a trivial solution of the  equation \(A\xv=\lambda \xv\)\,.
Consequently, we are only interested in those values of the \idx{eigenvalue}~\(\lambda\) when \emph{nonzero} solutions for the \idx{eigenvector}~\xv\ exist (as it is the directions which are of interest).

Rearranging the equation \(A\xv=\lambda\xv\) as the \idx{homogeneous} \idx{system} \((A-\lambda I)\xv=\ov\)\,, let's invoke properties of linear equations to solve the eigen-problem.
\begin{itemize}
\item \cref{pro:eeh} establishes that one way to find the eigenvalues~\(\lambda\) (albeit \emph{only} suitable for matrices of small size) is to solve the \idx{characteristic equation} \(\det(A-\lambda I)=0\)\,. 
\item Then for each eigenvalue, solving the homogeneous system \((A-\lambda I)\xv=\ov\) gives corresponding eigenvectors~\xv. 
\item The set of eigenvectors for a given eigenvalue forms a \idx{subspace} called the \idx{eigenspace}~\(\EE_\lambda\) (\cref{thm:espacedef}).

\end{itemize}






\paragraph{Three general difficulties in \idx{eigen-problem}s}
Recall that \cref{sec:iee} introduced one way to visually estimate eigenvectors and eigenvalues of a given \(2\times2\) matrix~\(A\) \cite[]{Schonefeld1995}. 
The graphical method is to plot many \idx{unit vector}s~\xv, and at the end of each~\xv\ to adjoin the vector~\(A\xv\).
Since eigenvectors satisfy \(A\xv=\lambda\xv\) for some \idx{scalar} eigenvalue~\(\lambda\), we visually identify eigenvectors as those~\xv\ for which~\(A\xv\) points in the same (or opposite) direction as~\xv.%
\footnote{Recall that the \script[1]\ function \index{eigshow()@\texttt{eigshow()}}\texttt{eigshow(A)} (download for R2017b etc.) provides an interactive alternative to the static views shown herein.}
Let's use this approach to identify three \text{general difficulties.}

\begin{enumerate}
\item 
\begin{figbox}{\eRose{1}1{0.125}{1}}%
In this first picture, for matrix
\( A=\begin{bmat} 1&1\\1/8&1 \end{bmat}\),
the eigenvectors appear to be in directions \(\xv_1\approx\pm(0.9,0.3)\) and \(\xv_2\approx\pm(0.9,-0.3)\) corresponding to eigenvalues \(\lambda_1\approx1.4\) and \(\lambda_2\approx 0.6\)\,.
(Recall that scalar multiples of an eigenvector are always also eigenvectors, \cref{sec:iee}, so we always see \(\pm\)~pairs of eigenvectors in these pictures.)
The eigenvectors \(\pm(0.9,0.3)\) are not orthogonal to the other 
eigenvectors \(\pm(0.9,-0.3)\), not at right-angles---as happens for symmetric matrices (\cref{thm:orthoevec}).
This lack of orthogonality in general means we soon generalize the concept of orthogonal sets of vectors to a new concept of \idx{linearly independent} sets (\cref{sec:lisb}).
\end{figbox}

\item 
\begin{figbox}{\eRose01{-1}{0.5}}%
In this second case, for
\( A=\begin{bmat} 0&1\\-1&1/2 \end{bmat}\),
there appears to be no (red) vector~\(A\xv\) in the same direction as the corresponding (blue) vector~\xv.
Thus there appears to be no eigenvectors at all.
No eigenvectors and eigenvalues is the answer if we require real answers.
However, in most applications we find it sensible to have non-real \idx{complex valued} eigenvalues and eigenvectors (\cref{sec:eennm})\index{complex eigenvalue}\index{complex eigenvector}, written using \index{i@$\i$}\(\i=\sqrt{-1}\).
So although we cannot see them graphically, for this matrix there are two \idx{complex eigenvalue}s and two families of \idx{complex eigenvector}s (analogous to those found in \cref{eg:ccevals}).% 
\end{figbox}

In this second case the vectors~\(A\xv\) all appear to be pointing clockwise.  
Such a consistent `\idx{rotation}' in~\(A\xv\) is characteristic of matrices with non-real complex valued eigenvalues and eigenvectors.

\item 
\begin{figbox}{\eRose1101}%
In this third case, for
\( A=\begin{bmat} 1&1\\0&1 \end{bmat}\),
there appear to be only the vectors \(\xv=\pm(1,0)\), aligned along the horizontal axis, for which \(A\xv=\lambda\xv\).
Whereas for symmetric matrices there were always two pairs, here we only appear to have one pair of eigenvectors (\cref{thm:dimee}).
Such degeneracy occurs for matrices on the border between reality \text{and complexity.}
\vspace{1\baselineskip}
\end{figbox}

\end{enumerate}


The first problem of the general lack of orthogonality of the eigenvectors is most clearly seen in the case of triangular matrices (\cref{def:trim}).
The reason is  linked to \cref{thm:rpdet:vi} that the \idx{determinant} of a \idx{triangular matrix} is simply the product of its \idx{diagonal entries}.

\begin{example} 
Find by algebra the eigenvalues and eigenvectors of the \idx{triangular matrix} \(A=\begin{bmat} 2&1\\0&3 \end{bmat}\).
\begin{solution} 
Recall \cref{pro:eeh}. 
\begin{enumerate}
\item Find all eigenvalues by solving the characteristic equation \(\det(A-\lambda I)=0\)\,. 
Here \(\det(A-\lambda I)=\det\begin{bmat} 2-\lambda&1\\0&3-\lambda \end{bmat}\) which, being a triangular matrix, has determinant that is the product of the diagonals, namely \(\det(A-\lambda I)=(2-\lambda)(3-\lambda)\)\,.
This determinant is zero only for eigenvalues \(\lambda=2\) or~\(3\).
These are the diagonal entries in the \text{triangular matrix.}
\item For each eigenvalue, find corresponding eigenvectors by solving the system \((A-\lambda I)\xv=\ov\)\,.
\begin{itemize}
\item For \(\lambda=2\), the system is \(\begin{bmat} 0&1\\0&1 \end{bmat}\xv=\ov\) which requires \(x_2=0\)\,.  That is, all eigenvectors are~\(x_1(1,0)\).
\item For \(\lambda=3\), the system is \(\begin{bmat} -1&1\\0&0 \end{bmat}\xv=\ov\) which requires \(x_1=x_2\)\,.  That is, all eigenvectors are~\(x_2(1,1)\).
\end{itemize}
\end{enumerate}
%\marginpar{\eRose{2}1{0}{3}}%
The eigenvectors corresponding to the different eigenvalues are not orthogonal as their dot product \((1,0)\cdot(1,1)=1+0=1\neq 0\)\,. 
Instead the different eigenvectors are at~\(45^\circ\) to each other.
\end{solution}
\end{example}



\begin{theorem}[triangular matrices] \label{thm:trieig} 
The \idx{diagonal entries} of a \idx{triangular matrix} are the only \idx{eigenvalue}s of the matrix.  
The corresponding \idx{eigenvector}s of distinct \idx{eigenvalue}s are \emph{generally} not \idx{orthogonal}.
\end{theorem}

\begin{proof} 
We detail only the case of upper triangular matrices, as the argument is similar for lower triangular matrices. 
First establish that the diagonal entries are eigenvalues, and second prove there are no others.
Let \(\lambda\) be any value in the diagonal of the matrix~\(A\), and let \(k\)~be the smallest index such that \(a_{k,k}=\lambda\) (this `smallest' caters for duplicated diagonal values).
Let's construct an eigenvector in the form \(\xv=(\hlist x{k-1},1,0,\ldots,0)\).  
Set \(x_k=1\) and \(x_j=0\) for \(j>k\)\,.
Then set 
\begin{eqnarray*}
x_{k-1}&=&-a_{k-1,k}x_k/(a_{k-1,k-1}-\lambda),
\\x_{k-2}&=&-(a_{k-2,k}x_k+a_{k-2,k-1}x_{k-1})/(a_{k-2,k-2}-\lambda),
\\&\vdots&
\\x_1&=&-(a_{1,k}x_k+a_{1,k-1}x_{k-1}+\cdots+a_{1,2}x_2)/(a_{1,1}-\lambda).
\end{eqnarray*}
Since \(k\)~is the smallest index for which \(\lambda=a_{k,k}\) none of the above expressions involve divisions by zero, and so all are well-defined.
Rearranging the above equations shows that this vector~\xv\ satisfies, for  \(\lambda=a_{k,k}\)\,,
\begin{equation*}
\begin{array}{l@{}l@{}l@{}l@{}l@{}l}
(a_{1,1}-\lambda)x_1
&{}+a_{1,2}x_2
&{}+\cdots+a_{1,k-1}x_{k-1}
&{}+a_{1,k}x_k
&{}=0\,,
\\&\hspace{-1em}
(a_{2,2}-\lambda)x_2
&{}+\cdots+a_{2,k-1}x_{k-1}
&{}+a_{1,k}x_k
&{}=0\,,
\\&&\ddots&\vdots
\\&&\hspace{-1em}
(a_{k-1,k-1}-\lambda)x_{k-1}
&{}+a_{k-1,k}x_k
&{}=0\,,
\\&&&\hspace{-1em}
(a_{k,k}-\lambda)x_k
&{}=0\,;
\end{array}
\end{equation*}
that is, \((A-\lambda I)\xv=\ov\)\,.
Rearranging gives \(A\xv=\lambda\xv\) for nonzero eigenvector~\xv\ (nonzero since at least \(x_k\neq 0\)) and corresponding eigenvalue~\(\lambda=a_{k,k}\)\,.

Second, there can be no other eigenvalues.
Every eigenvalue has to have non-trivial solutions, nonzero eigenvectors~\xv, to \((A-\lambda I)\xv=\ov\)\,, which by \cref{thm:detinv} requires \(\det(A-\lambda I)=0\)\,.
But as \(A\)~is upper triangular, matrix \((A-\lambda I)\) is upper triangular and so \cref{thm:rpdet:vi} asserts the determinant
\begin{equation*}
\det(A-\lambda I)=(a_{1,1}-\lambda)(a_{2,2}-\lambda)
\cdots(a_{n,n}-\lambda).
\end{equation*}
This expression is zero iff the eigenvalue~\(\lambda\) is one of the diagonal elements of~\(A\).

As an example of the non-orthogonality of eigenvectors, consider the two eigenvalues \(\lambda_1=a_{1,1}\) and \(\lambda_2=a_{2,2}\), assume distinct, with  corresponding eigenvectors \(\xv_1=(1,0,0,\ldots,0)\) and \(\xv_2=(-a_{1,2}/(a_{1,1}-a_{2,2}),1,0,\ldots,0)\).
Then the dot product \(\xv_1\cdot\xv_2=-a_{1,2}/(a_{1,1}-a_{2,2})\neq 0\) in general (the dot product is zero only in the special case when \(a_{1,2}=0\)).
Since the dot product is generally nonzero, \(\xv_1\) and~\(\xv_2\) are generally not orthogonal.
Similarly for other pairs of eigenvectors corresponding to \text{distinct eigenvalues.}
\end{proof}



\begin{example} \label{eg:trieig}
Use \cref{thm:trieig} to find the eigenvalues, corresponding eigenvectors, and corresponding eigenspaces, of the following triangular matrices.\index{triangular matrix}
\begin{enumerate}
\item \(\eAii=\begin{bmatrix}-3&2&0
\\0&-4&2
\\0&0&4\end{bmatrix}\)
\begin{solution} 
Matrix~\(\eAii\)\ is \idx{upper triangular} so read off the eigenvalues from the diagonal to be \(-3\) and~\(\pm 4\).
\begin{itemize}
\item For \(\lambda=-3\)\,, and by inspection, all eigenvectors are proportional to \((1,0,0)\).
Hence eigenspace \(\EE_{-3}=\Span\{(1,0,0)\}\).
\item For \(\lambda=-4\) we need to solve
\begin{equation*}
(\eAii+4I)\xv=
\begin{bmatrix}1&2&0
\\0&0&2
\\0&0&8\end{bmatrix}\xv=\ov\,.
\end{equation*}
By inspection an eigenvector must be of the form \((x_1,1,0)\).
Then the first line of the system asserts \(x_1+2=0\)\,.
Hence eigenvectors are proportional to \((-2,1,0)\).
That is, the eigenspace \(\EE_{-4}=\Span\{(-2,1,0)\}\).
\item For \(\lambda=+4\) we need to solve
\begin{equation*}
(\eAii-4I)\xv=
\begin{bmatrix}-7&2&0
\\0&-8&2
\\0&0&0\end{bmatrix}\xv=\ov\,.
\end{equation*}
Consider eigenvectors of the form \((x_1,x_2,1)\).
The second line asserts \(-8x_2+2=0\), that is \(x_2=\frac14\).
The first line asserts \(-7x_1+2x_2=0\), that is \(x_1=\frac27x_2=\frac1{14}\).
Hence eigenvectors are proportional to \((\frac1{14},\frac14,1)\).
That is, the eigenspace \(\EE_{4}=\Span\{(\frac1{14},\frac14,1)\}\).
\end{itemize}
\end{solution}


\item \(\eAii=\begin{bmatrix}3&0&0&0
\\-2&-4&0&0
\\-3&1&0&0
\\0&0&-3&1 \end{bmatrix}\)
\begin{solution} 
Matrix~\(\eAii\)\ is lower triangular so read the eigenvalues from the diagonal to be~\(3\), \(-4\), \(0\), and~\(1\).
\begin{itemize}
\item For \(\lambda=1\)\,, by inspection all eigenvectors are of the form \((0,0,0,1)\).
Hence eigenspace \(\EE_{1}=\Span\{(0,0,0,1)\}\).

\item For \(\lambda=0\)\,, seek an eigenvector \((0,0,1,x_4)\) then the last line of the system
\begin{equation*}
(\eAii-0I)\xv=
\begin{bmatrix}3&0&0&0
\\-2&-4&0&0
\\-3&1&0&0
\\0&0&-3&1 \end{bmatrix}\xv=\ov
\end{equation*}
requires \(-3+x_4=0\)\,.  
Hence eigenvectors are proportional to \((0,0,1,3)\).
That is, the eigenspace \(\EE_{0}=\Span\{(0,0,1,3)\}\).

\item For \(\lambda=-4\), seek an eigenvector \((0,1,x_3,x_4)\) then the third line of the system
\begin{equation*}
(\eAii+4I)\xv=
\begin{bmatrix}7&0&0&0
\\-2&0&0&0
\\-3&1&4&0
\\0&0&-3&5 \end{bmatrix}\xv=\ov
\end{equation*}
requires \(1+4x_3=0\)\,, that is \(x_3=-\frac14\).
Then the last line of the system requires \(\frac34+5x_4=0\)\,, that is \(x_4=-\frac3{20}\)\,. 
Hence eigenvectors are proportional to \((0,1,-\frac14,-\frac3{20})\).
That is, the eigenspace \(\EE_{-4}=\Span\{(0,1,-\frac14,-\frac3{20})\}\).

\item For \(\lambda=3\), seek an eigenvector \((1,x_2,x_3,x_4)\) then the second line of the system
\begin{equation*}
(\eAii-3I)\xv=
\begin{bmatrix}0&0&0&0
\\-2&-7&0&0
\\-3&1&-3&0
\\0&0&-3&-2 \end{bmatrix}\xv=\ov
\end{equation*}
requires \(-2-7x_2=0\)\,, that is \(x_2=-\frac27\).
Then the third line of the system requires \(-3-\frac27-3x_3=0\)\,, that is \(x_3=-\frac{23}{21}\)\,. 
Lastly, the last line requires \(\frac{23}7-2x_4=0\)\,, that is \(x_4=\frac{23}{14}\).
Hence eigenvectors are proportional to \((1,-\frac27,-\frac{23}{21},\frac{23}{14})\).
That is, the eigenspace \(\EE_{3}=\Span\{(1,-\frac27,-\frac{23}{21},\frac{23}{14})\}\).

\end{itemize}
\end{solution}


\item \(\eAii=\begin{bmatrix} -1&1&-8&-5&5
\\-3&6&4&-3&0
\\1&-3&1&0&0
\\-7&1&0&0&0
\\-1&0&0&0&0 \end{bmatrix}\)
\begin{solution} 
Matrix~\(\eAii\) is not a triangular matrix (\cref{def:trim}), so \cref{thm:trieig} does not apply.
Row or column swaps could transform~\(C\) to be triangular, but we have not investigated the effect of such swaps on eigenvalues \text{and eigenvectors.}
\end{solution}
\end{enumerate}
\end{example}




\begin{activity}
What are all the eigenvalues of the matrix
% 0+round(toeplitz(1:5,[1 0 0 0 0]).*rand(5))
\begin{equation*}
\begin{bmatrix} 1&0&0&0&0
\\0&1&0&0&0
\\2&2&1&0&0
\\3&3&1&0&0
\\2&2&2&0&1
 \end{bmatrix}?
\end{equation*}
\actposs[4]{\(0,1\)}{\(1\)}{\(0,1,2\)}{\(0,1,2,3\)}
\end{activity}






One consequence of the second part of the proof of \cref{thm:trieig} is that, when counted according to multiplicity, there are precisely \(n\)~eigenvalues of an \(n\times n\) \idx{triangular matrix}.
Correspondingly, the next \cref{sec:eennm} establishes that there are precisely \(n\)~eigenvalues of general \(n\times n\) matrices, provided we count the eigenvalues according to multiplicity and allow \idx{complex eigenvalue}s.









\needlines{16}
\sectionExercises

\begin{exercise}  
Each of the following pictures applies to some specific real matrix, say called~\(A\).
The pictures plot~\(A\xv\) (red) adjoined to the end of \idx{unit vector}s~\xv\ (blue).
By inspection decide whether the matrix, in each case, has real eigenvalues or non-real \idx{complex eigenvalue}s.
%\begin{verbatim}
%a=round(eye(2)*5+randn(2)*5)/10,eig(a),num2str(a(:)')
%\end{verbatim}

\begin{Parts}
\begin{OmitV1}
\item \eRose{0.2}{-1.1}{ 1.3}{ 0.9}
\answer{complex}

\item \eRose{0.9}{0.9}{0.8}{0.2}
\answer{real}

\item \eRose{0.5}{ 0.3}{-0.8}{ 1.3}
\answer{complex}

\item \eRose{-0.7}{-0.2}{ 0.4}{-0.3}
\answer{complex}
\end{OmitV1}

\item \eRose{0.3}{-0.3}{-0.2}{1}
\answer{real}

\item \eRose{0.1}{-0.4}{ 0.8}{ 0.7}
\answer{complex}

\item \eRose{0.8}{-0.7}{-0.3}{ 0.8}
\answer{real}

\item \eRose{0.8}{ 0.3}{-0.6}{1}
\answer{complex}

\end{Parts}
\end{exercise}








\begin{exercise}  
For each of the following \index{triangular matrix}triangular matrices, write down all \idx{eigenvalue}s and then find the corresponding \idx{eigenspace}s.
Show your working.
%\begin{verbatim}
%format rat, n=3
%A=0+tril(round(randn(n)*3)), if rand>0.5, A=A', end, [V,D]=eig(A); DVt=[diag(D) diag(1./max(abs(V)))*V']
%\end{verbatim}
\begin{Parts}
\item \(\begin{bmatrix} 2 & 0
\\-1 & 4 \end{bmatrix}\)
\answer{\(\EE_{4}=\Span\{(0,1)\}\), 
\(\EE_{2}=\Span\{(1,\frac12)\}\).}

\item \(\begin{bmatrix} 2 & 0
\\-3 & 2 \end{bmatrix}\)
\answer{\(\EE_{2}=\Span\{(0,1)\}\).}

\item \(\begin{bmatrix} -1 & 0 & 3
\\0 & -1 & 2
\\0 & 0 & -5 \end{bmatrix}\)
\answer{\(\EE_{-1}=\Span\{(1,0,0),(0,1,0)\}\), 
\(\EE_{-5}=\Span\{(-3,-2,4)\}\).}

\item \(\begin{bmatrix} 0 & 0 & 0
\\-3 & -4 & 0
\\1 & 5 & -2 \end{bmatrix}\)
\answer{\(\EE_{-2}=\Span\{(0,0,1)\}\), 
\(\EE_{-4}=\Span\{(0,2,-5)\}\), 
\(\EE_{0}=\Span\{(8,-6,-11)\}\).}

\begin{OmitV1}
\item \(\begin{bmatrix} 1 & -5 & 0
\\0 & 0 & -4
\\0 & 0 & 0 \end{bmatrix}\)
\answer{\(\EE_{1}=\Span\{(1,0,0)\}\), 
\(\EE_{0}=\Span\{(5,1,0)\}\).}

\item \(\begin{bmatrix} 0 & 0 & -2
\\0 & 4 & 1
\\-3 & -3 & -1 \end{bmatrix}\)
\answer{This matrix is not triangular so we cannot answer (as yet).}

\item \(\begin{bmatrix} -1 & 0 & 0
\\-2 & 2 & 0
\\-2 & -1 & -1 \end{bmatrix}\)
\answer{\(\EE_{-1}=\Span\{(0,0,1)\}\), 
\(\EE_{2}=\Span\{(0,3,-1)\}\).}

\item \(\begin{bmatrix} -2 & 4 & -2 & -2
\\0 & -2 & 1 & 7
\\0 & 0 & -3 & 1
\\0 & 0 & 0 & 2 \end{bmatrix}\)
\answer{\(\EE_{-2}=\Span\{(1,0,0,0)\}\), 
\(\EE_{-3}=\Span\{(6,-1,1,0)\}\), 
\(\EE_{2}=\Span\{(6,9,1,5)\}\).}
\end{OmitV1}

\item \(\begin{bmatrix} 8 & -2 & 3 & 2
\\0 & -6 & 1 & -2
\\0 & 0 & 3 & -2
\\0 & 0 & 0 & 0 \end{bmatrix}\)
\answer{\(\EE_{8}=\Span\{(1,0,0,0)\}\), 
\(\EE_{-6}=\Span\{(1,7,0,0)\}\), 
\(\EE_{3}=\Span\{(-5,1,9,0)\}\), 
\(\EE_{0}=\Span\{(-5,-2,6,9)\}\).}

\item \(\begin{bmatrix} 0 & -2 & -5 & 2
\\0 & 7 & -1 & 2
\\0 & 0 & 3 & -4
\\0 & 0 & 0 & 3 \end{bmatrix}\)
\answer{\(\EE_{0}=\Span\{(1,0,0,0)\}\), 
\(\EE_{7}=\Span\{(-2,7,0,0)\}\), 
\(\EE_{3}=\Span\{(-22,3,12,0)\}\).}

\end{Parts}
\end{exercise}


\index{eigenvalue|)}
\index{eigenvector|)}
\index{eigenspace|)}





%%!TEX root = ../larxxia.tex


\section{Find eigenvalues and eigenvectors of matrices}
\label{sec:eennm}

\secttoc
\begin{comment}
\pooliv{\S4.3}  \layiv{\S5.2} \holti{\S6.1}  \cite[Ch.~9]{Chartier2015}.
Need exercises for the fitting exponentials subsection??
\end{comment}


Given the additional determinant methods of \autoref{ch:ddm}, this section begins exploring the properties and some applications of the eigen-problem \(A\xv=\lambda\xv\) for general matrices~\(A\).
We establish that there are generally \(n\)~eigenvalues of an \(n\times n\) matrix, albeit possibly complex valued, and that repeated eigenvalues are sensitive to errors.
Applications include population modelling, connecting to the computation of \svd{}s, and fitting exponentials to real data.




\subsection{A characteristic equation gives eigenvalues}
\label{sec:cege}

\begin{quoted}{J.~H. Wilkinson, 1984 \cite[p.103]{Higham1996}}
The Fundamental Theorem of Algebra asserts that every polynomial equation over the complex field has a root.  
It is almost beneath the dignity of such a majestic theorem to mention that in fact it has precisely \(n\)~roots.
\end{quoted}


Recall that eigenvalues~\(\lambda\) and non-zero eigenvectors~\xv\ of a square matrix~\(A\) must satisfy \((A-\lambda I)\xv=\ov\)\,.
\autoref{thm:detinv} then implies the \idx{eigenvalue}s of a \idx{square matrix} are precisely the solutions of the \bfidx{characteristic equation} \(\det(A-\lambda I)=0\)\,.  

\begin{theorem} \label{thm:geecp}
For every \(n\times n\) \idx{square matrix}~\(A\) we call \(\det(A-\lambda I)\) the \bfidx{characteristic polynomial} of~\(A\):
\footnote{Alternatively, many call \(\det(\lambda I-A)\) the characteristic polynomial, as does \script.  
The distinction is immaterial as,  for an \(n\times n\) matrix~\(A\) and by \autoref{thm:basicdet:iii} with multiplicative factor \(k=-1\)\,, the only difference in the determinant is a factor of~\((-1)^n\).
In \script, \index{poly()@\texttt{poly()}}\texttt{poly(A)} computes the characteristic polynomial of the matrix,  \(\det(\lambda I-A)\), which might be useful for exercises, but is rarely useful in practice due to poor conditioning.}
\begin{itemize}
\item the characteristic polynomial of~\(A\) is a polynomial of \(n\)th~degree in~\(\lambda\);
\item  there are at most \(n\)~\idx{distinct eigenvalues} of~\(A\).
\end{itemize}
\end{theorem}

\begin{proof} 
% Do not do this as we want the induction for the next theorem's proof.
%Consider the characteristic polynomial
%\begin{equation*}
%\det(A-\lambda I)=
%\det\begin{bmatrix} a_{11}-\lambda&a_{12}&\cdots&a_{1n}
%\\a_{21}&a_{22}-\lambda&&a_{2n}
%\\\vdots&&\ddots&\vdots
%\\a_{n1}&a_{n2}&\cdots&a_{nn}-\lambda \end{bmatrix}.
%\end{equation*}
%Recall \autoref{thm:detnfac} asserts that the determinant of an \(n\times n\) matrix is a sum of \(n!\)~terms where each term is a product of \(n\)~factors where each factor in a term comes from a different row and column of the matrix.
%Being a sum of products of these entries, \(\det(A-\lambda I)\) is a polynomial in~\(\lambda\).
%Since each term is a product of \(n\)~factors, and each factor is an entry in the matrix~\(A-\lambda I\)\,, and each entry is at most linear in~\(\lambda\), then the polynomial is of degree at most~\(n\) in~\(\lambda\).
%One of the \(n!\)~terms in the sum is the product along the diagonal \((a_{11}-\lambda)(a_{22}-\lambda)\cdots(a_{nn}-\lambda)\): consequently the polynomial has a term in~\(\lambda^n\) and so is of degree~\(n\).
Use induction on the size of the matrix.
First, for \(1\times 1\) matrix \(A=\begin{bmatrix} a_{11} \end{bmatrix}\), the determinant \(\det(A-\lambda I)=\det \begin{bmatrix} a_{11}-\lambda \end{bmatrix}=a_{11}-\lambda\) which is of  degree one in~\(\lambda\).
Second, assume that for all \((n-1)\times (n-1)\) matrices~\(A\), \(\det(A-\lambda I)\) is a polynomial of degree~\(n-1\) in~\(\lambda\).
Then use the Laplace Expansion \autoref{thm:letdet} to give the first row expansion, in terms of minors of~\(A\) and~\(I\),
\begin{eqnarray*}
\det(A-\lambda I)
&=&(a_{11}-\lambda)\det (A_{11}-\lambda I_{11})
\\&&{}
-a_{12}\det (A_{12}-\lambda I_{12})
\\&&{}
+\cdots-(-1)^{n}a_{1n}\det (A_{1n}-\lambda I_{1n}).
\end{eqnarray*}
Now the minor~\(I_{11}\) is precisely the \((n-1)\times(n-1)\) identity, and hence by assumption \(\det(A_{11}-\lambda I_{11})\) is a polynomial of degree~\((n-1)\) in~\(\lambda\).
But differently, the minors~\(I_{12},\ldots,I_{1n}\) have two of the ones removed from the \(n\times n\)~identity and hence \(\lambda I_{12},\ldots,\lambda I_{1n}\) each have only \((n-2)\)~\(\lambda\)s: since each term in a determinant is a product of distinct entries of the matrix (\autoref{thm:detnfac}) it follows that for \(j\geq2\) the determinant \(\det(A_{1j}-\lambda I_{1j})\) is a polynomial in~\(\lambda\) of degree\({}\leq n-2\)\,.
Consequently, the first row expansion of 
\begin{eqnarray*}
\det(A-\lambda I)
&=&(a_{11}-\lambda)(\text{poly degree }n-1)
\\&&{}
-a_{12}(\text{poly degree }\leq n-2)
\\&&{}
+\cdots-(-1)^{n}a_{1n}(\text{poly degree }\leq n-2)
\\&=&(a_{11}-\lambda)(\text{poly degree }n-1)
\\&&{}
+(\text{poly degree }\leq n-2)
\\&=&(\text{poly degree }n)
\end{eqnarray*}
as the highest power of~\(\lambda\) cannot be cancelled by any term.
Induction thus implies the characteristic polynomial of an \(n\times n\) matrix~\(A\) is a polynomial of \(n\)th~degree in~\(\lambda\).

Lastly, because the characteristic polynomial of~\(A\) is of \(n\)th~degree in~\(\lambda\), the Fundamental Theorem of Algebra asserts that the polynomial has at most \(n\)~roots (possibly complex).  
Hence there are at most \(n\)~distinct eigenvalues.
\end{proof}



\begin{activity}
A given matrix has eigenvalues of~\(-7\), \(-1\), \(3\), \(4\)  and~\(6\).
The matrix must be of size \(n\times n\) for~\(n\) at least which of the following? (Select the smallest valid answer.)
\actposs[4]5467
%\partswidth=5em
%\begin{parts}
%\item \(4\)
%\item \(5\)
%\item \(6\)
%\item \(7\)
%\end{parts}
\end{activity}



\begin{example} \label{eg:}
Find the \idx{characteristic polynomial} of each of the following matrices.  
Where in the coefficients of the polynomial can you see the determinant? and the sum of the diagonal elements?
\begin{enumerate}
\item \(\eAii=\begin{bmatrix} 1&-1\\-2&4 \end{bmatrix}\)
\begin{solution} 
The characteristic polynomial is \(\det(\eAii-\lambda I)
=\det\begin{bmatrix} 1-\lambda&-1\\-2&4-\lambda \end{bmatrix}
=(1-\lambda)(4-\lambda)-2
=\lambda^2-5\lambda+2\)\,.
Now \(\det \eAii=4-2=2\) which is the coefficient of the constant term in this polynomial.
Whereas the sum of the diagonal of~\(\eAii\) is \(1+4=5\) which is the negative of the \(\lambda\)~coefficient in the polynomial.
\end{solution}

\item \(\eAii=\begin{bmatrix}4&-2&1
\\1&-2&0
\\8&2&6 \end{bmatrix}\)
\begin{solution} 
The characteristic polynomial is 
\begin{eqnarray*}
\det(\eAii-\lambda I)
&=&\det\begin{bmatrix} 4-\lambda&-2&1
\\1&-2-\lambda&0
\\8&2&6-\lambda \end{bmatrix}
\\&=&(4-\lambda)(-2-\lambda)(6-\lambda)+0+2
\\&&{}
-(-2-\lambda)8-0-(-2)(6-\lambda)
\\&=&-48+4\lambda+8\lambda^2-\lambda^3+2
\\&&{}
+16+8\lambda+12+2\lambda
\\&=&-\lambda^3+8\lambda^2+2\lambda-18\,.
\end{eqnarray*}
Now \(\det \eAii=4(-2)6+0+2-(-2)8-0-(-2)6=-48+2+16+12=-18\) which is the coefficient of the constant term in the polynomial.
Whereas the sum of the diagonal of~\(\eAii\) is \(4-2+6=8\) which is the \(\lambda^2\)~coefficient in the polynomial.
\end{solution}

\end{enumerate}
\end{example}




These observations about the coefficients in the characteristic polynomials leads to the next theorem.

\begin{theorem} \label{thm:charpolyc}
For every \(n\times n\) matrix~\(A\), the product of the \idx{eigenvalue}s equals~\(\det A\) and equals the \idx{constant term} in the \idx{characteristic polynomial}.  
\begin{aside}
This optional theorem helps establish the nature of a characteristic polynomial.
\end{aside}
The sum of the \idx{eigenvalue}s equals \((-1)^{n-1}\)~times the \idx{coefficient} of~\(\lambda^{n-1}\) in the \idx{characteristic polynomial} and equals the \bfidx{trace} of the matrix, defined as the sum of the diagonal elements \(a_{11}+a_{22}+\cdots+a_{nn}\)\,.
% that is, \(\sum_{j=1}^n a_{jj}\).
\end{theorem}

\begin{proof} 
\autoref{thm:geecp}, and its proof, establishes that the characteristic polynomial has the form
\begin{eqnarray}
\det(A-\lambda I)&=&c_0+c_1\lambda+\cdots+c_{n-1}\lambda^{n-1}+(-1)^n\lambda^n
\nonumber\\&=&(\lambda_1-\lambda)(\lambda_2-\lambda)\cdots(\lambda_n-\lambda),
\label{eq:charpolyc}
\end{eqnarray}
where the second equality follows from the Fundamental Theorem of Algebra that an \(n\)th~degree polynomial factors into \(n\)~linear factors (albeit possibly complex).
First, substitute \(\lambda=0\) and this equation~\eqref{eq:charpolyc} gives \(\det A=c_0=\lambda_1\lambda_2\cdots\lambda_n\)\,, as required.

Second, consider the term~\(c_{n-1}\lambda^{n-1}\) in equation~\eqref{eq:charpolyc}.
From the factorisation on the right-hand side, the \(\lambda^{n-1}\)~term arises as
\begin{eqnarray*}
c_{n-1}\lambda^{n-1}&=&\lambda_1(-\lambda)^{n-1}
+(-\lambda)\lambda_2(-\lambda)^{n-2}
+\cdots+(-\lambda)^{n-1}\lambda_n
\\&=&(-1)^{n-1}(\lambda_1+\lambda_2+\cdots+\lambda_n)\lambda^{n-1},
\end{eqnarray*}
and hence the coefficient \(c_{n-1}=(-1)^{n-1}(\lambda_1+\lambda_2+\cdots+\lambda_n)\), as required.
Now use induction to prove the trace formula.
Recall that the proof of \autoref{thm:geecp} establishes that
\begin{eqnarray*}
\det(A-\lambda I)
&=&(a_{11}-\lambda)\det (A_{11}-\lambda I_{11})
\\&&{}
+(\text{poly degree }\leq n-2).
\end{eqnarray*}
Assume the trace formula holds for \((n-1)\times(n-1)\) matrices such as the minor~\(A_{11}\). Then the previous identity gives
\begin{eqnarray*}
\det(A-\lambda I)
&=&(a_{11}-\lambda)\big[(-1)^{n-1}\lambda^{n-1}
\\&&\quad{}
+(-1)^{n-2}(a_{22}+\cdots+a_{nn})\lambda^{n-2} 
\\&&\quad{}
+(\text{poly degree }\leq n-3)\big]
\\&&{}
+(\text{poly degree }\leq n-2)
\\&=&(-1)^{n}\lambda^{n}
%\\&&{}
+(-1)^{n-1}(a_{11}+a_{22}+\cdots+a_{nn})\lambda^{n-1} 
\\&&{}
+(\text{poly degree }\leq n-2).
\end{eqnarray*}
Hence the coefficient \(c_{n-1}=(-1)^{n-1}(a_{11}+a_{22}+\cdots+a_{nn})\).
Since the formula holds for the basic case \(n=1\), namely \(c_0=+a_{11}\)\,, induction implies the sum of the eigenvalues
\(\lambda_1+\lambda_2+\cdots+\lambda_n=a_{11}+a_{22}+\cdots+a_{nn}\)\,, the trace of the matrix~\(A\).
\end{proof}


\begin{activity}
What is the trace of the matrix
\begin{equation*}
\begin{bmatrix} 4&5&-4&3
\\-2&2&-5&-1
\\-1&2&2&-6
\\-13&4&3&-1 \end{bmatrix}?
\end{equation*}
\actposs[4]{\(7\)}{\(-13\)}{\(-12\)}{\(8\)}
%\partswidth=5em
%\begin{parts}
%\item \(-13\)
%\item \(-12\)
%\item \(7\)
%\item \(8\)
%\end{parts}
\end{activity}



\begin{example} \label{eg:}
\begin{enumerate}
\item What are the two highest order terms and the constant term in the \idx{characteristic polynomial} of the matrix
\begin{equation*}
\eAii=\begin{bmatrix}-2&-1&3&-2
\\-1&3&-2&2
\\2&-3&0&1
\\0&1&0&-3  \end{bmatrix}.
\end{equation*}
\begin{solution} 
First compute the determinant using the Laplace expansion (\autoref{thm:letdet}).  
The two zeros in the last row suggest a last row expansion:
\begin{eqnarray*}
\det \eAii&=&(-1)^61\det\begin{bmatrix}-2&3&-2
\\-1&-2&2
\\2&0&1 \end{bmatrix}
\\&&{}
+(-1)^8(-3)\det\begin{bmatrix}-2&-1&3
\\-1&3&-2
\\2&-3&0\end{bmatrix}
\\&=&(4+12+0-8-0+3)
\\&&{}-3(0+4+9-18+12-0)
=-10\,.
\end{eqnarray*}
This is the constant term in the characteristic polynomial.
Second, the trace of~\(\eAii\) is \(-2+3+0-3=-2\) so the cubic coefficient in the characteristic polynomial is~\((-1)^3(-2)=2\)\,.
That is, the characteristic polynomial of~\(\eAii\) is of the form \(\lambda^4+2\lambda^3+\cdots-10\)\,.
\end{solution}

\item After laborious calculation you find the \idx{characteristic polynomial} of the matrix
\begin{equation*}
\eAii=\begin{bmatrix}-2&5&-3&-1&2
\\-2&-5&-1&-1&3
\\1&4&-2&1&-7
\\1&-5&1&4&-5
\\-1&0&3&-3&1\end{bmatrix}
\end{equation*}
is \(-\lambda^5+2\lambda^4-3\lambda^3+234\lambda^2+884\lambda+1564\)\,.  
Could this polynomial be correct?
\begin{solution} 
No, because the trace of~\(\eAii\) is \(-2-5-2+4+1=-4\) so the coefficient of the \(\lambda^4\)~term must be~\((-1)^4(-4)=-4\) instead of the calculated~\(2\).
\end{solution}

\item After much calculation you find the \idx{characteristic polynomial} of the matrix
\begin{equation*}
\eAii=\begin{bmatrix}0&0&3&1&0&0
\\0&0&0&0&0&1
\\0&0&-4&0&3&0
\\-5&0&3&0&0&0
\\0&0&0&0&0&2
\\0&0&0&-6&0&0\end{bmatrix}
\end{equation*}
is \(\lambda^6 +4\lambda^5 +5\lambda^4 +20\lambda^3 +108\lambda^2 -540\lambda+668\)\,.  
Could this polynomial be correct?
\begin{solution} 
No. By the column of zeros in~\(\eAii\), \(\det \eAii\) must be zero instead of the calculated~\(668\).
\end{solution}

\item What are the two highest order terms and the constant term in the \idx{characteristic polynomial} of the matrix
\begin{equation*}
\eAii=\begin{bmatrix}0&4&0&0&3&0
\\-2&0&0&1&0&-2
\\0&0&0&-1&0&0
\\0&0&-5&0&-4&3
\\0&2&-3&0&-4&0
\\0&-3&0&0&0&0\end{bmatrix}.
\end{equation*}
\begin{solution} 
First compute the determinant using the Laplace expansion (\autoref{thm:letdet}).  
The nearly all zeros in the last row suggests starting with a last row expansion (although others are just as good):
\begin{eqnarray*}
\det \eAii&=&(-1)^8(-3)\det\begin{bmatrix}0&0&0&3&0
\\-2&0&1&0&-2
\\0&0&-1&0&0
\\0&-5&0&-4&3
\\0&-3&0&-4&0\end{bmatrix}
\\&&(\text{using a 3rd row expansion})
\\&=&-3(-1)^6(-1)\det\begin{bmatrix}0&0&3&0
\\-2&0&0&-2
\\0&-5&-4&3
\\0&-3&-4&0\end{bmatrix}
\\&&(\text{using a 1st column expansion})
\\&=&3(-1)^3(-2)\det\begin{bmatrix}0&3&0
\\-5&-4&3
\\-3&-4&0\end{bmatrix}
\\&&(\text{using a 3rd column expansion})
\\&=&6(-1)^53\det\begin{bmatrix}0&3
\\-3&-4\end{bmatrix}
\\&=&-18(0+9)=-162\,.
\end{eqnarray*}
This is the constant term in the characteristic polynomial.
Second, the trace of~\(\eAii\) is \(0+0+0+0-4+0=-4\) so the quintic coefficient in the characteristic polynomial is~\((-1)^5(-4)=4\)\,.
That is, the characteristic polynomial of~\(\eAii\) is of the form \(\lambda^6+4\lambda^5+\cdots-162\)\,.
\end{solution}

\end{enumerate}
\end{example}


Recall that an important characteristic of an eigenvalue is their multiplicity.
The following definition of \emph{multiplicity} generalises to all matrices the somewhat different \autoref{def:eigsymult} that applies to only symmetric matrices.
For symmetric matrices the definitions are equivalent.


\begin{definition} \label{def:eigmult}
An \idx{eigenvalue}~\(\lambda_0\) of a matrix~\(A\) is said to have \bfidx{multiplicity}~\(m\) if the \idx{characteristic polynomial} factorises to \(\det(A-\lambda I)=(\lambda-\lambda_0)^mg(\lambda)\) with \(g(\lambda_0)\neq0\)\,, and \(g(\lambda)\)~is a polynomial of degree~\(n-m\)\,.
Every eigenvalue of multiplicity \(m\geq2\) is also called a \bfidx{repeated eigenvalue}.
\end{definition}



\begin{activity}
A given matrix~\(A\) has characteristic polynomial \(\det(A-\lambda I)=(\lambda+2)\lambda^2(\lambda-2)^3(\lambda-3)^4\).
The eigenvalue~\(2\) has what multiplicity?
\actposs[4]{three}{one}{two}{four}
%\partswidth=5em
%\begin{parts}
%\item one
%\item two
%\item three\actans
%\item four
%\end{parts}
\end{activity}



\begin{example} \label{eg:faem}
Use the \idx{characteristic polynomial}s for each of the following matrices to find all \idx{eigenvalue}s and their \idx{multiplicity}.
\begin{enumerate}
\item\label{eg:faem:a} \(\eAii=\begin{bmatrix} 3&1\\0&3 \end{bmatrix}\)
\begin{solution}The characteristic equation is \(\det(\eAii-\lambda I)
=\det{\small\begin{bmatrix} 3-\lambda&1\\0&3-\lambda \end{bmatrix}}
=(3-\lambda)^2-0\cdot1=(\lambda-3)^2=0\). 
The only eigenvalue is \(\lambda=3\) with multiplicity two.
(Since this is an upper triangular matrix, \autoref{thm:trieig} asserts the eigenvalues are simply the diagonal elements, namely~\(3\) twice.)
\end{solution}

\item\label{eg:faem:b} \(\eAii=\begin{bmatrix}-1&1&-2
\\-1&0&-1
\\0&-3&1 \end{bmatrix}\)
\begin{solution}The characteristic equation is 
\begin{eqnarray*}
\det(\eAii-\lambda I)
&=&\det\begin{bmatrix}-1-\lambda&1&-2
\\-1&-\lambda&-1
\\0&-3&1-\lambda\end{bmatrix}
\\&=&(1+\lambda)\lambda(1-\lambda)+0-6
\\&&{}
-0+3(1+\lambda)+(1-\lambda)
\\&=&-\lambda^3+3\lambda-2
\\&=&-(\lambda-1)^2(\lambda+2)=0\,.
\end{eqnarray*}
Eigenvalues are \(\lambda=1\) with multiplicity two, and \(\lambda=-2\) with multiplicity one.
\end{solution}


\item\label{eg:faem:c} \(\eAii=\begin{bmatrix}-1&0&-2
\\0&-3&2
\\0&-2&1\end{bmatrix}\)
\begin{solution}The characteristic equation is 
\begin{eqnarray*}
\det(\eAii-\lambda I)
&=&\det\begin{bmatrix}-1-\lambda&0&-2
\\0&-3-\lambda&2
\\0&-2&1-\lambda\end{bmatrix}
\\&=&(-1-\lambda)\det\begin{bmatrix} -3-\lambda&2
\\-2&1-\lambda \end{bmatrix}
\\&=&-(1+\lambda)\big[(-3-\lambda)(1-\lambda)+4\big]
\\&=&-(\lambda+1)\big[\lambda^2+2\lambda+1\big]
\\&=&-(\lambda+1)^3=0\,.
\end{eqnarray*}
The only eigenvalue is \(\lambda=-1\) with multiplicity three.
\end{solution}

\item \(\eAii=\begin{bmatrix}2&0&-1
\\-5&3&-5
\\5&-2&-2\end{bmatrix}\)
\begin{solution}The characteristic equation is 
\begin{eqnarray*}
\det(\eAii-\lambda I)
&=&\det\begin{bmatrix}2-\lambda&0&-1
\\-5&3-\lambda&-5
\\5&-2&-2-\lambda\end{bmatrix}
\\&=&(2-\lambda)(3-\lambda)(-2-\lambda)+0-10
\\&&{}
+5(3-\lambda)-10(2-\lambda)-0
\\&=&-\lambda^3+3\lambda^2+9\lambda-27
\\&=&-(\lambda-3)^2(\lambda+3)=0\,.
\end{eqnarray*}
Eigenvalues are \(\lambda=3\) with multiplicity two, and \(\lambda=-3\) with multiplicity one.
\end{solution}

\item \(E=\begin{bmatrix} 0&1\\-1&1 \end{bmatrix}\)
\begin{solution} 
The characteristic equation is
\begin{eqnarray*}
\det(E-\lambda I)
&=&\det\begin{bmatrix} -\lambda&1\\-1&1-\lambda \end{bmatrix}
\\&=&-\lambda(1-\lambda)+1
\\&=&\lambda^2-\lambda+1=0\,.
\end{eqnarray*}
This quadratic equation does not factor easily so use the formula 
\begin{equation*}
\lambda=\frac{-b\pm\sqrt{b^2-4ac}}{2a}
=\frac{1\pm\sqrt{1-4}}{2}
%=\tfrac12(1\pm\sqrt{-3})
=\frac12\pm\frac{\sqrt3}2\i\,,
\end{equation*}
for \(\i=\sqrt{-1}\), are two eigenvalues (complex valued) each of multiplicity one.
\end{solution}

\end{enumerate}
\end{example}


%\begin{verbatim}
%for i=1:99999,a=round(randn(4)*3);l=eig(a);if min(abs(diff(l)))<1e-6,a=a,lam=l,break,end,end
%\end{verbatim}


\begin{example} \label{eg:multeval}
Use \index{eig()@\texttt{eig()}}\verb|eig()| in \script\  to find the \idx{eigenvalue}s and their multiplicity for the following matrices.
Recall (\autoref{tbl:mtlbeig}) that executing just \verb|eig(A)| gives a column vector of eigenvalues of~\(A\), repeated according to their \idx{multiplicity}. 
\begin{enumerate}
\item \(\begin{bmatrix}2&2&-1
\\0&1&-2
\\0&-1&0\end{bmatrix}\)
\begin{solution} 
Execute \verb|eig([2 2 -1;0 1 -2;0 -1 0])| to get
\setbox\ajrqrbox\hbox{\qrcode{% eigenvalue code
eig([2 2 -1;0 1 -2;0 -1 0])
}}%
\marginpar{\usebox{\ajrqrbox}}%
\begin{verbatim}
ans =
   2
   2
  -1
\end{verbatim}
So the eigenvalue \(\lambda=2\) has multiplicity two, and the eigenvalue \(\lambda=-1\) has multiplicity one.
\end{solution}

\item \(\begin{bmatrix}-2&-2&-5&0
\\0&-2&2&1
\\-1&1&0&-1
\\-2&1&4&0\end{bmatrix}\)
\begin{solution} 
In \script\ execute
\begin{verbatim}
eig([-2 -2 -5 0
  0 -2 2 1
 -1 1 0 -1
 -2 1 4 0])
\end{verbatim}
\setbox\ajrqrbox\hbox{\qrcode{% eigenvalue code
eig([-2 -2 -5 0
  0 -2 2 1
 -1 1 0 -1
 -2 1 4 0])
}}%
\marginpar{\usebox{\ajrqrbox}}%
to get
\begin{verbatim}
ans =
  -3.0000 + 0.0000i
  -3.0000 + 0.0000i
   1.0000 + 1.4142i
   1.0000 - 1.4142i
\end{verbatim}
There are two complex-valued eigenvalues, evidently \(1\pm\sqrt2\i\), each of multiplicity one, and also the (real) eigenvalue \(\lambda=-3\) which has multiplicity two.
\end{solution}

\item \(\begin{bmatrix}3&-1&-2&1&-2
\\0&0&-2&-2&0
\\2&1&1&1&-1
\\-1&-3&0&1&2
\\2&-2&1&0&3\end{bmatrix}\)
\begin{solution} 
In \script\ execute
\begin{verbatim}
eig([3 -1 -2 1 -2
  0 0 -2 -2 0
  2 1 1 1 -1
 -1 -3 0 1 2
  2 -2 1 0 3])
\end{verbatim}
\setbox\ajrqrbox\hbox{\qrcode{% eigenvalue code
eig([3 -1 -2 1 -2
  0 0 -2 -2 0
  2 1 1 1 -1
 -1 -3 0 1 2
  2 -2 1 0 3])
}}%
\marginpar{\usebox{\ajrqrbox}}%
to get
\begin{verbatim}
ans =
   2.0000 + 2.8284i
   2.0000 - 2.8284i
   4.0000 + 0.0000i
  -0.0000 + 0.0000i
  -0.0000 - 0.0000i
\end{verbatim}
There are three eigenvalues of multiplicity one, namely \(4\) and \(2\pm\sqrt8\i\)\,.  
The last two rows appear to be the eigenvalue \(\lambda=0\) with multiplicity two.
\end{solution}

\item \(\begin{bmatrix}-1&0&0&0
\\-1&2&-3&3
\\3&1&-1&0
\\0&3&-2&1\end{bmatrix}\)
\begin{solution} 
In \script\ execute
\begin{verbatim}
eig([-1 0 0 0
 -1 2 -3 3
  3 1 -1 0
  0 3 -2 1])
\end{verbatim}
\setbox\ajrqrbox\hbox{\qrcode{% eigenvalue code
eig([-1 0 0 0
 -1 2 -3 3
  3 1 -1 0
  0 3 -2 1])
}}%
\marginpar{\usebox{\ajrqrbox}}%
to get
\begin{verbatim}
ans =
   4.0000 + 0.0000i
  -1.0000 + 0.0000i
  -1.0000 - 0.0000i
  -1.0000 + 0.0000i
\end{verbatim}
There is one eigenvalue of multiplicity one, \(\lambda=4\).  
The last three rows appear to be the eigenvalue \(\lambda=-1\) with multiplicity three.
\end{solution}

\end{enumerate}
\end{example}





To find eigenvalues and eigenvectors, the following restates \autoref{pro:eeh} with a little more information, and now empowered to address larger matrices upon using the determinant tools from \autoref{ch:ddm}.

\begin{procedure}[eigenvalues and eigenvectors] \label{pro:geneig}
To find by hand \idx{eigenvalue}s and \idx{eigenvector}s of a (small) \idx{square matrix}~\(A\):
\begin{enumerate}
%\item determine the \idx{characteristic polynomial} of~\(A\), \(\det(A-\lambda I)\);
\item find all \idx{eigenvalue}s (possibly complex) by solving the \bfidx{characteristic equation} of~\(A\), \(\det(A-\lambda I)=0\)\,;
\item for each \idx{eigenvalue}~\(\lambda\), solve the \idx{homogeneous} linear equation \((A-\lambda I)\xv=\ov\) to find the \idx{eigenspace}~\(\EE_\lambda\) of all \idx{eigenvector}s (together with~\ov);
\item write each eigenspace as the \idx{span} of a few chosen \idx{eigenvector}s  (\autoref{def:basis} calls such a set a \idx{basis}).
\end{enumerate}
Since, for an \(n\times n\) matrix, the \idx{characteristic polynomial} is of \(n\)th~degree in~\(\lambda\) (\autoref{thm:geecp}), there are \(n\)~\idx{eigenvalue}s (when counted according to \idx{multiplicity} and allowing \idx{complex eigenvalue}s).
\end{procedure}

Correspondingly, the following restates the computational procedure of \autoref{sec:cee}, but slightly more generally: the extra generality caters for non-symmetric matrices.

\begin{compute}
For a given square matrix~\verb|A|, execute \verb|[V,D]=eig(A)|, then the diagonal entries of~\verb|D|, \verb|diag(D)|, are the \idx{eigenvalue}s of~\verb|A|. 
Corresponding to the eigenvalue~\verb|D(j,j)| is an  \idx{eigenvector} \(\vv_j=\verb|V(:,j)|\), the \(j\)th~column of~\verb|V|.  
\footnote{Be aware that \script\ does not use the determinant to find the eigenvalues, nor does it solve the linear equations to find eigenvectors.  
For any but the smallest matrices such a `by hand' approach takes far too long.  
Instead, to find eigenvalues and eigenvectors, just as for the \svd, \script\ uses an intriguing iteration based upon what is called the \idx{QR factorisation}\ifcsname ch:qrfuma\endcsname\ (\autoref{ch:qrfuma})\fi.}
If an \idx{eigenvalue} is repeated in the diagonal of~\verb|D| (multiplicity more than one), then the corresponding columns of~\verb|V| span the \idx{eigenspace} 
(and, as \autoref{sec:lisb} discusses, when the column vectors have a property called \idx{linear independence} then they form a so-called \idx{basis} for the eigenspace). 
\end{compute}



\begin{activity}
For the matrix \(A=\begin{bmatrix}2&0&-1
\\-5&3&-5
\\5&-2&-2\end{bmatrix}\), which one of the following vectors satisfy \((A-3I)\xv=\ov\) and hence is an eigenvector of~\(A\) corresponding to eigenvalue~\(3\)?
\actposs{\(\xv=(1,5,-1)\)}{\(\xv=(1,5,5)\)}{\(\xv=(0,1,0)\)}{\(\xv=(-1,0,1)\)}
%\begin{enumerate}
%\item \(\xv=(1,5,5)\)
%\item \(\xv=(0,1,0)\)
%\item \(\xv=(-1,0,1)\)
%\item \(\xv=(1,5,-1)\)
%\end{enumerate}
\end{activity}




\begin{example} \label{eg:faespm}
Find the eigenspaces corresponding to the eigenvalues found for the first three matrices of \autoref{eg:faem}.
\begin{itemize}
\item[\ref{eg:faem:a}.]
\(A=\begin{bmatrix} 3&1\\0&3 \end{bmatrix}\)
\begin{solution} 
The only eigenvalue is \(\lambda=3\) with multiplicity two.
Its eigenvectors~\xv\ satisfy
\begin{equation*}
(A-3 I)\xv=\begin{bmatrix} 0&1\\0&0 \end{bmatrix}\xv=\ov\,.
\end{equation*}
The second component of this equation is the trivial \(0=0\)\,.  The first component of the equation is \(0x_1+1x_2=0\), hence \(x_2=0\)\,.
All eigenvectors are of the form \(\xv=(1,0)x_1\)\,.  
That is the eigenspace  \(\EE_3=\Span\{(1,0)\}\).

In \script, executing \verb|[V,D]=eig([3 1;0 3])| gives us
\begin{verbatim}
V =
   1.0000  -1.0000
   0.0000   0.0000
D =
   3   0
   0   3
\end{verbatim}
Diagonal matrix~\verb|D| confirms the only eigenvalue is three (multiplicity two), whereas the two columns of~\verb|V| confirm the eigenspace \(\EE_1=\Span\{(1,0),(-1,0)\}=\Span\{(1,0)\}\).
\end{solution}


\item[\ref{eg:faem:b}.]
\(B=\begin{bmatrix}-1&1&-2
\\-1&0&-1
\\0&-3&1 \end{bmatrix}\)
\begin{solution} 
The eigenvalues are \(\lambda=-2\) (multiplicity one) and \(\lambda=1\) (multiplicity two).
\begin{itemize}
\item For \(\lambda=-2\) solve
\begin{equation*}
(B+2I)\xv=\begin{bmatrix}1&1&-2
\\-1&2&-1
\\0&-3&3 \end{bmatrix}\xv=\ov\,.
\end{equation*}
The third component of this equation requires \(-3x_2+3x_3=0\)\,, that is, \(x_2=x_3\)\,.
The second component requires \(-x_1+2x_2-x_3=0\)\,, that is, \(x_1=2x_2-x_3=2x_3-x_3=x_3\)\,.
The first component requires \(x_1+x_2-2x_3=0\) which is also satisfied by \(x_1=x_2=x_3\)\,.
All eigenvectors are of the form \(x_3(1,1,1)\).
That is, the eigenspace \(\EE_{-2}=\Span\{(1,1,1)\}\).

\item For \(\lambda=1\) solve
\begin{equation*}
(B-1I)\xv=\begin{bmatrix}-2&1&-2
\\-1&-1&-1
\\0&-3&0 \end{bmatrix}\xv=\ov\,.
\end{equation*}
The third component of this equation requires \(-3x_2=0\)\,, that is, \(x_2=0\)\,.
The second component requires \(-x_1-x_2-x_3=0\)\,, that is, \(x_1=-x_2-x_3=0-x_3=-x_3\)\,.
The first component requires \(-2x_1+x_2-2x_3=0\) which is also satisfied by \(x_1=-x_3\) and \(x_2=0\).
All eigenvectors are of the form \(x_3(-1,0,1)\).
That is, the eigenspace \(\EE_{1}=\Span\{(-1,0,1)\}\).
\end{itemize}

Alternatively, in \script, executing 
\begin{verbatim}
B=[-1 1 -2
 -1 0 -1
 0 -3 1]
[V,D]=eig(B)
\end{verbatim}
gives us
\begin{verbatim}
V =
    -0.5774     0.7071    -0.7071
    -0.5774     0.0000     0.0000
    -0.5774    -0.7071     0.7071
D =
         -2          0          0
          0          1          0
          0          0          1
\end{verbatim}
Diagonal matrix~\verb|D| confirms the eigenvalues are \(\lambda=-2\) (multiplicity one) and \(\lambda=1\) (multiplicity two). 
The first column of~\verb|V| confirms the eigenspace 
\begin{eqnarray*}
\EE_{-2}&=&\Span\{(-0.5774,-0.5774,-0.5774)\}
\\&=&\Span\{(1,1,1)\}.
\end{eqnarray*}
Whereas the last two columns of~\verb|V| confirm the eigenspace 
\begin{eqnarray*}
\EE_1&=&\Span\{(0.7071,0,-0.7071),(-0.7071,0,0.7071)\}
\\&=&\Span\{(-1,0,1)\}.
\end{eqnarray*}
\end{solution}

\item[\ref{eg:faem:c}.]
\(C=\begin{bmatrix}-1&0&-2
\\0&-3&2
\\0&-2&1\end{bmatrix}\)
\begin{solution} 
The only eigenvalue is \(\lambda=-1\) with multiplicity three.
Its eigenvectors~\xv\ satisfy
\begin{equation*}
(C+1 I)\xv=\begin{bmatrix} 0&0&-2
\\0&-2&2
\\0&-2&2 \end{bmatrix}\xv=\ov\,.
\end{equation*}
The first component of this equation requires \(x_3=0\)\,.  
The second and third components both requires \(-2x_2+2x_3=0\), hence \(x_2=x_3=0\)\,.
Since \(x_1\) is unconstrained, all eigenvectors are of the form \(\xv=x_1(1,0,0)\).  
That is the eigenspace  \(\EE_{-1}=\Span\{(1,0,0)\}\).

Alternatively, in \script, executing 
\begin{verbatim}
C=[-1 0 -2
 0 -3 2
 0 -2 1]
[V,D]=eig(C)
\end{verbatim}
gives us
\begin{verbatim}
V =
           1          -1          -1
           0           0           0
           0           0           0
D =
          -1           0           0
           0          -1           0
           0           0          -1
\end{verbatim}
Diagonal matrix~\verb|D| confirms the only eigenvalue is \(\lambda=-1\) with multiplicity three.  
The three columns of~\verb|V| confirm the eigenspace \(\EE_{-1}=\Span\{(1,0,0)\}\).
\end{solution}

\end{itemize}
\end{example}

The matrices in \autoref{eg:faespm} all have repeated eigenvalues.  
For these repeated eigenvalues the corresponding eigenspaces happen to be all one dimensional.
This contrasts with the case of symmetric matrices where the eigenspaces always have the same dimensionality as the multiplicity of the eigenvalue, as illustrated by Examples~\ref{eg:espace2d} and~\ref{eg:sier2eig}.
Subsequent sections work towards \autoref{thm:dimee} which establishes that for non-symmetric matrices an eigenspace has dimensionality between one and the multiplicity of the corresponding eigenvalue.


\begin{example} \label{eg:fiveev}
By hand, find the eigenvalues and eigenspaces of the matrix
\begin{equation*}
A=\begin{bmatrix}0&3&0&0&0
\\1&0&3&0&0
\\0&1&0&3&0
\\0&0&1&0&3
\\0&0&0&1&0\end{bmatrix}
\end{equation*}
(\autoref{eg:fiveevm} confirms the answer using \verb|eig()| in \script.)
\begin{solution} 
Adopt \autoref{pro:geneig}.
\begin{enumerate}
\item The characteristic equation is
\begin{eqnarray*}
&&\det(A-\lambda I)
\\&=&\begin{vmatrix}-\lambda&3&0&0&0
\\1&-\lambda&3&0&0
\\0&1&-\lambda&3&0
\\0&0&1&-\lambda&3
\\0&0&0&1&-\lambda\end{vmatrix}
\\&&(\text{by first row expansion~\eqref{eq:detlet}})
\\&=&(-\lambda)\begin{vmatrix}-\lambda&3&0&0
\\1&-\lambda&3&0
\\0&1&-\lambda&3
\\0&0&1&-\lambda\end{vmatrix}
-3\begin{vmatrix}1&3&0&0
\\0&-\lambda&3&0
\\0&1&-\lambda&3
\\0&0&1&-\lambda\end{vmatrix}
\\&&(\text{by first row and first column expansion, respectively})
\\&=&(-\lambda)^2\begin{vmatrix}-\lambda&3&0
\\1&-\lambda&3
\\0&1&-\lambda\end{vmatrix}
-(-\lambda)3\begin{vmatrix}1&3&0
\\0&-\lambda&3
\\0&1&-\lambda\end{vmatrix}
\\&&{}
-3\begin{vmatrix}-\lambda&3&0
\\1&-\lambda&3
\\0&1&-\lambda\end{vmatrix}
\\&&(\text{by common factor, and first column expansion})
\\&=&(\lambda^2-3)\begin{vmatrix}-\lambda&3&0
\\1&-\lambda&3
\\0&1&-\lambda\end{vmatrix}
+3\lambda\begin{vmatrix}-\lambda&3
\\1&-\lambda\end{vmatrix}
\\&&(\text{using~\eqref{eq:dets23b}})
\\&=&(\lambda^2-3)\big[(-\lambda)^3+0+0-0+3\lambda+3\lambda\big]
+3\lambda(\lambda^2-3)
\\&=&(\lambda^2-3)(-\lambda^3+9\lambda)
\\&=&-\lambda(\lambda^2-3)(\lambda^2-9)=0\,.
\end{eqnarray*}
The five eigenvalues are thus \(\lambda=0,\pm\sqrt3,\pm3\)\,, all of multiplicity one.

\item Consider each eigenvalue in turn.
\begin{itemize}
\item[\(\lambda=0\)]  Solve
\begin{equation*}
(A-0I)\xv=\begin{bmatrix}0&3&0&0&0
\\1&0&3&0&0
\\0&1&0&3&0
\\0&0&1&0&3
\\0&0&0&1&0\end{bmatrix}\xv=\ov\,.
\end{equation*}
The last row requires \(x_4=0\)\,.
The fourth row requires \(x_3+3x_5=0\)\,, that is, \(x_3=-3x_5\)\,.
The third row requires \(x_2+3x_4=0\)\,, that is, \(x_2=-3x_4=-3\cdot 0=0\)\,.
The second row requires \(x_1+3x_3=0\)\,, that is, \(x_1=-3x_3=9x_5\)\,.
The first row requires \(3x_2=0\)\,, which is satisfied as \(x_2=0\)\,.
Since all eigenvectors are of the form \((9x_5,0,-3x_5,0,x_5)\), the eigenspace \(\EE_0=\Span\{(9,0,-3,0,1)\}\).

\item[\(\lambda=\pm\sqrt3\)]  Being careful with the upper and lower signs, solve \((A\mp\sqrt3I)\xv=\ov\)\,, that is,
\begin{equation*}
\begin{bmatrix}\mp\sqrt3&3&0&0&0
\\1&\mp\sqrt3&3&0&0
\\0&1&\mp\sqrt3&3&0
\\0&0&1&\mp\sqrt3&3
\\0&0&0&1&\mp\sqrt3\end{bmatrix}\xv=\ov\,.
\end{equation*}
The last row requires \(x_4\mp\sqrt3x_5=0\)\,, that is, \(x_4=\pm\sqrt3x_5\)\,.
The fourth row requires \(x_3\mp\sqrt3x_4+3x_5=0\)\,, that is, \(x_3=\pm\sqrt3x_4-3x_5=3x_5-3x_5=0\)\,.
The third row requires \(x_2\mp\sqrt3x_3+3x_4=0\)\,, that is, \(x_2=\pm\sqrt3x_3-3x_4=\mp3\sqrt3x_5\)\,.
The second row requires \(x_1\mp\sqrt3x_2+3x_3=0\)\,, that is, \(x_1=\pm\sqrt3x_2-3x_3=-9x_5\)\,.
The first row requires \(\mp\sqrt3x_1+3x_2=0\)\,, which is satisfied as \(\mp\sqrt3(-9x_5)+3(\mp3\sqrt3x_5)=0\)\,.
Since all eigenvectors are of the form \((-9x_5, \mp3\sqrt3x_5, 0, \pm\sqrt3x_5, x_5)\), the eigenspaces \(\EE_{\pm\sqrt3}=\Span\{(-9,\mp3\sqrt3,0,\pm\sqrt3,1)\}\).

\item[\(\lambda=\pm3\)]  Being careful with the upper and lower signs, solve \((A\mp3I)\xv=\ov\)\,, that is,
\begin{equation*}
\begin{bmatrix}\mp3&3&0&0&0
\\1&\mp3&3&0&0
\\0&1&\mp3&3&0
\\0&0&1&\mp3&3
\\0&0&0&1&\mp3\end{bmatrix}\xv=\ov\,.
\end{equation*}
The last row requires \(x_4\mp3x_5=0\)\,, that is, \(x_4=\pm3x_5\)\,.
The fourth row requires \(x_3\mp3x_4+3x_5=0\)\,, that is, \(x_3=\pm3x_4-3x_5=9x_5-3x_5=6x_5\)\,.
The third row requires \(x_2\mp3x_3+3x_4=0\)\,, that is, \(x_2=\pm3x_3-3x_4=\pm9x_5\)\,.
The second row requires \(x_1\mp3x_2+3x_3=0\)\,, that is, \(x_1=\pm3x_2-3x_3=9x_5\)\,.
The first row requires \(\mp3x_1+3x_2=0\)\,, which is satisfied as \(\mp3(9x_5)+3(\pm9x_5)=0\)\,.
Since all eigenvectors are of the form \((9x_5, \pm9 x_5, 6x_5, \pm3x_5, x_5)\), the eigenspaces \(\EE_{\pm3}=\Span\{(9,\pm9,6,\pm3,1)\}\).

\end{itemize}
\end{enumerate}
\end{solution}
\end{example}


\begin{example} \label{eg:fiveevm}
Use \script\ to confirm the eigenvalues and eigenvectors found for the matrix of \autoref{eg:fiveev}.
\begin{solution} 
In \script\ execute
\begin{verbatim}
A=[0 3 0 0 0
 1 0 3 0 0
 0 1 0 3 0
 0 0 1 0 3
 0 0 0 1 0]
[V,D]=eig(A)
\end{verbatim}
\setbox\ajrqrbox\hbox{\qrcode{% eigenvectors
A=[0 3 0 0 0
 1 0 3 0 0
 0 1 0 3 0
 0 0 1 0 3
 0 0 0 1 0]
[V,D]=eig(A)
bsxfun(@rdivide,V,V(5,:))
}}%
\marginpar{\usebox{\ajrqrbox}}%
to obtain the report \twodp
\begin{verbatim}
V =
   0.62  -0.62   0.94  -0.85  -0.85
   0.62   0.62  -0.00   0.49  -0.49
   0.42  -0.42  -0.31  -0.00   0.00
   0.21   0.21  -0.00  -0.16   0.16
   0.07  -0.07   0.10   0.09   0.09
D =
   3.00      0      0      0      0
      0  -3.00      0      0      0
      0      0  -0.00      0      0
      0      0      0  -1.73      0
      0      0      0      0   1.73
\end{verbatim}
The eigenvalues in~\verb|D| agree with the hand calculations of \(\lambda=0,\pm\sqrt3,\pm3\).
To confirm the hand calculation of eigenvectors in \autoref{eg:fiveev}, here divide each column of~\verb|V| by its last element,~\verb|V(5,:)|, via \index{bsxfun()@\texttt{bsxfun()}}\verb|bsxfun(@rdivide,V,V(5,:))| which gives the more appealing matrix of eigenvectors \twodp
\begin{verbatim}
ans =
   9.00   9.00   9.00  -9.00  -9.00
   9.00  -9.00   0.00   5.20  -5.20
   6.00   6.00  -3.00   0.00   0.00
   3.00  -3.00   0.00  -1.73   1.73
   1.00   1.00   1.00   1.00   1.00
\end{verbatim}
These also agree with the hand calculation.
\end{solution}
\end{example}







\subsection{Repeated eigenvalues are sensitive}
\label{sec:reas}

\index{repeated eigenvalue|(}
\index{sensitivity|(}


Albeit hidden in \autoref{eg:multeval}, \idx{repeated eigenvalue}s are exquisitely sensitive to errors in either the matrix or the computation.
\begin{aside}
This optional subsection does not prove the sensitivity: it uses examples to introduce and illustrate.
\end{aside}
If the matrix or the computation has an error~\(\epsilon\)\index{error~$\epsilon$}\index{\(\epsilon\)}, then expect a repeated eigenvalue of multiplicity~\(m\) to appear as \(m\)~eigenvalues all within about~\(\epsilon^{1/m}\) of each other.
Consequently, when we find or compute \(m\)~eigenvalues all within about~\(\epsilon^{1/m}\), then \emph{suspect} them to be one eigenvalue of multiplicity~\(m\).

\begin{example} \label{eg:}
Explore the eigenvalues of the matrix
\(A=\begin{bmatrix} a&1\\0.0001&a \end{bmatrix}\) for every parameter~\(a\).
\begin{solution} 
By hand, the characteristic equation is
\begin{equation*}
\det\begin{bmatrix} a-\lambda&1\\0.0001&a-\lambda \end{bmatrix}
=(a-\lambda)^2-0.0001=0\,.
\end{equation*}
Rearranging gives \((\lambda-a)^2=0.0001\)\,.
Taking square roots, \(\lambda-a=\pm0.01\)\,; that is, the two eigenvalues are \(\lambda=a\pm0.01\)\,.

Alternatively, if we consider that the entry~\(0.0001\) in the matrix is an error, that the entry should really be zero, then the eigenvalues should really be one repeated eigenvalue \(\lambda=a\) of multiplicity two.
However, the `error'~\(0.0001\) splits the repeated eigenvalue into two by an amount \(\sqrt{0.0001}=0.01\)\,.
\end{solution}
\end{example}


Further, since computers work to a relative error of about~\(10^{-15}\), then expect a repeated eigenvalue of multiplicity~\(m\) to appear as \(m\)~eigenvalues within about~\(10^{-15/m}\) of each other---even when there are no experimental errors in the matrix.
Repeat some of the previous cases of \autoref{eg:multeval}, preceded by the \script\ command \index{format long@\texttt{format long}}\verb|format long|, to see that the repeated eigenvalues are sensitive to computational~errors.

%Similarly, repeated eigenvalues are sensitive to errors in the matrix. 
%The following examples show the sensitivity when we are uncertain about the components of a matrix. 


\begin{example} \label{eg:}
Use \script\ to compute eigenvalues of the following matrices and comment on the effect on repeated eigenvalues of errors in the matrix and/or the computation.
\begin{enumerate}

\item \(B=\begin{bmatrix} 3&0&-2&0&0
\\-1&5&0&-1&3
\\-1&2&4&0&1
\\5&-1&4&1&-1
\\3&2&1&-2&2\end{bmatrix}\)
\begin{solution} 
In \script\ execute
\begin{verbatim}
eig([3  0 -2  0  0
    -1  5  0 -1  3
    -1  2  4  0  1
     5 -1  4  1 -1
     3  2  1 -2  2])
\end{verbatim}
\setbox\ajrqrbox\hbox{\qrcode{% eigenvalue code
eig([3 0 -2 0 0
 -1 5 0 -1 3
 -1 2 4 0 1
 5 -1 4 1 -1
 3 2 1 -2 2])
}}%
\marginpar{\usebox{\ajrqrbox}}%
to get something like
\begin{verbatim}
ans =
  -3.5176e-08
   3.5176e-08
   6.4142
   5.0000
   3.5858
\end{verbatim}
There are three eigenvalues of multiplicity one, namely \(5\pm\sqrt2\) and~\(5\)\,.  
The two values \(\pm\verb|3.5176e-08|\) at first sight appear to be two eigenvalues, \(\pm3.5176\cdot10^{-8}\), each of multiplicity one.  
However, computers usually work to about 15~digits accuracy, that is, the typical error is about~\(10^{-15}\), so an eigenvalue of multiplicity two is generally computed as two eigenvalues separated by about \(\sqrt{10^{-15}}\approx3\cdot10^{-8}\).
Thus we suspect that the two values~\(\pm\verb|3.5176e-08|\) actually represent one eigenvalue \(\lambda=0\) with multiplicity~two.
\end{solution}


\item Suppose the above matrix~\(B\) is obtained from some experiment where there are \idx{experimental error}s in the entries with error about~\(0.0001\).
Randomly perturb the entries in matrix~\(B\) to see the effects of such errors on the eigenvalues (use \index{randn()@\texttt{randn()}}\verb|randn()|, \autoref{tbl:mtlbops}).
\begin{solution} 
In \script\ execute
\begin{verbatim}
B=[3  0 -2  0  0
  -1  5  0 -1  3
  -1  2  4  0  1
   5 -1  4  1 -1
   3  2  1 -2  2]
eig(B+0.0001*randn(5))
\end{verbatim}
\setbox\ajrqrbox\hbox{\qrcode{% eigenvalue code
B=[3 0 -2 0 0
 -1 5 0 -1 3
 -1 2 4 0 1
 5 -1 4 1 -1
 3 2 1 -2 2]
eig(B+0.0001*randn(5))
}}%
\marginpar{\usebox{\ajrqrbox}}%
to get something like
\begin{verbatim}
ans =
   -0.0226
    0.0225
    6.4145
    4.9999
    3.5860
\end{verbatim}
The repeated eigenvalue \(\lambda=0\) splits into two eigenvalues, \(\lambda=\pm0.0226\)\,, of size roughly \(\sqrt{0.0001}=0.01\).
The other eigenvalues are also perturbed by the errors but only by amounts of size roughly~\(0.0001\).

Depending upon the random numbers, other possible answers are like
\begin{verbatim}
ans =
   0.0001 + 0.0157i
   0.0001 - 0.0157i
   6.4146 + 0.0000i
   4.9993 + 0.0000i
   3.5860 + 0.0000i
\end{verbatim}
where the repeated eigenvalue of zero splits to be a pair of complex valued eigenvalues of roughly \(\pm \i\sqrt{0.0001}=\pm \i0.01\)\,.
\end{solution}



\item \(C=\begin{bmatrix}-1&0&0&0
\\-1&2&-3&3
\\3&1&-1&0
\\0&3&-2&1\end{bmatrix}\) perturbed by errors of size~\(10^{-6}\)
\begin{solution} 
In \script\ execute
\begin{verbatim}
C=[-1 0 0 0
 -1 2 -3 3
  3 1 -1 0
  0 3 -2 1]
eig(C+1e-6*randn(4))
\end{verbatim}
\setbox\ajrqrbox\hbox{\qrcode{% eigenvalue code
C=[-1 0 0 0
 -1 2 -3 3
  3 1 -1 0
  0 3 -2 1]
eig(C+1e-6*randn(4))
}}%
\marginpar{\usebox{\ajrqrbox}}%
to get something like
\begin{verbatim}
ans =
   4.0000 + 0.0000i
  -1.0156 + 0.0000i
  -0.9922 + 0.0139i
  -0.9922 - 0.0139i
\end{verbatim}
The eigenvalue~\(4\) of multiplicity one is not noticeably affected by the errors about~\(10^{-6}\).
However, the repeated eigenvalue of \(\lambda=-1\) with multiplicity three is split into three close eigenvalues (two complex-valued) all differing by roughly~\(0.01\) which is indeed the cube-root of the perturbation~\(10^{-6}\).
\end{solution}

\end{enumerate}
\end{example}






\begin{activity}
In an experiment measurements are made to three decimal place accuracy.
Then in analysing the results, a \(5\times5\)~matrix is formed from the measurements, and its eigenvalues computed by \script\ to be
%\begin{verbatim}
%a=tril(randn(5));a([1 7 13])=1;a=a+0.002*(rand(5)-0.5);eig(a)
%\end{verbatim}%
\begin{equation*}
-0.9851,\quad
0.1266,\quad
0.9954,\quad
1.0090,\quad
1.0850.
\end{equation*}
What should you suspect is the number of different eigenvalues?
\actposs[4]{three}{two}{four}{five}
%\partswidth=5em
%\begin{parts}
%\item two
%\item three
%\item four
%\item five
%\end{parts}
\end{activity}








%\begin{verbatim}
%for i=1:99999,a=randn(4)*2;a=round(a+a');l=eig(a);if min(abs(diff(l)))<1e-6,a=a,lam=l,break,end,end
%\end{verbatim}

\paragraph{But symmetric matrices are OK}
The \idx{eigenvalue}s of a \idx{symmetric matrix} are not so sensitive.
This lack of sensitivity is fortunate as many applications give rise to symmetric matrices (\autoref{ch:eesm}).
Such symmetry often reflects some symmetry in the natural world such as Newton's law of every action having an equal and opposite reaction.
For symmetric matrices, the eigenvalues and eigenvectors are reasonably robust to both computational perturbations and \idx{experimental error}s.




\begin{example} \label{eg:}
For perhaps the simplest example, consider the \(2\times2\) \idx{symmetric matrix} \(A=\begin{bmatrix} a&0\\0&a \end{bmatrix}\).
Being diagonal, matrix~\(A\) has eigenvalue \(\lambda=a\) (multiplicity two).
Now perturb the matrix by `experimental' error to \(B=\begin{bmatrix} a&10^{-4}\\ 10^{-4}&a \end{bmatrix}\).
The characteristic equation of~\(B\) is
\begin{equation*}
\det(B-\lambda I)=(a-\lambda)^2-10^{-8}=0\,.
\end{equation*}
Rearrange this equation to \((\lambda-a)^2=10^{-8}\).  
Taking square roots gives \(\lambda-a=\pm 10^{-4}\), that is, the eigenvalues of~\(B\) are \(\lambda=a\pm10^{-4}\).
Because a perturbation to the \idx{symmetric matrix} of size~\(10^{-4}\) only changes the eigenvalues by a similar amount, the eigenvalues are \emph{not} sensitive.
\end{example}





\begin{activity}
What are the eigenvalues of matrix \(\begin{bmatrix} a&0.01\\-0.01&a \end{bmatrix}\)?
\actposs[4]{\(a\pm 0.01\i\)}{\(a\pm 0.01\)}{\(a\pm 0.1\)}{\(a\pm 0.1\i\)}
%\partswidth=5em
%\begin{parts}
%\item \(a\pm 0.01\)
%\item \(a\pm 0.1\)
%\item \(a\pm 0.01\i\)\actans
%\item \(a\pm 0.1\i\)
%\end{parts}
\end{activity}





\begin{example} \label{eg:}
Compute the eigenvalues of the \idx{symmetric matrix}
\begin{equation*}
A=\begin{bmatrix}1&1&0&2
\\1&0&2&-1
\\0&2&1&4
\\2&-1&4&1\end{bmatrix}
\end{equation*}
and see matrix~\(A\) has an eigenvalue of multiplicity two.  
Explore the effects on the eigenvalues of errors in the matrix by perturbing the entries by random amounts of size~\(0.0001\).
\begin{solution} 
In \script\ execute
\begin{verbatim}
A=[1 1 0 2
 1 0 2 -1
 0 2 1 4
 2 -1 4 1]
eig(A)
eig(A+0.0001*randn(4))
\end{verbatim}
\setbox\ajrqrbox\hbox{\qrcode{% eigenvalue code
A=[1 1 0 2
 1 0 2 -1
 0 2 1 4
 2 -1 4 1]
eig(A)
eig(A+0.0001*randn(4))
}}%
\marginpar{\usebox{\ajrqrbox}}%
to firstly get
\begin{verbatim}
ans =
  -4.6235
   1.0000
   1.0000
   5.6235
\end{verbatim}
showing the eigenvalue \(\lambda=1\) has multiplicity two.
Whereas, secondly, the eigenvalues of the perturbed matrix depend upon the random numbers and so might be
\begin{verbatim}
ans =
  -4.6236
   5.6235
   0.9998
   0.9999
\end{verbatim}
or perhaps
\begin{verbatim}
ans =
  -4.6236 + 0.0000i
   5.6234 + 0.0000i
   1.0001 + 0.0000i
   1.0001 - 0.0000i
\end{verbatim}
In either case, for this symmetric matrix the perturbation by amounts~\(0.0001\) only change the eigenvalues, whether repeated or not, by an amount of about the same size.
\end{solution}
\end{example}




\index{repeated eigenvalue|)}
\index{sensitivity|)}







\subsection{Application: discrete dynamics of populations}
\label{sec:ddp}

\index{age structure|(}
\index{population|(}

Age structured \idx{population}s are one case where matrix properties and methods are crucial.
The approach of this section is also akin to much mathematical modelling of diseases and epidemics.
This section aims to show how to derive and use a matrix-vector model for the change in time~\(t\) of interesting properties~\(\yv\) of a population.
Specifically, this subsection derives and analyses the model \(\yv(t+1)=A\yv(t)\).


For a given species, let's define \begin{itemize}
\item \(y_1(t)\)~to be the number of \idx{juvenile}s (including \idx{infant}s), 
\item \(y_2(t)\)~the number of \idx{adolescent}s, and 
\item \(y_3(t)\)~the number of \idx{adult}s.
\end{itemize}
Mostly, biologists only count \idx{female}s as females are the determining sex for reproduction.
(Although some bacteria\slash algae have seven sexes!)
How do these numbers of females evolve over time? from generation to generation?  
First we need to choose a basic time interval (the unit of time): it could be one year, one month, one day, or maybe six months.
Whatever we choose as convenient, we then quantify the number of events that happen to the females in each time interval as shown schematically in the diagram below:
\begin{center}\small
\setlength{\unitlength}{0.011\linewidth}
\newcommand{\ta}[2]{\begin{tabular}{c}#1\\\(#2\)\end{tabular}}
\begin{picture}(90,20)
\put(0,0){\framebox(15,15){\ta{juveniles}{y_1(t)}}}
\put(30,5){\framebox(15,15){\ta{adolescent}{y_2(t)}}}
\put(60,0){\framebox(15,15){\ta{adults}{y_3(t)}}}
\put(15,10){\vector(3,1){15}} \put(22,10){\(k_1\)}
\put(45,15){\vector(3,-1){15}} \put(52,10){\(k_2\)}
\put(75,8){\vector(1,0){15}}  \put(82,10){\(k_3\)}
\put(60,3){\vector(-1,0){45}} \put(37,0){\(k_0\)}
\end{picture}
\end{center}
Over any one time interval, and only counting females: 
\begin{itemize}
\item a fraction~\(k_1\) of the juveniles become adolescents; 
\item a fraction~\(k_2\) of the adolescents become adults;
\item a fraction~\(k_3\) of the adults die; 
\item but adults also give birth to juveniles at rate~\(k_0\) per adult.
\end{itemize}
Model this scenario with a system of discrete dynamical equations which are of the form that the numbers at the next time, \(t+1\), depend upon the numbers at the time~\(t\):
\begin{eqnarray*}&&
{y_1(t+1)}=\cdots,
\\&&{y_2(t+1)}=\cdots,
\\&&{y_3(t+1)}=\cdots.
\end{eqnarray*}
Let's fill in the right-hand sides from the given information about the rate of particular events per time interval.
\begin{itemize}
\item A fraction~\(k_1\) of the \idx{juvenile}s~\(y_1(t)\) becoming adolescents also means a fraction \((1-k_1)\) of the juveniles remain juveniles, hence
\begin{eqnarray*}&&
{y_1(t+1)}=(1-k_1)y_1(t)+\cdots,
\\&&{y_2(t+1)}=+k_1y_1(t)+\cdots,
\\&&{y_3(t+1)}=\cdots.
\end{eqnarray*}
\item A fraction~\(k_2\) of the \idx{adolescent}s~\(y_2(t)\) becoming adults also means a fraction \((1-k_2)\)  of the adolescents remain adolescents, hence additionally
\begin{eqnarray*}&&
{y_1(t+1)}=(1-k_1)y_1(t)+\cdots,
\\&&{y_2(t+1)}=+k_1y_1(t)+(1-k_2)y_2(t),
\\&&{y_3(t+1)}=+k_2y_2(t)+\cdots.
\end{eqnarray*}
\item A fraction~\(k_3\) of the \idx{adult}s die mean that a fraction \((1-k_3)\) of the adults remain adults, hence
\begin{eqnarray*}&&
{y_1(t+1)}=(1-k_1)y_1(t)+\cdots,
\\&&{y_2(t+1)}=+k_1y_1(t)+(1-k_2)y_2(t),
\\&&{y_3(t+1)}=+k_2y_2(t)+(1-k_3)y_3(t).
\end{eqnarray*}
\item But adults also give birth to juveniles at rate~\(k_0\) per adult so the number of juveniles increases by~\(k_0y_3\) from births:
\begin{eqnarray*}&&
{y_1(t+1)}=(1-k_1)y_1(t)+k_0y_3(t),
\\&&{y_2(t+1)}=+k_1y_1(t)+(1-k_2)y_2(t),
\\&&{y_3(t+1)}=+k_2y_2(t)+(1-k_3)y_3(t).
\end{eqnarray*}
This is our mathematical model of the \idx{age structure} of the \idx{population}.
\end{itemize}

Finally, write the mathematical model as the matrix-vector system 
\begin{equation*}
\begin{bmatrix} y_1(t+1)\\y_2(t+1)\\y_3(t+1) \end{bmatrix}
=\begin{bmatrix} 1-k_1&0&k_0
\\k_1&1-k_2&0
\\0&k_2&1-k_3 \end{bmatrix}
\begin{bmatrix} y_1(t)\\y_2(t)\\y_3(t) \end{bmatrix},
\end{equation*}
that is, \(\yv(t+1)=A\yv(t)\)\,.
Such a model empowers predictions.



\Needspace{6\baselineskip}
\begin{example}[orangutans] \label{eg:orang}
From the following  extract of the \idx{Wikipedia} entry on \idx{orangutan}s [20 Mar 2014] derive a mathematical model for the \idx{age structure} of the orangutans from one year to the next.
\marginpar{\includegraphics{EigGen/P1030879}}
\begin{quote}
Gestation lasts for 9~months, with females giving birth to their first offspring between the ages of 14 and 15~years. Female orangutans have [seven to] eight-year intervals between births, the longest interbirth intervals among the great apes. \ldots\  Infant orangutans are completely dependent on their mothers for the first two years of their lives. The mother will carry the infant during travelling, as well as feed it and sleep with it in the same night nest. For the first four months, the infant is carried on its belly and never relieves physical contact. In the following months, the time an infant spends with its mother decreases. When an orangutan reaches the age of two, its climbing skills improve and it will travel through the canopy holding hands with other orangutans, a behaviour known as ``buddy travel". Orangutans are juveniles from about two to five years of age and will start to temporarily move away from their mothers.  Juveniles are usually weaned at about four years of age. Adolescent orangutans will socialize with their peers while still having contact with their mothers. Typically, orangutans live over 30~years in both the wild and captivity.
\end{quote}
Suppose the initial population of orangutans in some area at year zero of a study is that of \(30\)~adolescent females and \(15\)~adult females.
Use the mathematical model to predict the population for the next five years.

\begin{solution} 
\emph{Choose level:}  we choose a time unit of one year, and choose to model three age categories.%
\footnote{A coarse model considers just the total number in a species; a fine model might count the number of each age in years (here 30~years); a `micro' model might simulate each and every orangutan as individuals (thousands).}

\emph{Define:}
\begin{itemize}
\item \(y_1(t)\) is the number of \idx{juvenile} \idx{female}s (including infant females) at time~\(t\)~(years), say age\({}\leq5\)~years;
\item \(y_2(t)\) is the number of \idx{adolescent} females, say \(6\leq\text{age}\leq 14\)~years;
\item \(y_3(t)\) is the number of \idx{adult} females, say age\({}\geq15\)~years.
\end{itemize}
\emph{Quantify changes in a year:} from the Wikipedia information (with numbers slightly modified here to make the results numerically simpler):
\begin{itemize}
\item Orangutans are juvenile for say \(5\)~years, so in any one year a fraction~\(1/5\) of them grow to be adolescent and a fraction~\(4/5\) remain as juveniles.
That is,
\begin{equation*}
y_1(t+1)=\tfrac45y_1(t)+\cdots\,,\quad y_2(t+1)=\tfrac15y_1(t)+\cdots\,.
\end{equation*}
Say an adult female gives birth every \(7\)--\(8\)~years, so on average she gives birth to a juvenile female every \(15\)~years. 
Thus a fraction~\(1/15\) of adults give birth to a juvenile female in any year.
Consequently, model juveniles as 
\begin{equation*}
{y_1}(t+1)=\tfrac45y_1(t)+\tfrac1{15}y_3(t).
\end{equation*}
\item Adolescent orangutans become breeding adults after another \(9\)--\(10\)~years, so in any one year a fraction~\(1/10\) of them grow to adults and \(9/10\)ths remain adolescents.
That is, 
\begin{equation*}
{y_2(t+1)}=\tfrac15y_1(t)+\tfrac9{10}y_2(t),\quad
{y_3(t+1)}=\tfrac1{10}y_2(t)+\cdots\,.
\end{equation*}

\item Orangutans live to 30~years, about \(15\)~years of adulthood so in any one year a fraction~\(14/15\) of adult females live to the next year.  
Consequently, 
\begin{equation*}
{y_3(t+1)}=\tfrac1{10}y_2(t)+\tfrac{14}{15}y_3(t).
\end{equation*}

\end{itemize}

Our mathematical model of the \idx{age structure} is then, in terms of the vector \(\yv=(y_1,y_2,y_3)\),  
\begin{equation*}
\yv(t+1)=A\yv(t)
=\begin{bmatrix} \frac45&0&\frac1{15}
\\\frac15&\frac9{10}&0
\\0&\frac1{10}&\frac{14}{15} \end{bmatrix}\yv(t)
\end{equation*}

To predict the population we need to know the current population.
The given information is that there are initially \(30\)~adolescent females and \(15\)~adult females. 
This information specifies that at time zero the population vector \(\yv(0)=(0,30,15)\) in the study area.
\begin{enumerate}
\item Then the rule \(\yv(t+1)=A\yv(t)\) with time \(t=0\) gives
\begin{eqnarray*}
\yv(1)&=&A\yv(0)
\\&=&\begin{bmatrix} \frac45&0&\frac1{15}
\\\frac15&\frac9{10}&0
\\0&\frac1{10}&\frac{14}{15} \end{bmatrix}\begin{bmatrix} 0\\30\\15 \end{bmatrix}
\\&=&\begin{bmatrix} 1\\27\\17 \end{bmatrix}.
\end{eqnarray*}
That is, during the first year we predict that there is one birth of a female juvenile, three adolescents matured to adults, and one adult died. 

\item Then the rule \(\yv(t+1)=A\yv(t)\) with time \(t=1\)~year gives
\begin{eqnarray*}
\yv(2)&=&A\yv(1)
\\&=&\begin{bmatrix} \frac45&0&\frac1{15}
\\\frac15&\frac9{10}&0
\\0&\frac1{10}&\frac{14}{15} \end{bmatrix}\begin{bmatrix} 1\\27\\17 \end{bmatrix}
\\&=&\begin{bmatrix} \frac{29}{15}\\\frac{49}2\\\frac{557}{30} \end{bmatrix}=\begin{bmatrix} 1.93\\24.50\\18.57 \end{bmatrix}\ \twodp.
\end{eqnarray*}
In real life we cannot have such fractions of an orangutan.
These predictions are averages, or expectations on average, and must be interpreted as such. 
Thus after two years, we expect on average nearly two juveniles, \(24\) or~\(25\) adolescents, and \(18\) or~\(19\) adults.

\item Continuing on with the aid of \script\ or calculator, the rule \(\yv(t+1)=A\yv(t)\) with time \(t=2\)~years gives
\begin{eqnarray*}
\yv(3)&=&A\yv(2)
\\&=&\begin{bmatrix} \frac45&0&\frac1{15}
\\\frac15&\frac9{10}&0
\\0&\frac1{10}&\frac{14}{15} \end{bmatrix}
\begin{bmatrix} 1.93\\24.50\\18.57 \end{bmatrix}
\\&=&\begin{bmatrix} 2.78\\22.44\\19.78 \end{bmatrix}\ \twodp.
\end{eqnarray*}
In \script\ do this calculation via
\begin{verbatim}
A=[4/5 0 1/15;1/5 9/10 0;0 1/10 14/15]
y0=[0;30;15]
y1=A*y0
y2=A*y1
y3=A*y2
y4=A*y3
y5=A*y4
\end{verbatim}
\setbox\ajrqrbox\hbox{\qrcode{% orangutans
A=[4/5 0 1/15;1/5 9/10 0;0 1/10 14/15]
y0=[0;30;15]
y1=A*y0
y2=A*y1
y3=A*y2
y4=A*y3
y5=A*y4
}}%
\marginpar{\usebox{\ajrqrbox}}%

\item Consequently, the rule \(\yv(t+1)=A\yv(t)\) with time \(t=3\)~years gives
\begin{eqnarray*}
\yv(4)&=&A\yv(3)
\\&=&\begin{bmatrix} \frac45&0&\frac1{15}
\\\frac15&\frac9{10}&0
\\0&\frac1{10}&\frac{14}{15} \end{bmatrix}
\begin{bmatrix} 2.78\\22.44\\19.78 \end{bmatrix}
\\&=&\begin{bmatrix} 3.55\\20.75\\20.70 \end{bmatrix}\ \twodp.
\end{eqnarray*}
  
\item Lastly, the rule \(\yv(t+1)=A\yv(t)\) with time \(t=4\)~years gives (to complete the marginal plot)
\marginpar{\begin{tikzpicture}
\def\temp#1{\addplot+ table[x=t,y=y#1,row sep=\\] {
    t       y1     y2      y3 \\
    0.00    0.00   30.00   15.00 \\
    1.00    1.00   27.00   17.00 \\
    2.00    1.93   24.50   18.57 \\
    3.00    2.78   22.44   19.78 \\
    4.00    3.55   20.75   20.70 \\
    5.00    4.22   19.38   21.40 \\
};
\addlegendentry{$y_#1(t)$};
}
\begin{axis}[footnotesize,font=\footnotesize, 
xlabel={$t$\,yrs},ylabel={$y_j$},axis lines=middle
,only marks,xmax=5.4,legend pos=south west]
\temp1
\temp2
\temp3
\end{axis}
\end{tikzpicture}}%
\begin{eqnarray*}
\yv(5)&=&A\yv(4)
\\&=&\begin{bmatrix} \frac45&0&\frac1{15}
\\\frac15&\frac9{10}&0
\\0&\frac1{10}&\frac{14}{15} \end{bmatrix}
\begin{bmatrix} 3.55\\20.75\\20.70 \end{bmatrix}
\\&=&\begin{bmatrix} 4.22\\19.38\\21.40 \end{bmatrix}\ \twodp.
\end{eqnarray*}
Thus after five years the mathematical model predicts about \(4\)~juveniles, \(19\)~adolescents, and \(21\)~adults (on average).
\end{enumerate}
The five-year population of \(44\)~females (\(45\)~if you add all the fractions) is the same as the starting population.
This nearly static total population is no accident, as we next see, and is one reason why orangutans are critically endangered.
%\emph{Interpret} Two of the three modes decay like~\(e^{-t/5}\) so after ten years or so they are negligible. 
%What emerges is the remaining mode with age structure of  20\%~juveniles, 40\%~adolescents, and 40\%~adults.
%A further prediction is that this population remains constant.
%Although we fudged the numbers somewhat, this implies that orangutan survival is easily damaged by small influences. 
\end{solution}
\end{example}





The mathematical model \(\yv(t+1)=A\yv(t)\) does predict\slash forecast the future populations.
However, to make predictions for many years and for general initial populations we prefer the formula solution given by the upcoming \autoref{thm:dynsol} and introduced in the next example.





\begin{example} \label{eg:2x2map}
A vector \(\yv(t)\in\RR^2\) changes with time~\(t\) according to the model
\begin{equation*}
\yv(t+1)=A\yv(t)=\begin{bmatrix} 1&-1\\-4&1 \end{bmatrix}\yv(t).
\end{equation*}
First, what is \(\yv(3)\) if the \idx{initial value} \(\yv(0)=(0,1)\)?
Second, find a general formula for~\(\yv(t)\) from every initial~\(\yv(0)\).
\begin{solution} 
First, given \(\yv(0)=(0,1)\) just compute (as drawn in the margin)
\marginpar{\begin{tikzpicture}
\begin{axis}[footnotesize,font=\footnotesize,domain=0:3, 
xlabel={time $t$},ylabel={$y_j$},axis lines=middle
,only marks,samples=4,xmax=3.9]
\addplot+ {0.25*(-1)^x-0.25*3^x};
\addlegendentry{$y_1(t)$};
\addplot+ {0.25*(-1)^x*2+0.25*3^x*2};
\addlegendentry{$y_2(t)$};
\end{axis}
\end{tikzpicture}}%
\begin{eqnarray*}
\yv(1)&=&A\yv(0)=\begin{bmatrix} 1&-1\\-4&1 \end{bmatrix}
\begin{bmatrix} 0\\1 \end{bmatrix}
=\begin{bmatrix} -1\\1 \end{bmatrix},
\\\yv(2)&=&A\yv(1)=\begin{bmatrix} 1&-1\\-4&1 \end{bmatrix}
\begin{bmatrix} -1\\1 \end{bmatrix}
=\begin{bmatrix} -2\\5 \end{bmatrix},
\\\yv(3)&=&A\yv(2)=\begin{bmatrix} 1&-1\\-4&1 \end{bmatrix}
\begin{bmatrix} -2\\5 \end{bmatrix}
=\begin{bmatrix} -7\\13 \end{bmatrix}.
\end{eqnarray*}


Second, let's suppose there may be solutions in the form \(\yv=\lambda^t\vv\) for some non-zero scalar~\(\lambda\) and some vector \(\vv\in\RR^2\).
Substitute into the model \(\yv(t+1)=A\yv(t)\) to find \(\lambda^{t+1}\vv=A\lambda^t\vv\). 
Divide by (non-zero) \(\lambda^t\) to find we require \(A\vv=\lambda\vv\)\,: this is an \idx{eigen-problem}.
That is, for every eigenvalue~\(\lambda\) of matrix~\(A\), with corresponding eigenvector~\vv, there is a solution \(\yv=\lambda^t\vv\) of the model.

To find the eigenvalues here solve the characteristic equation 
\begin{equation*}
\det(A-\lambda I)=\begin{vmatrix} 1-\lambda&-1\\-4&1-\lambda \end{vmatrix}=(1-\lambda)^2-4=0\,.
\end{equation*}
Rearrange to \((\lambda-1)^2=4\) and take square roots to give \(\lambda-1=\pm2\)\,, that is, eigenvalues \(\lambda=1\pm2=-1,3\)\,.
\begin{itemize}
\item For eigenvalue \(\lambda_1=-1\) a corresponding eigenvector has to satisfy
\begin{equation*}
(A+I)\vv_1=\begin{bmatrix} 2&-1\\-4&2 \end{bmatrix}\vv_1=\ov\,.
\end{equation*}
That is, \(\vv_1\propto (1,2)\); let's take \(\vv_1=(1,2)\).

\item For eigenvalue \(\lambda_2=3\) a corresponding eigenvector has to satisfy
\begin{equation*}
(A-3I)\vv_2=\begin{bmatrix} -2&-1\\-4&-2 \end{bmatrix}\vv_2=\ov\,.
\end{equation*}
That is, \(\vv_2\propto (1,-2)\); let's take \(\vv_2=(-1,2)\).
\end{itemize}
Summarising so far: \(\yv'(t)=(-1)^t(1,2)\) and \(\yv''(t)=3^t(-1,2)\) are both solutions of the model \(\yv(t+1)=A\yv(t)\).

Further, the model \(\yv(t+1)=A\yv(t)\) is a linear equation in~\yv, so every linear combination of solutions is also a solution.
To see this, let's try \(\yv(t)=c_1\yv'(t)+c_2\yv''(t)\).
Substituting
\begin{eqnarray*}
\yv(t+1)&=&c_1\yv'(t+1)+c_2\yv''(t+1)
\\&=&c_1A\yv'(t)+c_2A\yv''(t)
\\&=&A\big(c_1\yv'(t)+c_2\yv''(t)\big)
\\&=&A\yv(t),
\end{eqnarray*}
as required.
So, \(\yv(t)=c_1(-1)^t(1,2)+c_23^t(-1,2)\) are solutions of the model for all constants~\(c_1\) and~\(c_2\), and over all time~\(t\).

What values of constants~\(c_1\) and~\(c_2\) should be chosen for any given initial~\(\yv(0)\)?
Substitute time \(t=0\) into \(\yv(t)=c_1(-1)^t(1,2)+c_23^t(-1,2)\) to require \(\yv(0)=c_1(-1)^0(1,2)+c_23^0(-1,2)\).
Since the zero powers \((-1)^0=3^0=1\)\,, this requires \(\yv(0)=c_1(1,2)+c_2(-1,2)\).
Write as the matrix-vector system
\begin{equation*}
\begin{bmatrix} 1&-1\\2&2 \end{bmatrix}
\begin{bmatrix} c_1\\c_2 \end{bmatrix}
=\yv(0)
\iff
\begin{bmatrix} c_1\\c_2 \end{bmatrix}
=\begin{bmatrix} \frac12&\frac14\\-\frac12&\frac14 \end{bmatrix}\yv(0)
\end{equation*}
upon invoking the inverse (\autoref{thm:2x2det}) of the matrix of eigenvectors, \(P=\begin{bmatrix} \vv_1&\vv_2 \end{bmatrix}\).
That is, because the matrix of eigenvectors is invertible, we find constants~\(c_1\) and~\(c_2\) to suit any initial~\(\yv(0)\).

For example, with initial \(\yv(0)=(0,1)\) the above formula gives \((c_1,c_2)=(\frac14,\frac14)\) and so the corresponding formula solution is \(\yv(t)=\frac14\cdot(-1)^t(1,2)+\frac14\cdot3^t(-1,2)\).
To check, evaluate at say \(t=3\) to find 
\(\yv(3)=-\frac14(1,2)+\frac{27}4(-1,2)=(-7,13)\),
as before.
In general, as here, as time~\(t\) increases, the solution~\(\yv(t)\) grows like~\(3^t\) with a little oscillation from the~\((-1)^t\) term.
\end{solution}
\end{example}




\begin{activity}
For \autoref{eg:2x2map}, what is the particular solution when \(\yv(0)=(1,1)\)?
\actposs[1]{\(\yv=\frac34\cdot(-1)^t(1,2)-\frac14\cdot3^t(-1,2)\)}
{\(\yv=-\frac14\cdot(-1)^t(1,2)+\frac34\cdot3^t(-1,2)\)}
{\(\yv=\frac34\cdot(-1)^t(-1,2)-\frac14\cdot3^t(1,2)\)}
{\(\yv=4\cdot3^t(-1,2)\)}
%\begin{enumerate}
%\item \(\yv=\frac34\cdot(-1)^t(1,2)-\frac14\cdot3^t(-1,2)\)\actans
%\item \(\yv=-\frac14\cdot(-1)^t(1,2)+\frac34\cdot3^t(-1,2)\)
%\item \(\yv=\frac34\cdot(-1)^t(-1,2)-\frac14\cdot3^t(1,2)\)
%\item \(\yv=4\cdot3^t(-1,2)\)
%\end{enumerate}
\end{activity}




Now we establish that the same sort of general solution occurs for all such models.


\begin{theorem} \label{thm:dynsol}
Suppose the \(n\times n\) \idx{square matrix}~\(A\) governs the dynamics of \(\yv(t)\in\RR^n\) according to \(\yv(t+1)=A\yv(t)\).
\begin{enumerate}
\item\label{thm:dynsola} Let \hlist\lambda m\ be \idx{eigenvalue}s of~\(A\) and \hlist\vv m\ be corresponding \idx{eigenvector}s, then a solution of \(\yv(t+1)=A\yv(t)\) is the \idx{linear combination}
\begin{equation}
\yv(t)=c_1\lambda_1^t\vv_1+c_2\lambda_2^t\vv_2+\cdots+c_m\lambda_m^t\vv_m
\label{eq:dynsol}
\end{equation}
for all constants \hlist cm.

\sloppy
\item\label{thm:dynsolb} Further, if 
%the number of eigenvectors \(m=n\) (the size of the matrix~\(A\)), and 
the matrix of eigenvectors \(P=\begin{bmatrix} \vv_1&\vv_2&\cdots&\vv_m \end{bmatrix}\) is \idx{invertible}, then the general linear combination~\eqref{eq:dynsol} is a \bfidx{general solution} in that unique constants \hlist cm\ may be found for every given \idx{initial value}~\(\yv(0)\).
\end{enumerate}
\end{theorem}

\begin{proof} \ 
\begin{itemize}
\item[\ref{thm:dynsola}] 
Just premultiply~\eqref{eq:dynsol} by matrix~\(A\) to find that
\begin{eqnarray*}
A\yv(t)&=&A(c_1\lambda_1^t\vv_1+c_2\lambda_2^t\vv_2+\cdots+c_m\lambda_m^t\vv_m)
\\&&\quad(\text{using distributivity \autoref{thm:pasm}})
\\&=&c_1\lambda_1^tA\vv_1+c_2\lambda_2^tA\vv_2+\cdots+c_m\lambda_m^tA\vv_m
\\&&\quad(\text{then as eigenvectors }A\vv_j=\lambda_j\vv_j)
\\&=&c_1\lambda_1^t\lambda_1\vv_1+c_2\lambda_2^t\lambda_2\vv_2+\cdots+c_m\lambda_m^t\lambda_m\vv_m
\\&=&c_1\lambda_1^{t+1}\vv_1+c_2\lambda_2^{t+1}\vv_2+\cdots+c_m\lambda_m^{t+1}\vv_m
\end{eqnarray*}
which is the given formula~\eqref{eq:dynsol} for~\(\yv(t+1)\).
Hence \eqref{eq:dynsol} is a solution of \(\yv(t+1)=A\yv(t)\) for all constants \hlist cm.

\item[\ref{thm:dynsolb}]
For every given initial value~\(\yv(0)\), the solution~\eqref{eq:dynsol} will hold if we can find constants \hlist cm\ such that the solution~\eqref{eq:dynsol} evaluates to~\(\yv(0)\) at time \(t=0\)\,.
Let's do this given the preconditions that 
%the number of eigenvectors \(m=n\) and 
the matrix \(P=\begin{bmatrix} \vv_1&\vv_2&\cdots&\vv_m \end{bmatrix}\) is {invertible}.
First, since matrix~\(P\) is invertible, it must be square, and hence \(m=n\) (that is, there must be \(n\)~eigenvectors and \(n\)~terms in~\eqref{eq:dynsol}).
Second, evaluating the solution~\eqref{eq:dynsol} at \(t=0\) gives, since the zeroth power \(\lambda_j^0=1\)\,,
\begin{equation*}
\yv(0)=\lincomb c\vv n\,,
\end{equation*}
as an equation to be solved.
Writing as a matrix-vector system this equation requires \(P\cv=\yv(0)\) for constant vector \(\cv=(\hlist cn)\).
Since matrix~\(P\) is invertible, \(P\cv=\yv(0)\) always has the unique solution \(\cv=P^{-1}\yv(0)\) (\autoref{thm:ftim2}) which determines the requisite constants.
\end{itemize}
\end{proof}




\begin{activity}
The matrix~\(A=\begin{bmatrix} 1&1\\a^2&1 \end{bmatrix}\) has eigenvectors \((1,a)\) and~\((1,-a)\).
For what value(s) of~\(a\) does \autoref{thm:dynsol} \emph{not} provide a general solution to \(\yv(t+1)=A\yv(t)\)?
\actposs[4]{\(a=0\)}{\(a=-1\)}{\(a=1\)}{\(a=\pm1\)}
%\partswidth=5em
%\begin{parts}
%\item \(a=-1\)
%\item \(a=0\)\actans
%\item \(a=1\)
%\item \(a=\pm1\)
%\end{parts}
\end{activity}




\begin{example} \label{eg:2x2ccgrow}
Consider the dynamics of \(\yv(t+1)=A\yv(t)\) for matrix \(A=\begin{bmatrix} 1&3\\-1&1 \end{bmatrix}\).
First, what is \(\yv(3)\) when the \idx{initial value} \(\yv(0)=(1,0)\)?
Second, find a \idx{general solution}.
% and interpret its behaviour in time~\(t\).
\begin{solution} 
First, just compute (as illustrated)
\begin{eqnarray*}
\yv(1)&=&A\yv(0)=\begin{bmatrix} 1&3\\-1&1 \end{bmatrix}
\begin{bmatrix} 1\\0 \end{bmatrix}
=\begin{bmatrix} 1\\-1 \end{bmatrix},
\\\yv(2)&=&A\yv(1)=\begin{bmatrix} 1&3\\-1&1 \end{bmatrix}
\begin{bmatrix} 1\\-1 \end{bmatrix}
=\begin{bmatrix} -2\\-2 \end{bmatrix},
\\\yv(3)&=&A\yv(2)=\begin{bmatrix} 1&3\\-1&1 \end{bmatrix}
\begin{bmatrix} -2\\-2 \end{bmatrix}
=\begin{bmatrix} -8\\0 \end{bmatrix}.
\end{eqnarray*}
\marginpar{\begin{tikzpicture}
\begin{axis}[footnotesize,font=\footnotesize, 
xlabel={time $t$},ylabel={$y_j$},axis lines=middle
,only marks,xmax=3.9,legend pos=south west]
\addplot+ coordinates {(0,1)(1,1)(2,-2)(3,-8)};
\addlegendentry{$y_1(t)$};
\addplot+ coordinates {(0,0)(1,-1)(2,-2)(3,0)};
\addlegendentry{$y_2(t)$};
\end{axis}
\end{tikzpicture}}%
Interestingly, after three steps in time~\(\yv(3)\) is~\((-8)\) times the initial~\(\yv(0)\).
This suggest after six steps in time \(\yv(6)\) will be \((-8)^2=64\) times the initial~\(\yv(0)\), and so on.
Perhaps the solution grows in size roughly like~\(2^t\) but in some irregular manner---let's see via a general solution.

Second, find a general solution via the eigenvalues and eigenvectors of the matrix~\(A\).  
Its characteristic equation is 
\begin{equation*}
\det(A-\lambda I)=\begin{vmatrix} 1-\lambda&3\\-1&1-\lambda \end{vmatrix}
=(1-\lambda)^2+3=0\,.
\end{equation*}
That is, \((\lambda-1)^2=-3\) which upon taking square roots gives the complex conjugate pair of eigenvalues \(\lambda=1\pm \i\sqrt3\).
\autoref{thm:dynsol} applies for \idx{complex eigenvalue}s and eigenvectors so we proceed.
\begin{itemize}
\item For eigenvalue \(\lambda_1=1+\i\sqrt3\) the corresponding eigenvectors~\(\vv_1\) satisfy
\begin{equation*}
\big(A-(1+\i\sqrt3)I\big)\vv_1=\begin{bmatrix} -\i\sqrt3&3\\-1&-\i\sqrt3 \end{bmatrix}\vv_1=\ov\,.
\end{equation*}
Solutions are proportional to \(\vv_1=(-\i\sqrt3,1)\).

\item For eigenvalue \(\lambda_2=1-\i\sqrt3\) the corresponding eigenvectors~\(\vv_2\) satisfy
\begin{equation*}
\big(A-(1-\i\sqrt3)I\big)\vv_2=\begin{bmatrix} \i\sqrt3&3\\-1&\i\sqrt3 \end{bmatrix}\vv_2=\ov\,.
\end{equation*}
Solutions are proportional to \(\vv_2=(+\i\sqrt3,1)\).
\end{itemize}
\autoref{thm:dynsol} then establishes that a solution to \(\yv(t+1)=A\yv(t)\) is
\begin{equation*}
\yv(t)=c_1(1+\i\sqrt3)^t\begin{bmatrix} -\i\sqrt3\\1 \end{bmatrix}
+c_2(1-\i\sqrt3)^t\begin{bmatrix} +\i\sqrt3\\1 \end{bmatrix}.
\end{equation*}
This is a general solution since the matrix of the two eigenvectors (albeit complex valued) is invertible:
\begin{equation*}
P=\begin{bmatrix} \vv_1&\vv_2 \end{bmatrix}
=\begin{bmatrix} -\i\sqrt3&\i\sqrt3\\1&1 \end{bmatrix}
\quad\text{has }
P^{-1}=\begin{bmatrix} \frac \i{2\sqrt3}&\frac12\\\frac{-\i}{2\sqrt3}&\frac12 \end{bmatrix}
\end{equation*}
as its inverse (\autoref{thm:2x2det}).

For example, if \(\yv(0)=(1,0)\), then the coefficient constants are \((c_1,c_2)=P^{-1}(1,0)=(\i,-\i)/(2\sqrt3)\).
Then the solution becomes
\begin{eqnarray*}
\yv(t)&=&\frac \i{2\sqrt3}(1+\i\sqrt3)^t\begin{bmatrix} -\i\sqrt3\\1 \end{bmatrix}
-\frac{\i}{2\sqrt3}(1-\i\sqrt3)^t\begin{bmatrix} +\i\sqrt3\\1 \end{bmatrix}
\\&=&\tfrac 12(1+\i\sqrt3)^t\begin{bmatrix} 1\\\frac \i{\sqrt3} \end{bmatrix}
+\tfrac12(1-\i\sqrt3)^t\begin{bmatrix} 1\\-\frac \i{\sqrt3} \end{bmatrix}.
\end{eqnarray*}
Through the magic of the complex conjugate form of the two terms in this expression, the complex parts cancel to always give a real result.
For example, this complex formula predicts at time step \(t=1\)
\begin{eqnarray*}
\yv(1)&=&\tfrac 12(1+\i\sqrt3)\begin{bmatrix} 1\\\frac \i{\sqrt3} \end{bmatrix}
+\tfrac12(1-\i\sqrt3)\begin{bmatrix} 1\\-\frac \i{\sqrt3} \end{bmatrix}
\\&=&\frac12\begin{bmatrix} 1+\i\sqrt3+1-\i\sqrt3\\
\frac \i{\sqrt3}-1-\frac \i{\sqrt3}-1 \end{bmatrix}
\\&=&\begin{bmatrix} 1\\-1 \end{bmatrix},
\end{eqnarray*}
as computed directly at the start of this example.
\end{solution}
\end{example}

One crucial qualitative aspect we need to know is whether components in the solution~\eqref{eq:dynsol} grow, decay, or stay the same size as time increases.
The growth or decay is determined by the eigenvalues:
the reason is that~\(\lambda_j^t\) is the only place that the time appears in the formula~\eqref{eq:dynsol}.
\begin{itemize}
\item For example, in the general solution for \autoref{eg:2x2map}, \(\yv(t)=c_1(-1)^t(1,2)+c_23^t(-1,2)\), the \(3^t\)~factor grows in time since \(3^1=3\), \(3^2=9\), \(3^3=27\), and so on.
Whereas the \((-1)^t\)~factor just oscillates in time since \((-1)^1=-1\), \((-1)^2=1\), \((-1)^3=-1\), and so on.
Thus for long times, large~\(t\), we know that the term involving the factor~\(3^t\) will dominate the solution as it grows.

\item In \autoref{eg:2x2ccgrow} with complex conjugate eigenvalues the situation is more complicated.
Let's write every given \idx{complex eigenvalue} in \idx{polar form} 
\(\lambda=r(\cos\theta+\i\sin\theta)\) where \idx{magnitude} \(r=|\lambda|\) and \idx{angle}~\(\theta\) is such that \(\tan\theta=(\Im\lambda)/(\Re\lambda)\).
For example, \(1-\i\) has magnitude \(r=|1-1\i|=\sqrt{1^2+(-1)^2}=\sqrt2\) and angle \(\theta=-\frac\pi4\) since \(\tan(-\frac\pi4)=-1/1\)\,.

Question: how does this help understand the solution which has \(\lambda_j^t\) in it?
Answer: \idx{De Moivre's theorem} says that if \(\lambda=r[\cos\theta+\i\sin\theta]\), then \(\lambda^t=r^t\big[\cos(\theta t)+\i\sin(\theta t)\big]\).
Since the magnitude \(|\cos(\theta t)+\i\sin(\theta t)|=\sqrt{\cos^2(\theta t)+\sin^2(\theta t)}=\sqrt1=1\)\,, the magnitude \(|\lambda^t|=r^t\).
For example, the magnitude \(|(1-\i)^2|=(\sqrt2)^2=2\) which we check by computing \((1-\i)^2=1^2-2\i+\i^2=-2\i\) and \(|-2\i|=2\)\,.

In \autoref{eg:2x2ccgrow}, the eigenvalue \(\lambda_1=1+\i\sqrt3\) so its magnitude is \(r_1=|\lambda_1|=|1+\i\sqrt3|=\sqrt{1+3}=2\)\,.
%Its angle \(\theta_1=\frac\pi3\) because \(\tan\frac\pi3=\sqrt3=\sqrt3/1\)\,.
Hence the magnitude \(|\lambda_1^t|=2^t\) at every time step~\(t\).
Similarly, the magnitude \(|\lambda_2^t|=2^t\) at every time step~\(t\).
Consequently, the general solution
\begin{equation*}
\yv(t)=c_1\lambda_1^t\begin{bmatrix} -\i\sqrt3\\1 \end{bmatrix}
+c_2\lambda_2^t\begin{bmatrix} +\i\sqrt3\\1 \end{bmatrix}
\end{equation*}
will grow in magnitude roughly like~\(2^t\) as both components grow like~\(2^t\).
It is a `rough' growth because the components \(\cos(\theta t)\) and~\(\sin(\theta t)\)  cause `oscillations' in time~\(t\).
Nonetheless the overall growth like \(|\lambda_1|^t=|\lambda_2|^t=2^t\) is inexorable---and seen previously in the particular solution where we observe \(\yv(3)\) is eight times the magnitude of~\(\yv(0)\).
\end{itemize}
In general, for both real or \idx{complex eigenvalue}s~\(\lambda\), a term involving the factor~\(\lambda^t\) will, as time~\(t\) increases,
\marginpar{\begin{tikzpicture}
\begin{axis}[footnotesize,font=\footnotesize,domain=0:10, 
xlabel={time $t$},ylabel={$|\lambda|^t$},axis lines=middle
,only marks,samples=11,ymin=0,xmax=11.5]
\addplot+ {1.16^x};
\addlegendentry{$|\lambda|>1$};
\addplot+ {0.8^x};
\addlegendentry{$|\lambda|<1$};
\addplot+ {1^x};
\addlegendentry{$|\lambda|=1$};
\end{axis}
\end{tikzpicture}}
\begin{itemize}
\item grow to infinity if \(|\lambda|>1\)\,,
\item decay to zero if \(|\lambda|<1\)\,, and
\item remain the same \idx{magnitude} if \(|\lambda|=1\)\,.
\end{itemize}




\begin{activity}
For which of the following values of~\(\lambda\), as time~\(t\) increases, will \(\lambda^t\) grow in an oscillatory fashion?
\actposs[4]{\(\lambda=-\frac45+\i\frac45\)}
{\(\lambda=-0.8\)}
{\(\lambda=1.5\)}
{\(\lambda=\frac35+\i\frac45\)}
%\partswidth=5em
%\begin{parts}
%\item \(\lambda=-0.8\)
%\item \(\lambda=1.5\)
%\item \(\lambda=\frac35+\i\frac45\)
%\item \(\lambda=-\frac45+\i\frac45\) \actans
%\end{parts}
\end{activity}




\begin{example}[orangutans over many years] \label{eg:orangs}
Extend the orangutan analysis of \autoref{eg:orang}. 
Use \autoref{thm:dynsol} to predict the population over many years: from an \idx{initial population} of \(30\)~adolescent females and \(15\)~adult females; and from a general initial population.
\begin{solution} 
\autoref{eg:orang} derived that the age structure population \(\yv=(y_1,y_2,y_3)\) satisfies \(\yv(t+1)=A\yv(t)\) for matrix
\begin{equation*}
A=\begin{bmatrix} \frac45&0&\frac1{15}
\\\frac15&\frac9{10}&0
\\0&\frac1{10}&\frac{14}{15} \end{bmatrix}.
\end{equation*}
Let's find the eigenvalues and eigenvectors of the matrix~\(A\) using \script\ via
\begin{verbatim}
A=[4/5 0 1/15;1/5 9/10 0;0 1/10 14/15]
[V,D]=eig(A)
\end{verbatim}
\setbox\ajrqrbox\hbox{\qrcode{% orangutans
A=[4/5 0 1/15;1/5 9/10 0;0 1/10 14/15]
[V,D]=eig(A)
y0=[0;30;15]
rcond(V)
c=V\slosh y0
}}%
\marginpar{\usebox{\ajrqrbox}}%
to find
\begin{verbatim}
V =
 -0.3077+0.2952i -0.3077-0.2952i  0.2673+0.0000i
  0.7385+0.0000i  0.7385+0.0000i  0.5345+0.0000i
 -0.4308-0.2952i -0.4308+0.2952i  0.8018+0.0000i
D =
  0.8167+0.0799i  0.0000+0.0000i  0.0000+0.0000i
  0.0000+0.0000i  0.8167-0.0799i  0.0000+0.0000i
  0.0000+0.0000i  0.0000+0.0000i  1.0000+0.0000i
\end{verbatim}
Evidently there is one real eigenvalue of \(\lambda_3=1\) and two complex conjugate eigenvalues \(\lambda_{1,2}=0.8167\pm \i0.0799\)\,.
Corresponding eigenvectors are the columns~\(\vv_j\) of~\(V\).
Thus a solution for the orangutan population is
\begin{equation*}
\yv(t)=c_1\lambda_1^t\vv_1+c_2\lambda_2^t\vv_2+c_3\lambda_3^t\vv_3\,.
\end{equation*}
\begin{itemize}
\item For the initial population \(\yv(0)=(0,30,15)\) we need to find constants \(\cv=(c_1,c_2,c_3)\) such that \(V\cv=\yv(0)\).
Solve this linear equation in \script\ with \autoref{pro:unisol}:
\begin{verbatim}
y0=[0;30;15]
rcond(V)
c=V\y0
\end{verbatim}
which gives the answer
\begin{verbatim}
ans =
    0.1963
c =
  10.1550+2.1175i
  10.1550-2.1175i
  28.0624+0.0000i
\end{verbatim}
The \verb|rcond| value of~\(0.1963\) indicates that matrix~\(V\) is invertible.
Then the backslash operator computes the above coefficients~\cv.
Via the magic of complex conjugates cancelling, the real population of orangutans is for all times predicted to be \twodp
\begin{eqnarray*}
\yv(t)&=&(10.16+2.12\i)(0.82+0.08\i)^t\begin{bmatrix} -0.31+0.30\i   \\0.74
\\-0.43-0.30\i \end{bmatrix}
\\&&{}
+(10.16-2.12\i)(0.82-0.08\i)^t\begin{bmatrix} -0.31-0.30\i   \\0.74
\\-0.43+0.30\i \end{bmatrix}
\\&&{}
+28.06\begin{bmatrix} 0.27
\\0.53
\\0.80 \end{bmatrix}
\end{eqnarray*}
since \(\lambda_3^t=1^t=1\)\,.

Since the magnitude \(|\lambda_1|=|\lambda_2|=0.82\) \twodp, 
the first two terms in this expression decay to zero as time~\(t\) increases.
For example, \(|\lambda_1^{12}|=|\lambda_2^{12}|=0.09\)\,.
Hence the model predicts that over long times the population
\begin{equation*}
\yv(t)\approx 28.06\begin{bmatrix} 0.27
\\0.53
\\0.80 \end{bmatrix}
=\begin{bmatrix} 7.5\\15.0\\22.5 \end{bmatrix}
\end{equation*}
Such a static population means that the orangutans are highly sensitive to disease, or deforestation, or chance events, and so on.

\item Such unfortunate sensitivity is typical for orangutans.  
It is not a quirk of the initial population.
Recall the general prediction for the orangutans is 
\begin{equation*}
\yv(t)=c_1\lambda_1^t\vv_1+c_2\lambda_2^t\vv_2+c_3\lambda_3^t\vv_3\,.
\end{equation*}
The initial population determines the constants~\(\cv\).
However, the long term population is always predicted to be static.
The reason is that the magnitude of the eigenvalues \(|\lambda_1|=|\lambda_2|=0.82\) and so the first two terms in this general solution will in time always decay to zero.
Further, the remaining third eigenvalue has magnitude \(|\lambda_3|=|1|=1\) and so the third term in the population prediction is always constant in time~\(t\).
That is, over long times the population is always
\begin{equation*}
\yv(t)\approx c_3\vv_3\,.
\end{equation*}
Such a static population means that the orangutans are always sensitive to  disease, or deforestation, or chance events, and so on.

\end{itemize}
\end{solution}
\end{example}




In \emph{Guidelines for Assessment and Instruction in Mathematics Modeling Education}, \cite{gaimme2016} discuss mathematical modelling.
\begin{itemize}
\item On page~23 they comment ``\emph{Modelling (like real life) is open-ended and messy}'': in our two examples here you have to extract the important factors from many unneeded details, and use them in the context of an imperfect model.
\item Also on p.23, modellers ``\emph{must be making genuine choices}'': in these problems, as in all modelling, there are choices that lead to different models---we have to operate and sensibly predict with such uncertainty.
\item Lastly, they recommend to ``\emph{focus on the process, not the product}'': depending upon your choices and interpretations you will develop alternative plausible models in these scenarios---it is the process of forming plausible models and interpreting the results that are important.
\end{itemize}



\begin{example}[servals grow] \label{eg:serval} 
The \idx{serval} is a member of the cat family that lives in Africa.
Given next is an extract from \idx{Wikipedia} of a serval's Reproduction and Life History.
\marginpar{\includegraphics{EigGen/P1040280}}
\begin{quote}
Kittens are born shortly before the peak breeding period of local rodent populations. A serval is able to give birth to multiple litters throughout the year, but commonly does so only if the earlier litters die shortly after birth. Gestation lasts from~66 to 77~days and commonly results in the birth of two kittens, although sometimes as few as one or as many as four have been recorded.

The kittens are born in dense vegetation or sheltered locations such as abandoned aardvark burrows. If such an ideal location is not available, a place beneath a shrub may be sufficient. The kittens weigh around~250\,gm at birth, and are initially blind and helpless, with a coat of greyish woolly hair. They open their eyes at~9 to 13~days of age, and begin to take solid food after around a month. At around six months, they acquire their permanent canine teeth and begin to hunt for themselves; they leave their mother at about 12~months of age. They may reach sexual maturity from~12 to 25~months of~age.

Life expectancy is about 10~years in the wild.
\end{quote}
From the information in this extract, create a plausible, \idx{age structure}d,  \idx{population model} of \idx{serval}s: give reasons for estimates of the coefficients in the model.
Choose three age categories of \idx{kitten}s, \idx{juvenile}s, sexually mature \idx{adult}s.
What does the model predict over long times?
\begin{solution} 
Recall we only model the number of \emph{female} servals as females are the limiting breeders.
Define 
\begin{itemize}
\item \(y_1(t)\) is the number of female kittens, less than 0.5~years old from when they ``begin to hunt for themselves'';
\item \(y_2(t)\) is the number of female juveniles, between 0.5~years and 1.5~years which is when they ``reach sexual maturity'' on average;
\item \(y_3(t)\) is the number of female breeding adults, older than 1.5~years, and dying at about the ``life expectancy'' of 10~years;
\item since servals transition from one age category to another in multiples of six months (0.5~years), let the unit of time be six~months, equivalently a half-year.  
Consequently, time~\(t+1\) is the time a half-year later than time~\(t\).
\end{itemize}
Modelling of the servals leads to the following equations.
\begin{itemize}
\item All kittens age to juveniles after 0.5~years, so none remain as kittens.  
Hence the model has 
\begin{equation*}
y_1(t+1)=0y_1(t)+\cdots,\quad
y_2(t+1)=1y_1(t)+\cdots\,.
\end{equation*}
Kittens are commonly born once a year to each female, and the common litter size is two, so \emph{on average} one female kitten is born per year per adult, that is, on average \(\frac12\)~female kitten is born per half-year per adult female: hence the kitten model is
\begin{equation*}
y_1(t+1)=0y_1(t)+\tfrac12y_3(t).
\end{equation*}

\item Juveniles mature from the kittens, and age to an adult after about one year: that is, on average half of them remain juveniles every half-year, and half become adults.
So the model has 
\begin{equation*}
y_2(t+1)=y_1(t)+\tfrac12y_2(t),\quad
y_3(t+1)=\tfrac12y_2(t)+\cdots\,.
\end{equation*}

\item Adults mature from the juveniles, and die after about 8.5~years which is about a rate \(1/8.5\)~per year: that is, a rate of \(\frac1{17}\)~per half-year leaving \(\tfrac{16}{17}\) of them to live into the next half-year.
So the adult model completes to 
\begin{equation*}
y_3(t+1)=\tfrac12y_2(t)+\tfrac{16}{17}y_3(t).
\end{equation*}

\end{itemize}

Bring these equations together, the age structure population \(\yv=(y_1,y_2,y_3)\) satisfies \(\yv(t+1)=A\yv(t)\) for matrix
\begin{equation*}
A=\begin{bmatrix} 0&0&\frac1{2}
\\1&\frac12&0
\\0&\frac12&\frac{16}{17} \end{bmatrix}.
\end{equation*}
Find the eigenvalues and eigenvectors of the matrix~\(A\) using \script\ via
\begin{verbatim}
A=[0 0 1/2;1 1/2 0;0 1/2 16/17]
[V,D]=eig(A)
\end{verbatim}
\setbox\ajrqrbox\hbox{\qrcode{% servals
A=[0 0 1/2;1 1/2 0;0 1/2 16/17]
[V,D]=eig(A)
}}%
\marginpar{\usebox{\ajrqrbox}}%
to find
\begin{verbatim}
V =
 -0.3066+0.3439i -0.3066-0.3439i  0.3352+0.0000i
  0.7838+0.0000i  0.7838+0.0000i  0.4633+0.0000i
 -0.3684-0.1942i -0.3684+0.1942i  0.8203+0.0000i
D =
  0.1088+0.4387i  0.0000+0.0000i  0.0000+0.0000i
  0.0000+0.0000i  0.1088-0.4387i  0.0000+0.0000i
  0.0000+0.0000i  0.0000+0.0000i  1.2236+0.0000i
\end{verbatim}
Evidently there is one real eigenvalue of \(\lambda_3=1.2236\) and two complex conjugate eigenvalues \(\lambda_{1,2}=0.1088\pm \i0.4387\)\,.
Corresponding eigenvectors are the columns~\(\vv_j\) of~\(V\).
Thus a general solution for the serval population is (\autoref{thm:dynsol})
\begin{equation*}
\yv(t)=c_1\lambda_1^t\vv_1+c_2\lambda_2^t\vv_2+c_3\lambda_3^t\vv_3\,.
\end{equation*}

In this general solution, the first two terms decay in time to zero.
The reason is that the magnitudes \(|\lambda_1|=|\lambda_2|=|0.1088\pm \i0.4387|=\sqrt{0.1088^2+0.4387^2}=0.4520\), and since this magnitude is less than one, then \(\lambda_1^t\) and~\(\lambda_2^t\) will decay to zero with increasing time~\(t\).
However, the third term increases in time as \(\lambda_3=1.2236>1\)\,.
The model predicts the serval population increases by about~22\% per half-year (about~50\% per year).
\end{solution} 
Predation, disease, and food shortages are just some processes not included in this model which act to limit the serval's population in ways not included in this model.
\end{example}





Crucial in this section---so that we find a solution for all \idx{initial value}s---is that the matrix of eigenvectors is invertible.
The next \autoref{sec:lisb} relates the invertibility of a matrix of eigenvectors to the new concept of `linear independence' (\autoref{thm:ftim3}).



\index{age structure|)}
\index{population|)}





\begin{comment}
Could have a subsection on Metastability in Markov Chains---probably not.
Ensure application-like.
\end{comment}








\subsection{Extension: SVDs connect to eigen-problems}

Recall that \autoref{ch:eesm} starts by illustrating the close connection between the \svd\ of a symmetric matrix and the eigenvalues and eigenvectors of that symmetric matrix.
\begin{aside}
This optional section connects the \svd\ of a general matrix to a symmetric eigen-problem, in principle.
\end{aside}%
This subsection establishes that an \svd\ of a general matrix is closely connected to the eigenvalues and eigenvectors of a specific matrix of double the size.   
The connection depends upon determinants and solving linear systems and so, in principle, is an approach to compute an \svd\ distinct from the inductive maximisation of \autoref{sec:psvdt}.



\begin{example} \label{eg:eigsvd}
Compute the eigenvalues and eigenvectors of the (symmetric) matrix
\begin{equation*}
B=\begin{bmatrix} 0&0&10&2\\
0&0&5&11\\
10&5&0&0\\
2&11&0&0 \end{bmatrix}.
\qquad \text{For matrix }A=\begin{bmatrix} 10&2\\5&11 \end{bmatrix},
\end{equation*}
compare with an \svd\ of~\(A\).
\begin{solution} In \script\ execute
\setbox\ajrqrbox\hbox{\qrcode{% eigenvalue/vector code
B=[0 0 10 2
   0 0 5 11
  10 5 0 0
  2 11 0 0]
[V,D]=eig(B)
}}%
\marginpar{\usebox{\ajrqrbox}}%
\begin{verbatim}
B=[0 0 10 2
   0 0 5 11
  10 5 0 0
  2 11 0 0]
[V,D]=eig(B)
\end{verbatim}
and obtain \twodp
\begin{verbatim}
V =
   0.42   0.57   0.57  -0.42
   0.57  -0.42  -0.42  -0.57
  -0.50  -0.50   0.50  -0.50
  -0.50   0.50  -0.50  -0.50
D =
  -14.14       0       0       0
       0   -7.07       0       0
       0       0    7.07       0
       0       0       0   14.14
\end{verbatim}
The eigenvalues are the pairs \(\pm7.07\) and \(\pm14.14\), with corresponding eigenvector pairs \((0.57,-0.42,\pm0.50,\mp0.50)\) and \((\mp0.42,\mp0.57,-0.50,-0.50)\).

%Whereas from \autoref{eg:2by2svd}, an \svd\ of matrix~\(A\) is
%\begin{verbatim}
%U =
%  -0.60  -0.80
%  -0.80   0.60
%S =
%   14.14       0
%       0    7.07
%V =
%  -0.71  -0.71
%  -0.71   0.71
%\end{verbatim}

These eigenvalues\slash vectors occur in \(\pm\)~pairs because this matrix has the form
\begin{equation*}
B=\begin{bmatrix} O_2& A\\ \tr A&O_2\end{bmatrix},
\quad\text{here for matrix }
A=\begin{bmatrix} 10&2\\5&11 \end{bmatrix}
\end{equation*}
from \autoref{eg:2by2svd}.
Observe that not only are the eigenvectors orthogonal, because \(B\)~is symmetric, but also the two parts of the eigenvectors are orthogonal:  
\begin{itemize}
\item the components \((0.57,-0.42)\) from the first pair is orthogonal to \((-0.42,-0.57)\) from the second pair; and 
\item the components \((0.50,-0.50)\) from the first pair is orthogonal to \((-0.50,-0.50)\) from the second pair.
\end{itemize}
The next \autoref{thm:eigsvd} establishes how these properties relate to an \svd\ for the matrix~\(A\).
\end{solution}
\end{example}


\autoref{pro:geneig} computes eigenvalues and eigenvectors by hand (in principle no matter how large the matrix). 
The procedure is independent of the \svd.
Let's now invoke this procedure to establish another method to find an \svd\ distinct from the inductive maximisation of the proof in \autoref{sec:psvdt}. 
The following \autoref{thm:eigsvd} is a step towards an efficient numerical computation of an \svd\ \cite[p.234]{Trefethen1997}.

% Trefethen has the transpose in the upper right---curses??


\begin{theorem}[\svd\ as an eigenproblem] \label{thm:eigsvd}
\index{singular value decomposition}\index{SVD}%??
For every real \(m\times n\) matrix~\(A\), the \idx{singular value}s of~\(A\) are the non-negative \idx{eigenvalue}s of the \((m+n)\times(m+n)\) \idx{symmetric matrix} \(B=\begin{bmatrix} O_m&A\\\tr A&O_n \end{bmatrix}\). 
Each corresponding \idx{eigenvector}~\(\wv\in\RR^{m+n}\)\ of~\(B\) gives corresponding  singular vectors\index{singular vector} of~\(A\), namely \(\wv=(\uv,\vv)\) for singular vectors \(\uv\in\RR^m\)\ and \(\vv\in\RR^n\). 
\end{theorem}

\begin{proof} 
First prove the \svd\ of~\(A\) gives eigenvalues\slash vectors of~\(B\), and second prove the converse.
For simplicity this proof addresses only the case \(m=n\); the case \(m\neq n\) is similar but the more intricate details are of little interest.

First, let \(n\times n\) matrix \(A=\usv\) be an \svd\ (\autoref{thm:svd}) for \(n\times n\) orthogonal \(U=\begin{bmatrix} \uv_1&\uv_2&\cdots&\uv_n \end{bmatrix}\),  orthogonal \(V=\begin{bmatrix} \vv_1&\vv_2&\cdots&\vv_n \end{bmatrix}\), and diagonal \(S=\diag(\hlist\sigma n)\).
Post-multiply the \svd\ by orthogonal~\(V\) gives \(AV=US\)\,.
% that is, \(A\vv_j=\sigma_j\uv_j\) for \(j=1,\ldots,n\)\,.
Also, transpose the \svd\ to \(\tr A=\tr{(\usv)}=V\tr S\tr U=VS\tr U\) and post-multiply by orthogonal~\(U\) to give \(\tr AU=VS\)\,.
% that is, \(\tr A\uv_j=\sigma_j\vv_j\) for \(j=1,\ldots,n\)\,.
Now consider each of the \(\pm\)~cases of
\begin{equation*}
B\begin{bmatrix} U\\\pm V \end{bmatrix}
=\begin{bmatrix} O_n&A\\\tr A&O_n \end{bmatrix}\begin{bmatrix} U\\\pm V \end{bmatrix}
=\begin{bmatrix} \pm AV\\\tr AU \end{bmatrix}
=\begin{bmatrix} \pm US\\VS \end{bmatrix}
=\begin{bmatrix} U\\\pm V \end{bmatrix}(\pm S).
\end{equation*}
Setting vector \(\wv=(\uv_j,\pm\vv_j)\neq\ov\)\,, then the \(j\)th~column of the above equation is \(B\wv=\pm\sigma_j\wv\) and hence \((\uv_j,\pm\vv_j)\) is an eigenvector of~\(B\) corresponding to eigenvalue~\(\pm\sigma_j\), respectively, for \(j=1,\ldots,n\) and each of the \(\pm\)~cases.
This completes the first part of the proof.

Second, let \(\wv\in\RR^{2n}\)\ be an eigenvector of~\(B\) corresponding to an eigenvalue~\(\lambda\) (real as \(B\)~is symmetric, \autoref{thm:realeigsym}) and normalised so that \(|\wv|=1\)\,.
Separate the components of~\wv\ into two: \(\wv=(\uv,\vv)\) for \(\uv,\vv\in\RR^n\).
Then the fundamental eigen-problem \(B\wv=\lambda\wv\)  (\autoref{def:evecval}) separates into
\begin{equation}
\begin{bmatrix} O& A\\\tr A&O \end{bmatrix}\begin{bmatrix} \uv\\\vv \end{bmatrix}=\lambda\begin{bmatrix} \uv\\\vv \end{bmatrix}
\iff  A\vv=\lambda\uv \text{ and }\tr A\uv=\lambda\vv\,.
\label{eq:eigsvd}
\end{equation}
Further, the eigenvalues and eigenvectors of~\(B\) come in pairs since the eigenvalue \(\lambda'=-\lambda\) corresponds to eigenvector \(\wv'=(\uv,-\vv)\): substitute into~\eqref{eq:eigsvd} to check, \(A(-\vv)=-A\vv=-\lambda\uv=\lambda'\uv\) and \(\tr A\uv=\lambda\vv=(-\lambda)(-\vv)=\lambda'(-\vv)\).
If the eigenvalue~\(\lambda\neq0\)\,, then corresponding to the distinct eigenvalues~\(\pm\lambda\) the eigenvectors \((\uv,\pm\vv)\) of symmetric~\(B\) are orthogonal (\autoref{thm:orthoevec}) and so \(0=(\uv,\vv)\cdot(\uv,-\vv)=\uv\cdot\uv-\vv\cdot\vv=|\uv|^2-|\vv|^2\) and hence \(|\uv|=|\vv|=\frac1{\sqrt2}\) as \(|\wv|^2=|\uv|^2+|\vv|^2=1\) (see \autoref{eg:eigsvd}).
Further, for any two distinct eigenvalues \(|\lambda_i|\neq|\lambda_j|\neq 0\), the eigenvectors \((\uv_i,\vv_i)\) and \((\uv_j,\pm\uv_j)\) are also orthogonal, hence \(0 =(\uv_i,\vv_i)\cdot(\uv_j,\pm\uv_j) =\uv_i\cdot\uv_j \pm\vv_i\cdot\vv_j\)\,.
Taking the sum and the difference of the \(\pm\)~cases of this equation gives \(\uv_i\cdot\uv_j=\vv_i\cdot\vv_j=0\)\,; that is,  (\autoref{def:orthoset}) \(\{\hlist\uv n\}\) and \(\{\hlist\vv n\}\) are both orthogonal sets.
In the cases when matrix~\(B\) has \(2n\)~distinct non-zero eigenvalues, choose \((\uv_j,\vv_j)\) to be a normalised eigenvector corresponding to positive eigenvalue~\(\lambda_j\), \(j=1,\ldots,n\)\,. 
Upon setting \(U=\sqrt2\begin{bmatrix} \uv_1&\uv_2&\cdots&\uv_n \end{bmatrix}\), \(V=\sqrt2\begin{bmatrix} \vv_1&\vv_2&\cdots&\vv_n \end{bmatrix}\), and \(S=\diag(\hlist\lambda n)\), equation~\eqref{eq:eigsvd} gives the columns of \(AV=US\) which since the columns of \(U\) and~\(V\) are orthonormal gives the \svd\ \(A=\usv\) for singular vectors \(\uv_j,\vv_j\in\RR^n\) and singular values~\(\lambda_j(>0)\).

Extensions of this proof cater for the case when zero is an eigenvalue and/or eigenvalues are repeated and/or the dimensions \(m\neq n\)\,.
\end{proof}










\subsection{Application: Exponential interpolation discovers dynamics}
\label{sec:eidd}

\index{exponential interpolation|(}

Many applications require identification of rates and frequencies \cite[e.g.]{Pereyra2010}:
\begin{aside}
This optional subsection develops a method useful in many modern applied disciplines.
\end{aside}%
as a played musical note decays, what is its frequency?
in the observed vibrations of a complicated bridge, what are its natural modes?
in measurements of complicated bio-chemical reactions, what rates can be identified?
All such tasks require fitting a sum of exponential functions to the data.



\begin{example} \label{eg:eidd1}
This example is the simplest case of fitting one exponential to two data points.
Suppose we take two measurements of some process: 
\begin{itemize}
\item at time \(t_1=1\) we measure the value \(f_1=5\)\,, and 
\item at time \(t_2=3\) we measure the value \(f_2=10\)\,.
\end{itemize}
Find an exponential fit to this data of the form \(f(t)=ce^{rt}\) for some as yet unknown coefficients~\(c\) and rate~\(r\).
\begin{solution} 
The classic approach is to take logarithms of both sides:
\begin{equation*}
f=ce^{rt} \iff (\ln f)=(\ln c)+rt\,,
\end{equation*}
and then determine \(\ln c\) and~\(r\) by fitting a straight line through the two data points \((t,\ln f)=(1,\ln5)\) and~\((3,\ln10)\) respectively
(recall that \index{ln@\(\ln\)}``\(\ln(x)\)'' denotes the \idx{natural logarithm} of~\(x\),  computed in \script\ with \index{log()@\texttt{log()}}\verb|log(x)|, see \autoref{tbl:mtlbmops}).
This approach has the great virtue that the approach generalises to fitting lots more noisy data (\autoref{sec:asie}).
But here we take a different approach---an approach that generalises to fitting multiple exponentials.


The following basic steps correspond to complicated steps in the general procedure developed next for fitting multiple exponentials.
\begin{enumerate}
\item We start with the question: is there a transformation that gives~\(f_2\) as a linear function of~\(f_1\)?  That is, can we write \(f_2=\lambda f_1\) for some constant~\(\lambda\)?
Answer: yes; since \(f_1=5\) and \(f_2=10\) we need \(10=\lambda5\) with solution \(\lambda=2\)\,.
\item  Since \(f_2=2f_1\) we presume extrapolation to future values is reasonable via \(f_3=2f_2=2\cdot2f_1=2^2f_1\), and \(f_4=2f_3=2\cdot2^2f_1=2^3f_1\), and so on.   
That is, \(f_n=2^{j-1}f_1\) for \(j=1,2,3,\ldots\)\,.
\item But we are given that \(f_1=5\), so the exponential fit is \(f_j=5\cdot2^{j-1}\).
\item Now these values occur at times \(t_1=1\), \(t_2=3\), and so on; that is, \(t_j=2j-1\)\,.
Rearranging, for any time~\(t\) the corresponding index \(j=(t+1)/2\)\,.
The corresponding~\(f_j\) then gives~\(f\) as a function of~\(t\), namely \(f(t) =5\cdot 2^{(t+1)/2-1} =5\cdot 2^{(t-1)/2} \).

\end{enumerate}
Equivalently, we write this~\(f(t)\) in terms of the exponential function.
Recall the rule \(2^x=e^{x\ln 2}\).
Hence the fitted function \(f=5\cdot 2^{t/2-1/2} =(5/\sqrt2)e^{(t/2)\ln2}\). 
Since \(5/\sqrt2=3.5255\) and \(\frac12\ln 2=0.3466\), the exponential fit is \(f(t)=3.5255\,e^{0.3466\,t}\).
\end{solution}
\end{example}


\begin{activity}
Plotted in the margin is some points from a function~\(f(t)\).  
\marginpar{\begin{tikzpicture}[] 
\begin{axis}[footnotesize,axis lines=middle,xlabel={$t$},ylabel={$f(t)$}
  ,ymin=0,xmax=3.6,ymax=4.9,grid,xtick={0,1,2,3}]
\addplot+[domain=0:3,samples=4,only marks] {4/3^x};
\end{axis}\end{tikzpicture}}
Which of the following exponentials best represents the data plotted?
\actposs{$f\propto1/3^t$}{$f\propto1/2^t$}{$f\propto e^{-3t}$}{$f\propto e^{-t/2}$}
\end{activity}


Now let's develop the approach of \autoref{eg:eidd1} to the more complicated and interesting example of fitting the linear combination of two exponentials to four data points.


\begin{example} \label{eg:2expfit}
\newcommand{\mytemp}[1]{\begin{tikzpicture}
\begin{axis}[footnotesize,font=\footnotesize
, xlabel={$t\ $},ylabel={$f$},axis lines=middle
,ymin=0,ymax=1.3,xmax=3.7,xtick={0,1,2,3} ]
\addplot+[mark=*,only marks] coordinates { (0,1) (1,1) (2,2/3) (3,7/18)};
\ifnum1=#1\addplot+[thick,no marks,domain=0:3.99] {-3/3^x+4/2^x};\fi
\end{axis}
\end{tikzpicture}}%
Suppose in some chemical or biochemical experiment you measure the concentration of a key chemical at four times (as illustrated): 
\marginpar{\mytemp0}%
at the start of the experiment, time \(t_1=0\) you measure concentration \(f_1=1\) (in some units); 
at time \(t_2=1\) the measurement is \(f_2=1\) (again); at \(t_3=2\) the measurement is \(f_2=\frac23=0.6667\)\,; and at \(t_4=3\) the measurement is \(f_3=\frac7{18}=0.3889\) (as plotted in the margin).
We generally expect chemical reactions to decay exponentially in time. 
So our task is to find a function of the form
\begin{equation*}
f(t)=c_1e^{r_1t}+c_2e^{r_2t}.
\end{equation*}
Our aim is for this function to fit the four data points, as plotted.
The four unknown coefficients~\(c_1\) and~\(c_2\) and rates~\(r_1\) and~\(r_2\) need to be determined from the four data points.
That is, let's find these four unknowns from the data that
\begin{eqnarray*}
f(0)=1&\iff&c_1+c_2=1\,,\\
f(1)=1&\iff&c_1e^{r_1}+c_2e^{r_2}=1\,,\\
f(2)=\tfrac23&\iff&c_1e^{r_12}+c_2e^{r_22}=\tfrac23\,,\\
f(3)=\tfrac7{18}&\iff&c_1e^{r_13}+c_2e^{r_23}=\tfrac7{18}\,.
\end{eqnarray*}
These are nonlinear equations, but some algebraic tricks empower us to use our beautiful linear algebra to find the solution.
\marginpar{\mytemp1}%
After solving these four \idx{nonlinear equation}s,  we ultimately plot the function~\(f(t)\) that interpolates between the data as also shown in the margin. 


\begin{solution} 
This solution has several twists and several new ideas that turns these four nonlinear equations into two linear algebra problems!
\begin{enumerate}
\item First form the data into vectors of consecutive values:
\begin{equation*}
\fv_1=\begin{bmatrix} f_1\\f_2 \end{bmatrix}
=\begin{bmatrix} 1\\1 \end{bmatrix},\ 
\fv_2=\begin{bmatrix} f_2\\f_3 \end{bmatrix}
=\begin{bmatrix} 1\\\frac23 \end{bmatrix},\ 
\fv_3=\begin{bmatrix} f_3\\f_4 \end{bmatrix}
=\begin{bmatrix} \frac23\\\frac7{18} \end{bmatrix},\ 
\end{equation*}
and, if only we had more data, notionally on to~\(\fv_4,\fv_5,\ldots\)\,.
\item\label{eg:2expfitb}
Question: is there a \idx{linear transformation} (\autoref{sec:ilt}) that gives both \(\fv_2\) as a function of~\(\fv_1\) and \(\fv_3\) as a function of~\(\fv_2\)?  
That is, is there a \(2\times2\) matrix~\(K\) such that both \(\fv_2=K\fv_1\) and \(\fv_3=K\fv_2\)?

Answer: if there is such a matrix~\(K\), then these two equations require both
\begin{eqnarray*}
\begin{bmatrix} f_2\\f_3 \end{bmatrix}
=K\begin{bmatrix} f_1\\f_2 \end{bmatrix}
&\leftrightarrow& \begin{bmatrix} 1\\\frac23 \end{bmatrix}
=K\begin{bmatrix} 1\\1 \end{bmatrix},
\\
\begin{bmatrix} f_3\\f_4 \end{bmatrix}
=K\begin{bmatrix} f_2\\f_3 \end{bmatrix}
&\leftrightarrow& \begin{bmatrix} \frac23\\\frac7{18} \end{bmatrix}
=K\begin{bmatrix} 1\\\frac23 \end{bmatrix}.
\end{eqnarray*}
Put the two columns of these two equations side-by-side, we require the following combined equation to hold:
\begin{equation*}
\begin{bmatrix} f_2&f_3\\f_3&f_4 \end{bmatrix}
=K\begin{bmatrix} f_1&f_2\\f_2&f_3 \end{bmatrix}
\leftrightarrow 
\underbrace{\begin{bmatrix} 1&\frac23\\\frac23&\frac7{18} \end{bmatrix}}_A
=K\underbrace{\begin{bmatrix} 1&1\\1&\frac23 \end{bmatrix}}_B.
\end{equation*}
This last equation has the form \(A=KB\) for known matrices~\(A\) and~\(B\) (named~\(A\) for `After' values, and~\(B\) for `Before' values).

Post-multiplying \(A=KB\) by the inverse of matrix~\(B\) gives the matrix of the transformation as \(K=AB^{-1}\).
Here one readily calculates (\autoref{thm:2x2det})
\begin{equation*}
B^{-1}=\begin{bmatrix} -2&3\\3&-3 \end{bmatrix},
\quad\text{so}\quad
K=AB^{-1}=\begin{bmatrix} 0&1\\-\frac16& \frac56\end{bmatrix}.
\end{equation*}

\item\label{eg:2expfitc}
Since, for this matrix~\(K\), we know that \(\fv_2=K\fv_1\)\,, \(\fv_3=K\fv_2\)\,, we extrapolate into the future to predict \(\fv_4=K\fv_3\)\,, \(\fv_5=K\gv_4\)\,, and so on, where the vectors \(\fv_j=(f_j,f_{j+1})\).
That is, our model through the four data points is summarised by \(\fv_{j+1}=K\fv_j\) for \(j=1,2,3,\ldots\,\).
This is the same class of dynamic models as explored by \autoref{sec:ddp}, albeit here not for populations.
Consequently, \autoref{thm:dynsol} gives us a general algebraic solution of the model \(\fv_{j+1}=K\fv_j\) in terms of the eigenvalues and eigenvectors of matrix~\(K\).
Such an algebraic solution interpolates and extrapolates (with caution) the data.

Let's find the general solution here.
\begin{enumerate}
\item The eigenvalues of matrix~\(K\) satisfy the characteristic equation
\begin{eqnarray*}
\det(K-\lambda I)&=&\begin{vmatrix} -\lambda&1\\-\frac16& \frac56-\lambda\end{vmatrix}
=\lambda^2-\tfrac56\lambda+\tfrac16
\\&
=&(\lambda-\tfrac12)(\lambda-\tfrac13)
=0\,.
\end{eqnarray*}
So the eigenvalues are \(\lambda_1=\tfrac12\) and \(\lambda_2=\tfrac13\)\,.
\item For eigenvalue \(\lambda_1=\tfrac12\) every corresponding eigenvector satisfies
\begin{equation*}
(K-\tfrac12I)\vv_1=\begin{bmatrix} -\frac12&1\\-\frac16&\frac13 \end{bmatrix}\vv_1=\ov\,.
\end{equation*}
That is, \(\vv_1\propto (2,1)\); let's take \(\vv_1=(2,1)\).
\item For eigenvalue \(\lambda_2=\tfrac13\) every corresponding eigenvector satisfies
\begin{equation*}
(K-\tfrac13I)\vv_2=\begin{bmatrix} -\frac13&1\\-\frac16&\frac12 \end{bmatrix}\vv_2=\ov\,.
\end{equation*}
That is, \(\vv_2\propto (3,1)\); let's take \(\vv_2=(3,1)\).
\end{enumerate}
By \autoref{thm:dynsol}, a general solution of \(\fv_{j+1}=K\fv_j\) is 
\begin{equation*}
\fv_j=c_1\vv_1\lambda_1^j+c_2\vv_2\lambda_2^j
=c_1\begin{bmatrix} 2\\1 \end{bmatrix}\frac1{2^j} +c_2\begin{bmatrix} 3\\1 \end{bmatrix}\frac1{3^j}\,.
\end{equation*}


\item\label{eg:2expfitd}
The data values also determine the constants~\(c_1\) and~\(c_2\).
Recall that \(\fv_1=(1,1)\) so substituting \(j=1\) in the above general solution gives
\begin{equation*}
\begin{bmatrix} 1\\1 \end{bmatrix}
=c_1\begin{bmatrix} 1\\\tfrac12 \end{bmatrix}
+c_2\begin{bmatrix} 1\\\tfrac13 \end{bmatrix}
=\begin{bmatrix} 1&1\\\tfrac12&\tfrac13 \end{bmatrix}\begin{bmatrix} c_1\\c_2 \end{bmatrix}.
\end{equation*}
Solving this system of equations (for example, by subtracting twice the second row from the first, and subtracting three times the second row from the first) gives solution \(c_1=4\) and \(c_2=-3\).
Substitute these coefficients into the general solution to obtain the specific model
\begin{equation*}
\fv_j =\begin{bmatrix} 8\\4 \end{bmatrix}\frac1{2^j} +\begin{bmatrix} -9\\-3 \end{bmatrix}\frac1{3^j}\,.
\end{equation*}

\item\label{eg:2expfite}
Recall that this example's quest is to fit a function~\(f(t)\) to the data.
Here \(\fv_j=(f_j,f_{j+1})\) and so the first component of the above specific model gives \(f_j=8(\frac12)^j -9(\frac13)^j\)\,.
Then recall that \(t_1=0\)\,, \(t_2=1\)\,, \(t_3=2\)\,, and so on; 
that is, \(t_j=j-1\)\,.
Reverting this relation gives the index \(j=t+1\) and so in terms of time~\(t\) the model is \(f(t)=8(\frac12)^{t+1} -9(\frac13)^{t+1}\)\,.
\end{enumerate}
%\begin{comment}%%%%%%%%%%%%%%
%\begin{enumerate}
%\item \label{eg:2expfita}
%First, rather than finding the rates directly, find the multipliers \(\lambda_1=e^{r_1}\) and \(\lambda_2=e^{r_2}\) instead (then the rates \(r_1=\ln\lambda_1\) and \(r_2=\ln\lambda_2\)).
%The four equations to solve become (still nonlinear)
%\begin{eqnarray*}
%&&c_1+c_2=1\,,\\
%&&c_1\lambda_1+c_2\lambda_2=1\,,\\
%&&c_1\lambda_1^2+c_2\lambda_2^2=\tfrac23\,,\\
%&&c_1\lambda_1^3+c_2\lambda_2^3=\tfrac7{18}\,.
%\end{eqnarray*}
%
%\item \label{eg:2expfitb}
%Second, introduce the vector and two matrices
%\begin{equation*}
%\cv=\begin{bmatrix} c_1\\c_2 \end{bmatrix},\quad
%D=\begin{bmatrix} \lambda_1&0\\0&\lambda_2 \end{bmatrix},\quad
%U=\begin{bmatrix} 1&1\\\lambda_1&\lambda_2 \end{bmatrix}.
%\end{equation*}
%\begin{enumerate}
%\item Then write the first pair of the four equations as
%\begin{equation*}
%\begin{bmatrix} 1\\1 \end{bmatrix}
%=\begin{bmatrix} c_1+c_2\\ c_1\lambda_1+c_2\lambda_2\end{bmatrix}
%=\begin{bmatrix} 1&1\\\lambda_1&\lambda_2 \end{bmatrix}\begin{bmatrix} c_1\\c_2 \end{bmatrix}
%=U\cv\,.
%\end{equation*}
%
%\item Also write the second pair of the four equations, the second and third equations, as
%\begin{eqnarray*}
%\begin{bmatrix} 1\\\frac23 \end{bmatrix}
%&=&\begin{bmatrix} c_1\lambda_1+c_2\lambda_2\\c_1\lambda_1^2+c_2\lambda_2^2\end{bmatrix}
%=\begin{bmatrix} \lambda_1&\lambda_2\\\lambda_1^2&\lambda_2^2 \end{bmatrix}\begin{bmatrix} c_1\\c_2 \end{bmatrix}
%\\&=&\begin{bmatrix} 1&1\\\lambda_1&\lambda_2 \end{bmatrix}
%\begin{bmatrix} \lambda_1&0\\0&\lambda_2 \end{bmatrix}\cv
%=UD\cv\,.
%\end{eqnarray*}
%
%\item And write the third pair of the four equations, the third and fourth equations, as
%\begin{eqnarray*}
%\begin{bmatrix} \frac23\\\frac7{18} \end{bmatrix}
%&=&\begin{bmatrix} c_1\lambda_1^2+c_2\lambda_2^2\\ c_1\lambda_1^3+c_2\lambda_2^3\end{bmatrix}
%=\begin{bmatrix} \lambda_1^2&\lambda_2^2\\\lambda_1^3&\lambda_2^3 \end{bmatrix}\begin{bmatrix} c_1\\c_2 \end{bmatrix}
%\\&=&\begin{bmatrix} 1&1\\\lambda_1&\lambda_2 \end{bmatrix}
%\begin{bmatrix} \lambda_1^2&0\\0&\lambda_2^2 \end{bmatrix}\cv
%=UD^2\cv\,.
%\end{eqnarray*}
%\end{enumerate}
%
%\item \label{eg:2expfitc}
%Third, form two matrices: by adjoining side-by-side the first two of the above equalities in step~\ref{eg:2expfitb},
%\begin{equation*}
%\begin{bmatrix} 1&1\\1&\tfrac23 \end{bmatrix}
%=\begin{bmatrix} U\cv& UD\cv \end{bmatrix}
%=U\begin{bmatrix} \cv&D\cv \end{bmatrix};
%\end{equation*}
%and by adjoining side-by-side the last two of the previous equalities in step~\ref{eg:2expfitb}, 
%\begin{equation*}
%\begin{bmatrix} 1&\tfrac23\\\tfrac23&\frac7{18} \end{bmatrix}
%=\begin{bmatrix} UD\cv&UD^2\cv \end{bmatrix}
%=UD\begin{bmatrix} \cv&D\cv \end{bmatrix}.
%\end{equation*}
%Common to both of these last two equations is the matrix \(\begin{bmatrix} \cv&D\cv \end{bmatrix}\) involving the coefficients~\cv: eliminate the matrix by writing the first equation as \(\begin{bmatrix} \cv&D\cv \end{bmatrix}=U^{-1}\begin{bmatrix} 1&1\\1&\tfrac23 \end{bmatrix}\) and substituting into \(U^{-1}\)~times the second to get 
%\begin{equation*}
%U^{-1}\begin{bmatrix} 1&\tfrac23\\\tfrac23&\frac7{18} \end{bmatrix}
%=DU^{-1}\begin{bmatrix} 1&1\\1&\tfrac23 \end{bmatrix}.
%\end{equation*}
%This matrix equation only involves the unknown multipliers~\(\lambda_1\) and~\(\lambda_2\)  (withinin \(U\) and~\(D\)), the coefficients~\(c_1\) and~\(c_2\) are eliminated.
%
%\item \label{eg:2expfitd}
%Fourth, discover a new eigenvalue problem by taking the transpose of the above equation. 
%Since all matrices except~\(U^{-1}\) are symmetric the transposition gives
%\begin{equation*}
%\begin{bmatrix} 1&\tfrac23\\\tfrac23&\frac7{18} \end{bmatrix}\tr{(U^{-1})}
%=\begin{bmatrix} 1&1\\1&\tfrac23 \end{bmatrix}\tr{(U^{-1})}D.
%\end{equation*}
%Denote the two columns of~\(\tr{(U^{-1})}\) by \(\vv_1\) and~\(\vv_2\), that is, \(\tr{(U^{-1})}=\begin{bmatrix} \vv_1&\vv_2 \end{bmatrix}\).
%Then, recalling \(D=\diag(\lambda_1,\lambda_2)\), the equation becomes
%\begin{eqnarray*}
%&&\begin{bmatrix} 1&\tfrac23\\\tfrac23&\frac7{18} \end{bmatrix}\begin{bmatrix} \vv_1&\vv_2 \end{bmatrix}
%=\begin{bmatrix} 1&1\\1&\tfrac23 \end{bmatrix}\begin{bmatrix} \vv_1&\vv_2 \end{bmatrix}
%\begin{bmatrix} \lambda_1&0\\0&\lambda_2 \end{bmatrix}
%\\\iff&&
%\begin{bmatrix} 1&\tfrac23\\\tfrac23&\frac7{18} \end{bmatrix}\begin{bmatrix} \vv_1&\vv_2 \end{bmatrix}
%=\begin{bmatrix} 1&1\\1&\tfrac23 \end{bmatrix}\begin{bmatrix} \lambda_1\vv_1&\lambda_2\vv_2 \end{bmatrix}
%%\\\iff&&
%%\begin{bmatrix} \begin{bmatrix} 1&\tfrac23\\\tfrac23&\frac7{18} \end{bmatrix}\vv_1&\begin{bmatrix} 1&\tfrac23\\\tfrac23&\frac7{18} \end{bmatrix}\vv_2 \end{bmatrix}
%%\\&&{}
%%=\begin{bmatrix} \begin{bmatrix} 1&1\\1&\tfrac23 \end{bmatrix}\lambda_1\vv_1&\begin{bmatrix} 1&1\\1&\tfrac23 \end{bmatrix}\lambda_2\vv_2 \end{bmatrix}
%\\&&(\text{then equate the columns from either side})
%\\\iff&&\begin{bmatrix} 1&\tfrac23\\\tfrac23&\frac7{18} \end{bmatrix}\vv_1=\lambda_1\begin{bmatrix} 1&1\\1&\tfrac23 \end{bmatrix}\vv_1
%\\\text{and}&&\begin{bmatrix} 1&\tfrac23\\\tfrac23&\frac7{18} \end{bmatrix}\vv_2=\lambda_2\begin{bmatrix} 1&1\\1&\tfrac23 \end{bmatrix}\vv_2.
%\end{eqnarray*}
%These last two equations have the same form.
%That is, if we can find two distinct solutions to
%\begin{equation*}
%\begin{bmatrix} 1&\tfrac23\\\tfrac23&\frac7{18} \end{bmatrix}\vv=\lambda\begin{bmatrix} 1&1\\1&\tfrac23 \end{bmatrix}\vv\,,
%\end{equation*}
%then one solution gives~\(\lambda_1,\vv_1\) and the other~\(\lambda_2,\vv_2\).
%This looks like an eigenvalue problem, \(A\vv=\lambda\vv\)\,, but with a matrix~\(B\) inserted in the right-hand side as an extra.
%The form \(A\vv=\lambda B\vv\) is called a `\idx{generalised eigen-problem}'.
%
%\item \label{eg:2expfite}
%Fifth, solve for the `eigenvalues'~\(\lambda\) by rearranging the equation to
%\begin{equation*}
%\begin{bmatrix} 1&\tfrac23\\\tfrac23&\frac7{18} \end{bmatrix}\vv-\lambda\begin{bmatrix} 1&1\\1&\tfrac23 \end{bmatrix}\vv=\ov
%\iff
%\begin{bmatrix} 1-\lambda&\tfrac23-\lambda\\\tfrac23-\lambda&\frac7{18}-\tfrac23\lambda\end{bmatrix}\vv=\ov\,.
%\end{equation*}
%Being a homogeneous linear equation, this only has nontrivial solutions~\vv\ when the determinant is zero:
%\begin{eqnarray*}
%\det\begin{bmatrix} 1-\lambda&\tfrac23-\lambda\\\tfrac23-\lambda&\frac7{18}-\tfrac23\lambda\end{bmatrix}
%&=&(1-\lambda)(\tfrac7{18}-\tfrac23\lambda)-(\tfrac23-\lambda)^2
%\\&=&\tfrac23\lambda^2-\tfrac{19}{18}\lambda+\tfrac7{18}
%-\lambda^2+\tfrac43\lambda-\tfrac49
%\\&=&-\tfrac13\lambda^2+\tfrac5{18}\lambda-\tfrac1{18}
%\\&=&-\tfrac1{18}(6\lambda^2-5\lambda+1)
%\\&=&-\tfrac1{18}(2\lambda-1)(3\lambda-1).
%\end{eqnarray*}
%This determinant is zero only for multipliers \(\lambda_1=\tfrac12\) and \(\lambda_2=\tfrac13\)\,. 
%Consequently, the corresponding rates \(r_1=\ln\tfrac12=-0.6932\) and \(r_2=\ln\tfrac13=-1.0986\)\,.
%We now know two of the four unknowns in the exponential~fit.
%
%\item \label{eg:2expfitf}
%Finally, to solve for the corresponding coefficients, just use any pair of data points.
%For example, using the first two data points we know
%\begin{equation*}
%c_1+c_2=1,\quad\text{and}\quad
%c_1\lambda_1+c_2\lambda_2=\tfrac12c_1+\tfrac13c_2=1\,.
%\end{equation*}
%Subtracting three times the second from the first gives \(-\tfrac12c_1=-2\)\,, that is \(c_1=4\).
%Then the first equation gives \(c_2=1-c_1=1-4=-3\)\,.
%\end{enumerate}
%\end{comment}%%%%%%%%%%%%%%%%%%%
Consequently, the ultimate exponential fit to the data is (as previously plotted)
\begin{equation*}
f(t)=4(\tfrac12)^t-3(\tfrac13)^t
=4e^{-0.6932\,t}-3e^{-1.0986\,t}
\end{equation*}
as \(\ln\frac12=-0.6932\) and \(\ln \frac13=-1.0986\)\,.
\end{solution}
\end{example}



\autoref{eg:2expfit} shows one way that fitting exponentials to data can be done with eigenvalues and eigenvectors.
But one undesirable attribute of the example is the need to invert the matrix~\(B\) to form matrix \(K=AB^{-1}\).
We avoid this inversion by generalising \idx{eigen-problem}s as introduced by the following reworking of parts of \autoref{eg:2expfit}.

\begin{example}[two short-cuts] \label{eg:2expfit2}
Recall that \autoref{sec:ddp} derived general solutions of dynamic equations such as \(\fv_{j+1}=K\fv_j\) by seeking solutions of the form \(\fv_j=\vv\lambda^j\).
For the previous \autoref{eg:2expfit} let's instead seek solutions of the form \(\fv_j=B\wv\lambda^j\).
Substituting this form, the dynamic equation \(\fv_{j+1}=K\fv_j\) becomes \(B\wv\lambda^{j+1}=KB\wv\lambda^j\); then factoring~\(\lambda^j\), recognising that \(KB=A\), and swapping sides, this equation becomes \(A\wv=\lambda B\wv\)\,.
This \(A\wv=\lambda B\wv\) forms a \idx{generalised eigen-problem} because it reduces to the standard eigen-problem in cases when the matrix \(B=I\).
Rework parts of \autoref{eg:2expfit} via this generalised eigen-problem.

\begin{solution} 
Restart the analysis in \autoref{eg:2expfitc}: instead of the eigen-problem \(K\vv=\lambda\vv\), let's solve the generalised eigen-problem \(A\wv=\lambda B\wv\)\,.
Here matrices
\begin{equation*}
A=\begin{bmatrix} 1&\frac23\\\frac23&\frac7{18} \end{bmatrix},\quad
B=\begin{bmatrix} 1&1\\1&\frac23 \end{bmatrix}.
\end{equation*}
So the equation \(A\wv=\lambda B\wv\) becomes
\begin{equation*}
\begin{bmatrix} 1&\frac23\\\frac23&\frac7{18} \end{bmatrix}\wv
=\lambda \begin{bmatrix} 1&1\\1&\frac23 \end{bmatrix}\wv
= \begin{bmatrix} \lambda&\lambda\\\lambda&\frac23\lambda \end{bmatrix}\wv\,.
\end{equation*}
Subtract the matrix on the right-hand side from both sides to obtain
\begin{equation*}
\begin{bmatrix} 1&\frac23\\\frac23&\frac7{18} \end{bmatrix}\wv
- \begin{bmatrix} \lambda&\lambda\\\lambda&\frac23\lambda \end{bmatrix}\wv 
=\begin{bmatrix} 1-\lambda&\frac23-\lambda\\\frac23-\lambda&\frac7{18}-\frac23\lambda \end{bmatrix}\wv
=\ov
\end{equation*}
Being a homogeneous linear equation, this last equation only has nontrivial solutions~\wv\ when the determinant is zero:
\begin{eqnarray*}
\det\begin{bmatrix} 1-\lambda&\tfrac23-\lambda\\\tfrac23-\lambda&\frac7{18}-\tfrac23\lambda\end{bmatrix}
&=&(1-\lambda)(\tfrac7{18}-\tfrac23\lambda)-(\tfrac23-\lambda)^2
\\&=&\tfrac23\lambda^2-\tfrac{19}{18}\lambda+\tfrac7{18}
-\lambda^2+\tfrac43\lambda-\tfrac49
\\&=&-\tfrac13\lambda^2+\tfrac5{18}\lambda-\tfrac1{18}
\\&=&-\tfrac1{18}(6\lambda^2-5\lambda+1)
\\&=&-\tfrac1{18}(2\lambda-1)(3\lambda-1).
\end{eqnarray*}
This determinant is zero only for multipliers \(\lambda_1=\tfrac12\) and \(\lambda_2=\tfrac13\)\,. 
These are precisely the same eigenvalues as in \autoref{eg:2expfitc}.

Another simplification is to avoid forming the full vector solution (\autoref{eg:2expfitd}--\ref{eg:2expfite}):
we only use just the first vector-component in the final answer for~\(f(t)\), so the full vector is redundant.
Instead, given the multipliers \(\lambda_1=\tfrac12\) and \(\lambda_2=\tfrac13\), just seek the penultimate form \(f_j=c_1(\tfrac12)^j+c_2(\tfrac13)^j\).
Given the data values, evaluating this form of solution at \(j=1,2\) gives two equations:
\begin{eqnarray*}
j=1,&&c_1\tfrac12+c_2\tfrac13=f_1=f(0)=1\,;
\\
j=2,&&c_1\tfrac14+c_2\tfrac19=f_2=f(1)=1\,.
\end{eqnarray*}
Subtracting twice the second from the first, and subtracting three times the second from the first gives the coefficients \(c_1=8\) and \(c_2=-9\).
Hence the fit to the data is \(f_j=8(\frac12)^j -9(\frac13)^j\) as in \autoref{eg:2expfite}.
Then, as before, index \(j=t+1\) so the function fit becomes \(f(t)=4(\frac12)^{t} -3(\frac13)^{t}\) to match the conclusion of the previous \autoref{eg:2expfit}.
\end{solution}
\end{example}




\subsubsection{Generalised eigen-problem}

As introduced by \autoref{eg:2expfit2}, let's generalise the \autoref{def:evecval} of eigenvalues and eigenvectors.
Such generalised eigen-problems also occur in the design analysis of complicated structures, such as buildings and bridges, where the second matrix~\(B\) represents the various masses of the various elements making up a structure.

\begin{definition} \label{def:geneig}
Let \(A\) and~\(B\) be \(n\times n\) square matrices\index{square matrix}. 
The \bfidx{generalised eigen-problem} is to find scalar \idx{eigenvalue}s~\(\lambda\) and corresponding nonzero \idx{eigenvector}s~\vv\ such that \(A\vv=\lambda B\vv\).
\end{definition}

\index{generalised eigen-problem|(}

\begin{example} \label{eg:}
Given \(A=\begin{bmatrix} -2&2\\3&1 \end{bmatrix}\) and \(B=\begin{bmatrix} 3&-3\\2&0 \end{bmatrix}\), what eigenvalue corresponds to the eigenvector \(\vv_1=(1,1)\) of the generalised eigen-problem \(A\vv=\lambda B\vv\)\,?  
Also answer for \(\vv_2=(-3,13)\).
\begin{solution} 
\begin{itemize}
\item To test \(\vv_1=(1,1)\), multiply by the matrices:
\begin{eqnarray*}
&&A\vv_1=\begin{bmatrix} -2&2\\3&1 \end{bmatrix}\begin{bmatrix} 1\\1 \end{bmatrix}
=\begin{bmatrix} 0\\4 \end{bmatrix};
\\&&B\vv_1=\begin{bmatrix} 3&-3\\2&0 \end{bmatrix}\begin{bmatrix} 1\\1 \end{bmatrix}
=\begin{bmatrix} 0\\2 \end{bmatrix}.
\end{eqnarray*}
Since \(A\vv_1\) is twice \(B\vv_1\) then \(\vv_1\)~is an eigenvector corresponding to eigenvalue \(\lambda_1=2\)\,.
\item To test \(\vv_1=(-3,13)\), multiply by the matrices:
\begin{eqnarray*}
&&A\vv_1=\begin{bmatrix} -2&2\\3&1 \end{bmatrix}\begin{bmatrix} -3\\13 \end{bmatrix}
=\begin{bmatrix} 32\\4 \end{bmatrix};
\\&&B\vv_1=\begin{bmatrix} 3&-3\\2&0 \end{bmatrix}\begin{bmatrix} -3\\13 \end{bmatrix}
=\begin{bmatrix} -48\\-6 \end{bmatrix}.
\end{eqnarray*}
Since \(32=-\tfrac23(-48)\) and \(4=-\tfrac23(-6)\), then \(A\vv_2\) is~\(-2/3\) times~\(B\vv_2\), and so \(\vv_2\)~is an eigenvector corresponding to eigenvalue \(\lambda_2=-2/3\)\,.
\end{itemize}
\end{solution}
\end{example}


\begin{activity}
Which of the following vectors is an eigenvector of the generalised eigen-problem
\begin{equation*}
\begin{bmatrix} -2&0\\1&0 \end{bmatrix}\vv=\lambda\begin{bmatrix} 3&-1\\-1&1 \end{bmatrix}\vv\,?
\end{equation*}
\actposs[4]{$(1,-1)$}{$(0,0)$}{$(0,1)$}{$(2,-1)$}
\end{activity}


The standard eigen-problem is the case when matrix \(B=I\)\,.
Many of the properties for standard eigenvalues and eigenvectors also hold for generalised eigen-problems, although there are some differences.
Most importantly here, albeit without proof, provided matrix~\(B\) is invertible, then counted according to \idx{multiplicity} there are \(n\)~eigenvalues of a \idx{generalised eigen-problem} in \(n\times n\) matrices.

\begin{example} \label{eg:}
Find all \idx{eigenvalue}s and corresponding \idx{eigenvector}s of the \idx{generalised eigen-problem} \(A\vv=\lambda B\vv\) for matrices
\begin{equation*}
A=\begin{bmatrix} -1&1\\-4&0 \end{bmatrix},\quad
B=\begin{bmatrix} -1&1\\-1&2 \end{bmatrix}
\end{equation*}
\begin{solution} 
Now solving \(A\vv=\lambda B\vv\) is the same as solving \((A-\lambda B)\vv=\ov\), which requires \(\det(A-\lambda B)=0\) (\autoref{thm:detinv}).
Hence determine the eigenvalues by finding the zeros of the characteristic polynomial
\begin{eqnarray*}
\det(A-\lambda B)
&=&\det\begin{bmatrix} -1+\lambda&1-\lambda\\-4+\lambda&0-2\lambda \end{bmatrix}
\\&=&(-1+\lambda)(-2\lambda)-(-4+\lambda)(1-\lambda)
\\&=&-2\lambda^2+2\lambda+\lambda^2-5\lambda+4
\\&=&-\lambda^2-3\lambda+4
\\&=&-(\lambda+4)(\lambda-1).
\end{eqnarray*}
Hence the generalised eigenvalues are \(\lambda_1=-4\) and \(\lambda_2=1\)\,.
\begin{itemize}
\item For the case \(\lambda_1=-4\) the corresponding eigenvectors satisfy
\begin{equation*}
(A+4B)\vv=\begin{bmatrix} -5&5\\-8&8 \end{bmatrix}\vv=\ov.
\end{equation*}
That is, the two components of~\vv\ must be the same: corresponding eigenvectors are \(\vv_1\propto(1,1)\).

\item For the case \(\lambda_2=1\) the corresponding eigenvectors satisfy
\begin{equation*}
(A-B)\vv=\begin{bmatrix} 0&0\\-3&-2 \end{bmatrix}\vv=\ov.
\end{equation*}
That is, \(-3v_1-2v_2=0\) which rearranged is \(v_1=-\tfrac23v_2\): corresponding eigenvectors are \(\vv_2\propto(-\tfrac23,1)\).
\end{itemize}
\end{solution}
\end{example}




\begin{example} \label{eg:}
Find all \idx{eigenvalue}s and corresponding \idx{eigenvector}s of the \idx{generalised eigen-problem} \(A\vv=\lambda B\vv\) for matrices
\begin{equation*}
A=\begin{bmatrix} 2&2\\-1&0 \end{bmatrix},\quad
B=\begin{bmatrix} -1&1\\1&-1 \end{bmatrix}
\end{equation*}
\begin{solution} 
Now solving \(A\vv=\lambda B\vv\) is the same as solving \((A-\lambda B)\vv=\ov\), which requires \(\det(A-\lambda B)=0\) (\autoref{thm:detinv}).
Hence determine the eigenvalues by finding the zeros of the characteristic polynomial
\begin{eqnarray*}
\det(A-\lambda B)
&=&\det\begin{bmatrix} 2+\lambda&2-\lambda\\-1-\lambda&0+\lambda \end{bmatrix}
\\&=&(2+\lambda)(\lambda)-(2-\lambda)(-1-\lambda)
\\&=&\lambda^2+2\lambda-2-\lambda^2+\lambda+2
\\&=&3\lambda+2\,.
\end{eqnarray*}
Hence the only eigenvalue is \(\lambda=-2/3\)\,.
Here the matrix~\(B\) is not invertible---its determinant is zero---and consequently in this example we do not get the typical full complement of two eigenvalues.

For the only eigenvalue \(\lambda=-2/3\) the corresponding eigenvectors satisfy
\begin{equation*}
(A+\tfrac23B)\vv=\begin{bmatrix} \tfrac43&\tfrac83\\-\tfrac13&-\tfrac23 \end{bmatrix}\vv=\ov.
\end{equation*}
That is, the two components of~\vv\ must satisfy \(\tfrac13v_1+\tfrac23v_2=0\) which rearranged is \(v_1=-2v_2\): corresponding eigenvectors are \(\vv\propto(-2,1)\).
\end{solution}
\end{example}


\index{generalised eigen-problem|)}





\subsubsection{General fitting of exponentials}

Suppose that some experiment or other observation has given us \(2n\)~data values \hlist f{2n}\ at equi-spaced `times' \hlist t{2n}, where the spacing \(t_{j+1}-t_j=h\)\,.
The general aim is to fit the multi-exponential function \cite[\S2.6, e.g.]{Cuyt2015}
\begin{equation}
f(t)=c_1e^{r_1t}+c_2e^{r_2t}+\cdots+c_ne^{r_nt}.
%=\sum_{j=1}^n c_je^{r_jt}
\label{eq:eiddse}
\end{equation}
for some coefficients~\hlist cn\ and some rates~\hlist rn\ to be determined (possibly \idx{complex valued} for \idx{oscillation}s).
In general, finding the coefficients and rates is a delicate nonlinear task outside the remit of this book.
However, as the previous two examples illustrate, in these circumstances we instead invoke our powerful linear algebra methods.

Because the data is sampled at equi-spaced times, \(h\)~apart, then instead of seeking~\(r_k\) we seek multipliers \(\lambda_k=e^{r_kh}\).

\begin{procedure}[\bfidx{exponential interpolation}] \label{pro:ei}
Given measured data \hlist f{2n}\ at \(2n\)~equi-spaced times \hlist t{2n} where time \(t_j=(j-1)h\) for time-spacing~\(h\) (and starting from time \(t_1=0\) without loss of applicability), we seek to fit the data with a sum of exponentials~\eqref{eq:eiddse}.
\begin{enumerate}
\item\label{pro:eia} From the \(2n\)~data points, form two \(n\times n\) (symmetric) \index{Hankel matrix}Hankel matrices 
\begin{eqnarray*}&&
A=\begin{bmatrix} f_2&f_3&\cdots&f_{n+1}
\\f_3&f_4&\cdots&f_{n+2}
\\\vdots&\vdots&&\vdots
\\f_{n+1}&f_{n+2}&\cdots&f_{2n} \end{bmatrix},
\\&&
B=\begin{bmatrix} f_1&f_2&\cdots&f_{n}
\\f_2&f_3&\cdots&f_{n+1}
\\\vdots&\vdots&&\vdots
\\f_{n}&f_{n+1}&\cdots&f_{2n-1} \end{bmatrix}.
\end{eqnarray*}
In \script\ \index{hankel()@\texttt{hankel()}}\verb|A=hankel(f(2:n+1),f(n+1:2*n))| and \verb|B=hankel(f(1:n),f(n:2*n-1))| forms these two matrices (this Hankel function is also invoked in exploring El~Nino, \autoref{eg:orthbapp}).

\item\label{pro:eib} Find the \idx{eigenvalue}s of the \idx{generalised eigen-problem} \(A\vv=\lambda B\vv\)\,: 
\begin{itemize}
\item by hand on small problems solve \(\det(A-\lambda B)=0\)\,;
\item in \script\ invoke \index{eig()@\texttt{eig()}}\verb|lambda=eig(A,B)|\,, and then \index{log()@\texttt{log()}}\verb|r=log(lambda)/h|\,.
\end{itemize}
This eigen-problem typically determines \(n\)~multipliers \hlist\lambda n, and thence the \(n\)~rates \(r_k=(\ln\lambda_k)/h\)\,.

\item\label{pro:eic} Determine the corresponding \(n\)~coefficients \hlist cn\ from any \(n\)~point subset of the \(2n\)~data points.
For example, the first \(n\)~data points give the linear system 
\begin{equation*}
\begin{bmatrix} 1&1&\cdots&1
\\\lambda_1&\lambda_2&\cdots&\lambda_n
\\\lambda_1^2&\lambda_2^2&\cdots&\lambda_n^2
\\\vdots&\vdots&&\vdots
\\\lambda_1^{n-1}&\lambda_2^{n-1}&\cdots&\lambda_n^{n-1}
 \end{bmatrix}\begin{bmatrix} c_1\\c_2\\\vdots\\c_n \end{bmatrix}
 =\begin{bmatrix} f_1\\f_2\\f_3\\\vdots\\f_n \end{bmatrix}
\end{equation*}
In \script\ one may construct the matrix~\(U\) appearing here with \index{bsxfun()@\texttt{bsxfun()}}\verb|U=bsxfun(@power,lambda,0:n-1).'| when \verb|lambda| is a column vector of the eigenvalues.
Since the eigenvalues~\(\lambda\) may be \idx{complex valued} we need the transpose~``\verb|.'|'' not the complex conjugate transpose~``\verb|'|'' (\autoref{tbl:mtlbops}).
\end{enumerate}
\end{procedure}




\begin{table}
\caption{As well as the \script\ commands and operations listed in Tables~\ref{tbl:mtlbpre}, \ref{tbl:mtlbbasics}, \ref{tbl:mtlbops}, \ref{tbl:mtlbmops}, \ref{tbl:mtlbsvd}, \ref{tbl:mtlbimag}, and~\ref{tbl:mtlbnorm} this section invokes these functions.\index{Matlab@\textsc{Matlab}|textbf}\index{Octave|textbf}} \label{tbl:mtlbexpf}
\smallskip\hrule
\begin{minipage}{\linewidth}
%\parbox{\linewidth}{
\begin{itemize}
\item \index{hankel()@\texttt{hankel()}}\verb|hankel(x,y)| for two vectors of the same length, \(\verb|x|=(\hlist xn)\) and \(\verb|y|=(\hlist yn)\), forms the \(n\times n\) matrix
\def\adots{\reflectbox{$\ddots$}}
\begin{equation*}
\begin{bmatrix} 
x_1&x_2&x_3&\cdots&x_n \\
x_2&x_3&\cdots&x_n&y_2\\
x_3&&\adots&\adots&\vdots\\
\vdots&\adots&\adots&&y_{n-1}\\
x_n&y_2&\cdots&y_{n-1}&y_n
\end{bmatrix}.
\end{equation*}

\item \index{eig()@\texttt{eig()}}\verb|eig(A,B)| for \(n\times n\) matrices~\(A\) and~\(B\) computes a vector in~\(\RR^n\) of \idx{eigenvalue}s~\(\lambda\) such that \(\det(A-\lambda B)=0\)\,.
Some of the computed eigenvalues in the vector may be \(\pm\verb|Inf|\)\index{Inf@\texttt{Inf}} (depending upon the nature of~\(B\)) which denotes that a corresponding eigenvalue does \text{not exist}.

The command \verb|[V,D]=eig(A,B)| solves the \idx{generalised eigen-problem} \(A\xv=\lambda B\xv\) for eigenvalues~\(\lambda\) returned in the diagonal of matrix~\(D\) (some may be~\(\pm\verb|Inf|\)), and corresponding \idx{eigenvector}s returned in the corresponding column of~\(V\).

%\item \index{meshgrid()@\texttt{meshgrid()}}\verb|[X,Y]=meshgrid(x,y)|  for vectors \(\verb|x|=(\hlist xn)\) and \(\verb|y|=(\hlist ym)\) forms two \(m\times n\) matrices:
%\begin{equation*}
%\verb|X|=\begin{bmatrix} x_1&x_2&\cdots&x_n\\
% x_1&x_2&\cdots&x_n\\
% \vdots&\vdots&&\vdots\\
% x_1&x_2&\cdots&x_n \end{bmatrix},
% \quad
%\verb|Y|=\begin{bmatrix} y_1&y_1&\cdots&y_1\\
% y_2&y_2&\cdots&y_2\\
% \vdots&\vdots&&\vdots\\
% y_m&y_m&\cdots&y_m \end{bmatrix}.
%\end{equation*}
\end{itemize}
\end{minipage}
%}
\hrule
\end{table}





\begin{comment}
Link to singular spectrum analysis of \S\ref{sec:ssa} when it is written.
\end{comment}

\begin{proof} 
This proof establishes the generic case.
Denote the columns of matrices~\(A\) and~\(B\) by vectors of consecutive data values \(\fv_j=(f_j, f_{j+1},\ldots,f_{j+n-1})\in\RR^n\) for \(j=1,2,\ldots,n+1\) (\autoref{pro:eia} in \autoref{pro:ei}).
\begin{itemize}
\item First seek a linear transformation, by some \(n\times n\) matrix~\(K\), from one data vector to the next in the form \(\fv_{j+1}=K\fv_j\) for \(j=1,2,\ldots,n\)\,.
Arrange these \(n\)~equations \(\fv_{j+1}=K\fv_j\) into the one matrix equation
\begin{eqnarray*}
\begin{bmatrix} \fv_2&\fv_3&\cdots&\fv_4 \end{bmatrix}
&=&\begin{bmatrix} K\fv_1&K\fv_2&\cdots&K\fv_3 \end{bmatrix}
\\&=&K\begin{bmatrix} \fv_1&\fv_2&\cdots&\fv_3 \end{bmatrix},
\end{eqnarray*}
after factoring the matrix~\(K\).
That is, matrix~\(K\) satisfies \(A=KB\) (which in principle we could typically solve to obtain \(K=AB^{-1}\)).

\item Second, view \(\fv_{j+1}=K\fv_j\) as a dynamic system in~\(j\) and modify the approach of \autoref{sec:ddp}.
Seek solutions in the form \(\fv_j=B\vv \lambda^j\) for some non-zero vector~\vv\ and multiplier~\(\lambda\).
Then \(\fv_{j+1}=K\fv_j\) becomes \(B\vv \lambda^{j+1}=KB\vv \lambda^j\).
This equation becomes the generalised eigen-problem \(A\vv=\lambda B\vv\) after factoring out~\(\lambda^j\) (provided non-zero), and recalling that \(A=KB\)\,.

\item Typically, \autoref{pro:eib} in \autoref{pro:ei}, the eigen-problem \(A\vv=\lambda B\vv\) has \(n\)~distinct eigenvalues \hlist\lambda n\ (possibly complex valued) with some corresponding eigenvectors \hlist\vv n.
By the linearity of \(\fv_{j+1}=K\fv_j\)\,, a general solution is then the linear combination \(\fv_j=C_1B\vv_1\lambda_1^{j} +C_2B\vv_2\lambda_2^{j} +\cdots +C_nB\vv_n\lambda_n^{j}\).
We only need one component of this general solution, say the first, and so the data values \(f_j=\lincomb c{\lambda^{j-1}}n\) for some constants \hlist cn\ to be found (\(c_i=C_i\lambda_i(B\vv_i)_1\) in terms of~\(C_i\)).

\item Determine the constants \hlist cn\ by matching the general solution to \(n\)~data values.
Using the first~\(n\) data values gives the \(n\)~equations:
\begin{eqnarray*}
f_1&=&c_1+c_2+\cdots+c_n\,,
\\f_2&=&\lincomb c{\lambda^{1}}n,
\\f_3&=&\lincomb c{\lambda^{2}}n,
\\&\vdots&
\\f_n&=&\lincomb c{\lambda^{n-1}}n.
\end{eqnarray*}
Form as a matrix-vector system for the unknown vector \((\hlist cn)\) to give the system solved by \autoref{pro:eic}.
\end{itemize}
Lastly, rearrange \(t_j=(j-1)h\) to \(j-1=t_j/h\), then \(\lambda_k^{j-1}=e^{(\log \lambda_k)t_j/h}=e^{r_kt_j}\) for rates \(r_k=(\log\lambda_k)/h\) (\autoref{pro:eib}).
So the solution \(f_j=\lincomb c{\lambda^{j-1}}n\) becomes the exponential interpolation \(f(t)=c_1e^{r_1t}+c_2e^{r_2t}+\cdots+c_ne^{r_nt}\) as required.
\end{proof}





\begin{example} \label{eg:eipiano}
\newcommand{\mytemp}[1]{\begin{tikzpicture}
\def\aone{#1}
\def\ugh{\ifnum2>\aone xmax=19,xtick={5,10,15},domain=0:19
    \else domain=0:159,samples=100,xtick={50,100,150}\fi}
\begin{axis}[footnotesize,font=\footnotesize,width=4.7cm
, xlabel={$t$\,(ms)},ylabel={$f$\,(mm)},axis lines=middle
, ymin=-0.9,ymax=1.3, \ugh ]
\addplot+[mark=*,only marks] coordinates { (0,1) (5,-0.3766) (10,-0.5352) (15,0.7114)};
\ifnum0<\aone\addplot+[smooth,thick,no marks] {cos(deg(2*x/5))*exp(-x/50)};\fi
\end{axis}
\end{tikzpicture}}%
A damped \idx{piano string} is struck and the sideways displacement of the string is measured at four times, \(5\)\,ms~apart. 
\marginpar{\mytemp0}%
The measurements (in~mm) are \(f_1=1.0000\), \(f_2=-0.3766\), \(f_3=-0.5352\) and \(f_4=0.7114\) (as illustrated).  
Determine, by hand calculation, the \idx{frequency} and damping of the string.%
\footnote{Some of you will know that identification of frequencies is most commonly done by what is called a \idx{Fourier transform}.  
However, with a limited amount of good data, or for decaying oscillations, this approach may be better.}

Recall \idx{Euler's formula} that \(e^{\i\theta}=\cos\theta+\i\sin\theta\) so \idx{oscillation}s are here captured by \idx{complex valued} exponentials.

\begin{solution} 
Follow \autoref{pro:ei} by hand. 
\begin{enumerate}
\item 
Form the two Hankel matrices from the data values~\hlist f4:
\begin{eqnarray*}&&
A=\begin{bmatrix} f_2&f_3\\f_3&f_4 \end{bmatrix}
=\begin{bmatrix} -0.3766 & -0.5352
\\  -0.5352&   0.7114 \end{bmatrix},
\\&&
B=\begin{bmatrix} f_1&f_2\\f_2&f_3 \end{bmatrix}
=\begin{bmatrix} 1.0000&  -0.3766
\\  -0.3766&  -0.5352 \end{bmatrix}.
\end{eqnarray*}
\item 
With some arithmetic, the determinant 
\begin{eqnarray*}
&&\det(A-\lambda B)
\\&=&\det\begin{bmatrix} -0.3766-\lambda & -0.5352+0.3766\lambda
\\  -0.5352+0.3766\lambda&   0.7114+0.5352\lambda \end{bmatrix} 
\\&=&(-0.3766-\lambda)(0.7114+0.5352\lambda)
\\&&{}
-(-0.5352+0.3766\lambda)^2
\\&\vdots& 
\\&=& -0.6770\lambda^2  -0.5098\lambda  -0.5544\,.
\end{eqnarray*}
More arithmetic with the well known formula for solving quadratic equations finds that this determinant is zero for complex conjugate eigenvalues \(\lambda=-0.3765 \pm 0.8228\i\)\,.
Such \idx{complex eigenvalue}s are characteristic of oscillations.
The complex logarithm of these complex values, divided by the time step of \(0.005\,\text{s}=5\,\)ms, gives the two complex rates \(r=-20\pm400\i\) (to two decimal places).
\item 
To find the corresponding coefficients, solve the complex linear equations
\begin{eqnarray*}&&
c_1+c_2=1\,, 
\\&&
(-0.3765 + 0.8228\i)c_1+(-0.3765 - 0.8228\i)c_2
\\&&{}
=-0.3766\,.
\end{eqnarray*}
By inspection the solution is \(c_1=c_2=\tfrac12\) (to three decimal places).
\end{enumerate}
\marginpar{\mytemp1}%
We conclude the exponential fit is, as plotted in the margin,
\begin{eqnarray*}
f(t)&=&\tfrac12e^{-20t+400\i t}+\tfrac12e^{-20t-400\i t}
\\&=&\left(\tfrac12e^{400\i t}+\tfrac12e^{-400\i t}\right)e^{-20t}
\\&=&\cos(400t)e^{-20t}.
\end{eqnarray*}
Interpreting, the cosine factor indicates the piano string oscillates at \(400\)~radians per second which is \(400/(2\pi)=63.66\)~cycles/sec.
\marginpar{\mytemp2}%
However, the piano string is damped by the factor~\(e^{-20t}\) so that in just a fraction of a second the oscillations, and the sound, stop.
\end{solution}
\end{example}




\begin{example} \label{eg:eipianu}
For the data of the previous \autoref{eg:eipiano}, determine the \idx{frequency} and damping of the \idx{piano string} using \script.

\begin{solution} 
Follow \autoref{pro:ei} using \script. 
\begin{enumerate}
\item Form the Hankel matrices from the data with commands
\begin{verbatim}
f=[1.0000
  -0.3766
  -0.5352
   0.7114]
A=hankel(f(2:3),f(3:4))
B=hankel(f(1:2),f(2:3))
\end{verbatim}
\setbox\ajrqrbox\hbox{\qrcode{% exp fit
f=[1.0000
  -0.3766
  -0.5352
   0.7114]
A=hankel(f(2:3),f(3:4))
B=hankel(f(1:2),f(2:3))
lambda=eig(A,B)
r=log(lambda)/0.005
U=[1 1;lambda(1) lambda(2)]
rcond(U)
c=U\slosh f(1:2)
}}%
\marginpar{\usebox{\ajrqrbox}}%

\item Compute the multipliers as the eigenvalues of the generalised problem, and then determine the rates with
\begin{verbatim}
lambda=eig(A,B)
r=log(lambda)/0.005
\end{verbatim}
giving the results
\begin{verbatim}
lambda =
  -0.3765 + 0.8228i
  -0.3765 - 0.8228i
r =
   -19.99 + 399.99i
   -19.99 - 399.99i
\end{verbatim}


\item Compute the coefficients in the exponential fit with the following, as the value \verb|rcond=0.43| is good (\autoref{pro:unisol}),
\begin{verbatim}
U=[1 1; lambda(1) lambda(2)]
rcond(U)
c=U\f(1:2)
\end{verbatim}
giving results
\begin{verbatim}
U =
   1.0000 + 0.0000i   1.0000 + 0.0000i
  -0.3765 + 0.8228i  -0.3765 - 0.8228i
ans =  0.4319
c =
   0.5000 + 0.0000i
   0.5000 - 0.0000i
\end{verbatim}
\end{enumerate}
This gives the same answer as by hand, namely 
\begin{equation*}
f(t)=\tfrac12e^{-20t+400\i t}+\tfrac12e^{-20t-400\i t}.
\end{equation*}
As before, Euler's formula transforms this to a real form.
\end{solution}
\end{example}




\begin{example} \label{eg:}
\newcommand{\mytemp}[1]{\begin{tikzpicture}
\def\aone{#1}
\begin{axis}[footnotesize,font=\footnotesize,width=4.7cm
, xlabel={$t$\,(secs)},ylabel={$f$},axis lines=middle
, ymin=-0.9,ymax=1.3, xmax=24,domain=0:24]
\addplot+[mark=*,only marks] coordinates { (0,0.0000)
(3,0.1000)
(6,0.2833)
(9,0.4639)
(12,0.6134)
(15,0.7277)
(18,0.8112)
(21,0.8705)};
\ifnum0<\aone\addplot+[smooth,thick,no marks] {1.01-2.28*exp(-0.13*x) +1.26*exp(-0.24*x)};\fi
\end{axis}
\end{tikzpicture}}%
In a \idx{biochemical experiment} every three seconds we measure the concentration of an output chemical as tabulated below (and illustrated in the margin).
Fit a sum of four exponentials to this data.
\marginpar{\mytemp0}%
\begin{equation*}
\begin{array}{rl}\hline
\text{secs}&\text{concentration}
\\\hline0&0.0000
\\3&0.1000
\\6&0.2833
\\9&0.4639
\\12&0.6134
\\15&0.7277
\\18&0.8112
\\21&0.8705
\\\hline
\end{array}
\end{equation*}

\begin{solution} 
Follow \autoref{pro:ei} using \script. 
\begin{enumerate}
\item Form the Hankel matrices from the data with commands
\begin{verbatim}
f=[0.0000
   0.1000
   0.2833
   0.4639
   0.6134
   0.7277
   0.8112
   0.8705]
A=hankel(f(2:5),f(5:8))
B=hankel(f(1:4),f(4:7))
\end{verbatim}
\setbox\ajrqrbox\hbox{\qrcode{% exp fit
f=[0.0000
   0.1000
   0.2833
   0.4639
   0.6134
   0.7277
   0.8112
   0.8705]
A=hankel(f(2:5),f(5:8))
B=hankel(f(1:4),f(4:7))
lambda=eig(A,B)
r=log(lambda)/3
U=bsxfun(@power,lambda,0:3).'
rcond(U)
c=U\slosh f(1:4)
}}%
\marginpar{\usebox{\ajrqrbox}}%

\item Compute the multipliers as the eigenvalues of the generalised eigen-problem, and then determine the rates with
\begin{verbatim}
lambda=eig(A,B)
r=log(lambda)/3
\end{verbatim}
giving the results
\begin{verbatim}
lambda =
   0.9990
  -0.4735
   0.6736
   0.4922
r =
  -0.0003 + 0.0000i
  -0.2492 + 1.0472i
  -0.1317 + 0.0000i
  -0.2363 + 0.0000i
\end{verbatim}

\item Compute the coefficients in the exponential fit with the following (\autoref{tbl:mtlbnorm} introduces \index{bsxfun()@\texttt{bsxfun()}}\verb|bsxfun()|)
\begin{verbatim}
U=bsxfun(@power,lambda,0:3).'
rcond(U)
c=U\f(1:4)
\end{verbatim}
giving results
\begin{verbatim}
U =
   1.0000   1.0000   1.0000   1.0000
   0.9990  -0.4735   0.6736   0.4922
   0.9981   0.2242   0.4538   0.2422
   0.9971  -0.1061   0.3057   0.1192
ans =  0.007656
c =
   1.0117
  -0.0001
  -2.2756
   1.2641
\end{verbatim}
However, the value \verb|rcond=0.008| is poor (\autoref{pro:unisol}).
This \verb|rcond| suggests that the results have two less significant digits than the original data (\autoref{thm:erramp}).
Since the original data is specified to four decimal places, the results are probably only accurate to two decimal~places.
\end{enumerate}
Consequently, this analysis fits the data with the exponential sum (as illustrated in the margin)
\marginpar{\mytemp1}%
\begin{eqnarray*}
f(t)&\approx&1.01\cdot1^{t/3}+0\cdot(-0.47)^{t/3}
-2.28\cdot0.67^{t/3}+1.26\cdot0.49^{t/3}
\\&\approx&1.01-2.28e^{-0.13t}+1.26e^{-0.24t}.
\end{eqnarray*}
\end{solution}
\end{example}


As with any data fitting, in practical applications be careful about the reliability of the results.
Sound statistical analysis needs to supplement \autoref{pro:ei} to inform us about expected errors and sensitivity.
This problem of fitting exponentials to data is often sensitive to errors.

The techniques and theory of this subsection generalise to cater for noisy data and for complex system interactions with vast amounts of data.
The generalisation, \emph{Dynamic Mode Decomposition} \cite[e.g.]{Kutz2016}, is applied across many areas such as fluid mechanics, video processing, epidemiology, neuroscience, and financial trading.



\index{exponential interpolation|)}








\sectionExercises





\begin{exercise} \label{ex:} 
For each of the following list of numbers, could the numbers be all the \idx{eigenvalue}s of a \(4\times4\) matrix? 
Justify your answer.
\begin{parts}
\item \(-1.2,-0.6,0.2,-1.4\)
\answer{yes}

\item \(\pm2,-3\)
\answer{yes}

\item \(0,3,\pm5,8\)
\answer{no}

\item \(0,3\pm5\i,8\)
\answer{yes}

\item \(-1.4\pm\sqrt7\i,-4,3\pm2\i\)
\answer{no}

\end{parts}
\end{exercise}





\begin{exercise} \label{ex:} 
Find the \idx{trace} of each of the following matrices.
%\begin{verbatim}
%a=0+round(randn(n)*20)/10,trace(a)
%\end{verbatim}
\begin{parts}
\item \(\begin{bmatrix} 1&3
\\-1&-2 \end{bmatrix}\)
\answer{\(-1\)}

\item \(\begin{bmatrix} -1&-2
\\6&1 \end{bmatrix}\)
\answer{\(0\)}

\item \(\begin{bmatrix} 3.2&-0.9&-4.3
\\0.8&-0.1&2.3
\\-0.9&0.8&-0.2 \end{bmatrix}\)
\answer{\(2.9\)}

\item \(\begin{bmatrix} 3&-1.4&2
\\0&-0.5&0.4
\\1.3&-0.2&-0.6 \end{bmatrix}\)
\answer{\(1.9\)}

\item \(\begin{bmatrix} 1.1&-1.9&1.8&-2.4
\\-4.3&1.1&-2.1&1.2
\\1.6&1.4&-0.6&0.9
\\-0.4&-3.7&-0.4&2.5 \end{bmatrix}\)
\answer{\(4.1\)}

\item \(\begin{bmatrix} -1.5&0.6&0.5&1.8
\\1.5&-1.9&-0.1&2.8
\\-2.1&-1.4&-3.3&-0.3
\\-1.4&-2.8&0.8&-2.5 \end{bmatrix}\)
\answer{\(-9.2\)}
  
\end{parts}
\end{exercise}







\begin{exercise} \label{ex:} 
For each of the following matrices, determine the two highest order terms and the constant term in the \idx{characteristic polynomial} of the matrix.
%\begin{verbatim}
%a=0+round(full(sprandn(n,n,2/n))*3),(-1)^n*poly(a)
%\end{verbatim}
\begin{parts}
\item \(\begin{bmatrix} -7 & 1
\\-2 & 2 \end{bmatrix}\)
\answer{\(\lambda^2+5\lambda-12\)}

\item \(\begin{bmatrix} 3 & -3
\\6 & 2 \end{bmatrix}\)
\answer{\(\lambda^2-5\lambda+24\)}

\item \(\begin{bmatrix} 0 & 0 & 2
\\1 & 0 & 2
\\4 & 2 & 0 \end{bmatrix}\)
\answer{\(-\lambda^3+0\lambda^2+\cdots+4\)}

\item \(\begin{bmatrix} 3 & -3 & 6
\\-1 & 4 & 0
\\0 & 0 & -4 \end{bmatrix}\)
\answer{\(-\lambda^3+3\lambda^2+\cdots-36\)}

\item \(\begin{bmatrix} -1 & 0 & -6
\\-7 & 0 & -4
\\0 & -6 & 0 \end{bmatrix}\)
\answer{\(-\lambda^3-\lambda^2+\cdots-228\)}

\item \(\begin{bmatrix} -1 & 0 & -2 & 0
\\0 & 0 & 0 & 0
\\-3 & -1 & 0 & -2
\\-4 & 0 & -5 & 0 \end{bmatrix}\)
\answer{\(\lambda^4+\lambda^3+\cdots+0\)}

\item \(\begin{bmatrix} -3 & 0 & 0 & 1
\\0 & 0 & 2 & 0
\\0 & 0 & 0 & 1
\\-4 & -4 & 0 & 4 \end{bmatrix}\)
\answer{\(\lambda^4-\lambda^3+\cdots+24\)}

\item \(\begin{bmatrix} 0 & 4 & 0 & 1
\\4 & -5 & 0 & -3
\\0 & -4 & 0 & 0
\\0 & 0 & -1 & 0 \end{bmatrix}\)
\answer{\(\lambda^4+5\lambda^3+\cdots-16\)}

\end{parts}
\end{exercise}








\begin{exercise} \label{ex:} 
For each of the following \idx{characteristic polynomial}s, write down the size of the corresponding matrix, and the matrix's \idx{trace} and \idx{determinant}.
%\begin{verbatim}
%n=2
%(-1)^n*poly(round(randn(n)*3))
%\end{verbatim}
\begin{parts}
\item \(\lambda^2+5\lambda-6\)
\answer{\(2\times2\), \(-5\), \(-6\)}

\item \(\lambda^2+2\lambda-10\)
\answer{\(2\times2\), \(-2\), \(-10\)}

\item \(\lambda^2-\lambda\)
\answer{\(2\times2\), \(1\), \(0\)}

\item \(-\lambda^3+5\lambda^2+7\lambda-20\)
\answer{\(3\times3\), \(5\), \(-20\)}

\item \(-\lambda^3+28\lambda-199\)
\answer{\(3\times3\), \(0\), \(-199\)}

\item \(-\lambda^3+8\lambda^2-5\lambda\)
\answer{\(3\times3\), \(8\), \(0\)}

\item \(\lambda^4+3\lambda^3+143\lambda-56\)
\answer{\(4\times4\), \(-3\), \(-56\)}

\item \(\lambda^4+5\lambda^2-41\lambda-5\)
\answer{\(4\times4\), \(0\), \(-5\)}

\end{parts}
\end{exercise}




\begin{exercise}
Some matrices have the following \idx{characteristic polynomial}s.
In each case, write down the eigenvalues and their \idx{multiplicity}.
%\begin{verbatim}
%lam=0+round(randn(1,5)*4),mul=ceil(rand(1,5)*4)
%\end{verbatim}

\begin{enumerate}
\item \((\lambda+1)\lambda^4(\lambda-2)^3\)
\answer{\(-1\), one; \(0\), four; \(2\), three}
\item \((\lambda+1)(\lambda+2)^2(\lambda-4)^3(\lambda-1)^2\)
\answer{\(-1\), one; \(-2\), two; \(1\), two; \(4\), three}
\item \((\lambda+6)(\lambda+2)^2\lambda(\lambda-2)^3\)
\answer{\(-6\), one; \(-2\), two; \(0\), one; \(2\), three}
\item \((\lambda+2.6)^4(\lambda-0.8)^3(\lambda+1.1)\)
\answer{\(-2.6\), four; \(0.8\), three; \(1.1\), one}
\item \((\lambda+0.8)^4(\lambda-2.6)^3(\lambda+0.2)(\lambda-0.7)^2\)
\answer{\(-0.8\), four; \(2.6\), three; \(-0.2\), one; \(0.7\), two}

\end{enumerate}
\end{exercise}






\begin{exercise} \label{ex:} 
For each the following matrices, determine the \idx{characteristic polynomial} by hand, and hence find all \idx{eigenvalue}s of the matrix and their \idx{multiplicity}.
Show your working.
%\begin{verbatim}
%format rat
%lams=round(randn(1,n)*3)
%for j=1:9999,p=round(randn(n)*2);if abs(det(p))==1,break,end,end
%[j,i]=meshgrid(lams,lams);[i,j]=find(i==j);k=min(find(i<j));i=i(k),j=j(k)
%d=diag(lams),if length(i)>0,d(i,j)=round(ceil(3*rand)),if rand>0.7,d(j,i)=-d(i,j),end;end
%a=p*d*inv(p),(-1)^n*poly(a),eig(a)
%\end{verbatim}

\begin{parts}
\item \(\begin{bmatrix} 0 & -3
\\-1 & -2 \end{bmatrix}\)
\answer{\(1\) once, \(-3\) once}

\item \(\begin{bmatrix} 0 & 5
\\-2 & 2 \end{bmatrix}\)
\answer{\(1\pm 3\i\) once}

\item \(\begin{bmatrix} 0 & -3
\\3 & 6 \end{bmatrix}\)
\answer{\(3\) twice}

\item \(\begin{bmatrix} 3 & -4
\\2 & -3 \end{bmatrix}\)
\answer{\(\pm1\) once}

\item \(\begin{bmatrix} 4.5 & 16
\\-1 & -3.5 \end{bmatrix}\)
\answer{\(1/2\) twice}

\item \(\begin{bmatrix} -1 & 1 & -1
\\-6 & -6 & 2
\\-5 & -3 & -3 \end{bmatrix}\)
\answer{\(-4\) once, \(-3\pm \i\) once}

\item \(\begin{bmatrix} -2 & -5 & -1
\\0 & 3 & 1
\\0 & -6 & -2 \end{bmatrix}\)
\answer{\(0\) once, \(-2\) once, \(1\) once}

\item \(\begin{bmatrix} 9 & 3 & 0
\\-12 & -3 & 0
\\2 & -4 & -2 \end{bmatrix}\)
\answer{\(-2\) once, \(3\) twice}

\item \(\begin{bmatrix} -14 & 24 & 52
\\-4 & 8 & 18
\\-2 & 3 & 6 \end{bmatrix}\)
\answer{\(2\) once, \(-1\pm \i\) once}

\item \(\begin{bmatrix} -1 & -2 & 2
\\7 & 18 & -12
\\7 & 17 & -11 \end{bmatrix}\)
\answer{\(-1\) once, \(1\) once, \(6\) once}

\item \(\begin{bmatrix} -10 & -10 & -16
\\4 & 4 & 6
\\3 & 3 & 5 \end{bmatrix}\)
\answer{\(-1\) once, \(0\) twice}

\item \(\begin{bmatrix} 1 & -15 & 7
\\-1 & -1 & -1
\\-5 & -15 & -1 \end{bmatrix}\)
\answer{\(-1\) once, \(\pm2\i\) once}

\end{parts}
\end{exercise}






\begin{exercise} \label{ex:} 
For each the following matrices, use \script\ to find all \idx{eigenvalue}s of the matrix and their \idx{multiplicity}.
%\begin{verbatim}
%lams=round(randn(1,n)*20)/10; lams=lams(ceil(n*rand(1,n)))
%for j=1:99999,p=round(randn(n)*3);if abs(det(p))==1,break,end,end
%[j,i]=meshgrid(lams,lams);[i,j]=find(i==j);k=min(find(i<j));i=i(k),j=j(k)
%d=diag(lams),if length(i)>0,d(i,j)=round(ceil(rand*20))/10,if rand>0.7,d(j,i)=-d(i,j),end;end
%a=p*d*inv(p),eig(a)
%\end{verbatim}

\begin{enumerate}
\item \(\begin{bmatrix} -2.7 & 0 & 1.6
\\6.3 & -1 & -27.8
\\-0.1 & 0 & -1.9 \end{bmatrix}\)
\setbox\ajrqrbox\hbox{\qrcode{% evals
A=[-2.7 0 1.6
6.3 -1 -27.8
-0.1 0 -1.9]
}}%
\marginpar{\usebox{\ajrqrbox}}%
\answer{\(-1\) once, \(-2.3\) twice}

\item \(\begin{bmatrix} 4.9 & -8.1 & 5.4
\\8 & -11.2 & 8
\\3.9 & -3.9 & 3.4 \end{bmatrix}\)
\setbox\ajrqrbox\hbox{\qrcode{% evals
A=[4.9 & -8.1 & 5.4
8 & -11.2 & 8
3.9 & -3.9 & 3.4]
}}%
\marginpar{\usebox{\ajrqrbox}}%
\answer{\(-3.2\) once, \(-0.5\) once, \(0.8\) once}

\item \(\begin{bmatrix} -6.7 & -0.6 & -6.6 & 3.6
\\3 & 0.1 & 3 & -2
\\2.8 & 0.6 & 2.7 & -1.6
\\-6 & 0 & -6 & 3.1 \end{bmatrix}\)
\setbox\ajrqrbox\hbox{\qrcode{% evals
A=[-6.7 -0.6 -6.6 3.6
3 0.1 3 -2
2.8 0.6 2.7 -1.6
-6 0 -6 3.1]
}}%
\marginpar{\usebox{\ajrqrbox}}%
\answer{\(-0.9\) once, \(-0.1\) once, \(0.1\) twice}

\item \(\begin{bmatrix} 11 & 17.9 & -33.4 & 46.4
\\1.2 & 0.9 & 2.8 & 2.2
\\-12.8 & -21 & 37.2 & -54.8
\\-12.3 & -19.7 & 33.7 & -51.3 \end{bmatrix}\)
\setbox\ajrqrbox\hbox{\qrcode{% evals
A=[11 17.9 -33.4 46.4
1.2 0.9 2.8 2.2
-12.8 -21 37.2 -54.8
-12.3 -19.7 33.7 -51.3]
}}%
\marginpar{\usebox{\ajrqrbox}}%
\answer{\(-0.7\pm1.8\i\) once, \(-1.2\) once, \(0.4\) once}

\item \(\begin{bmatrix} 0.6 & 0 & 0 & 0
\\9.6 & -1 & 33.6 & 17.6
\\-9.6 & 1.6 & -33 & -17.6
\\19.2 & -3.2 & 67.2 & 35.8 \end{bmatrix}\)
\setbox\ajrqrbox\hbox{\qrcode{% evals
A=[0.6 0 0 0
9.6 -1 33.6 17.6
-9.6 1.6 -33 -17.6
19.2 -3.2 67.2 35.8]
}}%
\marginpar{\usebox{\ajrqrbox}}%
\answer{\(0.6\) four times}

\item \(\begin{bmatrix} 52.8 & -93.3 & 73.4 & 57.1 & 104
\\18.3 & -30.7 & 28.5 & 20.6 & 36.6
\\-20 & 34.9 & -29.4 & -22.3 & -40
\\18.2 & -31.8 & 27.4 & 20.8 & 36.4
\\-6.8 & 13.6 & -6.8 & -6.8 & -12.8 \end{bmatrix}\)
\setbox\ajrqrbox\hbox{\qrcode{% evals
A=[52.8 -93.3 73.4 57.1 104
18.3 -30.7 28.5 20.6 36.6
-20 34.9 -29.4 -22.3 -40
18.2 -31.8 27.4 20.8 36.4
-6.8 13.6 -6.8 -6.8 -12.8]
}}%
\marginpar{\usebox{\ajrqrbox}}%
\answer{\(-2\) once, \(0.3\) once, \(0.8\) thrice}

\item \(\begin{bmatrix} -356.4 & 264.6 & -880.8 & -689.9 & 455.5
\\-58.5 & 41.7 & -142.8 & -111.4 & 77
\\157 & -115.6 & 392.9 & 311.4 & -201
\\-78.5 & 57.8 & -198.4 & -159.6 & 100.5
\\-58.5 & 45.6 & -142.8 & -111.4 & 73.1 \end{bmatrix}\)
\setbox\ajrqrbox\hbox{\qrcode{% evals
A=[-356.4 264.6 -880.8 -689.9 455.5
-58.5 41.7 -142.8 -111.4 77
157 -115.6 392.9 311.4 -201
-78.5 57.8 -198.4 -159.6 100.5
-58.5 45.6 -142.8 -111.4 73.1]
}}%
\marginpar{\usebox{\ajrqrbox}}%
\answer{\(1.7\pm1.3\i\) once, \(-3.9\) thrice}

\item \(\begin{bmatrix} 309.4 & -29.7 & 451.3 & 337.3 & 305.9
\\-217.6 & 20.3 & -313.5 & -236.1 & -215.9
\\3 & 0 & 0.7 & 3 & 3
\\-232.6 & 24.1 & -336 & -254.9 & -230.9
\\-83.6 & 5.6 & -119.8 & -89.2 & -81.8 \end{bmatrix}\)
\setbox\ajrqrbox\hbox{\qrcode{% evals
A=[309.4 -29.7 451.3 337.3 305.9
-217.6 20.3 -313.5 -236.1 -215.9
3 0 0.7 3 3
-232.6 24.1 -336 -254.9 -230.9
-83.6 5.6 -119.8 -89.2 -81.8]
}}%
\marginpar{\usebox{\ajrqrbox}}%
\answer{\(1.8\) twice, \(-2.3\) once, \(-3.8\) twice}

\end{enumerate}
\end{exercise}






\begin{exercise} \label{ex:} 
For each of the following matrices, find by hand the \idx{eigenspace} of the nominated eigenvalue.  
Confirm your answer with \script.  
Show your working.  
%\begin{verbatim}
%format rat
%lams=round(randn(1,n)*3); lams=lams(ceil(n*rand(1,n)))
%for j=1:9999,p=round(randn(n)*3);if abs(det(p))==1,break,end,end
%[j,i]=meshgrid(lams,lams);[i,j]=find(i==j);k=min(find(i<j));i=i(k),j=j(k)
%d=diag(lams),if length(i)>0,if rand>0.99,d(i,j)=round(ceil(3*rand)),end;end
%a=p*d*inv(p),[v,d]=eig(a), v=v*diag(1./max(abs(v)))
%\end{verbatim}
\begin{enumerate}
\item \(\begin{bmatrix} -12 & 10
\\-15 & 13 \end{bmatrix}\), \(\lambda=3\)
\answer{\(\EE_{3}=\Span\{(2,3)\}\)}

\item \(\begin{bmatrix} -1 & 9
\\-1 & 5 \end{bmatrix}\), \(\lambda=2\)
\answer{\(\EE_{2}=\Span\{(3,1)\}\)}

\item \(\begin{bmatrix} -1 & 0
\\-2 & 1 \end{bmatrix}\), \(\lambda=-1\)
\answer{\(\EE_{-1}=\Span\{(1,1)\}\)}

\item \(\begin{bmatrix} 11 & -4 & -12
\\-27 & 10 & 27
\\19 & -7 & -20 \end{bmatrix}\), \(\lambda=1\)
\answer{\(\EE_{1}=\Span\{(0,-3,1)\}\)}

\item \(\begin{bmatrix} -1 & -7 & -2
\\8 & 14 & 2
\\0 & 0 & 7 \end{bmatrix}\), \(\lambda=7\)
\answer{\(\EE_{7}=\Span\{(7,-8,0),(-2,2,1)\}\)}

\item \(\begin{bmatrix} -12 & -82 & -17
\\3 & 18 & 3
\\-6 & -26 & -1 \end{bmatrix}\), \(\lambda=0\)
\answer{\(\EE_{0}=\Span\{(-4,1,-2)\}\)}

\item \(\begin{bmatrix} -4 & 0 & -4
\\-2 & -4 & 0
\\8 & 4 & 6 \end{bmatrix}\), \(\lambda=-2\)
\answer{\(\EE_{-2}=\Span\{(-2,2,1)\}\)}

\item \(\begin{bmatrix} -3 & 2 & -6
\\-4 & 3 & -6
\\2 & -1 & 4 \end{bmatrix}\), \(\lambda=1\)
\answer{\(\EE_{1}=\Span\{(2,1,-1),(-1,1,1)\}\)}

\end{enumerate}
\end{exercise}







\begin{exercise} \label{ex:} 
For each of the following matrices, Use \script\ to find their eigenvalues, with \idx{multiplicity}, and to find eigenvectors corresponding to each eigenvalue \twodp.  
(The next \autoref{sec:lisb} discusses that for \idx{repeated eigenvalue}s we generally want to record the so-called `linearly independent' eigenvectors.)
%\begin{verbatim}
%lams=round(randn(1,n)*3); lams=lams(ceil(n*rand(1,n)))
%for j=1:9999,p=round(randn(n)*3);if ismember(abs(det(p)),[1]),break,end,end
%[j,i]=meshgrid(lams,lams);[i,j]=find(i==j);k=min(find(i<j));i=i(k),j=j(k)
%d=diag(lams),if length(i)>0,d(i,j)=round(ceil(3*rand)),if rand>0.7,d(j,i)=-d(i,j),end;end
%a=0+round(p*d*inv(p)),[v,d]=eig(a),vt=num2str(conj(v'),2)
%\end{verbatim}
\begin{enumerate}
\item \(\begin{bmatrix} 14 & 3 & -6
\\14 & -2 & -2
\\31 & 4 & -11 \end{bmatrix}\)
\setbox\ajrqrbox\hbox{\qrcode{% eig
A=[14 3 -6
14 -2 -2
31 4 -11]
}}\marginpar{\usebox{\ajrqrbox}}%
\answer{\(\lambda=-3\), once, \((0,-0.89,-0.45)\); 
\(\lambda=2\), twice, \((0.27,0.53,0.80)\).}

\item \(\begin{bmatrix} -1 & 0 & 0
\\5 & 5 & 6
\\-5 & -6 & -7 \end{bmatrix}\)
\setbox\ajrqrbox\hbox{\qrcode{% eig
A=[-1 0 0
5 5 6
-5 -6 -7]
}}\marginpar{\usebox{\ajrqrbox}}%
\answer{\(\lambda=-1\), thrice, \((0,0.71,-0.71)\), \((0.53,-0.78,0.34)\).}

\item \(\begin{bmatrix} -144 & -374 & 316 & 18
\\21 & 45 & -42 & 0
\\-49 & -138 & 112 & 9
\\134 & 336 & -286 & -13 \end{bmatrix}\)
\setbox\ajrqrbox\hbox{\qrcode{% eig
A=[-144 -374 316 18
21 45 -42 0
-49 -138 112 9
134 336 -286 -13]
}}\marginpar{\usebox{\ajrqrbox}}%
\answer{\(\lambda=5\), once, \((-0.82,0,-0.41,0.41)\);
\(\lambda=3\), once, \((0.55,0.28,0.55,0.55)\);
\(\lambda=-4\), twice, \((0.67,0,0.33,-0.67)\).}

\item \(\begin{bmatrix} -50 & 30 & 46 & -62
\\0 & 2 & 0 & 0
\\-104 & 62 & 104 & -142
\\-39 & 24 & 42 & -58 \end{bmatrix}\)
\setbox\ajrqrbox\hbox{\qrcode{% eig
A=[-50 30 46 -62
0 2 0 0
-104 62 104 -142
-39 24 42 -58]
}}\marginpar{\usebox{\ajrqrbox}}%
\answer{\(\lambda=2\), twice, \((0.19,0,0.86,0.48)\), \((0.37,0.92,-0.15,0.018)\);
\(\lambda=-3\pm \i\), once, \((0.14\mp0.28\i,0,-0.69,-0.62\pm0.21\i)\).}

\item \(\begin{bmatrix} 4 & -11 & -12 & 9 & -19
\\0 & 0 & 0 & 3 & -2
\\4 & 16 & 19 & -12 & 23
\\6 & 16 & 20 & -7 & 31
\\-1 & -3 & -3 & 3 & 1 \end{bmatrix}\)
\setbox\ajrqrbox\hbox{\qrcode{% eig
A=[4 -11 -12 9 -19
0 0 0 3 -2
4 16 19 -12 23
6 16 20 -7 31
-1 -3 -3 3 1]
}}\marginpar{\usebox{\ajrqrbox}}%
\answer{\(\lambda=-3.94\), once, \((-0.33,-0.41,0.72,0.44,-0.15)\);
\(\lambda=1\), once, \((0,0.8,-0.53,0.27,0)\);
\(\lambda=7.49\pm0.39\i\), once, \((-0.24\pm0.55\i,0.23\pm0.02\i,0.043\pm0.01\i,0.72,0.24\mp0.11\i)\);
\(\lambda=4.95\), once, \((-0.14,0.45,0.62,0.46,-0.43)\).}

\item \(\begin{bmatrix} 75 & 7 & -13 & -51 & -129
\\120 & 12 & -24 & -84 & -208
\\62 & 6 & -12 & -42 & -106
\\-48 & -5 & 9 & 32 & 83
\\62 & 6 & -11 & -42 & -107 \end{bmatrix}\)
\setbox\ajrqrbox\hbox{\qrcode{% eig
A=[75 7 -13 -51 -129
120 12 -24 -84 -208
62 6 -12 -42 -106
-48 -5 9 32 83
62 6 -11 -42 -107]
}}\marginpar{\usebox{\ajrqrbox}}%
\answer{\(\lambda=2\), twice, \((-0.64,-0.43,-0.43,0.21,-0.43)\);
\(\lambda=-2\), once, \((0.35,0.71,0.35,-0.35,0.35)\);
\(\lambda=-1\), twice, \((-0.33,-0.89,-0.22 0,-0.22)\), \((-0.18,-0.81,-0.46,0.25,-0.20)\).}


\end{enumerate}
\end{exercise}






\begin{exercise} \label{ex:} 
Consider the following three matrices which are perturbed versions of the matrix
\(\begin{bmatrix} 2&1\\0&2 \end{bmatrix}\).
Which perturbations show the high \idx{sensitivity} of the \idx{repeated eigenvalue}?  
Give reasons.
\begin{parts}
\item \(\eAii=\begin{bmatrix} 2&1\\-0.0001&2 \end{bmatrix}\)
\answer{Sensitive.}

\item \(\eAii=\begin{bmatrix} 2&1.0001\\0&2 \end{bmatrix}\)
\answer{Insensitive.}

\item \(\eAii=\begin{bmatrix} 2.0001&1\\0&2 \end{bmatrix}\)
\answer{Not sensitive.}

\end{parts}
\end{exercise}





\begin{exercise} \label{ex:} 
For each of the following matrices, use \script\ to compute the eigenvalues, and to compute the eigenvalues of matrices obtained by adding \idx{random perturbation}s of size~\(0.0001\) (use \verb|randn|).
Give reasons for which eigenvalues appear \index{sensitivity}sensitive and which appear to be not sensitive.
\begin{enumerate}
\item \(\eAii=\begin{bmatrix} 0&-3\\-1&-2 \end{bmatrix}\)
\answer{Both not sensitive.}

\item \(\eAii=\begin{bmatrix} 0&-3\\3&6 \end{bmatrix}\)
\answer{\(\lambda=3\) sensitive.}

\item \(\eAii=\begin{bmatrix} -10 & -10 & -16
\\4 & 4 & 6
\\3 & 3 & 5 \end{bmatrix}\)
\answer{\(-1\) not sensitive, \(0\) sensitive.}

\item \(\eAii=\begin{bmatrix} 2&1&1
\\1&2&1
\\1&1&2 \end{bmatrix}\)
\answer{neither sensitive, symmetric.}

\item \(\eAii=\begin{bmatrix} -1 & 1 & -1
\\-6 & -6 & 2
\\-5 & -3 & -3 \end{bmatrix}\)
\answer{all not sensitive.}

\item \(\eAii=\begin{bmatrix} -6.7 & -0.6 & -6.6 & 3.6
\\3 & 0.1 & 3 & -2
\\2.8 & 0.6 & 2.7 & -1.6
\\-6 & 0 & -6 & 3.1 \end{bmatrix}\)
\setbox\ajrqrbox\hbox{\qrcode{% evals
[-6.7 -0.6 -6.6 3.6
3 0.1 3 -2
2.8 0.6 2.7 -1.6
-6 0 -6 3.1]
}}\marginpar{\usebox{\ajrqrbox}}%
\answer{\(-0.9\) and \(-0.1\) not sensitive, \(0.1\) sensitive}

\item \(\eAii=\begin{bmatrix} 1.4 & -7.1 & -0.7 & 6.2
\\ -7.1 & -1.0 & -2.2 & -2.5
\\ -0.7 & -2.2 & -3.4 & -4.1
\\ 6.2 & -2.5 & -4.1 & -1.0 \end{bmatrix}\)
\setbox\ajrqrbox\hbox{\qrcode{% evals
[1.4 -7.1 -0.7 6.2
-7.1 -1.0 -2.2 -2.5
-0.7 -2.2 -3.4 -4.1
6.2 -2.5 -4.1 -1.0]
}}\marginpar{\usebox{\ajrqrbox}}%
\answer{one sensitive, symmetric.}

\item \(\eAii=\begin{bmatrix} 0.6 & 0 & 0 & 0
\\9.6 & -1 & 33.6 & 17.6
\\-9.6 & 1.6 & -33 & -17.6
\\19.2 & -3.2 & 67.2 & 35.8 \end{bmatrix}\)
\setbox\ajrqrbox\hbox{\qrcode{% evals
[0.6 0 0 0
9.6 -1 33.6 17.6
-9.6 1.6 -33 -17.6
19.2 -3.2 67.2 35.8]
}}\marginpar{\usebox{\ajrqrbox}}%
\answer{\(0.6\) mixed sensitivity!}

\end{enumerate}
\end{exercise}






\begin{exercise} \label{ex:yay} 
Consider the evolving system \(\yv(t+1)=A\yv(t)\) for each of the following cases. 
Predict \(\yv(1)\), \(\yv(2)\) and~\(\yv(3)\), for the given initial~\(\yv(0)\).
%\begin{verbatim}
%lams=round(randn(1,n)*2)
%for j=1:9999,p=round(randn(n)*3);if ismember(abs(det(p)),[1]),break,end,end
%[j,i]=meshgrid(lams,lams);[i,j]=find(i==j);k=min(find(i<j));i=i(k),j=j(k)
%d=diag(lams),if length(i)>0,d(i,j)=round(ceil(3*rand)),if rand>0,d(j,i)=-d(i,j),end;end
%A=0+round(p*d*inv(p)*100)/100,[v,d]=eig(A)
%y0=round(randn(n,1)*2)+1, y1=y0'*A',y2=y1*A',y3=y2*A'
%\end{verbatim}
\begin{enumerate}
\item \(A=\begin{bmatrix} 3 & 0
\\3 & 2 \end{bmatrix}\), 
\(\yv(0)=\begin{bmatrix} 6
\\-1 \end{bmatrix}\)
\answer{\((18,16)\), \((54,86)\) and \((162,334)\)}

\item \(A=\begin{bmatrix} 0 & -1
\\-4 & 2 \end{bmatrix}\), 
\(\yv(0)=\begin{bmatrix} 3
\\0 \end{bmatrix}\)
\answer{\((0,-12)\), \((12,-24)\) and \((24,-96)\)}

\item \(A=\begin{bmatrix} 26 & 21
\\-28 & -23 \end{bmatrix}\), 
\(\yv(0)=\begin{bmatrix} 3
\\-4 \end{bmatrix}\)
\answer{\((-6,8)\), \((12,-16)\) and \((-24,32)\)}

\item \(A=\begin{bmatrix} -2 & 5
\\-2 & 4 \end{bmatrix}\), 
\(\yv(0)=\begin{bmatrix} 1
\\1 \end{bmatrix}\)
\answer{\((3,2)\), \((4,2)\) and \((2,0)\)}

\item \(A=\begin{bmatrix} 11 & -14 & -4
\\7 & -10 & -2
\\4 & -4 & -2 \end{bmatrix}\), 
\(\yv(0)=\begin{bmatrix} 0
\\2
\\-2 \end{bmatrix}\)
\answer{\((-20,-16,-4)\), \((20,28,-8)\) and \((-140,-124,-16)\)}

\item \(A=\begin{bmatrix} 9 & 7 & -5
\\-16 & -8 & 8
\\10 & 10 & -6 \end{bmatrix}\), 
\(\yv(0)=\begin{bmatrix} -1
\\1
\\1 \end{bmatrix}\)
\answer{\((-7,16,-6)\), \((79,-64,126)\) and \((-367,256,-606)\)}

\item \(A=\begin{bmatrix} 2 & -2 & 0
\\4 & -2 & 0
\\0 & -5 & -1 \end{bmatrix}\), 
\(\yv(0)=\begin{bmatrix} 2
\\0
\\1 \end{bmatrix}\)
\answer{\((4,8,-1)\), \((-8,0,-39)\) and \((-16,-32,39)\)}

\item \(A=\begin{bmatrix} -4 & 14 & 5
\\2 & -1 & 0
\\-13 & 26 & 8 \end{bmatrix}\), 
\(\yv(0)=\begin{bmatrix} 3
\\0
\\0 \end{bmatrix}\)
\answer{\((-12,6,-39)\), \((-63,-30,0)\) and \((-168,-96,39)\)}


\end{enumerate}
\end{exercise}






\begin{exercise} \label{ex:} 
For each of the matrices of the previous \autoref{ex:yay}, find a \idx{general solution} of \(\yv(t+1)=A\yv(t)\), if possible.
Then use the corresponding given initial~\(\yv(0)\) to find a formula for the specific~\(\yv(t)\).
Finally, check that the formula reproduces the values of \(\yv(1)\), \(\yv(2)\) and~\(\yv(3)\) found in \autoref{ex:yay}.
Show your working.
\end{exercise}





\begin{exercise} \label{ex:} 
Reconsider the mathematical modelling of the serval, \autoref{eg:serval}.
Derive the matrix~\(A\) in the model \(\yv(t+1)=A\yv(t)\) for the following cases:
\begin{enumerate}
\item choose the unit of time to be three months (\(1/4\)~year);
\answer{\(\begin{bmatrix} \frac12&0&\frac14
\\\frac12&\frac34&0
\\0&\frac14&\frac{33}{34} \end{bmatrix}\)}

\item choose the unit of time to be one month  (\(1/12\)~year).
\answer{\(\begin{bmatrix} \frac56&0&\frac1{12}
\\\frac16&\frac{11}{12}&0
\\0&\frac1{12}&\frac{101}{102} \end{bmatrix}\)}

\end{enumerate}
\end{exercise}







\begin{exercise} \label{ex:} 
From the following partial description of the \idx{Tasmanian Devil}, 
\marginpar{\includegraphics[width=\marginparwidth]{EigGen/TasmanianDevil}
\footnotesize\url{https://en.wikipedia.org/wiki/Tasmanian_devil}}
derive a mathematical model in the form \(\yv(t+1)=A\yv(t)\) for the age structure of the Tasmanian Devil.
By finding eigenvalues and an eigenvector, predict the long-term growth of the population, and predict the long-term relative numbers of Devils of various ages.
\begin{quoted}{Wikipedia, 2016}
Devils are not monogamous, and their reproductive process is very robust and competitive. Males fight one another for the females, and then guard their partners to prevent female infidelity. Females can ovulate three times in as many weeks during the mating season, and 80\%~of two-year-old females are seen to be pregnant during the annual mating season. Females average four breeding seasons in their life and give birth to 20--30 live young after three weeks' gestation. The newborn are pink, lack fur, have indistinct facial features and weigh around 0.20\,g (0.0071\,oz) at birth. As there are only four nipples in the pouch, competition is fierce and few newborns survive. The young grow rapidly and are ejected from the pouch after around 100~days, weighing roughly 200\,g (7.1\,oz). The young become independent after around nine months, so the female spends most of her year in activities related to birth and rearing.
\end{quoted}
\end{exercise}






\begin{exercise} \label{ex:} 
From the following partial description of the \idx{elephant}, 
\marginpar{\includegraphics[width=\marginparwidth]{EigGen/Loxodontacyclotis}
\footnotesize\url{https://en.wikipedia.org/wiki/Elephant}}
derive a mathematical model in the form \(\yv(t+1)=A\yv(t)\) for the age structure of the \idx{elephant}.
By finding eigenvalues and an eigenvector, predict the long-term growth of the population, and predict the long-term relative numbers of elephants of various ages.
\begin{quoted}{Wikipedia, 2016}
Gestation in elephants typically lasts around two years with interbirth intervals usually lasting four to five years. Births tend to take place during the wet season. Calves are born 85\,cm (33\,in) tall and weigh around 120\,kg (260\,lb). Typically, only a single young is born, but twins sometimes occur.  The relatively long pregnancy is maintained by five corpus luteums (as opposed to one in most mammals) and gives the foetus more time to develop, particularly the brain and trunk. As such, newborn elephants are precocial and quickly stand and walk to follow their mother and family herd.  A new calf is usually the centre of attention for herd members.  Adults and most of the other young will gather around the newborn, touching and caressing it with their trunks.  For the first few days, the mother is intolerant of other herd members near her young.  Alloparenting---where a calf is cared for by someone other than its mother---takes place in some family groups.  Allomothers are typically two to twelve years old. When a predator is near, the family group gathers together with the calves in the centre.

For the first few days, the newborn is unsteady on its feet, and needs the support of its mother. It relies on touch, smell and hearing, as its eyesight is poor. It has little precise control over its trunk, which wiggles around and may cause it to trip. By its second week of life, the calf can walk more firmly and has more control over its trunk. After its first month, a calf can pick up, hold and put objects in its mouth, but cannot suck water through the trunk and must drink directly through the mouth. It is still dependent on its mother and keeps close to her.

For its first three months, a calf relies entirely on milk from its mother for nutrition after which it begins to forage for vegetation and can use its trunk to collect water.  At the same time, improvements in lip and leg coordination occur.  Calves continue to suckle at the same rate as before until their sixth month, after which they become more independent when feeding.  By nine months, mouth, trunk and foot coordination is perfected.  After a year, a calf's abilities to groom, drink, and feed itself are fully developed.  It still needs its mother for nutrition and protection from predators for at least another year.  Suckling bouts tend to last 2--4\,min/hr for a calf younger than a year and it continues to suckle until it reaches three years of age or older.  Suckling after two years may serve to maintain growth rate, body condition and reproductive ability.  Play behaviour in calves differs between the sexes; females run or chase each other, while males play-fight. The former are sexually mature by the age of nine years while the latter become mature around 14--15 years.  Adulthood starts at about 18~years of age in both sexes.  Elephants have long lifespans, reaching 60--70~years of age.
\end{quoted}
\end{exercise}






\begin{exercise} \label{ex:} 
From the following partial description of the \idx{giant mouse lemur}, 
\marginpar{\includegraphics[width=\marginparwidth]{EigGen/Mirza_zaza}
\footnotesize\url{https://en.wikipedia.org/wiki/Giant_mouse_lemur}}
derive a mathematical model in the form \(\yv(t+1)=A\yv(t)\) for the age structure of the \idx{giant mouse lemur}.
By finding eigenvalues and an eigenvector, predict the long-term growth of the population, and predict the long-term relative numbers of \idx{giant mouse lemur}s of various ages.
\begin{quoted}{Wikipedia, 2016}
Reproduction starts in November for Coquerel's giant mouse lemur at Kirindy Forest; the estrous cycle runs approximately 22~days, while estrus lasts only a day or less. \ldots

One to three offspring (typically two) are born after 90~days of gestation, weighing approximately 12\,g (0.42\,oz). Because they are poorly developed, they initially remain in their mother's nest for up to three weeks, being transported by mouth between nests. Once they have grown sufficiently, typically after three weeks, the mother will park her offspring in vegetation while she forages nearby. After a month, the young begin to participate in social play and grooming with their mother, and between the first and second month, young males begin to exhibit early sexual behaviors (including mounting, neck biting, and pelvic thrusting). By the third month, the young forage independently, though they maintain vocal contact with their mother and use a small part of her range.

Females start reproducing after ten months, while males develop functional testicles by their second mating season. Testicle size in the northern giant mouse lemur does not appear to fluctuate by season, and is so large relative to the animal's body mass that it is the highest among all primates. This emphasis on sperm production in males, as well as the use of copulatory plugs, suggests a mating system best described as polygynandrous where males use scramble competition (roaming widely to find many females). In contrast, male Coquerel's giant mouse lemurs appear to fight for access to females (contest competition) during their breeding season. Males disperse from their natal range, and the age at which they leave varies from two years to several. Females reproduce every year, although postpartum estrus has been observed in captivity. In the wild, the lifespan of giant mouse lemurs is thought to rarely exceed five or six years
\end{quoted}
\end{exercise}




\begin{exercise} \label{ex:} 
From the following partial description of the \idx{dolphin} (Indo-Pacific bottlenose dolphin), 
\marginpar{\includegraphics[width=\marginparwidth]{EigGen/Tursiops_aduncus}
\footnotesize\url{https://en.wikipedia.org/wiki/Indo-Pacific_bottlenose_dolphin}}
derive a mathematical model in the form \(\yv(t+1)=A\yv(t)\) for the age structure of the \idx{dolphin}.
(Assume only one calf is born at a time.)
By finding eigenvalues and an eigenvector, predict the long-term growth of the population, and predict the long-term relative numbers of \idx{dolphin}s of various ages.
\begin{quoted}{Wikipedia, 2016}
Indo-Pacific bottlenose dolphins live in groups that can number in the hundreds, but groups of five to 15 dolphins are most common.  In some parts of their range, they associate with the common bottlenose dolphin and other dolphin species, such as the humpback dolphin.

The peak mating and calving seasons are in the spring and summer, although mating and calving occur throughout the year in some regions. Gestation period is about 12~months. Calves are between 0.84 and 1.5~metres (2.8 and 4.9\,ft) long, and weigh between 9 and 21~kilograms (20 and 46\,lb). The calves are weaned between 1.5 and two years, but can remain with their mothers for up to five years. The interbirth interval for females is typically four to six years.

In some parts of its range, this dolphin is subject to predation by sharks; its life span is more than 40 years.
\end{quoted}
\end{exercise}






\begin{exercise} \label{ex:} 
You are given that a mathematical model of the age structure of some animal population is
\begin{eqnarray*}
y_1(t+1)&=& 0.5y_1(t)+y_3(t),\\
y_2(t+1)&=& 0.5y_1(t)+0.7y_2(t),\\
y_3(t+1)&=& 0.3y_2(t)+0.9y_3(t).
\end{eqnarray*}
Invent an animal species, and time scale, and create details of a plausible scenario for the breeding and life cycle of the species that could lead to this mathematical model.  
Write a coherent paragraph about the breeding and life cycle of the species with enough information that someone could deduce this mathematical model from your description.  
Be creative.
\end{exercise}











\begin{exercise} \label{ex:} 
For each of the following matrices, say~\(A\) for instance, find by hand calculation the \idx{eigenvalue}s and \idx{eigenvector}s of the larger matrix~\(\begin{bmatrix} O&A\\\tr A&O \end{bmatrix}\).
Show your working.
Relate these to an \svd\ of the matrix~\(A\).
\begin{parts}
\item \(\eAii=\begin{bmatrix} 3\\4 \end{bmatrix}\)
\answer{\(\lambda=0,\pm5\), \((\frac45,-\frac35,0)\) and \((\frac35,\frac45,\pm1)\)}

\item \(\eAii=\begin{bmatrix} -5&12 \end{bmatrix}\)
\answer{\(\lambda=0,\pm13\), \((0,\frac{12}{13},\frac5{13})\) and \((\pm1,-\frac5{13},\frac{12}{13})\)}

\item \(\eAii=\begin{bmatrix} 1&0\\0&-2 \end{bmatrix}\)
\answer{\(\lambda=\pm1,\pm2\), \((\pm1,0,1,0)\) and \((0,1,0,\mp1)\)}

\item \(\eAii=\begin{bmatrix} 0&1\\-4&0 \end{bmatrix}\)
\answer{\(\lambda=\pm1,\pm4\), \((1,0,0,\pm1)\) and \((0,1,\mp1,0)\)}

\end{parts}
\end{exercise}




\begin{exercise} \label{ex:} 
Find by hand calculation all eigenvalues and corresponding eigenvectors of the \idx{generalised eigen-problem} \(A\vv=\lambda B\vv\) for the following pairs of matrices.
Check your calculations with \script.
%\begin{verbatim}
%n=2,A=0+round(randn(n)*2),B=0+round(randn(n)*2),[v,d]=eig(A,B), v*diag(1./v(1,:))
%\end{verbatim}
\begin{enumerate}
\item \(A=\begin{bmatrix} 1 & 3
\\ 0 & 0 \end{bmatrix}\),
\(B=\begin{bmatrix} 0 & 3
\\ 4 & 4 \end{bmatrix}\)
\answer{eigenvalues \(0,\tfrac23\); eigenvectors proportional to \((3,-1),(1,-1)\), respectively}

\item \(A=\begin{bmatrix} -1 & 0
\\ 5 & 1 \end{bmatrix}\),
\(B=\begin{bmatrix} 1 & 2
\\ -1 & -2 \end{bmatrix}\)
\answer{eigenvalue \(\tfrac17\); eigenvectors proportional to \((1,-4)\)}

\item \(A=\begin{bmatrix} 3 & -1
\\ -1 & -2 \end{bmatrix}\),
\(B=\begin{bmatrix} -3 & -3
\\ 1 & -1 \end{bmatrix}\)
\answer{eigenvalues \(-1,\tfrac76\); eigenvectors proportional to \((1,0),(1,-\tfrac{13}6\), respectively}

\item \(A=\begin{bmatrix} 1 & -1
\\ 1 & -4 \end{bmatrix}\),
\(B=\begin{bmatrix} 1 & -0
\\ 2 & -2 \end{bmatrix}\)
\answer{eigenvalues \(1\pm \i/\sqrt2\), eigenvectors proportional to \((1,\mp \i/\sqrt2)\), respectively}

\item \(A=\begin{bmatrix} 0 & 1 & -2
\\ -2 & -1 & -2
\\ 1 & 0 & 2 \end{bmatrix}\),
\(B=\begin{bmatrix} 1 & 0 & 1
\\ 0 & -2 & 1
\\ -1 & -3 & -1 \end{bmatrix}\)
\answer{eigenvalues \(\tfrac73,-2,0\); eigenvectors proportional to \((1,-\tfrac18,-\tfrac{59}{104}),(0,0,1),(1,-1,-\tfrac12)\), respectively}

\item \(A=\begin{bmatrix} 0 & -2 & 1
\\ -4 & 2 & 3
\\ 2 & -2 & -1 \end{bmatrix}\),
\(B=\begin{bmatrix} 0 & 0 & -1
\\ 0 & 0 & -1
\\ -2 & 0 & 0 \end{bmatrix}\)
\answer{eigenvalues \(0,-2\); eigenvectors proportional to \((1,\tfrac12,1),(0,-\tfrac12,1)\), respectively}

\item \(A=\begin{bmatrix} -1 & 1 & -2
\\ -4 & 0 & 3
\\ 0 & -3 & 0 \end{bmatrix}\),
\(B=\begin{bmatrix} -1 & 1 & 1
\\ 5 & -2 & 1
\\ 2 & -2 & -2 \end{bmatrix}\)
\answer{eigenvalue \(-\tfrac{11}2\); eigenvectors proportional to \((15,22,-13)\)}

\item \(A=\begin{bmatrix} -1 & 2 & -1
\\ 0 & -2 & -1
\\ 3 & 2 & 1 \end{bmatrix}\),
\(B=\begin{bmatrix} -1 & 2 & -1
\\ 0 & 2 & 1
\\ 3 & 1 & -2 \end{bmatrix}\)
\answer{eigenvalues \(1,\tfrac{12}{17},-1\); eigenvectors proportional to \((1,0,0),(4,1,-2),(1,-7,-15)\), respectively}

\end{enumerate}
\end{exercise}


%\begin{verbatim}
%n=4,A=0+round(randn(n)*2),B=0+round(randn(n)),[v,d]=eig(A,B), diag(d)'
%\end{verbatim}
\begin{exercise} \label{ex:} 
Use \script\ to find and describe all eigenvalues and corresponding eigenvectors of the \idx{generalised eigen-problem} \(A\vv=\lambda B\vv\) for the following pairs of matrices.
\begin{enumerate}
\item \(A=\begin{bmatrix} -2 & 1 & -2 & -2
\\ -2 & 2 & -1 & 0
\\ -1 & 1 & 2 & 1
\\ 1 & -2 & -1 & -1 \end{bmatrix}\),
\(B=\begin{bmatrix} -1 & 0 & -3 & 1
\\ -1 & 1 & -1 & 1
\\ -1 & 1 & 0 & 1
\\ 3 & 5 & 0 & -2 \end{bmatrix}\)
\answer{eigenvalues \(21.72, 1.50, -0.08, -2.15\) \twodp}

\item \(A=\begin{bmatrix} 3 & 1 & -2 & -4
\\ -2 & -3 & -3 & 2
\\ 1 & 1 & 1 & 2
\\ 1 & 2 & 0 & 0 \end{bmatrix}\),
\(B=\begin{bmatrix} -2 & 1 & 1 & 1
\\ 1 & -1 & 1 & -1
\\ 2 & 0 & -1 & 0
\\ 0 & 0 & 1 & 0 \end{bmatrix}\)
\answer{eigenvalues \(-5.26, 1.53, -0.98\) \twodp}

\item \(A=\begin{bmatrix} 0 & 2 & -1 & -1
\\ 0 & 0 & -1 & 0
\\ 3 & -3 & -1 & 1
\\ 2 & -3 & -2 & -3 \end{bmatrix}\),
\(B=\begin{bmatrix} 1 & 2 & 1 & 1
\\ 1 & 0 & 0 & -1
\\ 0 & 0 & 0 & 0
\\ 2 & -1 & -1 & 0 \end{bmatrix}\)
\answer{eigenvalues \(0.73\pm1.13\i, -1.37\) \twodp}

\item \(A=\begin{bmatrix} 4 & 4 & 0 & 3
\\ 1 & -2 & 0 & 1
\\ 0 & 2 & 4 & 2
\\ 0 & 2 & -1 & 1 \end{bmatrix}\),
\(B=\begin{bmatrix} -1 & 2 & -1 & 1
\\ -2 & 2 & 1 & 1
\\ -1 & 0 & 1 & -2
\\ 1 & 1 & 1 & -1 \end{bmatrix}\)
\answer{eigenvalues \(2.74, -2.20, -0.50\pm0.62\i\) \twodp}
  
\end{enumerate}
\end{exercise}


\begin{exercise} \label{ex:} 
Use the properties of determinants (\autoref{ch:ddm}), and that an \(n\)th~degree polynomial has exactly \(n\)~zeros (when counted according to multiplicity), to explain why the generalised eigen-problem \(A\vv=\lambda B\vv\), for real \(n\times n\) matrices~\(A\) and~\(B\), has \(n\)~eigenvalues iff matrix~\(B\) is invertible.
\end{exercise}


%\begin{verbatim}
%n=2,A=randn(2);A=0+round(A+A'),B=randn(n);B=0+round(B+B'),[v,d]=eig(A,B)
%\end{verbatim}
\begin{exercise} \label{ex:} 
Consider the \idx{generalised eigen-problem} \(A\vv=\lambda B\vv\) and let \(\lambda_1\neq\lambda_2\) be distinct eigenvalues with corresponding eigenvectors~\(\vv_1\) and~\(\vv_2\), respectively.
Invent two real \emph{symmetric} matrices~\(A\) and~\(B\) and record your working which demonstrates that although~\(\vv_1\) and~\(\vv_2\) are not orthogonal, nonetheless \(\tr{\vv_1}B\vv_2=0\).
\answer{one possibility is \(A=\begin{bmatrix} -1&-1\\-1&1 \end{bmatrix}\) and \(B=\begin{bmatrix} 1&1\\1&0 \end{bmatrix}\)} 
\end{exercise}


\begin{exercise} \label{ex:eennmov} 
Consider the \idx{generalised eigen-problem} \(A\vv=\lambda B\vv\) for real \emph{symmetric} \(n\times n\) matrices~\(A\) and~\(B\). 
Let \(\lambda_1\neq\lambda_2\) be distinct eigenvalues with corresponding eigenvectors~\(\vv_1\) and~\(\vv_2\), respectively.
Prove that \(\tr{\vv_1}B\vv_2=0\) always holds.
\end{exercise}


\begin{exercise} \label{ex:eennmre} 
Prove that if both~\(A\) and~\(B\) are real \emph{symmetric} \(n\times n\) matrices, then the eigenvalues of the \idx{generalised eigen-problem} \(A\vv=\lambda B\vv\) are all real.
Briefly explain why it is necessary to also have the additional proviso that the eigenvalues of~\(B\) are either all positive or all negative.
\end{exercise}


\begin{exercise} \label{ex:} 
In view of the preceding \autoref{ex:eennmre},
invent real symmetric matrices~\(A\) and~\(B\) such that the generalised eigen-problem \(A\vv=\lambda B\vv\) has complex valued eigenvalues.
% A=[0 1;1 0], B=[1 0;0 -1] will suffice, or vice versa.
\end{exercise}



\begin{exercise} \label{ex:} 
Explain briefly how the properties established in the the previous two Exercises~\ref{ex:eennmov} and~\ref{ex:eennmre} generalise important properties of the standard eigen-problem \(A\xv=\lambda\xv\) for symmetric~\(A\).
\end{exercise}






%\begin{verbatim}
%n=3,dt=0.5
%rs=[1 1/2 1/3 2/3 3/2 1/4 3/4 4/3]
%erat=rs(randperm(length(rs),n))
%[cn,cd]=rat(randn(1,n),0.2);
%[r,t]=meshgrid(erat,(0:2*n-1)*dt);
%f=(r.^t)*(cn./cd)'
%A=hankel(f(2:n+1),f(n+1:2*n))
%B=hankel(f(1:n),f(n:2*n-1))
%lambda=eig(A,B)
%r=log(lambda)/dt
%[U,P]=meshgrid(lambda,0:n-1);
%U=U.^P
%rcond(U)
%c=U\f(1:n)
%expr=exp(r)
%\end{verbatim}
\begin{exercise} \label{eg:}
Consider the specified data values~\fv\ at the specified times, and by hand or \script\ fit a sum of exponentials~\eqref{eq:eiddse}, \(f(t)=c_1e^{r_1t}+c_2e^{r_2t}+\cdots+c_ne^{r_nt}\).
Plot the data and the curve you have fitted.
\begin{enumerate}
\item For times \(0,1,2,3\) the data values are \[\fv=(3, 2.75, 2.5625, 2.4219).\]
\answer{\(f=(3/4)^t+2\)}

\item For times \(0,1,2,3\) the data values are \[\fv=(1.5833, 1.3333, 1.3056, 1.4954).\]
\answer{\(f=\tfrac54(2/3)^t+\tfrac13(3/2)^t\)}

\item For times \(0,2,4,6\) the data values are \[\fv=(-1, 0.222222, 0.172840, 0.085048).\]
\answer{\(f=-2(1/3)^t+(2/3)^t\)}

\item For times \(0,0.5,1,1.5\) the data values are \[\fv=(1.3, 0.75355, 0.45, 0.27678).\]
\answer{\(f=\tfrac45(1/4)^t+\tfrac12(1/2)^t\)}

\item For times \(0,1,2,\ldots,5\) the data values are 
\setbox\ajrqrbox\hbox{\qrcode{% f
f=[-1.00000
  0.12500
  0.71875
  1.03906
  1.21680
  1.31885]
}}\marginpar{\usebox{\ajrqrbox}}%
\[ \fv=\begin{bmatrix} -1.00000
\\ 0.12500
\\ 0.71875
\\ 1.03906
\\ 1.21680
\\ 1.31885 \end{bmatrix}.\]
\answer{\(f=-2(1/2)^t-\tfrac12(3/4)^t+\tfrac32\)}
  
\item For times \(0,1,2,\ldots,5\) the data values are 
\setbox\ajrqrbox\hbox{\qrcode{% f
f=[3.250000
   2.041667
   1.225694
   0.673900
   0.300178
   0.046636]
}}\marginpar{\usebox{\ajrqrbox}}%
\[ \fv=\begin{bmatrix} 3.250000
\\ 2.041667
\\ 1.225694
\\ 0.673900
\\ 0.300178
\\ 0.046636 \end{bmatrix}.\]
\answer{\(f=3.25(2/3)^t+\tfrac12(3/4)^t-\tfrac12\)}
  
\item For times \(0,2,4,\ldots,10\) the data values are 
\setbox\ajrqrbox\hbox{\qrcode{% f
f=[0.16667
    1.38889
    2.72984
    5.81259
   12.88104
   28.87004]
}}\marginpar{\usebox{\ajrqrbox}}%
\[ \fv=\begin{bmatrix} 0.16667
\\ 1.38889
\\ 2.72984
\\ 5.81259
\\ 12.88104
\\ 28.87004 \end{bmatrix}.\]
\answer{\(f=\tfrac12(3/2)^t-(1/3)^t+\tfrac23(3/4)^t\)}
  
\item For times \(0,0.5,1,\ldots,2.5\) the data values are 
\setbox\ajrqrbox\hbox{\qrcode{% f
f=[0.75000
  0.89846
  1.08333
  1.30323
  1.55903
  1.85343]
}}\marginpar{\usebox{\ajrqrbox}}%
\[ \fv=\begin{bmatrix} 0.75000
\\ 0.89846
\\ 1.08333
\\ 1.30323
\\ 1.55903
\\ 1.85343 \end{bmatrix}.\]
\answer{\(f=(4/3)^t+\tfrac14(1/2)^t-\tfrac12(3/4)^t\)}
  
\end{enumerate}
\end{exercise}



%\begin{verbatim}
%n=4,dt=2
%ws=1+rand(1,n/2),rs=-rand(1,n/2)
%ps=pi*rand(1,n/2),cs=randn(1,n/2)
%ts=(0:2*n-1)'*dt;
%f=0;for k=1:n/2,f=f+cs(k)*exp(rs(k)*ts).*cos(ws(k)*ts+ps(k));end
%f=f
%A=hankel(f(2:n+1),f(n+1:2*n));
%B=hankel(f(1:n),f(n:2*n-1));
%lambda=eig(A,B)
%r=log(lambda)/dt
%U=bsxfun(@power,lambda,0:n-1).';
%rcond(U)
%c=U\f(1:n);
%expr=exp(r);
%\end{verbatim}
\begin{exercise} \label{ex:} 
Consider the specified data values~\fv\ of some decaying oscillations at the specified times (in seconds).
Use \script\ to fit a sum of exponentials~\eqref{eq:eiddse}, \(f(t)=c_1e^{r_1t}+c_2e^{r_2t}+\cdots+c_ne^{r_nt}\).
Confirm your fit reproduces the data values.
What frequencies do you detect in the fitted constants?
\begin{enumerate}
\item  For times \(0,1,2,\ldots,7\) the data values are 
\setbox\ajrqrbox\hbox{\qrcode{% f
f=[-0.153161
   0.484787
  -0.124780
  -0.283690
   0.201896
   0.174832
  -0.200418
  -0.092107]
}}\marginpar{\usebox{\ajrqrbox}}%
\[ \fv=\begin{bmatrix} -0.153161
\\  0.484787
\\ -0.124780
\\ -0.283690
\\  0.201896
\\  0.174832
\\ -0.200418
\\ -0.092107 \end{bmatrix}.\]
\answer{frequencies \(1.20\) and \(1.70\) radians/sec \twodp}
  
\item  For times \(0,1,2,\ldots,7\) the data values are 
\setbox\ajrqrbox\hbox{\qrcode{% f
f=[1.225432
   0.499044
   0.034164
  -0.093971
  -0.103784
  -0.028558
   0.033825
   0.026439]
}}\marginpar{\usebox{\ajrqrbox}}%
\[ \fv=\begin{bmatrix} 1.225432
\\  0.499044
\\  0.034164
\\ -0.093971
\\ -0.103784
\\ -0.028558
\\  0.033825
\\  0.026439 \end{bmatrix}.\]
\answer{frequencies \(1.57\) and \(1.16\) radians/sec \twodp}
  
\item  For times \(0,0.5,1,\ldots,3.5\) the data values are 
\setbox\ajrqrbox\hbox{\qrcode{% f
f=[0.0060901
   0.1753701
   0.2385516
   0.1787653
   0.0404368
  -0.0998597
  -0.1742818
  -0.1552607]
}}\marginpar{\usebox{\ajrqrbox}}%
\[ \fv=\begin{bmatrix} 0.0060901
\\  0.1753701
\\  0.2385516
\\  0.1787653
\\  0.0404368
\\ -0.0998597
\\ -0.1742818
\\ -0.1552607 \end{bmatrix}.\]
\answer{frequencies \(1.94\) and \(1.46\) radians/sec \twodp}
  
\item  For times \(0,2,4,\ldots,14\) the data values are 
\setbox\ajrqrbox\hbox{\qrcode{% f
f=[1.297867
  -0.800035
   0.487305
  -0.285714
   0.161499
  -0.090411
   0.051244
  -0.028936]
}}\marginpar{\usebox{\ajrqrbox}}%
\[ \fv=\begin{bmatrix} 1.297867
\\ -0.800035
\\  0.487305
\\ -0.285714
\\  0.161499
\\ -0.090411
\\  0.051244
\\ -0.028936 \end{bmatrix}.\]
\answer{frequencies \(1.01\) and \(1.53\) radians/sec \twodp}
  
\end{enumerate}
\end{exercise}




% supernova CN1991T perhaps \cite[p.146, Ch.~8]{Pereyra2010}
% days from 2448300 and 20-B, digitised from SN1991T.pdf
%   64.2857    7.2833
%   65.0794    7.5333
%   66.6667    7.8500
%   73.8095    8.2500
%   81.7460    8.1333
%   88.0952    7.6000
%   95.2381    6.8667
%  104.7619    5.9167
%  121.4286    5.4500
%  132.5397    5.2500
%  146.0317    5.0333
%  159.5238    4.8667
% use data 3 to 8, say roughly 8 days apart, and mangle to get real
%\begin{verbatim}
%n=3,dt=8
%f=[0.6543
%   1.5783
%   1.6406
%   0.3678
%   0.05918
%   0.008227]
%A=hankel(f(2:n+1),f(n+1:2*n));
%B=hankel(f(1:n),f(n:2*n-1));
%lambda=eig(A,B)
%r=log(lambda)/dt
%U=bsxfun(@power,lambda,0:n-1).';
%rcond(U)
%c=U\f(1:n)
%\end{verbatim}
\begin{exercise} \label{ex:} 
In astronomy, Type~Ia supernova explode, peaking in luminosity in a few days, and then their luminosity declines over months. 
It is conjectured that the decline is powered by the radioactive decay from radioactive Nickel to Cobalt to Iron.
For the supernova \textsc{sn1991t} six measurements, starting on Julian day 2448366 and taken approximately 8~days apart, give the following luminosity values (in some units) \cite[p.146]{Pereyra2010}:
\setbox\ajrqrbox\hbox{\qrcode{% lumen
f=[0.6543
   1.5783
   1.6406
   0.3678
   0.05918
   0.008227]
}}\marginpar{\usebox{\ajrqrbox}}%
\begin{equation*}
\fv=\begin{bmatrix} 0.6543
\\ 1.5783
\\ 1.6406
\\ 0.3678
\\ 0.05918
\\ 0.008227 \end{bmatrix}
\end{equation*}
Detect the exponential decay in this supernova data by fitting a sum of exponentials~\eqref{eq:eiddse}, \(f(t)=c_1e^{r_1t}+c_2e^{r_2t}+c_3e^{r_3t}\).
Comment on the result.
\answer{\(f=1459e^{-0.36t}-2044e^{-0.33t}+585e^{-0.27t}\) where \(t\) is days post-2448366; but rcond is very poor!}
\end{exercise}




%\begin{verbatim}
%year=(1944:1993)';
%soi=[-0.03; 0.74; 6.37; -7.28; 0.44; -0.99; 1.32
%6.42; -6.51; 0.07; -1.96; 1.72; 6.49; -5.61
%-0.24; -2.90; 1.92; 6.54; -4.61; -0.47; -3.82
%1.94; 6.56; -3.53; -0.59; -4.69; 1.76; 6.53
%-2.38; -0.59; -5.48; 1.41; 6.41; -1.18; -0.45
%-6.19; 0.89; 6.19; 0.03; -0.16; -6.78; 0.21; 5.84
%1.23; 0.30; -7.22; -0.60; 5.33; 2.36; 0.91 ];
%n=length(year)/2, dt=1
%A=hankel(soi(2:n+1),soi(n+1:2*n));
%B=hankel(soi(1:n),soi(n:2*n-1));
%lambda=eig(A,B)
%r=log(lambda)/dt
%U=bsxfun(@power,lambda,0:n-1).';
%rcond(U)
%c=U\soi(1:n)
%expr=exp(r);
%\end{verbatim}
\begin{exercise} \label{ex:} 
Recall that \autoref{eg:orthbapp} introduced the phenomenon of \idx{El Nino} which makes a large impact on the world's weather.
El Nino is correlated significantly with the difference in atmospheric pressure between Darwin and Tahiti---the so-called \idx{Southern Oscillation Index} (\soi).
\autoref{fig:soiRoundData} plots a (`smoothed') yearly average \soi\ each year for fifty years up to 1993.
Here detect the cycles in the \soi\ by analysing the plotted data as a sum of exponentials~\eqref{eq:eiddse}, \(f(t)=c_1e^{r_1t}+c_2e^{r_2t}+\cdots+c_ne^{r_nt}\).
\begin{enumerate}
\item Enter the data into \script:
\setbox\ajrqrbox\hbox{\qrcode{% soi
year=(1944:1993)';
soi=[-0.03; 0.74; 6.37; -7.28; 0.44; -0.99; 1.32
6.42; -6.51; 0.07; -1.96; 1.72; 6.49; -5.61
-0.24; -2.90; 1.92; 6.54; -4.61; -0.47; -3.82
1.94; 6.56; -3.53; -0.59; -4.69; 1.76; 6.53
-2.38; -0.59; -5.48; 1.41; 6.41; -1.18; -0.45
-6.19; 0.89; 6.19; 0.03; -0.16; -6.78; 0.21; 5.84
1.23; 0.30; -7.22; -0.60; 5.33; 2.36; 0.91 ];
}}\marginpar{\usebox{\ajrqrbox}}%
\begin{verbatim}
year=(1944:1993)';
soi=[-0.03; 0.74; 6.37; -7.28; 0.44; -0.99; 1.32
6.42; -6.51; 0.07; -1.96; 1.72; 6.49; -5.61
-0.24; -2.90; 1.92; 6.54; -4.61; -0.47; -3.82
1.94; 6.56; -3.53; -0.59; -4.69; 1.76; 6.53
-2.38; -0.59; -5.48; 1.41; 6.41; -1.18; -0.45
-6.19; 0.89; 6.19; 0.03; -0.16; -6.78; 0.21; 5.84
1.23; 0.30; -7.22; -0.60; 5.33; 2.36; 0.91 ];
\end{verbatim}

\item Use \autoref{pro:ei} in \script\ to compute the twenty-five complex rates~\rv\ and twenty-five complex coefficients~\cv\ that fit the data.
% albeit poor rcond

\item Find the four relatively large coefficients, when compared to the relatively small coefficients of magnitude\({}<0.015\).
% should really look at size of c*exp(Re(r)t)---but almost the same

\item Explain why these four coefficients indicate that the \soi\ data appears to be dominantly composed of oscillations with periods about \(5.1\)~years and \(2.5\)~years.
% is the 2.5 year a harmonic of the 5 year oscillation?

\end{enumerate}
\end{exercise}











\begin{exercise} 
In a few sentences, answer\slash discuss each of the the following.
\begin{enumerate}
\item  Given an \(n\times n\) matrix, what leads to the characteristic polynomial being of \(n\)th~degree?

\item How does the trace of a matrix appear in its characteristic polynomial?

\item What is the importance of the multiplicity of eigenvalues of a matrix?

\item What is the relation between the characteristic polynomial of a matrix and its characteristic equation?

\item Why is it beautifully useful to cater for complex valued eigenvalues and eigenvectors of real matrices arising in real problems?

\item What is the evidence for repeated eigenvalues generally being sensitive to computational and experimental errors?

\item What causes \(\yv(t)=c_1\lambda_1^t\vv_1+c_2\lambda_2^t\vv_2+\cdots+c_m\lambda_m^t\vv_m\) to form a general solution to the evolving system \(\yv(t+1)=A\yv(t)\)?

\item How can the singular values of a matrix arise from an eigen-problem?

\item Describe some scenarios that require fitting a sum of exponentials to data.

\end{enumerate}
\end{exercise}

\begin{comment}%{ED498555.pdf}
why, what caused X?
how did X occur?
what-if? what-if-not?
how does X compare with Y?
what is the evidence for X?
why is X important?
\end{comment}





%%!TEX root = ../larxxia.tex

\section{Linearly independent vectors may form a basis}
\label{sec:lisb}

\secttoc
\begin{comment}
\pooliv{p.92--7,198--208} \holti{\S2.3}
\end{comment}

\index{linearly dependent|(}
\index{linearly independent|(}
\index{linear dependence|(}
\index{linear independence|(}


In \cref{ch:eesm} on symmetric matrices, the \idx{eigenvector}s from \idx{distinct eigenvalues} are proved to be always \idx{orthogonal}---because of matrix symmetry.  
For general matrices, the eigenvectors are not orthogonal---as introduced at the start of this \cref{ch:gee}.  
But the orthogonal property is extremely useful.
Question: is there some analogue of orthogonality that is similarly useful for general matrices?
Answer: yes. 
We now extend ``orthogonal'' to the more general concept of \text{``linearly independent''.}

One reason that \idx{orthogonal vectors} are useful is that they can form an \idx{orthonormal basis} and hence act as the unit vectors of an orthogonal coordinate systems.
Analogously, the concept of linearly independent vectors is closely connected to \idx{coordinate system}s that are not orthogonal.



\paragraph{Subspace \idx{coordinate system}s} 
In any given mathematical problem, an application wants two things from a \idx{general solution}: 
\begin{itemize}
\item firstly, the general solution must encompass every possibility (the solution must span the possibilities); and 
\item secondly, each possible solution should have a unique algebraic form in the general solution.
\end{itemize}
For an example of the need for a unique algebraic form, let's suppose we wanted to find solutions to the \idx{differential equation} \(d^2y/dt^2-y=0\)\,. 
You might find \(y=3e^x+2e^{-x}\), whereas I find \(y=5\cosh x+\sinh x\), and a friend finds \(y=e^x+4\cosh x\).
It appears from these disparate algebraic forms that we all disagree.
Should we all go and search for errors in the solution process?  No.
The reason is that all these solutions are the same.
The apparent differences arise only because you choose \idx{exponential}s to represent the solution, whereas I choose hyperbolic functions, and the friend a mixture; the solutions are the same, it is only the algebraic representation that appears different. 
In general, when we cannot immediately distinguish identical solutions, all algebraic manipulation becomes immensely more difficult due to algebraic ambiguity.
To avoid such immense difficulties we need to introduce, in both calculus and linear algebra, the concept of \text{linear independence.}

\needlines6
\begin{wrapfigure}r{0pt}
\begin{tikzpicture}
\newcommand{\ppoint}[2]{
    \pgfmathparse{#1*3+#2*1}\let\h\pgfmathresult
    \pgfmathparse{#1*1+#2*2}\let\v\pgfmathresult
    \addplot[red,mark=*,only marks] coordinates {(\h,\v)};
    \edef\tempe{\noexpand
    \node[left] at (axis cs:\h,\v) {$(#1,#2)$};
    }\tempe
    }
\begin{axis}[footnotesize,font=\footnotesize
  ,axis lines=none%,xlabel={$x$},ylabel={$y$}
  , axis equal image
  , view={0}{90}
  ,xmax=7.5,ymax=5.5,xmin=-7.6,ymin=-5.6
  ]
\addplot3[mesh,brown,samples=17,domain=-4:4,dotted] (3*x+y,x+2*y,0);
\addplot3[mesh,brown,samples=9,domain=-4:4] (3*x+y,x+2*y,0);
\addplot[blue,quiver={u=3,v=1},-stealth,thick] coordinates {(0,0)};
\node[right] at (axis cs:3,1) {$\vv_1$};
\addplot[blue,quiver={u=1,v=2},-stealth,thick] coordinates {(0,0)};
\node[above] at (axis cs:1,2) {$\vv_2$};
\ppoint21
\ppoint1{-3}
\ppoint{-2}3
\ppoint00
\end{axis}
\end{tikzpicture}
\end{wrapfigure}
Linear independence empowers us, often implicitly, to use a non-orthogonal \idx{coordinate system} in a \idx{subspace}.
We replace orthogonal \idx{unit vector}s by any suitable set of basis vectors.
For example, in the plane any two vectors at an angle to each other suffice to be able to describe uniquely every vector (point) in the plane.
As illustrated to the right, every point in the plane (end-point of a vector) is a unique \idx{linear combination} of the two drawn basis vectors~\(\vv_1\) and~\(\vv_2\).
Such a pair of basis vectors, termed a linearly independent pair, avoids the  difficulties of \text{algebraic ambiguity.}



\subsection{Linearly (in)dependent sets}

This section defines ``linearly dependent'' and ``linearly independent'', and then relates the concept to homogeneous linear equations, orthogonality, and sets of eigenvectors.


\begingroup
\def\temp#1{\begin{tikzpicture}% following top aligns tikzpicture
[baseline={([yshift={-\ht\strutbox}]current bounding box.north)}]
\begin{axis}[footnotesize,font=\footnotesize
  ,axis lines=middle, axis equal image
  ,xmax=2.3,ymax=2.3
  ]
\addplot[blue,quiver={u=2,v=1},-stealth,thick] coordinates {(0,0)};
\node[below] at (axis cs:2,1) {$\vv_1$};
\addplot[blue,quiver={u=1,v=2},-stealth,thick] coordinates {(0,0)};
\node[below] at (axis cs:1,2) {$\quad\vv_2$};
\ifcase#1
\or%1
\addplot[brown,quiver={u=-1.333,v=-0.667},-stealth] coordinates {(0,0)};
\addplot[brown,quiver={u=1.333,v=2.667},-stealth] coordinates {(-1.333,-0.667)};
\addplot[red,quiver={u=0,v=2},-stealth,thick] coordinates {(0,0)};
\node[right] at (axis cs:0,2) {$(0,2)$};
\or%2
\addplot[brown,quiver={u=1.333,v=0.667},-stealth] coordinates {(0,0)};
\addplot[brown,quiver={u=0.667,v=1.333},-stealth] coordinates {(1.333,0.667)};
\addplot[red,quiver={u=2,v=2},-stealth,thick] coordinates {(0,0)};
\node[left] at (axis cs:2,2) {$(2,2)$};
\or%3
\addplot[brown,quiver={u=2,v=1},-stealth] coordinates {(0,0)};
\addplot[brown,quiver={u=-1,v=-2},-stealth] coordinates {(2,1)};
\addplot[red,quiver={u=1,v=-1},-stealth,thick] coordinates {(0,0)};
\node[above] at (axis cs:1,-1) {\quad\qquad$(1,-1)$};
\or%4
\addplot[brown,quiver={u=-2,v=-1},-stealth] coordinates {(0,0)};
\addplot[brown,quiver={u=-1,v=-2},-stealth] coordinates {(-2,-1)};
\addplot[red,quiver={u=-3,v=-3},-stealth,thick] coordinates {(0,0)};
\node[above] at (axis cs:-3,-3) {\qquad\qquad$(-3,-3)$};
\else
\fi
\end{axis}
\end{tikzpicture}}%
\needlines7
\begin{wrapfigure}r{0pt} \temp0 \end{wrapfigure}
\begin{example}[2D \idx{non-orthogonal coordinates}] \label{eg:2dnoc}
Show that every vector in the plane~\(\RR^2\) can be written uniquely as a \idx{linear combination} of the two vectors \(\vv_1=(2,1)\) and \(\vv_2=(1,2)\) that are shown to the right.
\vspace*{3\baselineskip} % lines to clear the graphic??

\begin{solution} 
Let's start with some specific example vectors.
\def\tmpp#1#2{\parbox[t]{20em}{\raggedright #2} \temp{#1}}
\begin{enumerate}
\item \tmpp1{The vector \((0,2)\) may be written as the linear combination \((0,2)=-\frac23\vv_1+\frac43\vv_2\) as shown.} 

\item \tmpp2{The vector \((2,2)\) may be written as the linear combination \((2,2)=\frac23\vv_1+\frac23\vv_2\) as shown.}

\item \tmpp3{The vector \((1,-1)\) may be written as the linear combination \((1,-1)=\vv_1-\vv_2\) as shown.}

\item \tmpp4{The vector \((-3,-3)\) may be written as the linear combination \((-3,-3)=-\vv_1-\vv_2\) as shown.}

\end{enumerate}

\begin{wrapfigure}r{0pt}
\begin{tikzpicture}
\begin{axis}[small,font=\footnotesize
  ,axis lines=middle, axis equal image, view={0}{90}
  ,xmax=5.5,ymax=5.5,xmin=-5.5,ymin=-5.5
  ,xlabel={$x$},ylabel={$y$}
  ]
\addplot3[mesh,brown,samples=13,domain=-3:3,dotted] (2*x+y,x+2*y,0);
\addplot3[mesh,brown,samples=7,domain=-3:3] (2*x+y,x+2*y,0);
\addplot[blue,quiver={u=2,v=1},-stealth,thick] coordinates {(0,0)};
\node[right] at (axis cs:2,1) {$\vv_1$};
\addplot[blue,quiver={u=1,v=2},-stealth,thick] coordinates {(0,0)};
\node[above] at (axis cs:1,2) {$\vv_2$};
\end{axis}
\end{tikzpicture}
\end{wrapfigure}
Now proceed to consider a general vector \((x,y)\) and seek it as a linear combination of~\(\vv_1\) and~\(\vv_2\), namely \((x,y)=c_1\vv_1+c_2\vv_2\)\,.
That is, let's write each and every point in the plane as a linear combination of~\(\vv_1\) and~\(\vv_2\) as illustrated to the right.
Rewrite the equation in matrix-vector form as
\begin{equation*}
\begin{bmatrix} \vv_1&\vv_2 \end{bmatrix}
\begin{bmatrix} c_1\\c_2 \end{bmatrix}
=\begin{bmatrix} x\\y \end{bmatrix}
,\quad\text{that is,}\quad
V\cv=\begin{bmatrix} x\\y \end{bmatrix}
\text{ for }V=\begin{bmatrix} 2&1\\1&2 \end{bmatrix}.
\end{equation*}
For any given \((x,y)\), \(V\cv=(x,y)\) is a system of linear equations for the coefficients~\cv.
\cref{thm:ftim2} asserts that the system has a unique solution~\cv\ if and only if the matrix~\(V\) is invertible.
Here the unique solution is then that the vector of coefficients
\begin{equation*}
\cv=V^{-1}\begin{bmatrix} x\\y \end{bmatrix}
=\begin{bmatrix} \frac23&-\frac13\\-\frac13&\frac23 \end{bmatrix}
\begin{bmatrix} x\\y \end{bmatrix}.
\end{equation*}
Equivalently, \cref{thm:ftim2iii} asserts the system has a unique solution~\cv---unique coefficients~\cv---if and only if the {homogeneous} system \(V\cv=\ov\) has only the zero solution \(\cv=\ov\)\,.
It is this last statement that leads to the upcoming \cref{def:lindep} of \idx{linearly independent} vectors.
\end{solution}
\end{example}
\endgroup



{%group
\def\temp{\begin{tikzpicture}
\begin{axis}[footnotesize,font=\footnotesize
  ,axis lines=middle, axis equal image
  ,xmax=2.3,ymax=2.3
  ]
\addplot[blue,quiver={u=2,v=1},-stealth,thick] coordinates {(0,0)};
\node[below] at (axis cs:2,1) {$\vv_1$};
\addplot[blue,quiver={u=1,v=2},-stealth,thick] coordinates {(0,0)};
\node[below] at (axis cs:1,2) {$\quad\vv_2$};
\addplot[red,quiver={u=-2.5,v=1},-stealth,thick] coordinates {(0,0)};
%\node[right] at (axis cs:-2.5,1) {$(-2.5,1)$};
\end{axis}
\end{tikzpicture}}
\begin{activity}[\temp]
The vector shown to the right is which of the following \idx{linear combination}s of shown vectors~\(\vv_1\) and~\(\vv_2\)?
\actposs{\(-2\vv_1 +1.5\vv_2\)}
{\(-2.5\vv_1 +2\vv_2\)}
{\(-1.5\vv_1 +\vv_2\)}
{\(-\vv_1 +\vv_2\)}
\end{activity}
}





\needlines8
\begin{wrapfigure}r{0pt}
\qview{113}{117}{\begin{tikzpicture}
\begin{axis}[small,font=\small
  ,xlabel={$x$}, ylabel={$y$}, zlabel={$z$},label shift={-2ex}
  ,no marks ,domain=-2.5:2.5,axis equal image,zmax=2.5,zmin=-2.5
  ,view={\q}{30}
  ] 
\addplot3[surf,blue,samples=3,opacity=0.3]  {-x-y};
\addplot3[blue,thick,quiver={u=-1,v=1,w=0},-stealth] coordinates {(0,0,0)};
\node[right] at (axis cs:-1,1,0) {$\vv_1$};
\addplot3[blue,thick,quiver={u=1,v=-2,w=1},-stealth] coordinates {(0,0,0)};
\node[below] at (axis cs:1,-2,1) {$\vv_2$};
\addplot3[blue,thick,quiver={u=0,v=1,w=-1},-stealth] coordinates {(0,0,0)};
\node[below] at (axis cs:0,1,-1) {$\vv_3$};
\end{axis}
\end{tikzpicture}}
\end{wrapfigure}
\begin{example}[3D failure] 
Show that vectors in~\(\RR^3\) are not written uniquely as a \idx{linear combination} of \(\vv_1=(-1,1,0)\), \(\vv_2=(1,-2,1)\), and \(\vv_3=(0,1,-1)\). 

One reason for the failure is that these three vectors only span a plane, as shown to the right in stereo.  
The solution here looks at the different issue of unique representation.

\begin{solution} 
As one example, consider the vector~\((1,0,-1)\):
\begin{eqnarray*}
(1,0,-1)&=&-1\vv_1+0\vv_2+1\vv_3\,;
\\(1,0,-1)&=&1\vv_1+2\vv_2+3\vv_3\,;
\\(1,0,-1)&=&-2\vv_1-1\vv_2+0\vv_3\,;
\\(1,0,-1)&=&(-1+t)\vv_1+t\vv_2+(1+t)\vv_3\,,\quad\text{for every }t.
\end{eqnarray*}
This last combination shows that there are an infinite number of ways to write \((1,0,-1)\) as a linear combination of~\(\vv_1\), \(\vv_2\), and~\(\vv_3\).
Such an infinity of linear combinations means that~\(\vv_1\), \(\vv_2\), and~\(\vv_3\) cannot form the basis for a useful `coordinate system' because we cannot easily distinguish between the different combinations all describing the same vector.
The reason for the infinity of combinations is that there is a nontrivial linear combination of~\(\vv_1\), \(\vv_2\), and~\(\vv_3\) which is zero, namely here \(\vv_1+\vv_2+\vv_3=\ov\)\,.
It is this last statement that leads to the \cref{def:lindep} of \text{\idx{linearly dependent} vectors.}
\end{solution}
\end{example}





\begin{definition} \label{def:lindep} 
A set of vectors \(\{\hlist\vv k\}\) is \bfidx{linearly dependent} if there are \idx{scalar}s \hlist ck, at least one of which is nonzero, such that \(\lincomb c\vv k=\ov\)\,.
A set of vectors that is not linearly dependent is called \bfidx{linearly independent} (characterized by \emph{only} the linear combination with \(c_1=c_2=\cdots=c_k=0\) gives the \text{\idx{zero vector}).}
\end{definition}

When reading the terms ``linearly in/dependent'' be very careful: it is all too easy to misread the presence or absence of the crucial ``in''~syllable. 
The presence or absence of the ``in''~syllable makes all the difference between the property and its opposite.


\begin{example} \label{eg:lindep}
Are the following sets of vectors linearly dependent or linearly independent.  Give reasons.
\begin{enumerate}[ref=\ref{eg:lindep}(\alph*)]
\item\label[example]{eg:lindepa} \(\{(-1,1,0),\ (1,-2,1),\ (0,1,-1)\}\) 
\begin{solution} 
The set is linearly dependent as the linear combination \((-1,1,0)+(1,-2,1)+(0,1,-1)=(0,0,0)\). 
\end{solution}

\item\label[example]{eg:lindepb} \(\{(2,1),\ (1,2)\}\) 
\begin{solution} 
The set is linearly independent because the linear combination equation \(c_1(2,1)+c_2(1,2)=(0,0)\) is equivalent to the homogeneous matrix-vector system \(\begin{bmat} 2&1\\1&2 \end{bmat}\cv=\ov\) which has \emph{only} the zero solution \(\cv=\ov\)\,. 
\end{solution}

%\item \(\{(-1,1,0),\ (1,-2,1)\}\)
%\begin{solution} 
%The set is linearly independent: for the linear combination \(c_1(-1,1,0)+c_2(1,-2,1)
%=(-c_1+c_2,c_1-2c_2,c_2)=\ov\) requires the last component to be zero which requires \(c_2=0\);
%then either of the other components requires \(c_1=0\)\,.
%\end{solution}

\item\label[example]{eg:lindepc} \(\{(-2,4,1,-1,0)\}\)
\begin{solution} 
This set of one vector in~\(\RR^5\) is linearly independent as \(c_1(-2,4,1,-1,0)=\ov\) can only be satisfied with \(c_1=0\)\,. 
Indeed, any one nonzero vector~\vv\ in~\(\RR^n\) forms a linearly independent set,~\(\{\vv\}\), for the same reason.
\end{solution}



\item \(\{(2,1),\ (0,0)\}\) 
\begin{solution} 
The set is linearly dependent because the linear combination  \(0(2,1)+c_2(0,0)=(0,0)\) for every nonzero~\(c_2\). 
\end{solution}


\item \(\{\ov,\hlist[2]\vv k\}\)
\begin{solution} 
Every  set that includes the zero vector is linearly dependent as
\(c_1\ov+0\vv_2+\cdots+0\vv_k=\ov\) for every nonzero~\(c_1\).
\end{solution}

\item \index{e@$\ev_j$}\(\{\ev_1,\ev_2,\ev_3\}\), the set of \idx{standard unit vector}s in~\(\RR^3\).
\begin{solution} 
This set is linearly independent as
\(c_1\ev_1+c_2\ev_2+c_3\ev_3=(c_1,c_2,c_3)=\ov\)
\emph{only} when all three components are zero, \(c_1=c_2=c_3=0\)\,. 
\end{solution}

\item \(\{(\frac13,\frac23,\frac23),\ (\frac23,\frac13,-\frac23)\}\)
\begin{solution} 
This set is linearly independent.
Seek some linear combination \(c_1(\frac13,\frac23,\frac23) +c_2(\frac23,\frac13,-\frac23) =\ov\)\,.
Take the dot product of both sides of this equation with \((1,2,2)\):
\begin{eqnarray*}
&&c_1\begin{bmatrix}\frac13\\\frac23\\\frac23 \end{bmatrix}\cdot\begin{bmatrix} 1\\2\\2 \end{bmatrix} +c_2\begin{bmatrix} \frac23\\\frac13\\-\frac23 \end{bmatrix}\cdot\begin{bmatrix} 1\\2\\2 \end{bmatrix} =\ov\cdot\begin{bmatrix} 1\\2\\2 \end{bmatrix}
\\\implies&& c_13+c_20=0
\implies c_1=0\,.
\end{eqnarray*}
Also, take the dot product with \((2,1,-2)\): 
\begin{eqnarray*}
&&c_1\begin{bmatrix}\frac13\\\frac23\\\frac23 \end{bmatrix}\cdot\begin{bmatrix} 2\\1\\-2 \end{bmatrix} +c_2\begin{bmatrix} \frac23\\\frac13\\-\frac23 \end{bmatrix}\cdot\begin{bmatrix} 2\\1\\-2 \end{bmatrix} =\ov\cdot\begin{bmatrix} 2\\1\\-2 \end{bmatrix}
\\\implies&& c_10+c_23=0
\implies c_2=0\,.
\end{eqnarray*}
Hence \(c_1=c_2=0\) is the \emph{only} possibility and so the two given vectors are linearly independent.
\end{solution}

\end{enumerate}
These last two cases generalize to the next \cref{thm:ortholi} about the linear independence of every \idx{orthonormal set} of vectors.
\end{example}



\begin{activity}
Which of the following sets of vectors is linearly independent?
\actposs{\(\{(-1,1),(0,1)\}\)}
{\(\{(0,0),(-2,1)\}\)}
{\(\{(0,1),(0,-1)\}\)}
{\(\{(-1,2),(-2,4)\}\)}
\end{activity}




\begin{example}[calculus extension] 
In calculus the notion of a function corresponds to the notion of a vector in our linear algebra.  
For the purposes of this example, consider `vector' and `function' to be synonymous, and that `all \idx{components}' and `all~\(x\)' are synonymous. 
Show that the set \(\{e^x,e^{-x},\cosh x,\sinh x\}\) is linearly dependent.  
What is a subset that is \text{linearly independent?}
\begin{solution} 
The definition of the hyperbolic functions, namely that \(\cosh x=(e^x+e^{-x})/2\) and \(\sinh x=(e^x-e^{-x})/2\), immediately gives two nontrivial linear combinations that are zero for all~\(x\), namely \(2\cosh x-e^x-e^{-x}=0\) and \(2\sinh x-e^x+e^{-x}=0\) for all~\(x\).
Either one of these implies that the set \(\{e^x,e^{-x},\cosh x,\sinh x\}\) is \text{linearly dependent.}

Because \(e^x\) and~\(e^{-x}\) are not proportional to each other, there is no linear combination which is zero for all~\(x\), and hence the set \(\{e^x,e^{-x}\}\) is linearly independent (as are any other pairs of the four functions).
\end{solution}
\end{example}







\begin{theorem} \label{thm:ortholi}
Every \idx{orthonormal set} of vectors (\cref{def:orthoset}) is \idx{linearly independent}.
\end{theorem}
\begin{proof} %Take dot products with the defining equation.
Let \(\{\hlist\vv k\}\) be an orthonormal set of vectors in~\(\RR^n\).
Let's find all possible scalars \hlist ck such that \(\lincomb c\vv k=\ov\)\,.
Taking the dot product of this equation with~\(\vv_1\) requires
\(c_1\vv_1\cdot\vv_1+c_2\vv_2\cdot\vv_1+\cdots+c_k\vv_k\cdot\vv_1=\ov\cdot\vv_1\)\,;
by orthonormality this equation becomes
\(c_11+c_20+\cdots+c_k0=0\)\,; that is, \(c_1=0\)\,.
Similarly, taking the dot product with~\(\vv_2\) requires
\(c_1\vv_1\cdot\vv_2+c_2\vv_2\cdot\vv_2+\cdots+c_k\vv_k\cdot\vv_2=\ov\cdot\vv_2\)\,;
by orthonormality this equation becomes
\(c_10+c_21+\cdots+c_k0=0\)\,; that is, \(c_2=0\)\,.
And so on for every vector in the set, implying the coefficients \(c_1=c_2=\cdots=c_k=0\) is the only possibility.
By \cref{def:lindep}, the orthonormal set must be \text{linearly independent.}
\end{proof}


In contrast to orthonormal vectors which are always linearly independent, a set of two vectors proportional to each other is always linearly dependent as seen in the following examples.
This linear dependence of proportional vectors then generalizes in the forthcoming \cref{thm:lindeplc}.




\begin{example} 
Show that the following sets are linearly dependent.
\begin{enumerate}
\item \(\{(1,2),\ (3,6)\}\)
\begin{solution} 
Since \((3,6)=3(1,2)\) then the linear combination \(1(3,6)-3(1,2)=\ov\) and the set is linearly dependent. 
\end{solution}

\item \(\{(2.2,-2.1,0,1.5),\ (-8.8,8.4,0,-6)\}\)
\begin{solution} 
Since  \((-8.8,8.4,0,-6)=-4(2.2,-2.1,0,1.5)\) then the linear combination 
\begin{equation*}
(-8.8,8.4,0,-6)+4(2.2,-2.1,0,1.5)=\ov\,,
\end{equation*}
and so the set is linearly dependent.
\end{solution}
\end{enumerate}
\end{example}




\begin{activity}
For what value of~\(c\) is the set \(\{(-3c,-2+2c),(1,2)\}\) linearly dependent?
%v=round(randn(2)*2)+0,w=round(randn(2))+0,0-eig(v,w)
\actposs[4]{\(c=\frac14\)}{\(c=-\frac13\)}{\(c=0\)}{\(c=1\)}
\end{activity}




\begin{theorem} \label{thm:lindeplc} 
A set of vectors \(\{\hlist\vv m\}\) is \idx{linearly dependent} if and only if at least one of the vectors can be expressed as a \idx{linear combination} of the other vectors.
In particular, a set of two vectors \(\{\vv_1,\vv_2\}\) is linearly dependent if and only if one of the vectors is a scalar multiple of \text{the other.}
\end{theorem}

\begin{proof} 
\cref{ex:lindeplc} establishes the particular case of a set of two vectors.
%Start with the particular case of the set of two vectors \(\{\vv_1,\vv_2\}\).
%Say \(\vv_1\) is a multiple of~\(\vv_2\), \(\vv_1=c\vv_2\)\,, then the linear combination \(\vv_1-c\vv_2=\ov\)\,, and so the set is linearly dependent. 
%Similarly if \(\vv_2\) is a multiple of~\(\vv_1\).
%Conversely, if the set is linearly dependent then \(c_1\vv_1+c_2\vv_2=0\) where one or both of \(c_1,c_2\neq0\): if \(c_1\neq 0\) then rearrange to \(\vv_1=-(c_2/c_1)\vv_2\) and so \(\vv_1\) is a multiple of~\(\vv_2\); and similarly if \(c_2\neq 0\)\,.

In the general case of \(m\)~vectors, first establish that if one of the vectors can be expressed as a {linear combination} of the others, then the set is linearly dependent.
Let's label the set of vectors so that it is vector~\(\vv_1\) which is a linear combination of the others; that is, \(\vv_1=\lincomb[2]c\vv m\)\,.
Rearranging, \((-1)\vv_1+\lincomb[2]c\vv m=\ov\)\,; that is, there is a non-trivial (as at least \(c_1=-1\neq0\)) linear combination of the set of vectors which is zero.
Hence the set is \text{linearly dependent.}

Second, establish the converse.  
Given the set is linearly dependent, there exist coefficients, not all zero, such that \(\lincomb c\vv m=\ov\)\,.  
Suppose that we have labelled the vectors so that \(c_1\neq 0\)\,.  
Then rearranging the equation gives
\(c_1\vv_1=-c_2\vv_2-c_3\vv_3-\cdots-c_m\vv_m\)\,.
Divide by the nonzero~\(c_1\) to deduce
\(\vv_1=-(c_2/c_1)\vv_2-(c_3/c_1)\vv_3-\cdots-(c_m/c_1)\vv_m\)\,;
that is, \(\vv_1\)~is a linear combination of the \text{other vectors.}
\end{proof}


\begin{example} 
Invoke \cref{thm:lindeplc} to deduce whether the following sets are linearly independent or linearly dependent.
\begin{enumerate}
\item \(\{(-1,1,0),\ (1,-2,1),\ (0,1,-1)\}\)
\begin{solution} 
Since \((1,-2,1)=-(-1,1,0)-(0,1,-1)\) the set must be linearly dependent.
\end{solution}

\item 
\begin{figbox}{\begin{tikzpicture}
\begin{axis}[footnotesize,font=\footnotesize
  ,axis lines=none, axis equal image
  ]
\addplot[blue,quiver={u=3,v=2},-stealth,thick] coordinates {(0,0)};
\node[below] at (axis cs:3,2) {$\vv_2$};
\addplot[blue,quiver={u=1,v=2},-stealth,thick] coordinates {(0,0)};
\node[right] at (axis cs:1,2) {$\vv_1$};
\addplot[black,mark=*] coordinates {(0,0)};
\end{axis}
\end{tikzpicture}}%
The set of two vectors shown to the right.

\begin{solution} 
Since the two vectors are not proportional to each other, we cannot write either as a scalar multiple of the other, and so the pair are linearly independent. 

\end{solution}
\end{figbox}


\item 
\begin{figbox}{\begin{tikzpicture}
\begin{axis}[footnotesize,font=\footnotesize
  ,axis lines=none, axis equal image
  ]
\addplot[blue,quiver={u=3,v=2},-stealth,thick] coordinates {(0,0)};
\node[below] at (axis cs:3,2) {$\vv_2$};
\addplot[blue,quiver={u=-1,v=-0.667},-stealth,thick] coordinates {(0,0)};
\node[above] at (axis cs:-1,-0.667) {$\vv_1$};
\addplot[black,mark=*] coordinates {(0,0)};
\end{axis}
\end{tikzpicture}}%
The set of two vectors shown to the right.

\begin{solution} 
Since the two vectors appear proportional to each other, \(\vv_2\approx (-3)\vv_1\)\,,  so the pair appear linearly dependent. 

\end{solution}
\end{figbox}



\item \(\{(1,3,0,-1),\ (1,0,-4,2),\ (-2,3,0,-3),\ (0,6,-4,-2)\}\)
\begin{solution} 
Notice that the last vector is the sum of the first three, \((0,6,-4,-2)=(1,3,0,-1)+(1,0,-4,2)+(-2,3,0,-3)\), and so the set is linearly dependent. 
\end{solution}


\end{enumerate}
\end{example}




Recall that \cref{thm:orthoevec} established that for every two distinct {eigenvalue}s of a symmetric matrix~\(A\), any corresponding two {eigenvector}s are {orthogonal}.
Consequently, for a symmetric matrix~\(A\), a set of eigenvectors from distinct eigenvalues forms an \idx{orthogonal set}.
The following \cref{thm:indepev} generalizes this property to non-symmetric matrices using the concept of linear independence.
That the corresponding eigenvalues are all different \text{is crucial.}


\begin{theorem} \label{thm:indepev}
For every \(n\times n\) matrix~\(A\), let \hlist\lambda m\ be \idx{distinct eigenvalues} of~\(A\) with corresponding \idx{eigenvector}s \hlist\vv m.
Then the set \(\{\hlist \vv m\}\) is \idx{linearly independent}.
\end{theorem}
\begin{proof} 
Use \idx{contradiction}.
Assume the set \(\{\hlist \vv m\}\) is linearly dependent.
Choose~\(k<m\) such that the set \(\{\hlist\vv k\}\) is linearly independent, whereas the set \(\{\hlist\vv{k+1}\}\) is linearly dependent.
Hence there exists non-trivial coefficients such that 
\begin{equation*}
\lincomb c\vv{k}+c_{k+1}\vv_{k+1}=\ov\,;
\end{equation*}
further, \(c_{k+1}\neq0\) as \(\{\hlist\vv k\}\) is linearly independent.
Pre-multiply the linear combination by matrix~\(A\):
\begin{eqnarray*}&&
\lincomb c{A\vv}k+c_{k+1}A\vv_{k+1}=A\ov
\\\implies&&
c_1\lambda_1\vv_1+c_2\lambda_2\vv_2+\cdots+c_k\lambda_k\vv_k+c_{k+1}\lambda_{k+1}\vv_{k+1}=\ov\,.
\end{eqnarray*}
Now subtract \(\lambda_{k+1}\times\) the original linear combination:
\begin{eqnarray*}&&
\phantom{{}+(}c_1\lambda_1\vv_1+c_2\lambda_2\vv_2+\cdots+c_k\lambda_k\vv_k+c_{k+1}\lambda_{k+1}\vv_{k+1}
\\&&{}
-\left(\lincomb c{\lambda_{k+1}\vv}{k}
+c_{k+1}\lambda_{k+1}\vv_{k+1}\right)=\ov
\\\implies&&
c_1(\lambda_1-\lambda_{k+1})\vv_1+c_2(\lambda_2-\lambda_{k+1})\vv_2+\cdots
\\&&{}
+c_k(\lambda_k-\lambda_{k+1})\vv_k
+c_{k+1}\underbrace{(\lambda_{k+1}-\lambda_{k+1})}_{=0}\vv_{k+1}=\ov
\\\implies&&
\lincomb{c'}\vv k=\ov
\end{eqnarray*}
for coefficients \(c'_j=c_j(\lambda_j-\lambda_{k+1})\).
Since all the eigenvalues are distinct, \(\lambda_j-\lambda_{k+1}\neq0\)\,, and since the coefficients~\(c_j\) are not all zero, hence \(c'_j\)~are not all zero.
Thus we have created a non-trivial linear combination of \hlist \vv k\ which is zero, and so the set \(\{\hlist\vv k\}\) is linearly dependent.
This contradiction of the choice of~\(k\) proves that the assumption must be wrong.
Hence the set \(\{\hlist \vv m\}\) is linearly independent, \text{as required.}
\end{proof}



\begin{activity}
The matrix \(\begin{bmatrix} 2&1\\a^2&2 \end{bmatrix}\) has \idx{eigenvector}s proportional to~\((1,a)\), and proportional to~\((1,-a)\).
For what values of~\(a\) does the matrix have two \idx{distinct eigenvalues}?
\actposs[4]{\(a\neq0\)}{\(a\neq-1\)}{\(a\neq1\)}{\(a\neq2\)}
\end{activity}






\begin{example} \label{eg:indepev}
For each of the following matrices, show that the \idx{eigenvector}s from \idx{distinct eigenvalues} form linearly independent sets.
\begin{enumerate}[ref=\ref{eg:indepev}(\alph*)]
\item  Consider the matrix \(B=\begin{bmat}-1&1&-2
\\-1&0&-1
\\0&-3&1 \end{bmat}\) from \cref{eg:faespm}.
\begin{solution} \sloppy
In \script, executing 
\begin{verbatim}
B=[-1 1 -2
 -1 0 -1
 0 -3 1]
[V,D]=eig(B)
\end{verbatim}
\setbox\ajrqrbox\hbox{\qrcode{% eigenvectors
B=[-1 1 -2
 -1 0 -1
 0 -3 1]
[V,D]=eig(B)
}}%
\marginajrbox%
gives eigenvectors and corresponding eigenvalues in
\begin{verbatim}
V =
    -0.5774     0.7071    -0.7071
    -0.5774     0.0000     0.0000
    -0.5774    -0.7071     0.7071
D =
         -2          0          0
          0          1          0
          0          0          1
\end{verbatim} 
Recognizing \(0.7071=1/\sqrt2\)\,, the last two eigenvectors, \((1/\sqrt2,0,-1/\sqrt2)\) and  \((-1/\sqrt2,0,1/\sqrt2)\), form a linearly dependent set because they are proportional to each other.
This linear dependence does not confound  \cref{thm:linhomo} because the corresponding eigenvalues are the same, not distinct, namely \(\lambda=1\)\,.
The theorem only applies to eigenvectors of \text{distinct eigenvalues.}

Here the two distinct eigenvalues are \(\lambda=-2\) and \(\lambda=1\)\,.
Recognizing \(0.5774=1/\sqrt3\)\,, two corresponding eigenvectors are \((-1/\sqrt3,-1/\sqrt3,-1/\sqrt3)\) and \((1/\sqrt2,0,-1/\sqrt2)\).
Because the zero component in the second corresponds to a nonzero component in the first, these cannot be proportional to each other, and so the pair form a linearly \text{independent set.}
\end{solution}


\item\label[example]{eg:indepeva} \cref{eg:fiveev} found the eigenvalues and eigenvectors of matrix
\begin{equation*}
A=\begin{bmatrix}0&3&0&0&0
\\1&0&3&0&0
\\0&1&0&3&0
\\0&0&1&0&3
\\0&0&0&1&0\end{bmatrix}
\end{equation*}
In \script\ execute
\begin{verbatim}
A=[0 3 0 0 0
 1 0 3 0 0
 0 1 0 3 0
 0 0 1 0 3
 0 0 0 1 0]
[V,D]=eig(A)
\end{verbatim}
\setbox\ajrqrbox\hbox{\qrcode{% eigenvectors
A=[0 3 0 0 0
 1 0 3 0 0
 0 1 0 3 0
 0 0 1 0 3
 0 0 0 1 0]
[V,D]=eig(A)
svd(V)
}}%
\marginajrbox%
to obtain the report \twodp
\begin{verbatim}
V =
   0.62  -0.62   0.94  -0.85  -0.85
   0.62   0.62  -0.00   0.49  -0.49
   0.42  -0.42  -0.31  -0.00   0.00
   0.21   0.21  -0.00  -0.16   0.16
   0.07  -0.07   0.10   0.09   0.09
D =
   3.00      0      0      0      0
      0  -3.00      0      0      0
      0      0  -0.00      0      0
      0      0      0  -1.73      0
      0      0      0      0   1.73
\end{verbatim}
The five eigenvalues are all distinct, so \cref{thm:linhomo} asserts that a set of corresponding eigenvectors is linearly independent.
The five columns of~\(V\), call them \hlist\vv5, are a set of corresponding eigenvectors.

To confirm their linear independence let's seek a linear combination being zero, that is, \(\lincomb c\vv5=\ov\)\,.
Written as a matrix-vector system we seek \(\cv=(\hlist c5)\) such that \(V\cv=\ov\)\,.
Because the five singular values of square matrix~\(V\) are all nonzero,\footnote{One could alternatively compute the determinant \(\texttt{det(V)}=0.09090\) and because it is nonzero \cref{thm:ftim3} asserts that the equation has only the solution \(\cv=\ov\)\,.} obtained from \verb|svd(V)|~as
\begin{verbatim}
ans =
   1.7703
   1.1268
   0.6542
   0.3625
   0.1922
\end{verbatim}
consequently \cref{thm:ftim2} asserts that \(V\cv=\ov\) has only the zero solution.
Hence, by \cref{def:lindep} the set of eigenvectors in the columns of~\(V\) is linearly independent. 

\end{enumerate}
\end{example}


This last case of the preceding \cref{eg:indepeva} connects the concept of linear in/dependence to the existence or otherwise of nonzero solutions to a homogeneous system of \idx{linear equation}s, \(V\cv=\ov\)\,.
So does \cref{eg:lindepb}.
The great utility of this connection is that we understand a lot about homogeneous \idx{system}s of linear equations.
The next \cref{thm:linhomo} establishes this connection \text{in general.}


\begin{theorem} \label{thm:linhomo} 
Let \hlist\vv m\ be vectors in~\(\RR^n\),
and let the \(n\times m\) matrix \(V=\begin{bmatrix} \vv_1&\vv_2&\cdots&\vv_m \end{bmatrix}\).  
Then the set \(\{\hlist\vv m\}\) is \idx{linearly dependent} if and only if the \idx{homogeneous} \idx{system} \(V\cv=\ov\) has a nonzero solution~\cv.
\end{theorem}
\begin{proof} 
Now \(\{\hlist\vv m\}\) is \idx{linearly dependent} if and only if there are scalars, not all zero, such that the equation \(\lincomb c\vv m=\ov\) holds (\cref{def:lindep}).
Let the vector \(\cv=(\hlist cm)\), then this equation is equivalent to the statement \(V\cv=\ov\)\,.
That is, if and only if \(V\cv=\ov\) has a \text{nonzero solution.}
\end{proof}





Recall \cref{thm:orthcomp}, that in~\(\RR^n\) there can be no more that \(n\)~vectors in an \idx{orthogonal set} of vectors.
The following theorem is the generalization: in~\(\RR^n\) there can be no more than \(n\)~vectors in a linearly independent set of vectors.

\begin{activity}
Which of the following sets of vectors are linearly dependent?
\def\temp#1{\begin{tikzpicture}
\begin{axis}[footnotesize,font=\footnotesize
  ,axis lines=none, axis equal image
  ,xmax=2,ymax=1,xmin=-0.55,ymin=-1.5
  ]
\addplot[blue,quiver={u=2,v=1},-stealth,thick,mark=*] coordinates {(0,0)};
\ifnum#1>1\addplot[blue,quiver={u=-0.4,v=1},-stealth,thick] coordinates {(0,0)};\fi
\ifnum#1>2\addplot[blue,quiver={u=-0.5,v=-1.5},-stealth,thick] coordinates {(0,0)};\fi
\end{axis}
\end{tikzpicture}}
\actposs{\temp3}{\temp1}{\temp2}{None of these sets.}
\end{activity}



\begin{theorem} \label{thm:mgtnli} 
Every  set of \(m\)~vectors in~\(\RR^n\) is \idx{linearly dependent} when the number of vectors \(m>n\)\,.
\end{theorem}
\begin{proof} 
Form the \(m\)~vectors \(\hlist\vv m\in\RR^n\)\ into the \(n\times m\) matrix \(V=\begin{bmatrix} \vv_1&\vv_2&\cdots&\vv_m \end{bmatrix}\).
Consider the homogeneous system \(V\cv=\ov\)\,: 
%it has nonzero solutions iff the nullspace has nonzero dimension (\cref{thm:homosubsp}).
%Recall (\cref{def:nullity}) that the dimension of the nullspace is the nullity, and the Rank \cref{thm:rank} gives
%\begin{eqnarray*}
%\nullity A&=&m-\rank A
%\quad(\text{as \(m\) is the number of columns})
%\\&\geq&m-n
%\quad(\text{as }\rank A\leq \min(m,n)=n\text{ here})
%\\&>&0\quad(\text{again as }m>n\text{ here}).
%\end{eqnarray*}
as \(m>n\)\,, \cref{thm:feweqns} (with the meaning of \(m\) and~\(n\) swapped) asserts that \(V\cv=\ov\) has infinitely many solutions.
Thus \(V\cv=\ov\) has nonzero solutions, so \cref{thm:linhomo} implies that the set of eigenvectors \(\{\hlist\vv m\}\) is \text{linearly dependent.}
\end{proof}



\begin{example} 
%Use \cref{thm:linhomo} and \cref{thm:mgtnli}
Determine if the following sets of vectors are linearly dependent or independent.  Give reasons.

\begin{enumerate}
\item \(\{(-1,-2)\clb (-1,4)\clb (0,5)\clb (2,3)\}\)
\begin{solution} 
As there are four vectors in~\(\RR^2\) so \cref{thm:mgtnli} asserts the set is linearly dependent.
\end{solution}

\item \(\{(-6,-4,-1,-2)\clb (2,0,1,-2)\clb (2,-1,-1,1)\}\)
\begin{solution} 
In \script\ form the matrix with these vectors as columns
\begin{verbatim}
V=[-6   2   2
  -4   0  -1
  -1   1  -1
  -2  -2   1]
svd(V)
\end{verbatim}
\setbox\ajrqrbox\hbox{\qrcode{% lin indep
V=[-6   2   2
  -4   0  -1
  -1   1  -1
  -2  -2   1]
svd(V)
}}%
\marginajrbox%
and find that the three singular values are all nonzero (namely \(7.7568\), \(2.7474\), and~\(2.2988\)).
Hence there are \emph{no} free variables when solving \(V\cv=\ov\) (\cref{pro:gensol}), and consequently there is only the unique solution \(\cv=\ov\)\,.
By \cref{thm:linhomo}, the set of vectors is \text{linearly independent.}
\end{solution}

\item \(\{(-1,-2,2,-1)\clb (1,3,1,-1)\clb (-2,-4,4,-2)\}\)
\begin{solution} 
By inspection, the third vector is twice the first.
Hence the linear combination \(2(-1,-2,2,-1)+0(1,3,1,-1)-(-2,-4,4,-2)=\ov\) and so the set of vectors is linearly dependent. 
\end{solution}


\item \(\{(3,3,-1,-1)\clb (0,-3,-1,-7)\clb (1,2,0,2)\}\)
\begin{solution} 
In \script\ form the matrix with these vectors as columns
\begin{verbatim}
V=[3   0   1
   3  -3   2
  -1  -1   0
  -1  -7   2]
svd(V)
\end{verbatim}
\setbox\ajrqrbox\hbox{\qrcode{% lin indep
V=[3   0   1
   3  -3   2
  -1  -1   0
  -1  -7   2]
svd(V)
}}%
\marginajrbox%
and find that the three singular values are \(8.1393\), \(4.6638\), and~\(0.0000\).
The zero singular value implies that there is a free variable when solving \(V\cv=\ov\) (\cref{pro:gensol}), and consequently that there are infinitely many nonzero~\cv\ that solve \(V\cv=\ov\).
By \cref{thm:linhomo}, the set of vectors is \text{linearly dependent.}
\end{solution}

\item \(\{(10,3,3,1)\clb ( 2,-3,0,-1)\clb ( 1,-1,2,-1)\clb ( 2,-1,-3,0)\clb (-2,0,2,-1)\}\)
\begin{solution} 
There are five vectors in~\(\RR^4\) so \cref{thm:mgtnli} asserts that the set is linearly dependent.
\end{solution}




\item \(\sloppy\{(-0.4,-1.8,-0.2, 0.7,-0.2)\clb (-1.1, 2.8, 2.7,-3.0,-2.6)\clb (-2.3,-2.3, 4.1, 3.4,-1.6)\clb (-2.6,-5.3,-3.3,-1.3,-4.1)\clb ( 1.4, 5.2,-6.9,-0.7, 0.6)\}\)
\begin{solution} 
In \script\ form the matrix~\(V\) with these vectors as columns
\begin{verbatim}
V=[-0.4 -1.1 -2.3 -2.6 1.4
  -1.8  2.8 -2.3 -5.3  5.2
  -0.2  2.7  4.1 -3.3 -6.9
   0.7 -3.0  3.4 -1.3 -0.7
  -0.2 -2.6 -1.6 -4.1  0.6]
svd(V)
\end{verbatim}
\setbox\ajrqrbox\hbox{\qrcode{% lin indep
V=[-0.4 -1.1 -2.3 -2.6 1.4
  -1.8  2.8 -2.3 -5.3  5.2
  -0.2  2.7  4.1 -3.3 -6.9
   0.7 -3.0  3.4 -1.3 -0.7
  -0.2 -2.6 -1.6 -4.1  0.6]
svd(V)
}}%
\marginajrbox%
and find that the five singular values are \(10.6978\), \(8.0250\), \(5.5920\), \(3.0277\), and~\(0.0024\).
As the singular values are all nonzero, the homogeneous system \(V\cv=\ov\) has the unique solution \(\cv=\ov\) (\cref{pro:gensol}), and hence the set of five vectors are \text{linearly independent.}

However, the answer depends upon the context.  
In the strict mathematical context the vectors are unequivocally linearly independent.
But in the context of practical problems, where errors in matrix entries are likely, there are `shades of grey'. 
Here, one of the singular values is quite small, namely~\(0.0024\).
If the context informs us that the entries in the matrix had errors of say~\(0.01\), then this singular value is effectively zero (\cref{sec:rle}).
In the context of such errors, this set of five vectors would be  \text{linearly dependent.}
\end{solution}


\end{enumerate}
%for i=1:9999, a=round(randn(3,4)*3);if min(svd(a))<1e-7,null(a'),break,end,end
\end{example}













\subsection{Form a basis for subspaces}
\index{basis|(}
\index{subspace|(}


Recall that \cref{sec:lcss,sec:sbd} defined \idx{subspace}s and the \idx{span}, namely that a subspace is 
a set of vectors \idx{closed} under addition and \idx{scalar multiplication}, and a span gives a subspace as all \idx{linear combination}s of a set of vectors.
Also, \cref{def:orthobasis} defined an ``\idx{orthonormal basis}'' for a subspace to be a set of orthonormal vectors that span the subspace.
This section generalizes the concept of an ``orthonormal basis'' by relaxing the requirement of orthonormality to result in the concept of \text{a ``basis''.}


\begin{definition} \label{def:basis} 
A \bfidx{basis} for a \idx{subspace}~\WW\ of~\(\RR^n\) is a set of  vectors such that the set both \idx{span}s~\WW\ and is \idx{linearly independent}.
\end{definition}

\begin{example} 
\begin{enumerate}
\item Recall that \cref{eg:lindepb,eg:2dnoc} showed that the two vectors \((2,1)\) and \((1,2)\) are linearly independent and span~\(\RR^2\). 
Hence the set \(\{(2,1)\clb (1,2)\}\) is a basis of~\(\RR^2\).

\item Recall that \cref{eg:lindepa} showed that the set \(\{(-1,1,0)\clb (1,-2,1)\clb (0,1,-1)\}\) is linearly dependent, so this set cannot be a basis.

However, remove one vector from the set, such as the middle one, and consider the set \(\{(-1,1,0)\clb (0,1,-1)\}\).
As the two vectors are not proportional to each other, this set is linearly independent (\cref{thm:lindeplc}).
Also, the plane \(x+y+z=0\) is a subspace, say~\WW.
The plane~\WW\ is characterized by \(y=-x-z\)\,.
So every vector in~\WW\ can be written as \((x,-x-z,z)=(x,-x,0)+(0,-z,z)=(-x)(-1,1,0)+(-z)(0,1,-1)\). 
That is, \(\Span\{(-1,1,0)\clb (0,1,-1)\}=\WW\).
Hence \(\{(-1,1,0)\clb (0,1,-1)\}\) is a basis for the plane~\WW.

\item 
\begin{figbox}{\qview{30}{34} {\begin{tikzpicture}
\begin{axis}[small,font=\footnotesize, view={\q}{30}
  ,xlabel={$x$}, ylabel={$y$}, zlabel={$z$},label shift={-2ex}
  ,domain=-2:2,axis equal image
  ] 
\addplot3[blue,samples=3,no marks]  ({2.1*x},{1.3*x},{-1.1*x});
\addplot3[blue,thick,quiver={u=2.1,v=1.3,w=-1.1},-stealth,mark=*] coordinates {(0,0,0)};
\node[below] at (axis cs:2.1,1.3,-1.1) {$(2.1,1.3,-1.1)$};
\end{axis}
\end{tikzpicture}}}%
Find a basis for the line given parametrically as \(x=2.1t\)\,, \(y=1.3t\), and \(z=-1.1t\) (shown to the right in stereo).

\begin{solution} 
The vectors in the line may be written as \(\xv=(x,y,z) =(2.1t,1.3t,-1.1t) =(2.1,1.3,-1.1)t\)\,.
Since the parameter~\(t\) may vary over all values, vectors in the line form \(\Span\{(2.1,1.3,-1.1)\}\). 
Since \(\{(2.1,1.3,-1.1)\}\) is a linearly independent set of vectors (\cref{eg:lindepc}), it thus forms a basis for the vectors in the \text{given line.}

\end{solution}
\end{figbox}



\item 
\begin{figbox}{\begin{tikzpicture}
\begin{axis}[footnotesize,font=\footnotesize
  ,axis lines=middle, axis equal image
  ,xlabel={$x$}, ylabel={$y$}
  ,thick,no marks,samples=3 ,domain=-1:1 ] 
\addplot  ({5.7*x-0.6},{6.8*x+2.4});
\end{axis}
\end{tikzpicture}}%
Find a basis for the line given parametrically as \(x=5.7t-0.6\) and \(y=6.8t+2.4\)\,.

\begin{solution} 
The vectors in the line may be written as \(\xv=(5.7t-0.6,6.8t+2.4)\)\,.
But this does not form a subspace as it does not include the zero vector~\ov\ (as illustrated to the right). 
%the \(x\)-component is only zero for some positive~\(t\) whereas the \(y\)-component is only zero for some negative~\(t\) so they are never zero for the same value of parameter~\(t\).
Since this line is not a subspace, it cannot have a basis.


\end{solution}
\end{figbox}




\item 
\begin{figbox}{\qview{30}{34} {\begin{tikzpicture}
\begin{axis}[footnotesize,font=\footnotesize, view={\q}{30}
  ,xlabel={$x$}, ylabel={$y$}, zlabel={$z$},label shift={-2ex}
  ,no marks ,domain=-2:2 ] 
\addplot3[surf,blue,samples=3,opacity=0.3]  {-3*x+2*y};
\addplot3[blue,thick,quiver={u=1,v=0,w=-3},-stealth] coordinates {(0,0,0)};
\node[below] at (axis cs:1,0,-3) {$(1,0,-3)$};
\addplot3[blue,thick,quiver={u=0,v=1,w=2},-stealth] coordinates {(0,0,0)};
\node[above] at (axis cs:0,1,2) {$(0,1,2)$};
\end{axis}
\end{tikzpicture}}}%
Find a basis for the plane \(3x-2y+z=0\)\,.

\begin{solution} 
Writing the equation of the plane as \(z=-3x+2y\) we then write the plane parametrically (\cref{sec:nvep}) as the vectors \(\xv=(x,y,-3x+2y) =(x,0,-3x)+(0,y,2y) =x(1,0,-3) +y(0,1,2)\).
Since \(x\) and~\(y\) may vary over all values, the plane is the subspace \(\Span\{(1,0,-3)\clb (0,1,2)\}\) (as illustrated above-right in stereo).
Since \((1,0,-3)\) and \((0,1,2)\) are not proportional to each other, they form a linearly independent set.
Hence \(\{(1,0,-3)\clb (0,1,2)\}\) is a basis for \text{the plane.}

\end{solution}
\end{figbox}





\item %keep this as the last case of this example
Prove that every \idx{orthonormal basis} of a subspace~\WW\ is also a basis of~\WW.
\begin{solution} 
\cref{thm:ortholi} establishes that every orthonormal basis is linearly independent.
By \cref{def:orthobasis}, an orthonormal basis of~\WW\ spans~\WW.
Hence an orthonormal basis of~\WW\ is also a basis of~\WW.
\end{solution}
\end{enumerate}
\end{example}




\begin{activity}
Which of the following sets of vectors forms a basis for~\(\RR^2\), but is not an \idx{orthonormal basis} for~\(\RR^2\)?
\def\temp#1{\begin{tikzpicture}
\begin{axis}[footnotesize,font=\footnotesize
  ,axis lines=none, axis equal image
  ,xmax=1,ymax=2,xmin=-2,ymin=-1.3
  ]
\addplot[blue,quiver={u=1,v=2},-stealth,thick,mark=*] coordinates {(0,0)};
\ifnum#1=2\addplot[blue,quiver={u=-2,v=1},-stealth,thick] coordinates {(0,0)};\fi
\ifnum#1>2\addplot[blue,quiver={u=0.7,v=-1.3},-stealth,thick] coordinates {(0,0)};\fi
\ifnum#1>3\addplot[blue,quiver={u=-2,v=-0.5},-stealth,thick] coordinates {(0,0)};\fi
\end{axis}
\end{tikzpicture}}
\actposs[4]{\temp3}{\temp1}{\temp2}{\temp4}
\end{activity}






Recall that \cref{thm:sameD} establishes that an \idx{orthonormal basis} of a given subspace always has the same number of vectors.
The following theorem establishes that the same is true for general bases.
The proof directly generalizes that for \cref{thm:sameD}.

\begin{theorem} \label{thm:sameDii} 
Every basis for a given \idx{subspace} has the same number of vectors.
\end{theorem}
\begin{proof} 
Let \(\cU=\{\hlist\uv r\}\) and \(\cV=\{\hlist\vv s\}\) be any two  bases for a subspace in~\(\RR^n\).
We prove the number of vectors \(r=s\) by \idx{contradiction}.
In the first case, assume \(r<s\)\,.
Since \cU\ is a basis for the subspace every vector in the set~\cV\ can be written as a linear combination of vectors in~\cU\ with some coefficients~\(a_{ij}\):
\begin{eqnarray*}
  &&\vv_1=\uv_1a_{11}+\uv_2a_{21}+\cdots+\uv_ra_{r1}\,,
\\&&\vv_2=\uv_1a_{12}+\uv_2a_{22}+\cdots+\uv_ra_{r2}\,,
\\&&\quad\vdots
\\&&\vv_s=\uv_1a_{1s}+\uv_2a_{2s}+\cdots+\uv_ra_{rs}\,.
\end{eqnarray*}
Write each of these, such as the first one, in the form
\begin{equation*}
\vv_1=\begin{bmatrix} \uv_1&\uv_2&\cdots&\uv_r \end{bmatrix}
\begin{bmatrix} a_{11}\\a_{21}\\\vdots\\a_{r1} \end{bmatrix}
=U\av_1\,,
\end{equation*}
where \(n\times r\) matrix \(U=\begin{bmatrix} \uv_1&\uv_2&\cdots&\uv_r \end{bmatrix}\).
Similarly for the other equations \(\vv_2=\cdots=U\av_2\) through to \(\vv_s=\cdots=U\av_s\)\,.
Then the \(n\times s\) matrix
\begin{equation*}
V=\begin{bmatrix} \vv_1&\vv_2&\cdots&\vv_s \end{bmatrix}
=\begin{bmatrix} U\av_1& U\av_2&\cdots&U\av_s\end{bmatrix}
=U\begin{bmatrix} \av_1& \av_2&\cdots&\av_s\end{bmatrix}
=UA
\end{equation*}
for the \(r\times s\) matrix~\(A=\begin{bmatrix} \av_1& \av_2&\cdots&\av_s \end{bmatrix}\).
By assumption \(r<s\) and so \cref{thm:feweqns} assures us that the homogeneous system \(A\xv=\ov\) has infinitely many solutions; choose any non-trivial solution \(\xv\neq\ov\)\,.
Consider 
\begin{eqnarray*}
V\xv&=& UA\xv\quad(\text{from above})
\\&=&U\ov\quad(\text{since }A\xv=\ov)
\\&=&\ov\,.
\end{eqnarray*}
The identity \(V\xv=\ov\) implies that there is a nontrivial linear combination of the columns \hlist\vv s\ of~\(V\) which gives zero, hence the set~\cV\ is linearly dependent (\cref{thm:linhomo}).  
But this is a contradiction, so we cannot have \(r<s\)\,.

Second, a corresponding argument establishes we cannot have \(s<r\)\,.
Hence \(s=r\)\,: all bases of a given subspace must have the same number of vectors.
\end{proof}



\begin{example} \label{eg:samedi}
Consider the plane \(x+y+z=0\) in~\(\RR^3\).
Each of the following is a basis for the plane:
\begin{itemize}
\item \(\{(-1,1,0)\clb (1,-2,1)\}\);
\item \(\{(1,-2,1)\clb (0,1,-1)\}\); 
\item \(\{(0,1,-1)\clb (-1,1,0)\}\).
\end{itemize}
The reasons are that all three vectors involved are in the plane, and that each pair is linearly independent (in each pair, one is not proportional to the other).

However, consider the set \(\{(-1,1,0)\clb (1,-2,1)\clb (0,1,-1)\}\).
Although each of the three vectors is in the plane \(x+y+z=0\), this set is not a basis because the set is not linearly independent (\cref{eg:lindepa}).
Each individual vector, say \((-1,1,0)\), cannot form a basis for the plane because the span of one vector, such as \(\Span\{(-1,1,0)\}\), is a line not the \text{whole plane.}

The orthonormal basis \(\big\{(1,0,-1)/\sqrt2\clb (1,-2,1)/\sqrt6\big\}\) is another basis for the plane \(x+y+z=0\); both vectors satisfy \(x+y+z=0\) and are orthogonal and so linearly independent (\cref{thm:ortholi}).  
All these bases possess two vectors.
\end{example}

That all bases for a given subspace, including orthonormal bases, have the same number of vectors (\cref{thm:sameDii}) leads to the following theorem about the dimensionality.

\begin{theorem} \label{thm:dimii} 
For every \idx{subspace}~\WW\ of~\(\RR^n\),  
the \bfidx{dimension} of~\WW, denoted~\(\dim\WW\), is the number of vectors in any \idx{basis} for~\WW. 
\end{theorem}

\begin{proof} 
Recall that \cref{def:dim} defined \(\dim\WW\) to be the number of vectors in any orthonormal basis for~\WW.
\cref{thm:ortholi} certifies that all orthonormal bases are also bases (\cref{def:basis}), so \cref{thm:sameDii} implies that every basis of~\WW\ has \(\dim\WW\) vectors.
\end{proof}




\begin{activity}
Which of the following sets forms a basis for a subspace of \idx{dimension} two?
\actposs[2]{\(\{(1,-2,1)\clb (1,0,-1)\}\)}
{\(\{(1,2)\}\)}
{\(\{(1,1,-2)\clb (2,2,-4)\}\)}
{\(\{(-1,0,2)\clb (0,0,1)\clb (-1,2,0)\}\)}
\end{activity}





\begin{procedure}[basis for a span] \label{pro:bfs}
Find a \idx{basis} for the \idx{subspace} \(\AA=\Span\{\hlist\av n\}\) given that $\{\hlist\av n\}$ is a set of $n$~vectors in~\(\RR^m\).
Recall that \cref{pro:ospan} underpins finding an \idx{orthonormal basis} by the following.
\begin{enumerate}
\item Form \(m\times n\) matrix $A:= \begin{bmatrix} \av_1 & \av_2& \cdots&\av_n \end{bmatrix}$. 
\item Factorize~\(A\) into its \svd, $A=\usv$\,, and let \(r=\rank A\) be the number of nonzero \idx{singular value}s (or effectively nonzero when the matrix has \idx{experimental error}s, \cref{sec:rle}).
\item The set \(\{\hlist\uv r\}\)  (where \(\uv_j\)~denotes the columns of~$U$) is a \idx{basis}, specifically an \idx{orthonormal basis}, for the \(r\)-\idx{dimension}al subspace~\AA.
\end{enumerate}
Alternatively, if the rank \(r=n\)\,, then the set \(\{\hlist\av n\}\) is \idx{linearly independent} and spans the subspace~\AA, and so is also a \idx{basis} for the \(n\)-dimensional subspace~\AA.
\end{procedure}



\begin{example} 
Apply \cref{pro:bfs} to find a basis for the following sets.
\begin{enumerate}
\item Recall that \cref{eg:samedi} identified that every pair of vectors in the set \(\{(-1,1,0)\clb (1,-2,1)\clb (0,1,-1)\}\) forms a basis for the plane that they span.  
Find another basis for the plane.
\begin{solution} 
In \script\ form the matrix with these vectors as columns:
\begin{verbatim}
A=[-1 1 0
   1 -2 1
   0 1 -1]
[U,S,V]=svd(A)
\end{verbatim}
\setbox\ajrqrbox\hbox{\qrcode{% basis
A=[-1 1 0
   1 -2 1
   0 1 -1]
[U,S,V]=svd(A)
}}%
\marginajrbox%
Then the \svd\  obtains \twodp
\begin{verbatim}
U =
  -0.41  -0.71   0.58
   0.82   0.00   0.58
  -0.41   0.71   0.58
S =
   3.00      0      0
      0   1.00      0
      0      0   0.00
V = ...
\end{verbatim}
The two nonzero singular values determine \(\rank A=2\) and hence the first two columns of~\verb|U| form an (orthonormal) basis for \(\Span\{(-1,1,0)\clb (1,-2,1)\clb (0,1,-1)\}\).
That is, \(\{0.41(-1,2,-1)\clb 0.71(-1,0,1)\}\) is a (orthonormal) basis for the two-dimensional plane.
\end{solution}

\item Find a basis for the span of the three vectors
\begin{equation*}
(-2,0,-4,1,1),\ 
(7,1,2,-1,-5),\ 
(-5,-1,2,3,-2).
\end{equation*}
\begin{solution} 
In \script\ it is often easiest to enter such vectors as rows, and then \idx{transpose} with the dash operator to form the matrix with the vectors as columns:
\begin{verbatim}
A=[ -2 0 -4 1 1
 7 1 2 -1 -5
 -5 -1 2 3 -2]'
[U,S,V]=svd(A)
\end{verbatim}
\setbox\ajrqrbox\hbox{\qrcode{% basis
A=[ -2 0 -4 1 1
 7 1 2 -1 -5
 -5 -1 2 3 -2]'
[U,S,V]=svd(A)
}}%
\marginajrbox%
Then the \svd\ obtains \twodp
\begin{verbatim}
U =
   0.86  -0.32  -0.02   0.40   0.02
   0.12  -0.11  -0.06  -0.40   0.90
   0.22   0.65   0.72   0.07   0.13
  -0.23   0.35  -0.38   0.73   0.37
  -0.38  -0.59   0.58   0.37   0.19
S =
   10.07       0       0
       0    5.87       0
       0       0    3.01
       0       0       0
       0       0       0
V = ...
\end{verbatim}
The three nonzero singular values determine  \(\rank A=3\) and so the original three vectors are linearly independent.
Consequently, the original three vectors form a basis for their span.

If you prefer an orthonormal basis, then use the first three columns of~\verb|U| as an orthonormal basis. 
\end{solution}



\begin{OmitV1}
\item Find a basis for the span of the four vectors
\((1,0,3,-4,0)\), 
\((-1,-1,1,4,2)\), 
\((-3,2,2,2,1)\), 
\((3,-3,2,-2,1)\).
\begin{solution} 
In \script, enter these vectors as rows, and then transpose with the dash operator to form the matrix with these as columns:
\begin{verbatim}
A=[1 0 3 -4 0
 -1 -1 1 4 2
 -3 2 2 2 1
 3 -3 2 -2 1]'
[U,S,V]=svd(A)
\end{verbatim}
\setbox\ajrqrbox\hbox{\qrcode{% basis
A=[1 0 3 -4 0
 -1 -1 1 4 2
 -3 2 2 2 1
 3 -3 2 -2 1]'
[U,S,V]=svd(A)
}}%
\marginajrbox%
Then the \svd\ obtains \twodp
\begin{verbatim}
U =
  -0.52   0.11   0.46  -0.71   0.02
   0.28  -0.44  -0.55  -0.61   0.24
  -0.19   0.66  -0.60  -0.17  -0.37
   0.78   0.32   0.36  -0.30  -0.27
   0.10   0.51  -0.00   0.03   0.86
S =
   7.64      0      0      0
      0   4.59      0      0
      0      0   4.30      0
      0      0      0   0.00
      0      0      0      0
V = ...
\end{verbatim}
The three nonzero singular values determine  \(\rank A=3\).
Consequently, the first three columns of~\verb|U| form an orthonormal basis for the span. 

Alternatively, you might notice that the sum of the first two vectors is the sum of the last two vectors.
Consequently, given the rank is three, we obtain three linearly independent vectors by omitting any one.
That is, any three of the given vectors form a basis for the span of \text{the four.}
\end{solution}
\end{OmitV1}

\end{enumerate}
\end{example}


The procedure is different if the subspace of interest is defined by a \idx{system} of equations instead of the span of some vectors.

\begin{example} \label{eg:bas2sys}
Find a basis for the solutions of the system in~\(\RR^3\) of \(3x+y=0\) and \(3x+2y+3z=0\)\,.
\begin{solution} 
By hand manipulation, the first equation gives \(y=-3x\)\,; which when substituted into the second gives \(3x-6x+3z=0\)\,, namely \(z=x\)\,.  
That is, all solutions are of the form \((x,-3x,x)\), namely \(\Span\{(1,-3,1)\}\).
Thus a basis for the subspace of solutions is~\(\{(1,-3,1)\}\).

Other possible bases are \(\{(-1,3,-1)\}\),  \(\{(2,-6,2)\}\), and so on: there are an infinite number of possible answers.
\end{solution}
\end{example}




\begin{example} 
Find a basis for the solutions of \(-2x-y+3z=0\) in~\(\RR^3\).
\begin{solution} 
Rearrange the equation so that one variable is a function of the others, say \(y=-2x+3z\).
Then the vector form of solutions are \((x,y,z)=(x,-2x+3z,z)=(1,-2,0)x+(0,3,1)z\) in terms of \idx{free variable}s~\(x\) and~\(z\).
Since \((1,-2,0)\) and~\((0,3,1)\) are not proportional to each other, they are linearly independent, and so a basis for the solutions is \(\{(1,-2,0)\clb (0,3,1)\}\).
(Infinitely many other bases are \text{possible answers.)}
\end{solution}
\end{example}



\begin{activity}
Which of the following is \emph{not} a basis for the line \(3x+7y=0\)?
\actposs[4]{\(\{(3,7)\}\)}
{\(\{(-7,3)\}\)}
{\(\{(-\frac73,1)\}\)}
{\(\{(1,-\frac37)\}\)}
\end{activity}





\begin{procedure}[basis from equations] \label{pro:bfe}
Suppose we seek a \idx{basis} for a \idx{subspace}~\WW\ specified as the solutions of a system of equations.
\begin{enumerate}
\item Rewrite the \idx{system} of equations as the \idx{homogeneous} system \(A\xv=\ov\)\,. 
Then the subspace~\WW\ is the \idx{nullspace} of \(m\times n\) matrix~\(A\).
\item  Adapting \cref{pro:gensol} for the specific case of homogeneous systems, first find an \svd\ \idx{factorization} \(A=\usv\) and let \(r=\rank A\) be the number of nonzero \idx{singular value}s (or effectively nonzero when the matrix has \idx{experimental error}s, \cref{sec:rle}).
\item Then \(\yv=(0,\ldots,0,y_{r+1},\ldots,y_n)\) is a \idx{general solution} of \(S\yv=\zv=\ov\)\,.
Consequently, all possible solutions \(\xv=V\yv\) are spanned by the last \(n-r\) columns of~\(V\), which thus form an \idx{orthonormal basis} for the subspace~\WW.
\end{enumerate}
\end{procedure}


\begin{example} 
Find a \idx{basis} for all solutions to each of the following systems of equations.
\begin{enumerate}
\item \(3x+y=0\) and \(3x+2y+3z=0\) from \cref{eg:bas2sys}.
\begin{solution} 
Form matrix \(A=\begin{bmat} 3&1&0\\3&2&3 \end{bmat}\) and compute an \svd\ with \verb|[U,S,V]=svd(A)| to obtain \twodp
\begin{verbatim}
U = ...
S =
   5.34      0      0
      0   1.86      0
V =
   0.77   0.56   0.30
   0.42  -0.09  -0.90
   0.48  -0.82   0.30
\end{verbatim}
The two nonzero singular values determine \(\rank A=2\)\,.
Hence the solutions of the system are spanned by the last one column of~\(V\).  
That is, a basis for the solutions is \(\{(0.3,-0.9,0.3)\}\).
\end{solution}

\item \(7x=6y+z+3\) and \(4x+9y+2z+2=0\)\,.
\begin{solution} 
This system is not homogeneous (due to the constant terms, \cref{def:homosys}), therefore \(\xv=\ov\) is not a solution. Consequently, the solutions of the system cannot form a subspace (\cref{def:subspace}). 
Thus the concept of a basis does not apply (\cref{def:basis}). 
\end{solution}


\item \(w+x=z\)\,,
\(3w=x+y+5z\)\,,
\(4x+y+2z=0\)\,.
\begin{solution} 
Rearrange to the matrix-vector system \(A\xv=\ov\) for vector \(\xv=(w,x,y,z)\in\RR^4\) and matrix
\begin{verbatim}
A=[1  1  0 -1
   3 -1 -1 -5
   0  4  1  2]
\end{verbatim}
\setbox\ajrqrbox\hbox{\qrcode{% basis
A=[1  1  0 -1
   3 -1 -1 -5
   0  4  1  2]
[U,S,V]=svd(A)
}}%
\marginajrbox%
Enter into \script\ as above and then find an \svd\ with \verb|[U,S,V]=svd(A)| to obtain \twodp
\begin{verbatim}
U = ...
S =
   6.77      0      0      0
      0   3.76      0      0
      0      0   0.00      0
V =
  -0.40   0.45   0.09   0.80
   0.41   0.86  -0.19  -0.25
   0.20   0.10   0.97  -0.07
   0.80  -0.24  -0.10   0.54
\end{verbatim}
There are two nonzero singular values, so \(\rank A=2\)\,.
There are thus \(4-2=2\) free variables in solving \(A\xv=\ov\) leading to a 2D subspace with orthonormal basis of the last two columns of~\(V\).
That is, an orthonormal basis for the subspace of all solutions is \(\{(0.09,-0.19,0.97,-0.10)\clb (0.80,-0.25,-0.07,0.54)\}\). 
\end{solution}

\end{enumerate}
%for i=1:9999,A=round(randn(3,4)*4); if min(svd(A))<1e-7, A=A, break, end, end
\end{example}







Recall that this \cref{sec:lisb} started by discussing the need to have a unique representation of solutions to problems.
If we do not have uniqueness, then the ambiguity in algebraic representation ruins basic algebra.
The forthcoming theorem assures us that the linear independence of a basis ensures the unique representation that we need.
In essence it says that every basis, whether orthogonal or not, can be used to form a \text{\idx{coordinate system}}.



\needlines6
\begin{wrapfigure}[6]r{0pt}
\begin{tikzpicture}
\newcommand{\ppoint}[3]{
    \pgfmathparse{#1*3+#2*1}\let\h\pgfmathresult
    \pgfmathparse{#1*1+#2*2}\let\v\pgfmathresult
    \addplot[blue,mark=*,thick,quiver={u=\h,v=\v},-stealth] coordinates {(0,0)};
    \edef\tempe{
    \noexpand\node[right] at (axis cs:\h,\v) {$#3$};
    }\tempe
    }
\begin{axis}[footnotesize,font=\footnotesize
  , axis lines=none%,xlabel={$x$},ylabel={$y$}
  , axis equal image
  , view={0}{90}
  ,xmax=8.5,ymax=5.5,xmin=-6.5,ymin=-5.5
  ]
\ppoint21{\noexpand\av}
\ppoint1{-3}{\noexpand\bv}
\ppoint{-2}3{\noexpand\cv}
\end{axis}
\end{tikzpicture}
\end{wrapfigure}
\begin{example}[a tale of two \idx{coordinate system}s] 
To the right are plotted three vectors and the origin.
Take the view that these are fixed physically meaningful vectors; the issue of this example is how do we code such vectors in mathematics?
\\\ \\\ 

\needlines8
\begin{wrapfigure}r{0pt}
\begin{tikzpicture}%
\newcommand{\ppoint}[3]{
    \pgfmathparse{#1*3+#2*1}\let\h\pgfmathresult
    \pgfmathparse{#1*1+#2*2}\let\v\pgfmathresult
    \addplot[blue,mark=*,only marks] coordinates {(\h,\v)};
    \edef\tempe{%
    \noexpand\node[left] at (axis cs:\h,\v) {$(\h,\v)$};
    \noexpand\node[right] at (axis cs:\h,\v) {$#3$};
    }\tempe
    }
\begin{axis}[small,font=\footnotesize
  , axis lines=middle,xlabel={$x$},ylabel={$y$},grid
  , axis equal image
  , view={0}{90}
  ,xmax=7.9,ymax=5.5,xmin=-7.4,ymin=-5.5
  ]
\ppoint21{\noexpand\av}
\ppoint1{-3}{\noexpand\bv}
\ppoint{-2}3{\noexpand\cv}
\ppoint00{}
\end{axis}
\end{tikzpicture}
\end{wrapfigure}
In the orthogonal \idx{standard coordinate system} these three vectors and the origin have coordinates as plotted on the right by their end-points.
Consequently, we write \(\av=(7,4)\), \(\bv=(0,-5)\), and \(\cv=(-3,4)\).
\\\ \\\ \\\ \\\

\needlines7
\begin{wrapfigure}r{0pt}
\begin{tikzpicture}%
%[baseline={([yshift={-\ht\strutbox}]current bounding box.north)}]
\newcommand{\ppoint}[3]{
    \pgfmathparse{#1*3+#2*1}\let\h\pgfmathresult
    \pgfmathparse{#1*1+#2*2}\let\v\pgfmathresult
%    \addplot[blue,mark=*,thick,quiver={u=\h,v=\v},-stealth] coordinates {(0,0)};
    \addplot[blue,mark=*,only marks] coordinates {(\h,\v)};
    \edef\tempe{%
    \noexpand\node[left] at (axis cs:\h,\v) {$(#1,#2)$};
%    \node[left] at (axis cs:\h,\v) {$(\h,\v)$};
    \noexpand\node[right] at (axis cs:\h,\v) {$#3$};
    }\tempe
    }
\begin{axis}[small,font=\footnotesize
%  , axis lines=middle,xlabel={$x$},ylabel={$y$},grid
  , axis lines=none
  , axis equal image
  , view={0}{90}
  ,xmax=7.9,ymax=5.5,xmin=-7.4,ymin=-5.5
  ]
\addplot3[mesh,red,samples=9,domain=-4:4,dotted] (3*x+y,x+2*y,0);
\addplot[red,quiver={u=3,v=1},-stealth,thick] coordinates {(0,0)};
\node[right] at (axis cs:3,1) {$\vv_1$};
\addplot[red,quiver={u=1,v=2},-stealth,thick] coordinates {(0,0)};
\node[above] at (axis cs:1,2) {$\vv_2$};
\ppoint21{\noexpand\av}
\ppoint1{-3}{\noexpand\bv}
\ppoint{-2}3{\noexpand\cv}
\ppoint00{}
\end{axis}
\end{tikzpicture}
\end{wrapfigure}
Now use the (red) basis \(\cB=\{\vv_1,\vv_2\}\) to form a non-orthogonal \idx{coordinate system} (represented by the dotted grid).
Then in this system the three vectors have coordinates \(\av=(2,1)\), \(\bv=(1,-3)\), and \(\cv=(-2,3)\).


But surely we cannot say both \(\av=(7,4)\) and \(\av=(2,1)\); it appears nonsense.
The reason for the different coordinate numbers representing the one vector~\av\ is that the underlying coordinate systems are different.
For example, we can say both \index{e@$\ev_j$}\(\av=7\ev_1+4\ev_2\) 
and \(\av=2\vv_1+\vv_2\) without any contradiction: these 
two statements explicitly recognize the underlying \idx{standard unit 
vector}s in the first expression and the underlying \idx{non-orthogonal basis} vectors in \text{the second.}

Consequently, mathematicians, scientists, and engineers invented a more precise notation.
We write \([\av]_\cB=(2,1)\) to represent that the coordinates of vector~\av\ in the basis~\cB\ are~\((2,1)\).
Correspondingly, letting \(\cE=\{\ev_1,\ev_2\}\) denote the basis of the \idx{standard unit vector}s, we write \([\av]_\cE=(7,4)\) to represent that the coordinates of vector~\av\ in the \idx{standard basis}~\cE\ are~\((7,4)\).
Similarly, \([\bv]_\cE=(0,-5)\) and \([\bv]_\cB=(1,-3)\);
and \([\cv]_\cE=(-3,4)\) and \([\cv]_\cB=(-2,3)\).

The endemic practice of just writing \(\av=(7,4)\), \(\bv=(0,-5)\) and \(\cv=(-3,4)\) is rationalized in this more precise notation by the convention that if no basis is explicitly specified, then we assume the \idx{standard basis}~\cE.%
\footnote{Given that the values of the components in a vector change with the coordinate basis, some of you may wonder whether the same thing happens for matrices.  
The answer is yes:  for a given \idx{linear transformation} (\cref{sec:ilt}), the values of the \idx{components} in its matrix also depend upon the underlying coordinate basis. }
\end{example}



\begin{theorem} \label{thm:ssbc} 
For any given \idx{subspace}~\WW\ of~\(\RR^n\) let \(\cB=\{\hlist\vv k\}\) be a \idx{basis} for~\WW.  
Then there is exactly one way to write each and every vector \(\wv\in\WW\) as a \idx{linear combination} of the \idx{basis} vectors: \(\wv=\lincomb c\vv k\)\,.
The coefficients \(\hlist ck\) are called the \textbf{\bfidx{coordinates} of~\wv\ with respect to~\cB}, and the \idx{column vector} \([\wv]_{\cB}=(\hlist ck)\) is called the \textbf{\bfidx{coordinate vector} of~\wv\ with respect to~\cB}.
\end{theorem}
\begin{proof}  
Consider any vector \(\wv\in\WW\).
Since \(\{\hlist\vv k\}\) is a basis for the subspace~\WW, \wv~can be written as a linear combination of the basis vectors. 
Let two such linear combinations be
\begin{eqnarray*}&&
\wv=\lincomb c\vv k\,,
\\&&
\wv=\lincomb d\vv k\,.
\end{eqnarray*}
Subtract the second of these equations from the first, grouping common vectors:
\begin{equation*}
\ov=(c_1-d_1)\vv_1+(c_2-d_2)\vv_2+\cdots+(c_k-d_k)\vv_k\,.
\end{equation*}
Since \(\{\hlist\vv k\}\) is linearly independent, this equation implies all the coefficients in parentheses are zero:
\begin{equation*}
0=(c_1-d_1)=(c_2-d_2)=\cdots=(c_k-d_k).
\end{equation*}
That is, \(c_1=d_1\)\,, \(c_2=d_2\)\,, \ldots, \(c_k=d_k\)\,, and the two linear combinations are identical.
Consequently, there is exactly one way to write a vector \(\wv\in\WW\) as a {linear combination} of the {basis} vectors.
\end{proof}


\begin{example}  
\newcommand{\ppoint}[3]{%
    \pgfmathparse{#1*2-#2*1}\let\h\pgfmathresult
    \pgfmathparse{-#1*1+#2*2}\let\v\pgfmathresult
    \addplot[blue,mark=*,thick,quiver={u=\h,v=\v},-stealth] coordinates {(0,0)};
    \edef\tempe{%
    \noexpand\node[right] at (axis cs:\h,\v) {$\noexpand\vec #3$};
    }\tempe
    }%
\newcommand{\cbase}[1]{\begin{tikzpicture}%
[baseline={([yshift={-\ht\strutbox}]current bounding box.north)}]
\begin{axis}[small,font=\footnotesize
  , axis lines=none, axis equal image
  , view={0}{90}
  ,xmax=7.9,ymax=5.5,xmin=-7.1,ymin=-5.5
  ]
\ifcase#1
\or\addplot3[mesh,red,samples=9,domain=-4:4,dotted] (2*x-y,-x+2*y,0);
\fi
\addplot[red,quiver={u=2,v=-1},-stealth,thick] coordinates {(0,0)};
\node[right] at (axis cs:2,-1) {$\vv_1$};
\addplot[red,quiver={u=-1,v=2},-stealth,thick] coordinates {(0,0)};
\node[above] at (axis cs:-1,2) {$\vv_2$};
\ppoint32{a}
\ppoint1{-2}{b}
\ppoint{-2}{0.5}{c}
\ppoint{-1}{-1}{d}
\end{axis}
\end{tikzpicture}}%
\begin{enumerate}
\item 
\begin{figbox}{\cbase0}%
Consider the diagram to the right of six labelled vectors.
Estimate the coordinates of the four shown vectors~\av, \bv, \cv, and~\dv\ in the shown basis \(\cB=\{\vv_1,\vv_2\}\).
\vspace{3\baselineskip}
\end{figbox}

\begin{solution}\\
\begin{figbox}{\cbase1}%
Draw in a grid corresponding to multiples of~\(\vv_1\) and~\(\vv_2\) in both directions, and parallel to~\(\vv_1\) and~\(\vv_2\),  as shown to the right.
Then from the grid, estimate that \(\av\approx3\vv_1+2\vv_2\) hence the coordinates \([\av]_\cB\approx(3,2)\).
Similarly, \(\bv\approx \vv_1-2\vv_2\) hence the coordinates \([\bv]_\cB\approx(1,-2)\).
Also, \(\cv\approx -2\vv_1+0.5\vv_2\) hence the coordinates \([\cv]_\cB\approx(-2,0.5)\).
And lastly,  \(\dv\approx -\vv_1-\vv_2\) hence the coordinates \([\dv]_\cB\approx(-1,-1)\).
\aqed

\end{figbox}
\end{solution}




\renewcommand{\ppoint}[3]{%
    \pgfmathparse{#1*1-#2*2}\let\h\pgfmathresult
    \pgfmathparse{-#1*2-#2*2}\let\v\pgfmathresult
    \addplot[blue,mark=*,thick,quiver={u=\h,v=\v},-stealth] coordinates {(0,0)};
    \edef\tempe{%
    \noexpand\node[right] at (axis cs:\h,\v) {$\noexpand\vec #3$};
    }\tempe
    }%
\renewcommand{\cbase}[1]{\begin{tikzpicture}%
[baseline={([yshift={-\ht\strutbox}]current bounding box.north)}]
\begin{axis}[small,font=\footnotesize
  , axis lines=none, axis equal image
  , view={0}{90}
  ,xmax=7.9,ymax=5.5,xmin=-7.1,ymin=-5.5
  ]
\ifcase#1
\or\addplot3[mesh,red,samples=9,domain=-4:4,dotted] (1*x-2*y,-2*x-2*y,0);
\fi
\addplot[red,quiver={u=1,v=-2},-stealth,thick] coordinates {(0,0)};
\node[below] at (axis cs:1,-2) {$\wv_1$};
\addplot[red,quiver={u=-2,v=-2},-stealth,thick] coordinates {(0,0)};
\node[below] at (axis cs:-2,-2) {$\wv_2$};
\ppoint1{-1.5}{a}
\ppoint3{-0.5}{b}
\ppoint{-2.5}{1}{c}
\ppoint{0}{0.5}{d}
\end{axis}
\end{tikzpicture}}%
\item 
\begin{figbox}{\cbase0}%
Consider the same four vectors but with a pair of different basis vectors:  let's see that although the vectors are the same, the coordinates in the different basis are different.
Estimate the coordinates of the four shown vectors~\av, \bv, \cv, and~\dv\ in the shown basis \(\cW=\{\wv_1,\wv_2\}\).
\vspace{2\baselineskip}
\end{figbox}

\begin{solution}\\
\begin{figbox}{\cbase1}%
Draw in a grid corresponding to multiples of~\(\wv_1\) and~\(\wv_2\) in both directions, and parallel to~\(\wv_1\) and~\(\wv_2\),  as shown to the right.
Then from the grid, estimate that \(\av\approx\wv_1-1.5\wv_2\) hence the coordinates \([\av]_\cW\approx(1,-1.5)\).
Similarly, \(\bv\approx 3\wv_1-0.5\wv_2\) hence the coordinates \([\bv]_\cW\approx(3,-0.5)\).
Also, \(\cv\approx -2.5\wv_1+\wv_2\) hence the coordinates \([\cv]_\cW\approx(-2.5,1)\).
And lastly,  \(\dv\approx 0.5\wv_2\) hence the coordinates \([\dv]_\cW\approx(0,0.5)\).
\aqed

\end{figbox}
\end{solution}

\end{enumerate}
\end{example}


\begingroup
\newcommand{\ppoint}[3]{%
    \pgfmathparse{#1*1-#2*1}\let\h\pgfmathresult
    \pgfmathparse{-#1*1+#2*3}\let\v\pgfmathresult
    \addplot[blue,mark=*,thick,quiver={u=\h,v=\v},-stealth] coordinates {(0,0)};
    \edef\tempe{%
    \noexpand\node[right] at (axis cs:\h,\v) {$\noexpand\vec #3$};
    }\tempe
    }%
\def\temp{\begin{tikzpicture}
\begin{axis}[footnotesize,font=\footnotesize
  , axis lines=none, axis equal image
  ,xmax=1.9,ymax=3.9,xmin=-1.9,ymin=-1.9
  ]
\addplot[red,quiver={u=1,v=-1},-stealth,thick] coordinates {(0,0)};
\node[right] at (axis cs:1,-1) {$\bv_1$};
\addplot[red,quiver={u=-1,v=3},-stealth,thick] coordinates {(0,0)};
\node[above] at (axis cs:-1,3) {$\bv_2$};
\ppoint32{x}
\end{axis}
\end{tikzpicture}}
\begin{activity}[\temp]
For the vector~\xv\ shown to the right, estimate the coordinates of~\xv\ in the shown basis \(\cB=\{\bv_1,\bv_2\}\).
\actposs{\([\xv]_\cB=(3,2)\)}
{\([\xv]_\cB=(1,4)\)}
{\([\xv]_\cB=(2,3)\)}
{\([\xv]_\cB=(4,1)\)}
\end{activity}
\endgroup



\begin{example} 
Let the basis \(\cB=\{\vv_1,\vv_2,\vv_3\}\) for the three given vectors \(\vv_1=(-1,1,-1)\), \(\vv_2=(1,-2,0)\), and \(\vv_3=(0,4,5)\) (each of these is specified in the \idx{standard basis}~\cE\ of the \idx{standard unit vector}s~\(\ev_1\), \(\ev_2\), and~\(\ev_3\)).
\begin{enumerate}
\item What is the vector with coordinates \([\av]_\cB=(3,-2,1)\)?
\begin{solution} 
Since a coordinate system is not specified for the answer, we answer with the default of the standard basis~\cE.
The vector \(\av=3\vv_1-2\vv_2+\vv_3\) which has standard coordinates 
\([\av]_\cE=3(-1,1,-1)-2(1,-2,0)+(0,4,5)=(-5,11,2)\).
\end{solution}

\item What is the vector with coordinates \([\bv]_\cB=(-1,1,1)\)?
\begin{solution} 
The vector \(\bv=-\vv_1+\vv_2+\vv_3\) which has standard coordinates (the default coordinates as the question does not specify) 
\([\bv]_\cE=-(-1,1,-1)+(1,-2,0)+(0,4,5)=(2,1,6)\).
\end{solution}

\item What are the coordinates in the basis~\cB\ of the vector~\cv\ where \([\cv]_\cE=(-1,3,3)\) in the \idx{standard basis}~\cE?
\begin{solution} 
We seek coordinate values \(c_1,c_2,c_3\) such that \(\cv=c_1\vv_1+c_2\vv_2+c_3\vv_3\)\,. 
Expressed in the standard basis~\cE\ this is the vector equation
\begin{equation*}
\begin{bmatrix} -1\\3\\3 \end{bmatrix}=
\begin{bmatrix} -1\\1\\-1 \end{bmatrix}c_1+
\begin{bmatrix} 1\\-2\\0 \end{bmatrix}c_2+
\begin{bmatrix} 0\\4\\5 \end{bmatrix}c_3\,.
\end{equation*}
A small system like this we solve by hand (\cref{sec:amss}): write in component form as
\begin{eqnarray*}
&&\left\{\begin{array}{r@{\ }r@{\ }r@{{}={}}l}  
-c_1&+c_2&&-1 \\
c_1&-2c_2&+4c_3&3 \\
-c_1&&+5c_3&3
\end{array}\right.
\\&&(\text{add 1st row to 2nd and take from 3rd})
\\&&\left\{\begin{array}{r@{\ }r@{\ }r@{{}={}}l}  
-c_1&+c_2&&-1 \\
&-c_2&+4c_3&2 \\
&-c_2&+5c_3&4
\end{array}\right.
\\&&(\text{subtract 2nd row from 3rd})
\\&&\left\{\begin{array}{r@{\ }r@{\ }r@{{}={}}l}  
-c_1&+c_2&&-1 \\
&-c_2&+4c_3&2 \\
&&c_3&2
\end{array}\right.
\end{eqnarray*}
Solving this triangular system gives \(c_3=2\)\,, \(c_2=4c_3-2=6\)\,, and \(c_1=c_2+1=7\)\,.
Thus the coordinates \([\cv]_\cB=(7,6,2)\) in the basis~\cB.
\end{solution}


\item What are the coordinates in the basis~\cB\ of the vector~\dv\ where \([\dv]_\cE=(-3,2,0)\) in the \idx{standard basis}~\cE?
\begin{solution} 
We seek coordinate values \(d_1,d_2,d_3\) such that \(\dv=d_1\vv_1+d_2\vv_2+d_3\vv_3\)\,. 
Expressed in the standard basis~\cE\ this equation is
\begin{equation*}
\begin{bmatrix} -3\\2\\0 \end{bmatrix}=
\begin{bmatrix} -1\\1\\-1 \end{bmatrix}d_1+
\begin{bmatrix} 1\\-2\\0 \end{bmatrix}d_2+
\begin{bmatrix} 0\\4\\5 \end{bmatrix}d_3\,.
\end{equation*}
To solve this system in \script\ (\cref{pro:unisol}), enter the matrix (easiest by transposing rows of~\(\vv_1\), \(\vv_2\), and~\(\vv_3\)) and the standard coordinates of~\dv:
\begin{verbatim}
A=[-1 1 -1
 1 -2 0
 0 4 5]'
d=[-3;2;0]
\end{verbatim}
\setbox\ajrqrbox\hbox{\qrcode{% new basis
A=[-1 1 -1
 1 -2 0
 0 4 5]'
d=[-3;2;0]
rcond(A)
dB=A\slosh d
}}%
\marginajrbox%
Then compute the coordinates~\([\dv]_\cB\) with \verb|dB=A\d| to determine \([\dv]_\cB=(20,17,4)\).

But remember, before using~\verb|A\| always check \verb|rcond(A)|, which here is the poor~\(0.0053\) (\cref{pro:unisol}).
Interestingly, such a poor small value of \verb|rcond| indicates that although the basis vectors in~\cB\ are linearly independent, they are `only just' linearly independent.
A small change, or error, might make the basis vectors linearly dependent and thus \cB~would be ruined as a basis for~\(\RR^3\).
The poor \verb|rcond| indicates that~\cB\ is a poor basis in \text{practical application.}
\end{solution}
\end{enumerate}
\end{example}





\begin{activity}
What are the coordinates in the \idx{basis}~\(\cB=\{(1,1)\clb (1,-1)\}\)\ of the vector~\dv\ where \([\dv]_\cE=(2,-4)\) in the \idx{standard basis}~\cE?
\actposs[4]{\([\dv]_\cB=(-1,3)\)}
{\([\dv]_\cB=(-2,6)\)}
{\([\dv]_\cB=(1,3)\)}
{\([\dv]_\cB=(3,-1)\)}
\end{activity}





\begin{example} 
You are given a \idx{basis} \(\cW=\{\wv_1,\wv_2,\wv_3\}\) for a 3D subspace~\WW\ of~\(\RR^5\), where the three basis vectors are
\(\wv_1=(1,3,-4,-3,3)\),
\(\wv_2=(-4,1,-2,-4,1)\), and
\(\wv_3=(-1,1,0,2,-3)\) (in the \idx{standard basis}~\cE).
\begin{enumerate}
\item What are the coordinates in the standard basis of the vector \(\av=2\wv_1+3\wv_2+\wv_3\)?
\begin{solution} 
In the standard basis
\begin{equation*}
[\av]_\cE=2\begin{bmatrix} 1\\3\\-4\\-3\\3 \end{bmatrix}
+3\begin{bmatrix} -4\\1\\-2\\-4\\1 \end{bmatrix}
+\begin{bmatrix} -1\\1\\0\\2\\-3 \end{bmatrix}
=\begin{bmatrix} -11\\ 10\\-14\\-16\\  6 \end{bmatrix}.
\end{equation*}
\end{solution}

\item What are the coordinates in the basis~\cW\ of the vector \(\sloppy\bv=(-1,2,-6,-11,10)\) (in the standard coordinates~\cE).
\begin{solution} 
We need to find coefficients \(c_1,c_2,c_3\) such that \(\bv=c_1\wv_1+c_2\wv_2+c_3\wv_3\)\,.
This forms the set of linear equations
\begin{equation*}
\begin{bmatrix}-1\\2\\-6\\-11\\10 \end{bmatrix}
=\begin{bmatrix} 1\\3\\-4\\-3\\3 \end{bmatrix}c_1
+\begin{bmatrix} -4\\1\\-2\\-4\\1 \end{bmatrix}c_2
+\begin{bmatrix} -1\\1\\0\\2\\-3 \end{bmatrix}c_3\,.
\end{equation*}
Form as the matrix-vector system
\begin{equation*}
\begin{bmatrix}1&-4&-1
\\ 3&1&1
\\-4&-2&0
\\-3&-4&2
\\ 3&1&-3\end{bmatrix}
\begin{bmatrix} c_1\\c_2\\c_3 \end{bmatrix}
=\begin{bmatrix}-1\\2\\-6\\-11\\10 \end{bmatrix},
\end{equation*}
and perhaps solve with \script.
Since there are more equations than unknowns, we should use an \svd\ in order to check that the system is consistent, namely, to check that \(\bv\in\WW\).
\begin{enumerate}
\item Code the matrix and the vector:
\begin{verbatim}
W=[1 -4 -1
 3 1 1
 -4 -2 0
 -3 -4 2
 3 1 -3]
b=[-1;2;-6;-11;10]
\end{verbatim}
\setbox\ajrqrbox\hbox{\qrcode{% new basis
W=[1 -4 -1
 3 1 1
 -4 -2 0
 -3 -4 2
 3 1 -3]
b=[-1;2;-6;-11;10]
[U,S,V]=svd(W)
z=U'*b
y=z(1:3)./diag(S)
bW=V*y
}}%
\marginajrbox%
\item Then obtain an \svd\ with \verb|[U,S,V]=svd(W)| \twodp
\begin{verbatim}
U =
   0.18  -0.88   0.04  -0.24  -0.37
  -0.32  -0.09   0.65   0.62  -0.28
   0.51   0.11  -0.46   0.56  -0.44
   0.64  -0.18   0.34   0.22   0.63
  -0.44  -0.42  -0.49   0.44   0.45
S =
   8.18      0      0
      0   4.52      0
      0      0   3.11
      0      0      0
      0      0      0
V =
  -0.74  -0.51   0.43
  -0.62   0.77  -0.15
   0.26   0.38   0.89
\end{verbatim}
The three nonzero singular values establish that the three vectors in the basis~\cW\ are indeed linearly independent (and since no singular value is small, then the vectors are robustly linearly independent).

\item Find \(\zv=\tr U\bv\) with \verb|z=U'*b| to get
\begin{verbatim}
z =
  -15.3413
   -2.2284
   -4.6562
   -0.0000
    0.0000
\end{verbatim}
The last two values of~\zv\ being zero confirm the system of equations is consistent, and so vector~\(\bv\) is in the range of~\cW, that is, \bv~is in the subspace~\WW.

\item Find \(y_j=z_j/\sigma_j\) with \verb|y=z(1:3)./diag(S)| to get
\begin{verbatim}
y =
  -1.8761
  -0.4929
  -1.4958
\end{verbatim}

\item Lastly, find the coefficients \([\bv]_\cW=V\yv\) with \verb|bW=V*y| to get
\begin{verbatim}
bW =
   1.00000
   1.00000
  -2.00000
\end{verbatim}
\end{enumerate}
That is, \([\bv]_\cW=(1,1,-2)\).

That \([\bv]_\cW\) has three components and \([\bv]_\cE\) has five components is not a contradiction.  
The difference in components occurs because the subspace~\WW\ is 3D but lies in~\(\RR^5\).  
Using the basis~\cW\ implicitly builds in the information that the vector~\bv\ is in a lower dimensional space, and so needs \text{fewer components.}
\end{solution}

\end{enumerate}
\end{example}


\index{subspace|)}



\begin{comment} 
Optionally: Gram--Schmidt (probably not here as would confuse the aims of lin indep, next section); and something on change of basis, recalling quadratics.
\end{comment}








\subsubsection{Revisit unique solutions}
Lastly, with all these extra concepts of determinants, eigenvalues,  linear independence, and a basis, we now revisit the issue of when there is a unique solution to a set of \idx{linear equation}s.


\begin{theorem}[Unique Solutions: version~3] \label{thm:ftim3} 
For every \(n\times n\) \idx{square matrix}~\(A\), and  
extending \cref{thm:ftim1,thm:ftim2}, the following statements are equivalent:
\begin{enumerate}[ref=\ref{thm:ftim3}(\alph*)]
\item\label[theorem]{thm:ftim3i} \(A\) is \idx{invertible};
\item\label[theorem]{thm:ftim3ii} \(A\xv=\bv\) has a \bfidx{unique solution} for every \(\bv\in\RR^n\);
\item\label[theorem]{thm:ftim3iii} \index{homogeneous}\(A\xv=\ov\) has only the zero solution;
\item\label[theorem]{thm:ftim3iv} all \(n\)~\idx{singular value}s of~\(A\) are nonzero;
\item\label[theorem]{thm:ftim3ivx} the \idx{condition number} of~\(A\) is finite (\(\verb|rcond|>0\));
\item\label[theorem]{thm:ftim3v} \index{rank}\(\rank A=n\)\,;
\item\label[theorem]{thm:ftim3vi} \index{nullity}\(\nullity A=0\)\,;
\item\label[theorem]{thm:ftim3vii} the \idx{column vector}s of~\(A\) span~\(\RR^n\);
\item\label[theorem]{thm:ftim3viii} the \idx{row vector}s of~\(A\) span~\(\RR^n\).
\item\label[theorem]{thm:ftim3ix} \(\det A\neq 0\)\,;
\item\label[theorem]{thm:ftim3x} \(0\) is not an \idx{eigenvalue} of~\(A\);
\item\label[theorem]{thm:ftim3xi} the \(n\) column vectors of~\(A\) are linearly independent;
\item\label[theorem]{thm:ftim3xii} the \(n\) \idx{row vector}s of~\(A\) are linearly independent.
\end{enumerate}
\end{theorem}
\begin{proof} 
Recall \cref{thm:ftim1,thm:ftim2} establish the equivalence of \ref{thm:ftim3i}--\ref{thm:ftim3viii}, and \cref{thm:detinv} proves the equivalence of \ref{thm:ftim3i} and~\ref{thm:ftim3ix}.
To establish that Property~\ref{thm:ftim3x} is equivalent to~\ref{thm:ftim3ix}, recall that \cref{thm:charpolyc} proved that \(\det A\) equals the product of the eigenvalues of~\(A\). 
Hence  \(\det A\) is not zero if and only if all the eigenvalues \text{are nonzero.}

Lastly, Property~\ref{thm:ftim3vii} says the \(n\)~column vectors span~\(\RR^n\), so they must be a basis, and hence linearly independent.
Conversely, if the \(n\)~columns of~\(A\) are linearly independent then they must span~\(\RR^n\).  
Hence Property~\ref{thm:ftim3xi} is equivalent to~\ref{thm:ftim3vii}.
Similarly  for Property~\ref{thm:ftim3viii} and the row vectors of~\ref{thm:ftim3xii}.
\end{proof}


\index{basis|)}







\sectionExercises



\begin{exercise}  
By inspection or basic arguments, decide whether the following sets of vectors are linearly dependent, or linearly independent.  Give reasons.
%\begin{verbatim}
%m=ceil(0.5+3*rand), n=ceil(1+2.5*rand), A=0+round(randn(m,n)*2), svd(A)
%\end{verbatim}

\begin{enumerate}
\item \(\{(-2,3,3)\clb (-1,2,-1)\}\)
\answer{lin.\ indep.}

\item \(\{(0,2)\clb (2,-2)\clb (0,-1)\}\)
\answer{lin.\ dep.}

\item \(\{(-3,0,-3)\clb (3,2,-2)\}\)
\answer{lin.\ indep.}

\item \(\{(0,2,2)\}\)
\answer{lin.\ indep.}

\begin{OmitV1}
\item \(\{(-2,0,-1)\clb (0,-2,2)\clb (2,0,1)\}\)
\answer{lin.\ dep.}

\item \(\{(0,3,-2)\clb (-2,-1,1)\clb (1,-2,-4)\}\)
\answer{lin.\ indep.}

\item \(\{(-2,1),(0,0)\}\)
\answer{lin.\ dep.}
\end{OmitV1}

\item \(\{(-2,-1,1)\clb (-2,-2,2)\clb (2,-1,-1)\clb (2,-2,-2)\}\)
\answer{lin.\ dep.}

\item \(\{(2,4)\clb (1,2)\}\)
\answer{lin.\ dep.}

\item \(\{( 1,-2,3)\clb (1,1,2)\}\)
\answer{lin.\ indep.}

\end{enumerate}
\end{exercise}







\begin{exercise}  
Compute an \svd\ to decide whether the following sets of vectors are linearly dependent, or linearly independent.  Give reasons.
%\begin{verbatim}
%m=ceil(1+3*rand), n=floor(m+2.5*rand), for j=1:(m*n)^2,A=0+round(randn(m,n)*3); sv=svd(A); if min(sv)<1e-7,break,end,end, A=A, sv=sv
%\end{verbatim}
\begin{Parts}
\item \(\begin{bmatrix} 5\\0\\-1 \end{bmatrix}, \begin{bmatrix}-2\\5\\-1 \end{bmatrix}\)
\answer{lin.\ indep.}

\item \(\begin{bmatrix} -2\\-1\\1 \end{bmatrix}, \begin{bmatrix} 3\\-4\\-1 \end{bmatrix}, \begin{bmatrix} 5\\-3\\-2 \end{bmatrix}\)
\answer{lin.\ dep.}

\item \(\begin{bmatrix} 2\\-4\\1\\6 \end{bmatrix}, \begin{bmatrix} -1\\1\\10\\-5 \end{bmatrix}, \begin{bmatrix}1\\6\\-2\\4 \end{bmatrix}\)
\answer{lin.\ indep.}

\item \(\begin{bmatrix} 1\\0\\-5\\2 \end{bmatrix}, \begin{bmatrix}-1\\0\\-2\\4 \end{bmatrix}, \begin{bmatrix}-1\\3\\1\\1 \end{bmatrix}, \begin{bmatrix} 3\\4\\-4\\-4 \end{bmatrix}\)
\answer{lin.\ dep.}

\begin{OmitV1}
\item \(\begin{bmatrix} 0\\-2\\2\\0 \end{bmatrix}, \begin{bmatrix}-1\\-4\\-6\\-1 \end{bmatrix}, \begin{bmatrix} 4\\3\\-5\\0 \end{bmatrix}\)
\answer{lin.\ indep.}

\item \(\begin{bmatrix} 0\\3\\2\\0 \end{bmatrix}, \begin{bmatrix}-3\\-3\\1\\-2 \end{bmatrix}, \begin{bmatrix}-3\\0\\3\\-2 \end{bmatrix}\)
\answer{lin.\ dep.}

\item \(\begin{bmatrix} -3\\2\\6\\3\\-2 \end{bmatrix}, \begin{bmatrix}-3\\-3\\1\\6\\-4 \end{bmatrix}, \begin{bmatrix} 3\\-2\\-2\\4\\-1 \end{bmatrix}, \begin{bmatrix} 1\\-2\\-2\\3\\2 \end{bmatrix}\)
\answer{lin.\ indep.}

\item \(\begin{bmatrix} 3\\0\\4\\1\\2 \end{bmatrix}, \begin{bmatrix}-2\\1\\2\\2\\2 \end{bmatrix}, \begin{bmatrix} 2\\1\\0\\0\\0 \end{bmatrix}, \begin{bmatrix}-9\\-1\\-2\\1\\0 \end{bmatrix}\)
\answer{lin.\ dep.}
\end{OmitV1}

\end{Parts}
\end{exercise}


\begin{exercise} 
Prove that every \idx{orthogonal set} of vectors is also a linearly independent set.
\end{exercise}




\begin{exercise} \label{ex:lindeplc} 
Prove the particular case of \cref{thm:lindeplc}, namely that a 
set of two vectors \(\{\vv_1,\vv_2\}\) is linearly dependent if and 
only if one of the vectors is a \index{scalar multiplication}scalar multiple of the other.
\end{exercise}




\begin{exercise}  
Prove that every (non-empty) subset of a linearly independent set is also linearly independent.  (Perhaps use contradiction.)
\end{exercise}





\begin{exercise}  
Let \(\{\hlist\vv m\}\) be a linearly independent set of vectors in~\(\RR^n\).
Given that a vector \(\uv=\lincomb c\vv m\) with coefficient \(c_1\neq0\)\,, prove that the set \(\{\hlist[2]\vv m,\uv\}\) is linearly independent.
\end{exercise}




\begin{exercise}  
For each of the following systems of equations find by hand two different bases for their solution set (among the infinitely many bases that are possible).  
Show your working.
% format rat
% A=0+round(randn(2,3)*3), v=null(A); v=v*diag(1./max(abs(v)))
\begin{enumerate}
\item \(-x-5y=0\) and \(y-3z=0\)
\answer{Two possibilities are 
\(\{(-1,1/5,1/15)\}\) and 
\(\{(15,-3,-1)\}\).}

\item \(6x+4y+2z=0\) and \(-2x-y-2z=0\)
\answer{Two possibilities are 
\(\{(-3/4,1,1/4)\}\) and 
\(\{(3,-4,-1)\}\).}

\begin{OmitV1}
\item \(-2y-z+2=0\) and \(3x+4z=0\)
\answer{Does not have a solution subspace.}

\item \(-7x+y-z=0\) and \(-3x+2y-2z=0\)
\answer{Two possibilities are 
\(\{(0,1,1)\}\) and 
\(\{(0,-2,-2)\}\).}

% A=quads(ceil(6*rand),1:3).*sign(randn(1,3)), v=null(A); v=v*diag(1./max(abs(v))), v*round(randn(2)*2)

\item \(x+2y+2z=0\)
\answer{Two possibilities are 
\(\{(-2,2,-1),(-2,-1,2)\}\) and 
\(\{(0,1,-1),(-8,11,-7)\}\).}

\item \(2x+0y-4z=0\)
\answer{Two possibilities are 
\(\{(0,1,0),(2,0,1)\}\) and 
\(\{(2,-4,1),(4,1,2)\}\).}
\end{OmitV1}

\item \(-2x+3y-6z=0\)
\answer{Two possibilities are 
\(\{(1/2,1,1/3),(-1,1/3,1/2)\}\) and 
\(\{(3,-8,-5),(12,10,1)\}\).}

\item \(9x+4y-9z=6\)
\answer{Does not have a solution subspace.}

\end{enumerate}
\end{exercise}







\begin{exercise}  
Use \cref{pro:bfe} to compute in \script\ an \idx{orthonormal basis} for all solutions to each of the following systems of equations.
%\begin{verbatim}
%format short,for i=1:9999,A=0+round(randn(3,3+ceil(2*rand))*4); if min(svd(A))<1e-7, A=A, break, end, end
%format bank, [u,s,v]=svd(A); sv=diag(s), vt=v'
%\end{verbatim}

\begin{enumerate}
\item \(2x-6y-9z=0\), \(2x-2z=0\), \(-x+z=0\)
\answer{\(\{(0.55,-0.64,0.55)\}\) \twodp}

\item \(3x-3y+8z=0\), \(-2x-4y+2z=0\), \(-4x+y-7z=0\)
\answer{\(\{(0.67,-0.57,-0.47)\}\) \twodp}

\begin{OmitV1}
\item \(-2x+2y-2z=0\), \(x+3y-z=0\), \(3x+3y=0\)
\answer{\(\{(-0.41,0.41,0.82)\}\) \twodp}

\item \(2w+x+2y+z=0\), \(w+4x-4y-z=0\), \(3w-2x+5y=0\), \(2w-x+y-2z=0\)
\answer{\(\{(-0.50,0.50,0.50,-0.50)\}\) \twodp}

\item \(5w+y-3z=0\), \(-5x-5y=0\), \(-3x-y+4z=0\), \(3x+y-4z=0\)
\answer{\(\{(-0.32,-0.63,0.63,-0.32)\}\) \twodp}

\item \(-w-2y+4z=0\), \(2w+2y+2z=0\), \(-2w+3x+y+z=0\), \(-w+x-y+5z=0\)
\answer{\(\{(0.61,0.61,-0.51,-0.10)\}\) \twodp}
\end{OmitV1}

\item \(-2w+x+2y-6z=0\), \(-2w+3x+4y=0\), \(-2w+2x+3y-3z=0\)
\answer{\(\{(0.76,-0.27,0.58,-0.10)\clb (-0.51,-0.78,0.33,0.15)\}\) \twodp}

\item \(-w-2x-3y+2z=0\), \(2x+2y-2z=0\), \(-w-3x-4y+3z=0\)
\answer{\(\{(0.62,0.45,-0.62,-0.17)\clb (-0.13,0.63,0.13,0.76)\}\) \twodp}

\end{enumerate}
\end{exercise}






\begin{exercise} \label{ex:lenlam} 
Recall that \cref{thm:lenlam} establishes that there are at most \(n\)~\idx{eigenvalue}s of an \(n\times n\) symmetric matrix.
Adapt the proof of that theorem, using linear independence, to prove that there are at most \(n\)~eigenvalues of an \(n\times n\) \idx{non-symmetric matrix}.
(This is an alternative to the given proof of \cref{thm:geecp}.)
%\begin{proof} 
%Invoke the \idx{pigeonhole principle} and \idx{contradiction}. 
%Assume there are more than \(n\)~distinct eigenvalues, then there would be more than \(n\)~eigenvectors corresponding to distinct eigenvalues.
%\cref{thm:indepev} asserts all such eigenvectors are linearly independent. 
%But there cannot be more than \(n\)~vectors in a linearly independent set in~\(\RR^n\) (\cref{thm:mgtnli}).
%Hence the assumption is wrong and there cannot be any more than \(n\)~distinct eigenvalues.
%\end{proof}
\end{exercise}



\needlines{10}
\begin{exercise}  
For each diagram, estimate roughly the \idx{components} of each of the four vectors~\av, \bv, \cv, and~\dv\ in the basis \(\cP=\{\pv_1,\pv_2\}\).
\newcommand{\ppoint}[3]{%
    \pgfmathparse{#1*\px+#2*\qx}\let\h\pgfmathresult
    \pgfmathparse{#1*\py+#2*\qy}\let\v\pgfmathresult
    \addplot[blue,mark=*,thick,quiver={u=\h,v=\v},-stealth] coordinates {(0,0)};
    \edef\tempe{%
    \noexpand\node[right] at (axis cs:\h,\v) {$\noexpand\vec #3$};
    }\tempe
    }%
\newcommand{\cbase}[9]{\begin{tikzpicture}%
%[baseline={([yshift={-\ht\strutbox}]current bounding box.north)}]
\begin{axis}[small,font=\footnotesize
  , axis lines=none, unit vector ratio=1 1 1
  , view={0}{90}
  ]
\ifcase#1
\or\addplot3[mesh,red,samples=9,domain=-4:4,dotted] (\px*x+\qx*y,\py*x+\qy*y,0);
\fi
\addplot[red,quiver={u=\px,v=\py},-stealth,thick] coordinates {(0,0)};
\node[right] at (axis cs:\px,\py) {$\pv_1$};
\addplot[red,quiver={u=\qx,v=\qy},-stealth,thick] coordinates {(0,0)};
\node[right] at (axis cs:\qx,\qy) {$\pv_2$};
\ppoint{#2}{#3}{a}
\ppoint{#4}{#5}{b}
\ppoint{#6}{#7}{c}
\ppoint{#8}{#9}{d}
\end{axis}
\end{tikzpicture}
\answer{$(#2,#3)$, $(#4,#5)$, $(#6,#7)$, $(#8,#9)$}}%
\begin{Parts}
\item \def\px{2} \def\py{-1} \def\qx{1} \def\qy{2} 
\cbase0{-2}{-3}{ 1}{-3}{-3}{-1}{ 3}{ 2}

\item \def\px{2} \def\py{-1} \def\qx{-1} \def\qy{3} 
\cbase0{-1.5}{ 0}{ 0.5}{-0.5}{-1}{ 2.5}{ 2}{ 0.5}

\item \def\px{-2} \def\py{1} \def\qx{1} \def\qy{-2} 
\cbase0{2.5}{ 0.5}{-1.5}{ 0}{ 1}{ 0.5}{ 1}{ 1.5}

\item \def\px{-3} \def\py{1} \def\qx{-1} \def\qy{3} 
\cbase0{-1}{-0.5}{ 1}{ 1}{-1}{0.5}{ 0.5}{0}

\begin{OmitV1}
\item \def\px{0} \def\py{1} \def\qx{1} \def\qy{0} 
\cbase0{-0.9}{ 0.3}{ 0.6}{ 0.2}{-0.2}{-1.1}{ 0.7}{-0.6}

\item \def\px{-1} \def\py{-1} \def\qx{1} \def\qy{0} 
\cbase0{-0.9}{ 0.3}{ 0.6}{ 0.2}{-0.2}{-1.1}{ 0.7}{-0.6}

\item \def\px{1} \def\py{-3} \def\qx{1} \def\qy{1} 
\cbase0{0.7}{ 1.6}{ 1.2}{ 0.5}{-0.4}{-0.3}{ 1.4}{2}
\end{OmitV1}

\end{Parts}
\end{exercise}



\begin{exercise} \label{ex:3basisv} 
Let the three given vectors \(\bv_1=(-1,1,-1)\), \(\bv_2=(1,-2,0)\), and \(\bv_3=(0,4,5)\) form a \idx{basis} \(\cB=\{\bv_1,\bv_2,\bv_3\}\) (where these vectors are specified in the \idx{standard basis}~\cE\ of the \idx{standard unit vector}s~\(\ev_1\), \(\ev_2\), and~\(\ev_3\)).
For each of the following vectors with specified coordinates in basis~\cB, what is the vector when written in the \text{standard basis?}
%\begin{verbatim}
%p=0+round(randn(1,3)*2)/2, pe=p*[-1 1 -1;1 -2 0;0 4 5]
%\end{verbatim}

\def\temp#1{% 
  \pgfmathparse{\value{enumii}}%
  \pgfmathprintnumberto[assume math mode=true]{\pgfmathresult}{\tempx}%
  \answer{\noexpand\setcounter{enumii}{\tempx}%
  $[\noexpand\evii]_{\noexpand\cE}=#1$}%
  }

\begin{Parts}
\item \([\evii]_\cB=(1,-1,2)\)
\temp{(-2,11,9)}

\item \([\evii]_\cB=(0,2,3)\)
\temp{(2,8,15)}

\item \([\evii]_\cB=(1,-3,-2)\)
\temp{(-4,-1,-11)}

\item \([\evii]_\cB=(1,2,1)\)
\temp{(1,1,4)}

\item \([\evii]_\cB=(1/2,-1/2,1)\)
\temp{(-1,11/2,9/2)}

\item \([\evii]_\cB=(-1/2,1/2,-1/2)\)
\temp{(1,-7/2,-2)}

\begin{OmitV1}
\item \([\evii]_\cB=(0,1/2,-1/2)\)
\temp{(1/2,-3,-5/2)}

\item \([\evii]_\cB=(-0.7,0.5,1.1)\)
\temp{(1.2,2.7,6.2)}

\item \([\evii]_\cB=(0.2,-0.1,0.9)\)
\temp{(-0.3,4.0,4.3)}

\item \([\evii]_\cB=(2.1,-0.2,0.1)\)
\temp{(-2.3,2.9,-1.6)}
\end{OmitV1}

\end{Parts}
\end{exercise}





\begin{exercise}  
Repeat \cref{ex:3basisv} but with the three basis vectors \(\bv_1=(6,2,1)\), \(\bv_2=(-2,-1,-2)\), and \(\bv_3=(-3,-1,5)\).
% V=0+round(randn(3)*3),sv=svd(V)
\end{exercise}





\begin{exercise} \label{ex:2basev} 
Let the two given vectors \(\bv_1=(1,-2,2)\) and \(\bv_2=(1,-1,-1)\) form a basis \(\cB=\{\bv_1,\bv_2\}\) for the \idx{subspace}~\BB\ of~\(\RR^3\) (specified in the \idx{standard basis}~\cE\ of the \idx{standard unit vector}s~\(\ev_1\), \(\ev_2\), and~\(\ev_3\)).
For each of the following vectors, specified in the standard basis~\cE, what is the vector when written in the basis~\cB, if it is possible?
%\begin{verbatim}
%B=[1 -2 2;1 -1 -1]'
%pb=0+round(randn(1,2)*3)/10, pe=pb*B'+(rand<0.33)*round(randn(1,3)), pbt=B\pe'
%\end{verbatim}

\def\temp#1{% 
  \pgfmathparse{\value{enumii}}%
  \pgfmathprintnumberto[assume math mode=true]{\pgfmathresult}{\tempx}%
  \answer{\noexpand\setcounter{enumii}{\tempx}%
  $[\noexpand\evii]_{\noexpand\cB}=#1$}%
  }

\begin{Parts}
\item \([\evii]_{\cE}=(0,1,3)\)
\temp{(-1,1)}

\item \([\evii]_{\cE}=(-4,9,-11)\)
\temp{(-5,1)}

\item \([\evii]_{\cE}=(-4,7,-5)\)
\temp{(-3,-1)}

\item \([\evii]_{\cE}=(0,2,-4)\)
\answer{not in \cB}

\item \([\evii]_{\cE}=(0,-2,5)\)
\answer{not in \cB}

\item \([\evii]_{\cE}=(-2/3,0,8/3)\)
\temp{(2/3,-4/3)}

\begin{OmitV1}
\item \([\evii]_{\cE}=( -8/3,19/3,-19/3)\)
\answer{not in \cB}

\item \([\evii]_{\cE}=(-10/3,5,-5/3)\)
\temp{(-5/3,-5/3)}

\item \([\evii]_{\cE}=(-0.3,0.4,0)\)
\temp{(-0.1,-0.2)}

\item \([\evii]_{\cE}=(0.5,-0.7,0.1)\)
\temp{(0.2,0.3)}
\end{OmitV1}

\end{Parts}
\end{exercise}





\begin{exercise}  
Repeat \cref{ex:2basev} but with the two basis vectors \(\bv_1=(-2,3,-1)\) and \(\bv_2=(0,-1,3)\).
\end{exercise}






\begin{exercise} \label{ex:3base5v} 
Let the three vectors \(\bv_1=(-1,-2,5,3,-2)\), \(\bv_2=(-2,-2,2,-1,-1)\), and \(\bv_3=(-4,6,-4,2,-1)\) form a \idx{basis} \(\cB=\{\bv_1,\bv_2,\bv_3\}\) for the \idx{subspace}~\BB\ in~\(\RR^5\) (each basis vector specified in the \idx{standard basis}~\cE).
\setbox\ajrqrbox\hbox{\qrcode{% basis
B=[-1 -2 -4
 -2 -2 6
 5 2 -4
 3 -1 2
 -2 -1 -1]
}}\marginajrbox%
For each of the following vectors, use \script\ to find the requested coordinates, if possible.
%\begin{verbatim}
%B=0+round(randn(5,3)*3), svb=svd(B)
%pb=0+round(randn(1,3)*3)/1, pe=pb*B'+(rand<0.33)*round(randn(1,5)), pbt=B\pe'
%\end{verbatim}
\def\temp#1#2{% 
  \pgfmathparse{\value{enumii}}%
  \pgfmathprintnumberto[assume math mode=true]{\pgfmathresult}{\tempx}%
  \answer{\noexpand\setcounter{enumii}{\tempx}%
  $[\noexpand\evii]_{\noexpand#1}=#2$}%
  }
\begin{enumerate}
\item Find \([\evii]_\cE\) when \([\evii]_\cB=(2,2,-1)\).
\temp{\cE}{(-2,-14,18,2,-5)}

\item Find \([\evii]_\cE\) when \([\evii]_\cB=(-5,0,-2)\).
\temp{\cE}{(13,-2,-17,-19,12)}

\item Find \([\evii]_\cE\) when \([\evii]_\cB=(0,3,3)\).
\temp{\cE}{(-18,12,-6,3,-6)}

\item Find \([\evii]_\cE\) when \([\evii]_\cB=(-1,5,0)\).
\temp{\cE}{(-9,-8,5,-8,-3)}

\item Find \([\evii]_\cB\) when \([\evii]_\cE=(-31,26,-5,19,-14)\).
\temp{\cB}{(3,2,6)}

\item Find \([\evii]_\cB\) when \([\evii]_\cE=(-1,6,4,14,-5)\).
\answer{not in \BB}

\begin{OmitV1}
\item Find \([\evii]_\cB\) when \([\evii]_\cE=(-21,18,-7,9,-8)\).
\temp{\cB}{(1,2,4)}

\item Find \([\evii]_\cB\) when \([\evii]_\cE=(-0.2,-0.6,1.8,1.3,-0.7)\).
\temp{\cB}{(0.4,-0.1,0)}

\item Find \([\evii]_\cB\) when \([\evii]_\cE=(0.7,-0.4,-0.3,-0.3,1.2)\).
\answer{not in \BB}

\item Find \([\evii]_\cB\) when \([\evii]_\cE=(4.8,-3.8,0.8,-2.6,2.1)\).
\temp{\cB}{(-0.4,-0.4,-0.9)}
\end{OmitV1}

\end{enumerate}
\end{exercise}




\begin{exercise}  
Repeat \cref{ex:3base5v} but with basis vectors \(\bv_1=(-3,8,-9,-1,1)\), \(\bv_2=(10,-20,14,-7,2)\), and \(\bv_3=(-1,-2,5,3,-2)\) (specified in the \idx{standard basis}~\cE).
%for i=1:9999,P=round(randn(3)*2); if abs(det(P))==1, P=P, break, end, end
\setbox\ajrqrbox\hbox{\qrcode{% basis
B=[-3 10 -1
 8 -20 -2
 -9 14 5
 -1 -7 3
 1 2 -2]
}}\marginajrbox%
\end{exercise}




\begin{exercise} 
In a few sentences, answer\slash discuss each of the following.
\begin{enumerate}
\item  Why is it important to establish a mathematical framework in which solutions have a unique algebraic form?

\item How is \cref{def:lindep} of linear independence linked to a unique algebraic form?

\item How does linear independence compare to orthogonality of a set of vectors?

\item For distinct eigenvalues of a matrix~\(A\), how does the linear independence of the corresponding eigenvectors arise?

\item How can an \svd\ be useful in testing whether a set of vectors is linearly dependent or not?

\item How does a ``basis'' compare to an ``orthonormal basis''?

\item What problems might arise if, for a given subspace, there were bases for the subspace with different numbers of vectors?

\item How does an \svd\ help us to find bases for subspaces?

\item Why is new notation needed for vectors when we invoke the possibility of different bases?

\end{enumerate}
\end{exercise}

\begin{comment}%{ED498555.pdf}
why, what caused X?
how did X occur?
what-if? what-if-not?
how does X compare with Y?
what is the evidence for X?
why is X important?
\end{comment}




\index{linearly dependent|)}
\index{linearly independent|)}
\index{linear dependence|)}
\index{linear independence|)}



%%!TEX root = ../larxxia.tex


\section{Diagonalization identifies the transformation}
\label{sec:dit}

\secttoc
\begin{comment}
\pooliv{\S4.4} \layiv{\S5.3} \holti{\S6.4} \cite[\S09]{Davis99a}
\end{comment}

\index{diagonalization|(}
\index{diagonalizable|(}

\paragraph{Population modelling} \index{population modelling}
Recall that this \cref{ch:gee} started by introducing the dynamics of two interacting \idx{species} of animals.
Recall we let \(y(t)\) and~\(z(t)\) be the number of \idx{female} animals in each of the two species at time~\(t\) (years). 
Modelling might deduce the populations interact according to the rule that the population one year later is \(y(t+1)=2y(t)-4z(t)\) and \(z(t+1)=-y(t)+2z(t)\).
Then seeking solutions proportional to~\(\lambda^t\) led to the \idx{eigen-problem}
\begin{equation*}
\begin{bmatrix} 2&-4\\-1&2 \end{bmatrix}\xv=\lambda\xv\,.
\end{equation*}
This section introduces an alternate equivalent approach.

The alternate approach invokes \idx{non-orthogonal coordinates}.
Start by writing the \idx{population model} as a \idx{system} in terms of vector \(\yv(t)=(y(t),z(t))\), namely
\begin{align*}&
\begin{bmatrix} y(t+1)\\z(t+1) \end{bmatrix}
=\begin{bmatrix} 2y(t)-4z(t)\\-y(t)+2z(t) \end{bmatrix}
,&& \text{that is, }
\yv(t+1)=\begin{bmatrix} 2&-4\\-1&2 \end{bmatrix}\yv(t).
\end{align*}
Now let's ask if there is a basis \(\cP=\{\pv_1,\pv_2\}\) for the \(yz\)-plane that simplifies this matrix-vector system?
In such a basis every vector may be written as \(\yv=Y_1\pv_1+Y_2\pv_2\) for some scalar components~\(Y_1\) and~\(Y_2\).
That is,  \([\yv]_\cP=(Y_1,Y_2)\).
But to simplify writing we use the vector symbol~\(\Yv=(Y_1,Y_2)\)\ in place of~\([\yv]_\cP\).
Write the relation \(\yv=Y_1\pv_1+Y_2\pv_2\) as the \idx{matrix-vector product} \(\yv=P\Yv\) where matrix \(P=\begin{bmatrix} \pv_1&\pv_2 \end{bmatrix}\) and vector \(\Yv=(Y_1,Y_2)\).
The populations~\(\yv\) depends upon time~\(t\), and hence so does~\(\Yv\) since \(\Yv=[\yv]_\cP\); that is, \(\yv(t)=P\Yv(t)\).
Substitute this identity into the system \text{of equations:}
\begin{equation*}
\yv(t+1)=P\Yv(t+1)=\begin{bmatrix} 2&-4\\-1&2 \end{bmatrix}P\Yv(t).
\end{equation*}
Multiply both sides by~\(P^{-1}\) (which exists by linear independence of the columns, \cref{thm:ftim3xi}) to give
\begin{equation*}
\Yv(t+1)=\underbrace{P^{-1}\begin{bmatrix} 2&-4\\-1&2 \end{bmatrix}P}_{P^{-1}AP}\Yv\,.
\end{equation*}
The question then becomes, for a given \idx{square matrix}~\(A\), such as this, can we find a matrix~\(P\) such that \(P^{-1}AP\) is somehow simple?
The answer is yes, in most cases: using \idx{eigenvalue}s and \idx{eigenvector}s the product~\(P^{-1}AP\) usually can be made into a simple \text{\idx{diagonal matrix}.}



\subsection{Eigenvectors achieve diagonalization}
\label{sec:ditead}


Recall that (\cref{sec:smod}) for a \idx{symmetric matrix}~\(A\) we could always factor \(A=VD\tr V=VDV^{-1}\) for \idx{orthogonal matrix}~\(V\) and \idx{diagonal matrix}~\(D\): thus a symmetric matrix is always \idx{orthogonally diagonalizable} (\cref{def:odsble}).
For non-symmetric matrices, a \idx{diagonalization} mostly can be done  (although not always):
the difference being we need an invertible matrix, typically called~\(P\), instead of an orthogonal matrix~\(V\).
Such a matrix~\(A\) is generally termed `diagonalizable' instead of the more specific `\idx{orthogonally diagonalizable}'.


\begin{definition} \label{def:diagonalise} 
An \(n\times n\) \idx{square matrix}~\(A\) is \bfidx{diagonalizable} if there exists a \idx{diagonal matrix}~\(D\) and an \idx{invertible} matrix~\(P\) such that \(A=PDP^{-1}\), equivalently \(AP=PD\) or \(P^{-1}AP=D\)\,.
\end{definition}


\begin{example} \label{eg:diagonalise}
\begin{enumerate}[ref=\ref{eg:diagonalise}(\alph*)]
\item\label[example]{eg:diagonalisea} 
Show that \(A=\begin{bmat} 0&1\\2&-1 \end{bmat}\) is diagonalizable by matrix \(P=\begin{bmat} 1&-1\\1&2 \end{bmat}\).
\begin{solution} 
First, find the \(2\times2\) inverse (\cref{thm:2x2det})
\begin{equation*}
P^{-1}=\frac1{3}\begin{bmatrix} 2&1\\-1&1 \end{bmatrix}.
\end{equation*}
Second, compute the product
\begin{equation*}
P^{-1}AP
=P^{-1}\begin{bmatrix} 1&2\\1&-4 \end{bmatrix}
=\frac13\begin{bmatrix} 3&0\\0&-6 \end{bmatrix}
=\begin{bmatrix} 1&0\\0&-2 \end{bmatrix}.
\end{equation*}
Since this product is diagonal, \(\diag(1,-2)\), the matrix~\(A\) is diagonalizable.
\end{solution}


\item\label[example]{eg:diagonaliseb} Use a contradiction to prove that \(B=\begin{bmat} 0&1\\0&0 \end{bmat}\) is not diagonalizable.
\begin{solution} 
Assume \(B\) is diagonalizable by an invertible matrix \(P=\begin{bmat} a&b\\c&d \end{bmat}\).
Being invertible, \(P\)~has inverse \(P^{-1}=\frac1{ad-bc}\begin{bmat} d&-b\\-c&a \end{bmat}\)  (\cref{thm:2x2det}).
Then the product
\begin{equation*}
P^{-1}BP=P^{-1}\begin{bmatrix} c&d\\0&0 \end{bmatrix}
=\frac1{ad-bc}\begin{bmatrix} cd&d^2\\-c^2&-cd \end{bmatrix}.
\end{equation*}
For the matrix~\(P^{-1}BP\) to be diagonal requires the off-diagonal elements to be zero: \(d^2=-c^2=0\)\,.
This requires both \(c=d=0\)\,, but then the determinant \(ad-bc=0-0=0\) and so matrix~\(P\) is not invertible (\cref{thm:2x2det}).
This \idx{contradiction} means that matrix~\(B\) is \text{not diagonalizable.}
\end{solution}

\item Is matrix \(C=\begin{bmat}1.2& 3.2& 2.3
\\   2.2&-0.5&-2.2\end{bmat}\) diagonalizable?
\begin{solution} 
No, as it is not a square matrix.
(Perhaps an \svd\ could answer the needs of whatever problem led to this question.) 
\end{solution}

\end{enumerate}
\end{example}


\begin{example} 
\cref{eg:diagonalisea} showed that matrix \(P=\begin{bmat} 1&-1\\1&2 \end{bmat}\) diagonalizes matrix \(A=\begin{bmat} 0&1\\2&-1 \end{bmat}\) to matrix \(D=\diag(1,-2)\). 
As a prelude to the next \cref{thm:gendiag}, show that the columns of~\(P\) are \idx{eigenvector}s of~\(A\).
\begin{solution} Invoke the original \cref{def:evecval} of an eigenvector for \text{a matrix.}
\begin{itemize}
\item The first column of~\(P\) is \(\pv_1=(1,1)\).
Multiplying \(A\pv_1=(0+1,2-1)=(1,1)=1\pv_1\) so the first column vector~\(\pv_1\) is an eigenvector of~\(A\) corresponding to the eigenvalue~\(1\).
Correspondingly, this eigenvalue~\(1\) is the first entry in the diagonal~\(D\).

\item The second column of~\(P\) is \(\pv_2=(-1,2)\).
Multiplying \(A\pv_2=(0+2,-2-2)=(2,-4)=-2\pv_2\) so the second column vector~\(\pv_2\) is an eigenvector of~\(A\) corresponding to the eigenvalue~\(-2\).
Correspondingly, this eigenvalue~\(-2\) is the second entry in the diagonal~\(D\).
\aqed

\end{itemize}
\end{solution}
\end{example}




\begin{activity}
%for it=1:999, p=0+round(randn(2)*3); if abs(det(p))==1,break,end, end, p=p, d=diag(0+round(randn(1,2)*3)), a=p*d*inv(p)
Given matrix~\(F=\begin{bmatrix} 5&8\\-4&-7 \end{bmatrix}\) has eigenvectors \((-1,1)\) and~\((2,-1)\) corresponding to respective {eigenvalue}s~\(-3\) and~\(1\), what matrix diagonalizes~\(F\) to \(D=\diag(-3,1)\)?
\actposs[4]{\(\begin{bmatrix} -1&2\\1&-1 \end{bmatrix}\)}
{\(\begin{bmatrix} -1&1\\2&-1 \end{bmatrix}\)}
{\(\begin{bmatrix} 2&-1\\-1&1 \end{bmatrix}\)}
{\(\begin{bmatrix} 2&-1\\1&-1 \end{bmatrix}\)}
\end{activity}




\begin{theorem} \label{thm:gendiag} 
For every \(n\times n\) \idx{square matrix}~\(A\), the matrix~\(A\) is \idx{diagonalizable} if and only if \(A\)~has \(n\)~\idx{linearly independent} \idx{eigenvector}s.  
If \(A\)~is \idx{diagonalizable}, with \idx{diagonal matrix} \(D=P^{-1}AP\), then  the \idx{diagonal entries} of~\(D\) are \idx{eigenvalue}s, and the columns of~\(P\) are corresponding \idx{eigenvector}s.
\end{theorem}
\begin{proof} 
First, let matrix~\(A\) be diagonalizable by invertible~\(P\) and diagonal~\(D\).
Write \(P=\begin{bmatrix} \pv_1&\pv_2&\cdots&\pv_n \end{bmatrix}\) in terms of its columns, and let \(D=\diag(\hlist\lambda n)\) in terms of its diagonal entries.
Then \(AP=PD\) becomes
\begin{equation*}
A\begin{bmatrix} \pv_1&\pv_2&\cdots&\pv_n \end{bmatrix}
=\begin{bmatrix} \pv_1&\pv_2&\cdots&\pv_n \end{bmatrix}
\begin{bmatrix} \lambda_1&0&\cdots&0\\
0&\lambda_2&&0\\
\vdots&&\ddots&\vdots\\
0&0&\cdots&\lambda_n \end{bmatrix}.
\end{equation*}
Multiplying the matrix-column products on both sides gives
\begin{equation*}
\begin{bmatrix} A\pv_1&A\pv_2&\cdots&A\pv_n \end{bmatrix}
=\begin{bmatrix} \lambda_1\pv_1& \lambda_2\pv_2&\cdots& \lambda_n\pv_n \end{bmatrix}.
\end{equation*}
Equating corresponding columns implies \(A\pv_1=\lambda_1\pv_1\)\,, \(A\pv_2=\lambda_2\pv_2\)\,, \ldots, \(A\pv_n=\lambda_n\pv_n\)\,.
As the matrix~\(P\) is invertible, all its columns must be nonzero (\cref{thm:ppdet:i}).
Hence \hlist\pv n\ are eigenvectors of matrix~\(A\) corresponding to eigenvalues \hlist\lambda n\,.
Since matrix~\(P\) is invertible, \cref{thm:ftim3xi} asserts the \(n\)~columns vectors \hlist\pv n\ are \text{linearly independent.}

Second, suppose matrix~\(A\) has \(n\)~{linearly independent} {eigenvector}s \hlist\pv n\ with corresponding eigenvalues \hlist\lambda n\,.  
Then follow the above argument backwards to deduce \(AP=PD\) for invertible matrix \(P=\begin{bmatrix} \pv_1&\pv_2&\cdots&\pv_n \end{bmatrix}\), and hence \(A\)~is diagonalizable.
Lastly, in these arguments, \(P\)~is the matrix of eigenvectors and the diagonal of~\(D\) are the corresponding eigenvalues, \text{as required.}
\end{proof}


\begin{example} \label{eg:trieigp}
Recall that \cref{eg:trieig} found the \idx{triangular matrix}
\(A=\begin{bmat}-3&2&0
\\0&-4&2
\\0&0&4\end{bmat}\)
has eigenvalues~\(-3\), \(-4\) and~\(4\) (from its diagonal) and corresponding eigenvectors are proportional to \((1,0,0)\), \((-2,1,0)\) and \((\frac1{14},\frac14,1)\).
Is matrix~\(A\) diagonalizable?
\begin{solution} 
These three eigenvectors are linearly independent as they correspond to distinct eigenvalues (\cref{thm:indepev}).
Hence the matrix is diagonalizable.

The previous statement answers the question. 
But further, forming these eigenvectors into the columns of matrix
\begin{equation*}
P=\begin{bmatrix} 1&-2&\frac1{14}
\\0&1&\frac14
\\0&0&1 \end{bmatrix},
\end{equation*}
we know that (\cref{thm:gendiag}) \(P^{-1}AP=\diag(-3,-4,4)\) where the eigenvalues appear in the same order as that of the eigenvectors in~\(P\).

One may check this by hand or with \script.
Enter the matrices with
\begin{verbatim}
A=[-3 2 0;0 -4 2;0 0 4]
P=[1 -2 1/14;0 1 1/4;0 0 1]
\end{verbatim}
\setbox\ajrqrbox\hbox{\qrcode{% check diag
A=[-3 2 0;0 -4 2;0 0 4]
P=[1 -2 1/14;0 1 1/4;0 0 1]
D=P\slosh A*P
}}\marginajrbox%
then compute \(D=P^{-1}AP\) with \verb|D=P\A*P| to find as required the following diagonal result
\begin{verbatim}
D =
  -3.0000   0.0000   0.0000
   0.0000  -4.0000   0.0000
   0.0000   0.0000   4.0000
\end{verbatim}
\end{solution}
\end{example}


\begin{wrapfigure}[4]r{0pt}
\begin{tikzpicture} 
\begin{axis}[footnotesize,font=\footnotesize
    ,axis equal image, axis lines=none ]
    \addplot[blue,mark=*,domain=40:400,samples=10
    ,nodes near coords={\quad\color{black}\ifnum\coordindex=0\else\coordindex\fi}] 
    ({cos(\x)+0.347*cos(2*\x)}
    ,{sin(\x)-0.347*sin(2*\x)});
    \addplot[blue,mark=*,samples at={40,80,160,200,280,320,40}] 
    ({cos(\x)+0.347*cos(2*\x)}
    ,{sin(\x)-0.347*sin(2*\x)});
\end{axis}
\end{tikzpicture}
\end{wrapfigure}
\begin{example} \label{eg:sier2eigp}
Recall the \idx{Sierpinski network} of \cref{eg:sier2eig} (shown to the right).
The following \(9\times9\) matrix~\(A\) encodes the network.  
Is the matrix diagonalizable?
\begin{equation*}
A=\begin{bmatrix}-3&1&1&0&0&0&0&0&1
\\1&-2&1&0&0&0&0&0&0
\\1&1&-3&1&0&0&0&0&0
\\0&0&1&-3&1&1&0&0&0
\\0&0&0&1&-2&1&0&0&0
\\0&0&0&1&1&-3&1&0&0
\\0&0&0&0&0&1&-3&1&1
\\0&0&0&0&0&0&1&-2&1
\\1&0&0&0&0&0&1&1&-3 \end{bmatrix}.
\end{equation*}

\begin{solution} 
In that example we used \script\ command \verb|[V,D]=eig(A)| to compute a matrix of eigenvectors~\verb|V| and the corresponding diagonal matrix of eigenvalues~\verb|D| where \twodp
\begin{verbatim}
V =
 -0.41  0.51 -0.16 -0.21 -0.45  0.18 -0.40  0.06  0.33
  0.00 -0.13  0.28  0.63  0.13 -0.18 -0.58 -0.08  0.33
  0.41 -0.20 -0.49 -0.42  0.32  0.01 -0.36 -0.17  0.33
 -0.41 -0.11  0.52 -0.42  0.32  0.01  0.14 -0.37  0.33
 -0.00 -0.18 -0.26  0.37 -0.22  0.51  0.36 -0.46  0.33
  0.41  0.53  0.07  0.05 -0.10 -0.51  0.33 -0.23  0.33
 -0.41 -0.39 -0.36  0.05 -0.10 -0.51  0.25  0.31  0.33
  0.00  0.31 -0.03  0.16  0.55  0.34  0.22  0.55  0.33
  0.41 -0.33  0.42 -0.21 -0.45  0.18  0.03  0.40  0.33
D =
 -5.00     0     0     0     0     0     0     0     0
     0 -4.30     0     0     0     0     0     0     0
     0     0 -4.30     0     0     0     0     0     0
     0     0     0 -3.00     0     0     0     0     0
     0     0     0     0 -3.00     0     0     0     0
     0     0     0     0     0 -3.00     0     0     0
     0     0     0     0     0     0 -0.70     0     0
     0     0     0     0     0     0     0 -0.70     0
     0     0     0     0     0     0     0     0 -0.00
\end{verbatim}
%From these we identify the following eigenspaces:
%\begin{eqnarray*}
%&&\EE_{-5}=\Span\{(-1,0,1,-1,0,1,-1,0,1)\};
%\\&&\EE_{-4.30}=\Span\left\{ \begin{bmatrix} 0.51\\-0.13\\-0.20\\-0.11\\-0.18\\0.53\\-0.39\\0.31\\-0.33 \end{bmatrix},\  \begin{bmatrix} -0.16\\0.28\\-0.49\\0.52\\-0.26\\0.07\\-0.36\\-0.03\\0.42 \end{bmatrix} \right\};
%\\&&\EE_{-3}=\Span\left\{ 
%\begin{bmatrix} -0.21\\0.63\\-0.42\\-0.42\\0.37\\0.05\\0.05\\0.16\\-0.21
%\end{bmatrix},\  
%\begin{bmatrix} -0.45\\0.13\\0.32\\0.32\\-0.22\\-0.10\\-0.10\\0.55\\-0.45
%\end{bmatrix},\ 
%\begin{bmatrix} 0.18\\-0.18\\0.01\\0.01\\0.51\\-0.51\\-0.51\\0.34\\0.18
% \end{bmatrix} \right\};
%\\&&\EE_{-0.70}=\Span\left\{ \begin{bmatrix} -0.40\\-0.58\\-0.36\\0.14\\0.36\\0.33\\0.25\\0.22\\0.03 \end{bmatrix},\  \begin{bmatrix} 0.06\\-0.08\\-0.17\\-0.37\\-0.46\\-0.23\\0.31\\0.55\\0.40 \end{bmatrix} \right\}; 
%\\&&\EE_{0}=\Span\{(1,1,1,1,1,1,1,1,1)\}.
%\end{eqnarray*}
Since matrix~\(A\) is symmetric, \script\ computes for us an orthogonal matrix~\(V\) with columns eigenvectors.
Since~\(V\) is orthogonal its column vectors are orthonormal and hence its columns are linearly independent (\cref{thm:ortholi}).
Since there exist nine linearly independent eigenvectors, the nine column vectors of~\(V\), the matrix~\(A\) is diagonalizable.
Further, the product \(V^{-1}AV=D\) for the above diagonal matrix~\(D\) of eigenvalues in the order of the eigenvectors in~\(V\).
(Also, since \(V\)~is orthogonal, \(\tr VAV=D\).)
\aqed

\end{solution}
\end{example}





\begin{example} \label{eg:faespmp}
Recall \cref{eg:faespm} found \idx{eigenvalue}s and corresponding \idx{eigenspace}s for various matrices.
Revisit these cases and show none of the matrices are diagonalizable.
\begin{enumerate}
\item Matrix \(A=\begin{bmat} 3&1\\0&3 \end{bmat}\) had one eigenvalue \(\lambda=3\) with \idx{multiplicity} two and corresponding eigenspace \(\EE_3=\Span\{(1,0)\}\).
This matrix is not diagonalizable as it has only one \idx{linearly independent} eigenvector, such as~\((1,0)\) or any nonzero multiple, and it needs two to \text{be diagonalizable.}

\item Matrix \(B=\begin{bmat}-1&1&-2
\\-1&0&-1
\\0&-3&1 \end{bmat}\)
has eigenvalues  \(\lambda=-2\) (multiplicity one) and \(\lambda=1\) (multiplicity two).
The corresponding eigenspaces are \(\EE_{-2}=\Span\{(1,1,1)\}\) and \(\EE_{1}=\Span\{(-1,0,1)\}\).
Thus the matrix has only two \idx{linearly independent} eigenvectors, one from each eigenspace, and it needs three to \text{be diagonalizable.}

\item Matrix \(C=\begin{bmat}-1&0&-2
\\0&-3&2
\\0&-2&1\end{bmat}\)
has only the eigenvalue \(\lambda=-1\) with multiplicity three.
The corresponding eigenspace \(\EE_{-1}=\Span\{(1,0,0)\}\).
With only one \idx{linearly independent} eigenvector, the matrix is not diagonalizable.

\end{enumerate}
\end{example}


\begin{example} \label{eg:fiveevp}
Use the results of \cref{eg:fiveev} to show the following matrix is \idx{diagonalizable}:
\begin{equation*}
A=\begin{bmatrix}0&3&0&0&0
\\1&0&3&0&0
\\0&1&0&3&0
\\0&0&1&0&3
\\0&0&0&1&0\end{bmatrix}
\end{equation*}
\begin{solution} 
\cref{eg:fiveev} derived the five eigenvalues are  \(\lambda=0,\pm\sqrt3,\pm3\)\,, all of multiplicity one.
Further, the corresponding eigenspaces are
\begin{eqnarray*}
&&\EE_0=\Span\{(9,0,-3,0,1)\},
\\&&\EE_{\pm\sqrt3}=\Span\{(-9,\mp3\sqrt3,0,\pm\sqrt3,1)\},
\\&&\EE_{\pm3}=\Span\{(9,\pm9,6,\pm3,1)\}.
\end{eqnarray*}
Here there are five linearly independent eigenvectors, one from each distinct eigenspace  (\cref{thm:indepev}).
Since \(A\) is a \(5\times5\) matrix it is thus diagonalizable.
Further, \cref{thm:gendiag} establishes that the matrix formed from the columns of the five eigenvectors is a possible matrix
\begin{equation*}
P=\begin{bmatrix} 9&-9&-9&9&9
\\0&-3\sqrt3&3\sqrt3&9&-9
\\-3&0&0&6&6
\\0&\sqrt3&-\sqrt3&3&-3
\\1&1&1&1&1
\end{bmatrix}.
\end{equation*}
(One could also scale each column of~\(P\) by a different arbitrary nonzero constant, and the diagonalization still holds.)
Lastly, \cref{thm:gendiag} establishes the diagonal matrix
\(P^{-1}AP=D=\diag(0,\sqrt3,-\sqrt3,3,-3)\) is that of the eigenvalues in the order corresponding to the eigenvectors in~\(P\).
\end{solution}
\end{example}





These examples illustrate a widely useful property.
The \(5\times 5\) matrix in \cref{eg:fiveevp} has five distinct eigenvalues whose corresponding \idx{eigenvector}s are necessarily \idx{linearly independent} (\cref{thm:indepev}) and so diagonalize the matrix (\cref{thm:gendiag}).
The \(3\times 3\) matrix in \cref{eg:trieigp} has three distinct eigenvalues whose corresponding eigenvectors are necessarily \idx{linearly independent} (\cref{thm:indepev}) and so diagonalize the matrix (\cref{thm:gendiag}).
However, the matrices of \cref{eg:sier2eigp,eg:faespmp} have \idx{repeated eigenvalue}s---eigenvalues of \idx{multiplicity} two or more---and these matrices may (\cref{eg:sier2eigp}) or may not (\cref{eg:faespmp}) be diagonalizable.
The following theorem confirms that matrices with as many \emph{distinct} eigenvalues as the \idx{size} of the matrix are \text{always diagonalizable.}






\begin{theorem} \label{thm:dlamd} 
For every \(n\times n\) \idx{square matrix}~\(A\), if~\(A\) has \(n\)~\idx{distinct eigenvalues}, then \(A\)~is \idx{diagonalizable}.
Consequently, and allowing nonreal \idx{complex eigenvalue}s, a real \idx{non-diagonalizable} matrix must be non-symmetric and must have at least one \idx{repeated eigenvalue} (an \idx{eigenvalue} with \idx{multiplicity} two or~more).
\end{theorem}
\begin{proof}  
First, let \hlist\vv n\ be eigenvectors corresponding to the \(n\)~distinct eigenvalues of matrix~\(A\).
(Recall that \cref{thm:geecp} establishes that there cannot be more than \(n\)~eigenvalues.)
As the corresponding eigenvalues are distinct, \cref{thm:indepev} establishes that these eigenvectors are linearly independent.
\cref{thm:gendiag} then establishes the matrix~\(A\) \text{is diagonalizable.}

Second, since an \(n\times n\) matrix has \(n\)~eigenvalues, when counted accordingly to multiplicity and allowing for nonreal complex-valued eigenvalues (\cref{pro:geneig}), a non-diagonalizable matrix must have at least one repeated eigenvalue.  
Further, by \cref{thm:symspec} a real symmetric matrix is always diagonalizable: hence a non-diagonalizable matrix must also \text{be non-symmetric.}
\end{proof}


\begin{example} 
From the given information, are the matrices diagonalizable?
\begin{enumerate}
\item The only \idx{eigenvalue}s of a \(4\times 4\) matrix are \(1.8\), \(-3\), \(0.4\) and~\(3.2\).
\begin{solution} 
\cref{thm:dlamd} implies the matrix must be diagonalizable.
\end{solution}

\item The only eigenvalues of a \(5\times 5\) matrix are \(1.8\), \(-3\), \(0.4\) and~\(3.2\).
\begin{solution} 
Here there are only four distinct eigenvalues of the \(5\times5\) matrix.
\cref{thm:dlamd} does not apply as the precondition that there be five distinct eigenvalues is not met: the matrix may or may not be diagonalizable---it is unknowable on this information.
\end{solution}

\item The only eigenvalues of a \(3\times 3\) matrix are \(1.8\), \(-3\), \(0.4\) and~\(3.2\).
\begin{solution} 
An error has been made in determining the eigenvalues as a \(3\times3\) matrix has at most three distinct eigenvalues (\cref{thm:geecp}).
Because of the error, we cannot answer.
\end{solution}

\end{enumerate}
\end{example}



\begin{activity}
A \(3\times3\) matrix~\(A\) depends upon a parameter~\(a\) and has \idx{eigenvalue}s \(6\), \(3-3a\) and \(2+a\)\,.
For which of the following values of parameter~\(a\) may the matrix be \emph{not} diagonalizable?
\actposs[4]{\(a=4\)}{\(a=1\)}{\(a=2\)}{\(a=3\)}
Extension: what are the two other values of~\(a\) for which the matrix \emph{may not} be diagonalizable?
\end{activity}





\begin{example} 
\script\ computes the \idx{eigenvalue}s of matrix
\begin{equation*}
A=\begin{bmatrix}-1&2&-2&1&-2
\\-3&-1&-2&5&6
\\3&1&6&-2&-1
\\1&1&2&1&-1
\\7&5&-3&0&0 \end{bmatrix}
\end{equation*}
via \index{eig()@\texttt{eig()}}\verb|eig(A)| and reports them to be \twodp
\begin{verbatim}
ans =
  -3.45 + 3.50i
  -3.45 - 3.50i
   5.00
   5.00
   1.91
\end{verbatim}
Is the matrix diagonalizable?
\begin{solution} 
The matrix appears to have only four distinct eigenvalues (two of them non-real \idx{complex valued}), and so on the given information \cref{thm:dlamd} cannot determine whether the matrix is diagonalizable \text{or not.}

However, upon reporting the eigenvalues to four \idx{decimal places} we find the two eigenvalues of~\verb|5.00| \twodp\ are more precisely two separate eigenvalues of~\verb|5.0000| and~\verb|4.9961|.
Hence this matrix has five distinct eigenvalues and so \cref{thm:dlamd} implies the matrix is diagonalizable.%
\footnote{Nonetheless, in an application where errors are significant then the matrix may be effectively \idx{non-diagonalizable}.
Such effective non-diagonalizability is indicated by poor conditioning of the matrix of eigenvectors which here has the poor \texttt{rcond} of~\(0.0004\) (\cref{pro:unisol}).}
\end{solution}
%\begin{verbatim}
%for i=1:99999,a=round(randn(5)*3); if abs(min(abs(diff(eig(a))))-3e-3)<2e-3, eig(a), break, end, end
%\end{verbatim}
\end{example}





Recall that \cref{def:eigmult} defined the \idx{multiplicity} of every eigenvalue in terms of the characteristic polynomial.
Also, the earlier \cref{def:eigsymult} identified that for every \index{symmetric matrix}\emph{symmetric} matrix the \idx{dimension} of an \idx{eigenspace}, \(\dim\EE_{\lambda_j}\), is equal to the \idx{multiplicity} of the corresponding eigenvalue~\(\lambda_j\).
However, for general non-symmetric matrices this equality between multiplicity and eigenspace dimension does not \text{necessarily hold.}


\begin{theorem} \label{thm:dimee} 
For every square matrix~\(A\), and for each \idx{eigenvalue}~\(\lambda_j\) of~\(A\), the corresponding \idx{eigenspace}~\(\EE_{\lambda_j}\) has \idx{dimension} less than or equal to the \idx{multiplicity} of~\(\lambda_j\);
that is, \(1\leq\dim\EE_{\lambda_j}\leq\text{\idx{multiplicity} of }\lambda_j\)\,.  
\end{theorem}
%\begin{comment}
%Maths I 2014 course omits any proof.
%\end{comment}
%\begin{proof} 
%Suppose \(\lambda_j\) is an eigenvalue of matrix~\(A\) and \(\dim\EE_{\lambda_j}=p<n\) (the case \(p=n\) is proved by \cref{ex:dimme}).
%Then to span \(\EE_{\lambda_j}\) there are \(p\) corresponding, linearly independent, eigenvectors \hlist\vv p\,.
%Let
%\begin{equation*}
%P=\begin{bmatrix} \vv_1&\cdots&\vv_p&\wv_{p+1}&\cdots&\wv_n \end{bmatrix}
%\end{equation*}
%be any invertible matrix with \hlist\vv p\ as its first \(p\)~columns.
%Equivalently, write as the partitioned matrix \(P=\begin{bmatrix} V&W \end{bmatrix}\) for corresponding \(n\times p\) and \(n\times(n-p)\) matrices.
%Since the columns of~\(V\) are eigenvectors corresponding to eigenvalue~\(\lambda_j\), \(AV=A\begin{bmatrix} \vv_1&\cdots&\vv_p \end{bmatrix}
%=\begin{bmatrix} A\vv_1&\cdots&A\vv_p \end{bmatrix}
%=\begin{bmatrix} \lambda_j\vv_1&\cdots&\lambda_j\vv_p \end{bmatrix}
%=\lambda_j\begin{bmatrix} \vv_1&\cdots&\vv_p \end{bmatrix}
%=\lambda_jV\).
%Partition the inverse of~\(P\) as \(P^{-1}=\begin{bmatrix} B\\C \end{bmatrix}\) for \(p\times n\) and \((n-p)\times n\) matrices~\(B\) and~\(C\), respectively.
%From this partitioning,
%\begin{eqnarray*}&&
%P^{-1}P=\begin{bmatrix} B\\C \end{bmatrix}\begin{bmatrix} V&W \end{bmatrix}
%=\begin{bmatrix} BV&BW\\CV&CW \end{bmatrix}
%\\\text{and}&&
%P^{-1}P=I_n=\begin{bmatrix} I_p&O\\O&I_{n-p} \end{bmatrix},
%\end{eqnarray*}
%so \(BV=I_p\)\,, \(BW=O\)\,, \(CV=O\) and \(CW=I_{n-p}\)\,.
%
%Now consider
%\begin{eqnarray*}
%P^{-1}AP
%&=&\begin{bmatrix} B\\C \end{bmatrix}A\begin{bmatrix} V&W \end{bmatrix}
%=\begin{bmatrix} BAV&BAW\\CAV&CAW \end{bmatrix}
%\\
%&=&\begin{bmatrix} \lambda_jBV&BAW\\\lambda_jCV&CAW \end{bmatrix}
%=\begin{bmatrix} \lambda_jI_p&BAW\\O&CAW \end{bmatrix}
%\end{eqnarray*}
%Then the characteristic polynomial of matrix~\(A\) becomes
%\begin{eqnarray*}
%&&\hspace{-1.5em}\det(A-\lambda I_n)
%\\&=&\det(PP^{-1}APP^{-1}-\lambda PP^{-1})
%\quad(\text{as }PP^{-1}=I_n)
%\\&=&\det[P(P^{-1}AP-\lambda I_n)P^{-1}]
%\\&=&\det P\det(P^{-1}AP-\lambda I_n)\det(P^{-1})
%\quad(\text{product \cref{thm:detprod}})
%\\&=&\det P\det(P^{-1}AP-\lambda I_n)\frac1{\det P}
%\quad(\text{inverse \cref{thm:detinv}})
%\\&=&\det(P^{-1}AP-\lambda I_n)
%\\&=&\det\begin{bmatrix} (\lambda_j-\lambda)I_p&BAW\\O&CAW -\lambda I_{n-p}\end{bmatrix}
%\quad(\text{by above }PAP^{-1})
%\\&=&(\lambda_j-\lambda)^p\det(CAW -\lambda I_{n-p})
%\end{eqnarray*}
%by \(p\)~successive first column expansions of the determinant (\cref{thm:letdet}).
%Because of the factor \((\lambda_j-\lambda)^p\) in the characteristic polynomial of~\(A\), the eigenvalue~\(\lambda_j\) must have multiplicity of at least~\(p=\dim\EE_{\lambda_j}\)
%(there may be more factors of \((\lambda_j-\lambda)\) hidden within \(\det(CAW -\lambda I_{n-p})\)).
%\end{proof}

\begin{comment}Reviewer comment:
It may be useful to point out that the computation actually establishes the fact that for any square matrix \(A\) and any same-size invertible matrix \(P\), the characteristic polynomials of \(A\) and \(P^{-1}AP\) are the same. (Change of basis formula, same underlying linear transformation, eigenvalues come from the lin. tranf. and not the matrix. But perhaps this is not the right place for this kind of explanation.)
\end{comment}
\begin{proof} 
Suppose \(\lambda_j\) is an eigenvalue of \(n\times n\) matrix~\(A\) and \(\dim\EE_{\lambda_j}=p<n\) (the case \(p=n\) is proved by \cref{ex:dimme}).
Because its dimension is~\(p\) the eigenspace may be spanned by \(p\)~vectors.
Then choose \(p\)~orthonormal vectors \hlist\vv p\ to span~\(\EE_{\lambda_j}\): these vectors \hlist\vv p\ are eigenvectors as they are in the eigenspace.
Let
\begin{equation*}
P=\begin{bmatrix} \vv_1&\cdots&\vv_p&\wv_{p+1}&\cdots&\wv_n \end{bmatrix}
\end{equation*}
be any orthogonal matrix with \hlist\vv p\ as its first \(p\)~columns.
Equivalently, write as the partitioned matrix \(P=\begin{bmatrix} V&W \end{bmatrix}\) for corresponding \(n\times p\) and \(n\times(n-p)\) matrices.
Since \(P\)~is orthogonal, its inverse \(P^{-1}=\tr P=\begin{bmat} \tr V\\\tr W \end{bmat}\).
Since the columns of~\(V\) are eigenvectors corresponding to eigenvalue~\(\lambda_j\), \(AV=A\begin{bmatrix} \vv_1&\cdots&\vv_p \end{bmatrix}
=\begin{bmatrix} A\vv_1&\cdots&A\vv_p \end{bmatrix}
=\begin{bmatrix} \lambda_j\vv_1&\cdots&\lambda_j\vv_p \end{bmatrix}
=\lambda_j\begin{bmatrix} \vv_1&\cdots&\vv_p \end{bmatrix}
=\lambda_jV\).

Now consider
\begin{eqnarray*}
P^{-1}AP
&=&\begin{bmatrix} \tr V\\\tr W \end{bmatrix}A\begin{bmatrix} V&W \end{bmatrix}
=\begin{bmatrix} \tr VAV&\tr VAW\\\tr WAV&\tr WAW \end{bmatrix}
\\
&=&\begin{bmatrix} \lambda_j\tr VV&\tr VAW\\\lambda_j\tr WV&\tr WAW \end{bmatrix}
=\begin{bmatrix} \lambda_jI_p&\tr VAW\\O&\tr WAW \end{bmatrix}
\end{eqnarray*}
where the last equality follows from the orthonormality of columns of \(P=\begin{bmatrix} V&W \end{bmatrix}\).
Then the characteristic polynomial of matrix~\(A\) becomes
\begin{eqnarray*}
&&\hspace{-1.5em}\det(A-\lambda I_n)
\\&=&\det(PP^{-1}APP^{-1}-\lambda PP^{-1})
\quad(\text{as }PP^{-1}=I_n)
\\&=&\det[P(P^{-1}AP-\lambda I_n)P^{-1}]
\\&=&\det P\det(P^{-1}AP-\lambda I_n)\det(P^{-1})
\quad(\text{product \cref{thm:detprod}})
\\&=&\det P\det(P^{-1}AP-\lambda I_n)\frac1{\det P}
\quad(\text{inverse \cref{thm:detinv}})
\\&=&\det(P^{-1}AP-\lambda I_n)
\\&=&\det\begin{bmatrix} (\lambda_j-\lambda)I_p&\tr VAW\\O&\tr WAW -\lambda I_{n-p}\end{bmatrix}
\quad(\text{by above }P^{-1}AP)
\\&=&(\lambda_j-\lambda)^p\det(\tr WAW -\lambda I_{n-p})
\end{eqnarray*}
by \(p\)~successive first column expansions of the determinant (\cref{thm:letdet}).
Because of the factor \((\lambda_j-\lambda)^p\) in the characteristic polynomial of~\(A\), the eigenvalue~\(\lambda_j\) must have multiplicity of at least~\(p=\dim\EE_{\lambda_j}\)---there may be more factors of \((\lambda_j-\lambda)\) hidden within \(\det(\tr WAW -\lambda I_{n-p})\).
\end{proof}



\begin{example} 
Show that the matrix \(A=\begin{bmat}0&5&6
\\-8&22&24
\\6&-15&-16 \end{bmat}\) has one \idx{eigenvalue} of \idx{multiplicity} three, and that the corresponding \idx{eigenspace} has \idx{dimension} two.
%\begin{equation*}
%A=\begin{bmatrix}0&5&6
%\\-8&22&24
%\\6&-15&-16 \end{bmatrix}
%\end{equation*}
\begin{solution} 
Find eigenvalues via the characteristic polynomial
\begin{eqnarray*}
\det(A-\lambda I)
&=&
\begin{vmatrix} -\lambda&5&6
\\-8&22-\lambda&24
\\6&-15&-16-\lambda \end{vmatrix}
\\&=&-\lambda(22-\lambda)(-16-\lambda)+5\cdot24\cdot6+6(-8)(-15)
\\&&{}-6(22-\lambda)6+\lambda24(-15)-5(-8)(-16-\lambda)
\\&\vdots&
\\&=&-\lambda^3+6\lambda^2-12\lambda+8
=-(\lambda-2)^3.
\end{eqnarray*}
This characteristic polynomial is zero only for eigenvalue \(\lambda=2\) which is of multiplicity three.

The corresponding eigenspace comes from solving \((A-\lambda I)\xv=\ov\) which here is
\begin{equation*}
\begin{bmatrix}-2&5&6
\\-8&20&24
\\6&-15&-18 \end{bmatrix}\xv=\ov.
\end{equation*}
Observe that the second row is just four times the first row, and the third row is \((-3)\times{}\)the first row, hence all three equations in this system are equivalent to just the one from the first row, namely \(-2x_1+5x_2+6x_3=0\)\,.
A general solution of this equation is \(x_1=\tfrac52x_2+3x_3\)\,.
That is, all solutions are \(\xv=(\tfrac52x_2+3x_3,x_2,x_3)=x_2(\tfrac52,1,0)+x_3(3,0,1)\). 
Hence all solutions form the two dimensional eigenspace 
\(\EE_2=\Span\{(\tfrac52,1,0)\clb (-3,0,1)\}\).
\end{solution}
\end{example}



\begin{example} 
Use \script\ to find the \idx{eigenvalue}s and the \idx{dimension} of the \idx{eigenspace}s of the matrix 
\begin{equation*}
B=\begin{bmatrix} 344&-1165&-149&-1031&1065&-2816
\\90&-306&-38&-272&280&-742
\\-45&140&12&117&-115&302
\\135&-470&-70&-421&445&-1175
\\-165&555&67&493&-506&1338
\\-105&360&48&322&-335&886 \end{bmatrix}.
\end{equation*}

\begin{solution} 
In \script\ enter the matrix with
\begin{verbatim}
B=[344 -1165  -149 -1031  1065 -2816
    90  -306   -38  -272   280  -742
   -45   140    12   117  -115   302
   135  -470   -70  -421   445 -1175
  -165   555    67   493  -506  1338
  -105   360    48   322  -335   886]
\end{verbatim}
\setbox\ajrqrbox\hbox{\qrcode{% eigenspaces
B=[344 -1165  -149 -1031  1065 -2816
    90  -306   -38  -272   280  -742
   -45   140    12   117  -115   302
   135  -470   -70  -421   445 -1175
  -165   555    67   493  -506  1338
  -105   360    48   322  -335   886]
[V,D]=eig(B)
svd(V(:,1:3))
svd(V(:,4:6))
}}%
\marginajrbox%
Then \verb|[V,D]=eig(B)| computes something like the following \twodp
\begin{verbatim}
V =
  -0.19   0.19  -0.45   0.75  -0.20   0.15
  -0.38   0.38   0.12   0.26  -0.03   0.00
  -0.58   0.58   0.08  -0.00  -0.54   0.56
  -0.00  -0.00  -0.83   0.28  -0.50   0.52
   0.58  -0.58  -0.29  -0.45  -0.64   0.63
   0.38  -0.38   0.08  -0.29  -0.04   0.03
D =
   4.00      0      0      0      0      0
      0   4.00      0      0      0      0
      0      0   4.00      0      0      0
      0      0      0  -1.00      0      0
      0      0      0      0  -1.00      0
      0      0      0      0      0  -1.00
\end{verbatim}
Evidently, the matrix~\(B\) has two eigenvalues, \(\lambda=4\) and \(\lambda=-1\)\,, both of multiplicity three.
Although due to round-off error \script\ reports these with errors of about~\(10^{-5}\) (\cref{sec:reas})---possibly complex errors in which case ignore small complex parts.
For each eigenvalue that \script\ reports, the three corresponding columns of~\verb|V| contain corresponding eigenvectors.
Each set of these three eigenvectors do span the corresponding eigenspace, but are not necessarily \text{linearly independent.}
\begin{itemize}
\item For eigenvalue \(\lambda=4\) the first two columns of~\verb|V| are clearly the negative of each other, and so are essentially the same eigenvector.  
The third column of~\verb|V| is clearly not proportional to the first two columns and so is linearly independent (\cref{thm:lindeplc}).
Thus \script\ has computed only two linearly independent eigenvectors, either the first and third column, or the second and third column.
Consequently, the dimension \(\dim\EE_4=2\)\,.

One can confirm this dimensionality by computing the singular values 
of the first three columns of~\verb|V| with \index{svd()@\texttt{svd()}}\verb|svd(V(:,1:3))| to find they are \(1.4278\), \(0.9806\) and \(0.0000\).
The two nonzero singular values indicate the dimension of the span is two (\cref{pro:ospan}).

\item For eigenvalue \(\lambda=-1\) the last three columns of~\verb|V| look linearly independent and so we suspect the eigenspace dimension \(\dim\EE_{-1}=3\)\,.

%If the dimensionality of an eigenspace is less than the multiplicity of an eigenvalue, then \script\ generally computes two columns as the same or the negative of each other (as for \(\lambda=-4\) here), so that the lower dimensionality is clearly apparent.
%However, this computed result is not guaranteed.
To confirm the dimension of this eigenspace via \cref{pro:ospan}, compute \verb|svd(V(:,4:6))| to find the three singular values are \(1.4136\), \(1.0001\) and \(0.0414\).
Since all three singular values are nonzero, \(\dim\EE_{-1}=3\)\,.
\end{itemize}
\end{solution}
%\begin{verbatim}
%d=diag([-1 -1 -1 4 4 4]);
%d(4,5)=1;%d(5,6)=1;;
%for i=1:999999,p=round(randn(6)*3); if min(abs(abs(det(p))-[1 2 5]))<1e-7, a=p*d*inv(p), break, end, end
%\end{verbatim}
The \script\ function \index{eig()@\texttt{eig()}}\verb|eig()| may produce for you a quite different matrix~\verb|V| of eigenvectors (possibly with complex parts). 
As discussed by \cref{sec:reas}, \idx{repeated eigenvalue}s are very sensitive and this \idx{sensitivity} means small variations in the hidden \script\ algorithm may produce quite large changes in the matrix~\verb|V| for repeated eigenvalues.
However, each eigenspace spanned by the appropriate columns of~\verb|V| \text{is robust.}
\end{example}

\index{diagonalizable|)}





\subsection{Solve systems of differential equations}
\label{sec:ssde}
\index{differential equation|(}
\index{system|(}

\paragraph{Population modelling}
The \idx{population modelling} seen so far (\cref{sec:ddp}) expressed the changes of the population over discrete intervals in time via discrete time equations such as \(y(t+1)=\cdots\) and \(z(t+1)=\cdots\)\,.
One such example is to describe the population numbers year by year.
The alternative is to model the changes in the population \emph{continuously} in time.
This alternative invokes and analyses differential equations.
Such continuous time, differential equation, models are common for exploring the interaction between different \idx{species}, such as between humans and viruses.
Such differential equations are also common throughout engineering \text{and science.}

Let's start with a continuous time version of the \idx{population modelling} discussed at the start of this \cref{ch:gee}.
Let two \idx{species} interact continuously in time with populations \(y(t)\) and~\(z(t)\) at time~\(t\) (years).
Suppose they interact according to \idx{differential equation}s \(dy/dt=y-4z\) and \(dz/dt=-y+z\) (instead of the discrete time equations \(y(t+1)=\cdots\) and \(z(t+1)=\cdots\)).
Analogous to the start of this \cref{sec:dit}, we now ask the following question: is there a matrix transformation to new variables, the vector~\(\Yv(t)\), such that \((y,z)=P\Yv\) for some as yet unknown matrix~\(P\), where the differential equations for~\Yv\ are simple?
Equivalently, is there a different basis \(\cP=\{\pv_1,\pv_2\}\) for the \(yz\)-plane in which the differential equations for \(\Yv=[\yv]_\cP\) \text{are simple?}
\begin{itemize}
\item First, form the differential equations into a matrix-vector system:
\begin{equation*}
\begin{bmatrix} dy/dt\\dz/dt \end{bmatrix}
=\begin{bmatrix} y-4z
\\ -y+z \end{bmatrix}
=\begin{bmatrix} 1&-4\\-1& 1 \end{bmatrix}\begin{bmatrix} y\\z \end{bmatrix}.
\end{equation*}
So using vector \(\yv=(y,z)\), this system is
\begin{equation*}
\D t\yv=A\yv
\quad\text{for matrix }A=\begin{bmatrix} 1&-4\\-1& 1 \end{bmatrix}.
\end{equation*}

\item Second, let's see what happens when we transform to some, as yet unknown, new vector variable~\(\Yv(t)\) such that \(\yv=P\Yv\) for some constant invertible matrix~\(P\).
Under such a transform: \(\D t\yv=\D t{}P\Yv=P\D t\Yv\); also \(A\yv=AP\Yv\).
Hence substituting such an assumed transformation into the matrix-vector form of the differential equation leads to
\begin{equation*}
P\D t\Yv=AP\Yv\,, 
\quad\text{that is}\quad
\D t\Yv=\left(P^{-1}AP\right)\Yv.
\end{equation*}
To simplify this system for~\Yv, we diagonalize the matrix on the right-hand side.
The procedure is to choose the columns of~\(P\) to be \idx{eigenvector}s of the matrix~\(A\) (\cref{thm:gendiag}).

\item Third, find the eigenvectors of~\(A\) by hand as it is a \(2\times2\) matrix. 
Here the matrix \(A=\begin{bmat} 1&-4\\-1&1 \end{bmat}\) has \idx{characteristic polynomial} \(\det(A-\lambda I)=(1-\lambda)^2-4\)\,.
This is zero for \((1-\lambda)^2=4\)\,, that is, \((1-\lambda)=\pm2\).
Hence the \idx{eigenvalue}s \(\lambda=1\pm2=3,-1\)\,.
\begin{itemize}
\item For eigenvalue \(\lambda_1=3\) the corresponding eigenvectors satisfy
\begin{equation*}
(A-\lambda_1I)\pv_1=\begin{bmatrix} -2&-4\\-1&-2 \end{bmatrix}\pv_1=\ov\,,
\end{equation*}
with general solution \(\pv_1\propto(2,-1)\).

\item For eigenvalue \(\lambda_2=-1\) the corresponding eigenvectors satisfy
\begin{equation*}
(A-\lambda_2I)\pv_2=\begin{bmatrix} 2&-4\\-1&2 \end{bmatrix}\pv_2=\ov\,,
\end{equation*}
with general solution \(\pv_2\propto(2,1)\).
\end{itemize}
Thus setting transformation matrix 
\begin{equation*}
P=\begin{bmatrix} \pv_1&\pv_2 \end{bmatrix}
=\begin{bmatrix} 2&2\\-1&1 \end{bmatrix}
\implies \D t\Yv=\begin{bmatrix} 3&0\\0&-1 \end{bmatrix}\Yv
\end{equation*}
(every non-zero scalar multiple of the two columns of~\(P\) also work).

\item Fourth, having diagonalized the matrix, expand this diagonalized set of differential equations to write this system in terms of \idx{components} \(\Yv=(Y_1,Y_2)\): 
\begin{equation*}
\D t{Y_1}=3Y_1 \quad\text{and}\quad \D t{Y_2}=-Y_2\,.
\end{equation*}
Each of these differential equations have well-known \idx{exponential} solutions, respectively \(Y_1=c_1e^{3t}\) and \(Y_2=c_2e^{-t}\),  for every value of the constants~\(c_1\) and~\(c_2\).

\item Lastly, what does this mean for the original problem?
Recall the relation
\begin{equation*}
\begin{bmatrix} y\\z \end{bmatrix}=\yv
=P\Yv=\begin{bmatrix} 2&2\\-1&1 \end{bmatrix}
\begin{bmatrix}c_1e^{3t}\\ c_2e^{-t} \end{bmatrix}
=\begin{bmatrix} 2c_1e^{3t}+2c_2e^{-t}
\\ -c_1e^{3t}+c_2e^{-t}\end{bmatrix}.
\end{equation*}
That is, a general solution of the original system of differential equations is \(y(t)=2c_1e^{3t}+2c_2e^{-t}\) and \(z(t)=-c_1e^{3t}+c_2e^{-t}\).

\end{itemize}
The diagonalization of the matrix empowers us to solve complicated systems of differential equations as a set of simple systems.


Such a \idx{general solution} makes predictions. 
For example, suppose at time zero (\(t=0\)) the \idx{initial population} (\idx{female}) of \(y\)-animals is~\(22\) and the population of \(z\)-animals is~\(9\).
From the above general solution we then know that at time \(t=0\)
\begin{equation*}
\begin{bmatrix} 22\\9 \end{bmatrix}=
\begin{bmatrix} y(0)\\z(0) \end{bmatrix}
=\begin{bmatrix} 2c_1e^{3\cdot0}+2c_2e^{-0}
\\ -c_1e^{3\cdot0}+c_2e^{-0} \end{bmatrix}
=\begin{bmatrix} 2c_1+2c_2
\\ -c_1+c_2 \end{bmatrix}
\end{equation*}

\begin{wrapfigure}r{0pt}
\begin{tikzpicture}
\begin{axis}[small,font=\footnotesize,domain=0:1, 
xlabel={time $t$},ylabel={females},axis lines=middle
,legend pos=outer north east]
\addplot [blue] {2*exp(3*x)+20*exp(-x)};
\addlegendentry{$y(t)$};
\addplot [red] {-exp(3*x)+10*exp(-x)};
\addlegendentry{$z(t)$};
\end{axis}
\end{tikzpicture}
\end{wrapfigure}
This determines the coefficients: \(2c_1+2c_2=22\)
and \(-c_1+c_2=9\)\,.
Adding the first to twice the second gives \(4c_2=40\)\,, that is, \(c_2=10\)\,.
Then either equation determines \(c_1=1\)\,.
Consequently, the \idx{particular solution} from this initial population is
\begin{equation*}
\begin{bmatrix} y\\z \end{bmatrix}
=\begin{bmatrix} 2\cdot1e^{3t}+2\cdot10e^{-t}
\\ -1e^{3t}+10e^{-t}\end{bmatrix}
=\begin{bmatrix} 2e^{3t}+20e^{-t}
\\ -e^{3t}+10e^{-t}\end{bmatrix}.
\end{equation*}
The graph, above-right, of this solution shows that the population of \(y\)-animals grows in time, whereas the population of \(z\)-animals crashes and becomes extinct at about time \(0.6\)~years.%
\footnote{After time \(0.6\)~years the differential equation model  and its predictions becomes meaningless as there is no biological meaning to a negative number of animals~\(z\).}


The forthcoming \cref{thm:ddtsol} confirms that the same approach solves general systems of differential equations: it is analogous to \cref{thm:dynsol} for discrete dynamics.



\begin{activity}
A given \idx{population model} is expressed as the differential equations \(dx/dt=x+y-3z\)\,, \(dy/dt=-2x+z\) and \(dz/dt=-2x+y+2z\)\,.
This may be written in \idx{matrix-vector form} \(d\xv/dt=A\xv\) for vector \(\xv(t)=(x,y,z)\) and which of the following matrices?
\actposs[4]{\(\begin{bmatrix} 1&1&-3
\\-2&0&1
\\-2&1&2 \end{bmatrix}\)}
{\(\begin{bmatrix} 1&-3&1
\\-2&1&0
\\-2&2&1 \end{bmatrix}\)}
{\(\begin{bmatrix} 1&-2&2
\\0&-2&1
\\1&1&-3 \end{bmatrix}\)}
{\(\begin{bmatrix} 1&1&-3
\\0&-2&1
\\1&-2&2 \end{bmatrix}\)}
\end{activity}




\begin{theorem} \label{thm:ddtsol}
Let \(n\times n\) \idx{square matrix}~\(A\) be \idx{diagonalizable} by matrix \(P=\begin{bmatrix} \pv_1&\pv_2&\cdots&\pv_n \end{bmatrix}\) whose columns are \idx{eigenvector}s corresponding to \idx{eigenvalue}s \hlist\lambda n\ of matrix~\(A\).  
Then a \idx{general solution}~\(\xv(t)\) to the \idx{differential equation} system \(d\xv/dt=A\xv\) is the \idx{linear combination}
\begin{equation}
\xv(t)=c_1\pv_1e^{\lambda_1t}+c_2\pv_2e^{\lambda_2t}+\cdots+c_n\pv_ne^{\lambda_nt}
\label{eq:ddtsol}
\end{equation}
for arbitrary constants \hlist cn.
\end{theorem}

\begin{proof} 
First, instead of finding solutions for~\(\xv(t)\) directly, let's write the differential equations in terms of the alternate basis for~\(\RR^n\), basis \(\cP=\{\hlist\pv n\}\) (as \hlist\pv n\ are linearly independent).
That is, solve for the coordinates \(\Xv(t)=[\xv(t)]_\cP\) with respect to basis~\cP.
From \cref{thm:ssbc} recall that  \(\Xv=[\xv]_\cP\) means that \(\xv=\lincomb X\pv n=P\Xv\)\,.
Substitute this into the differential equation \(d\xv/dt=A\xv\) requires \(\D t{}(P\Xv)=A(P\Xv)\) which is the same as \(P\D t\Xv=AP\Xv\)\,.
Since matrix~\(P\) is invertible, this equation is the same as \(\D t\Xv=P^{-1}AP\Xv\)\,.
Because the columns of matrix~\(P\) are eigenvectors, the product \(P^{-1}AP\) is the diagonal matrix \(D=\diag(\hlist\lambda n)\), hence the system becomes \(\D t\Xv=D\Xv\)\,.
Because matrix~\(D\) is diagonal, this is a much simpler system of differential equations.
The \(n\)~rows of the system \(\D t\Xv=D\Xv\) are
\begin{equation*}
\D t{X_1}=\lambda_1X_1\,,\quad
\D t{X_2}=\lambda_2X_2\,,\quad \ldots,\quad
\D t{X_n}=\lambda_nX_n\,.
\end{equation*}
Each of these have general solution
\begin{equation*}
X_1=c_1e^{\lambda_1t},\quad
X_2=c_2e^{\lambda_2t},\quad \ldots,\quad
X_n=c_ne^{\lambda_nt},
\end{equation*}
where \hlist cn\ are arbitrary constants.
To rewrite this solution for the original coordinates~\xv\ use
\begin{equation*}
\xv = P\Xv
=\begin{bmatrix} \pv_1&\pv_2&\cdots&\pv_n \end{bmatrix}
\begin{bmatrix} c_1e^{\lambda_1t}
\\c_2e^{\lambda_2t}
\\\vdots
\\c_ne^{\lambda_nt} \end{bmatrix}
= c_1\pv_1e^{\lambda_1t}+c_2\pv_2e^{\lambda_2t}+\cdots+c_n\pv_ne^{\lambda_nt}
\end{equation*}
to derive the solution~\eqref{eq:ddtsol}.

\begin{comment}
Straightforward substitution is simpler, but the above reinforces the simplification through diagonalization.
\end{comment}
%First show that formula~\eqref{eq:ddtsol} is a solution by substitution into the differential equation \(d\xv/dt=A\xv\)\,: recalling from calculus that \(\D t{}e^{at}=ae^{at}\), the left-hand side
%\begin{eqnarray*}
%\D t\xv&=&c_1\pv_1\D t{}e^{\lambda_1t}+c_2\pv_2\D t{}e^{\lambda_2t}+\cdots+c_n\pv_n\D t{}e^{\lambda_nt}
%\\&=&c_1\pv_1\lambda_1e^{\lambda_1t}+c_2\pv_2\lambda_2e^{\lambda_2t}+\cdots+c_n\pv_n\lambda_ne^{\lambda_nt}\,;
%\end{eqnarray*}
%whereas the right-hand side, as \(\pv_j\) are eigenvectors,
%\begin{eqnarray*}
%A\xv&=&c_1A\pv_1e^{\lambda_1t}+c_2A\pv_2e^{\lambda_2t}+\cdots+c_nA\pv_ne^{\lambda_nt}
%\\&=&c_1\lambda_1\pv_1e^{\lambda_1t}+c_2\lambda_2\pv_2e^{\lambda_2t}+\cdots+c_n\lambda_n\pv_ne^{\lambda_nt}.
%\end{eqnarray*}
%These expressions are the same and hence formula~\eqref{eq:ddtsol} solves \(d\xv/dt=A\xv\)\,.

Second, being able to use the constants \(\cv=(\hlist cn)\) to match every given initial condition shows formula~\eqref{eq:ddtsol} is a general solution.
Suppose the value of~\(\xv(0)\) is given.
Recalling \(e^0=1\)\,, formula~\eqref{eq:ddtsol} evaluated at \(t=0\)  requires
\begin{eqnarray*}
\xv(0)&=&c_1\pv_1e^{\lambda_10}+c_2\pv_2e^{\lambda_20}+\cdots+c_n\pv_ne^{\lambda_n0}
\\&=&c_1\pv_1+c_2\pv_2+\cdots+c_n\pv_n
=\begin{bmatrix} \pv_1&\pv_2&\cdots&\pv_n \end{bmatrix}\cv
=P\cv\,.
\end{eqnarray*}
Since matrix~\(P\) is invertible, choose constants \(\cv=P^{-1}\xv(0)\) for any given~\(\xv(0)\).
%\cref{thm:gendiag,thm:ftim3} establish that such an invertible~\(P\) exists for diagonalizable matrices~\(A\).
\end{proof}





\begin{activity}
Recall that the \idx{differential equation}s \(dy/dt=y-4z\) and \(dz/dt=-y+z\) have a \idx{general solution} \(y(t)=2c_1e^{3t}+2c_2e^{-t}\) and \(z(t)=-c_1e^{3t}+c_2e^{-t}\).
What are the values of these constants given that \(y(0)=2\) and \(z(0)=3\)?
\actposs[4]{\(c_1=-1\), \(c_2=2\)}
{\(c_1=c_2=1\)}
{\(c_1=0\), \(c_2=-1\)}
{\(c_1=-2\), \(c_2=0\)}
\end{activity}




\begin{example} \label{eg:de3t} 
Use matrix analysis to find (by hand) a \idx{general solution} to the system of differential equations \(\D tu=-2u+2v\)\,, \(\D tv=u-2v+w\)\,, and \(\D tw=2v-2w\)\,.

\begin{solution} 
Let vector~\(\uv=(u,v,w)\), and then form the differential equations into the matrix-vector system
\begin{equation*}
\D t\uv=\begin{bmatrix} \D tu\\\D tv\\\D t w \end{bmatrix}
=\begin{bmatrix} -2u+2v\\u-2v+w\\2v-2w \end{bmatrix}
=\begin{bmatrix} -2&2&0\\1&-2&1\\0&2&-2 \end{bmatrix}
\begin{bmatrix} u\\v\\w \end{bmatrix}
=\underbrace{\begin{bmatrix} -2&2&0\\1&-2&1\\0&2&-2 \end{bmatrix}}_A\uv\,.
\end{equation*}
To use \cref{thm:ddtsol} we need eigenvalues and eigenvectors of the matrix~\(A\).
Here the characteristic polynomial of~\(A\) is, using~\eqref{eq:dets23b},
\begin{eqnarray*}
\det(A-\lambda I)&=&\det\begin{bmatrix} -2-\lambda&2&0\\1&-2-\lambda&1\\0&2&-2-\lambda \end{bmatrix}
\\&=&-(2+\lambda)^3+0+0-0+2(2+\lambda)+2(2+\lambda)
\\&=&(2+\lambda)\left[-(2+\lambda)^2+4\right]
\\&=&(2+\lambda)[-\lambda^2-4\lambda]
=-\lambda(\lambda+2)(\lambda+4).
\end{eqnarray*}
This determinant is only zero for eigenvalues \(\lambda=0,-2,-4\)\,.
\begin{itemize}
\item For eigenvalue \(\lambda=0\)\,, corresponding eigenvectors~\pv\ satisfy
\begin{equation*}
(A-0I)\pv=\begin{bmatrix} -2&2&0\\1&-2&1\\0&2&-2 \end{bmatrix}\pv=\ov\,.
\end{equation*}
The last row of this equation requires \(p_3=p_2\), and the first row requires \(p_1=p_2\).
Hence all solutions may be written as \(\pv=(p_2,p_2,p_2)\).
Choose any nonzero one, say \(\pv=(1,1,1)\).

\item For eigenvalue \(\lambda=-2\)\,, corresponding eigenvectors~\pv\ satisfy
\begin{equation*}
(A+2I)\pv=\begin{bmatrix} 0&2&0\\1&0&1\\0&2&0 \end{bmatrix}\pv=\ov\,.
\end{equation*}
The first and last rows of this equation require \(p_2=0\), and the second row requires \(p_3=-p_1\).
Hence all solutions may be written as \(\pv=(p_1,0,-p_1)\).
Choose any nonzero one, say \(\pv=(1,0,-1)\).

\item For eigenvalue \(\lambda=-4\)\,, corresponding eigenvectors~\pv\ satisfy
\begin{equation*}
(A+4I)\pv=\begin{bmatrix} 2&2&0\\1&2&1\\0&2&2 \end{bmatrix}\pv=\ov\,.
\end{equation*}
The last row of this equation requires \(p_3=-p_2\), and the first row requires \(p_1=-p_2\).
Hence all solutions may be written as \(\pv=(-p_2,p_2,-p_2)\).
Choose any nonzero one, say \(\pv=(-1,1,-1)\).

\end{itemize}
With these three distinct eigenvalues, corresponding eigenvectors are linearly independent, and so \cref{thm:ddtsol} gives a general solution of the differential equations as
\begin{equation*}
\begin{bmatrix} u\\v\\w \end{bmatrix}
=c_1\begin{bmatrix} 1\\1\\1 \end{bmatrix}e^{0t}
+c_2\begin{bmatrix} 1\\0\\-1 \end{bmatrix}e^{-2t}
+c_3\begin{bmatrix} -1\\1\\-1 \end{bmatrix}e^{-4t}.
\end{equation*}
That is, \(u(t)=c_1+c_2e^{-2t}-c_3e^{-4t}\), \(v(t)=c_1+c_3e^{-4t}\), and \(w(t)=c_1-c_2e^{-2t}-c_3e^{-4t}\) for every values of the constants~\(c_1\), \(c_2\) and~\(c_3\).
\end{solution}
\end{example}





\begin{example} 
Use the general solution derived in \cref{eg:de3t} to predict the solution of the differential equations \(\D tu=-2u+2v\)\,, \(\D tv=u-2v+w\)\,, and \(\D tw=2v-2w\) given the \idx{initial condition}s that \(u(0)=v(0)=0\) and \(w(0)=4\)\,.
\begin{solution} 
Evaluating the general solution
\begin{equation*}
\begin{bmatrix} u\\v\\w \end{bmatrix}
=c_1\begin{bmatrix} 1\\1\\1 \end{bmatrix}
+c_2\begin{bmatrix} 1\\0\\-1 \end{bmatrix}e^{-2t}
+c_3\begin{bmatrix} -1\\1\\-1 \end{bmatrix}e^{-4t}
\end{equation*}
at time \(t=0\) gives, using the initial conditions and \(e^0=1\)\,,
\begin{equation*}
\begin{bmatrix} 0\\0\\4 \end{bmatrix}
=c_1\begin{bmatrix} 1\\1\\1 \end{bmatrix}
+c_2\begin{bmatrix} 1\\0\\-1 \end{bmatrix}
+c_3\begin{bmatrix} -1\\1\\-1 \end{bmatrix}
=\begin{bmatrix} c_1+c_2-c_3\\c_1+c_3\\c_1-c_2-c_3 \end{bmatrix}.
\end{equation*}


\begin{wrapfigure}r{0pt}
\begin{tikzpicture}
\begin{axis}[footnotesize,font=\footnotesize
,axis lines=middle,xlabel={$t$},no marks,thick
,domain=0:2.2,smooth,legend pos=north east]
\addplot+[]{1-2*exp(-2*x)+exp(-4*x)};
\addlegendentry{$u(t)$};
\addplot+[]{1-exp(-4*x)};
\addlegendentry{$v(t)$};
\addplot+[]{1+2*exp(-2*x)+exp(-4*x)};
\addlegendentry{$w(t)$};
\end{axis}
\end{tikzpicture}
\end{wrapfigure}
Solving by hand, the second row requires \(c_3=-c_1\)\,, so the first row then requires \(c_1+c_2+c_1=0\)\,, that is, \(c_2=-2c_1\)\,.
Putting both of these into the third row requires \(c_1+2c_1+c_1=4\)\,, that is, \(c_1=1\)\,.
Then \(c_2=-2\) and \(c_3=-1\)\,.
Consequently, as drawn to the right, the particular solution is 
\begin{equation*}
\begin{bmatrix} u\\v\\w \end{bmatrix}
=\begin{bmatrix} 1\\1\\1 \end{bmatrix}
-2\begin{bmatrix} 1\\0\\-1 \end{bmatrix}e^{-2t}
-\begin{bmatrix} -1\\1\\-1 \end{bmatrix}e^{-4t}
=\begin{bmatrix} 1-2e^{-2t}+e^{-4t}\\1-e^{-4t}\\1+2e^{-2t}+e^{-4t} \end{bmatrix}.
\end{equation*}
\end{solution}
\end{example}





\begin{example} 
Use \script\ to find a \idx{general solution} to the system of differential equations
\begin{eqnarray*}
&&dx_1/dt=-\tfrac12x_1-\tfrac12x_2+x_3+2x_4\,, 
\\&&dx_2/dt=-\tfrac12x_1-\tfrac12x_2+2x_3+x_4\,, 
\\&&dx_3/dt=\phantom{+2} x_1+2x_2-\tfrac12x_3-\tfrac12x_4\,,  
\\&&dx_4/dt=\phantom{+}2x_1+x_2-\tfrac12x_3-\tfrac12x_4\,.
\end{eqnarray*}
What is the \idx{particular solution} that satisfies the \idx{initial condition}s \(x_1(0)=-5\)\,, \(x_2(0)=-1\) and \(x_3(0)=x_4(0)=0\)\,?
Record your commands and give reasons.
%a=0+round(randn(1,4)*3),A=toeplitz(a([1 end:-1:2]),a),[v,d]=eig(A)
%P=[1 1 1 1;1 1 -1 -1;1 -1 -1 1;1 -1 1 -1]/2, A=P*diag(round(randn(1,4)*2))*P', [V,D]=eig(A)
\begin{solution} 
Write the system in matrix-vector form \(\D t\xv=A\xv\) for vector 
\begin{equation*}
\xv=\begin{bmatrix} x_1\\x_2\\x_3\\x_4 \end{bmatrix}
\quad\text{and matrix }
A=\begin{bmatrix} -\tfrac12 & -\tfrac12 & 1 & 2
\\-\tfrac12 & -\tfrac12 & 2 & 1
\\1 & 2 & -\tfrac12 & -\tfrac12
\\2 & 1 & -\tfrac12 & -\tfrac12 \end{bmatrix}.
\end{equation*}
Enter the matrix into \script\ and then find its eigenvalues and eigenvectors as follows
\begin{verbatim}
A=[-1/2 -1/2 1 2
 -1/2 -1/2 2 1
 1 2 -1/2 -1/2
 2 1 -1/2 -1/2]
[V,D]=eig(A)
\end{verbatim}
\setbox\ajrqrbox\hbox{\qrcode{% DEs
A=[-1/2 -1/2 1 2
 -1/2 -1/2 2 1
 1 2 -1/2 -1/2
 2 1 -1/2 -1/2]
[V,D]=eig(A)
c=V\slosh[-5;-1;0;0]
}}\marginajrbox%
\script\ tells us the eigenvectors and eigenvalues:
\begin{verbatim}
V =
  -0.5000   0.5000  -0.5000  -0.5000
  -0.5000  -0.5000   0.5000  -0.5000
   0.5000   0.5000   0.5000  -0.5000
   0.5000  -0.5000  -0.5000  -0.5000
D =
  -4.0000        0        0        0
        0  -1.0000        0        0
        0        0   1.0000        0
        0        0        0   2.0000
\end{verbatim}
Then \cref{thm:ddtsol} gives that a general solution of the differential equations is
\begin{equation*}
\xv=
c_1\begin{bmatrix} -\tfrac12\\-\tfrac12\\\tfrac12\\\tfrac12 \end{bmatrix}e^{-4t}
+c_2\begin{bmatrix} \tfrac12\\-\tfrac12\\\tfrac12\\-\tfrac12 \end{bmatrix}e^{-t}
+c_3\begin{bmatrix} -\tfrac12\\\tfrac12\\\tfrac12\\-\tfrac12 \end{bmatrix}e^{t}
+c_4\begin{bmatrix} -\tfrac12\\-\tfrac12\\-\tfrac12\\-\tfrac12 \end{bmatrix}e^{2t}.
\end{equation*}
Given the specified initial conditions at \(t=0\)\,, when all the above exponentials reduce to \(e^0=1\)\,, we just need to find the linear combination of the eigenvectors that equals the initial vector \(\xv(0)=(-5,-1,0,0)\); that is, solve \(V\cv=\xv(0)\).
In \script, after checking \verb|rcond(V)| is the good value~\(0.25\), compute \verb|c=V\[-5;-1;0;0]| to find the vector of coefficients is \(\cv=(3,-2,2,3)\).
Hence the particular solution is
\begin{equation*}
\xv=
\begin{bmatrix} -\tfrac32\\-\tfrac32\\\tfrac32\\\tfrac32 \end{bmatrix}e^{-4t}
+\begin{bmatrix} -1\\1\\-1\\1 \end{bmatrix}e^{-t}
+\begin{bmatrix} -1\\1\\1\\-1 \end{bmatrix}e^{t}
+\begin{bmatrix} -\tfrac32\\-\tfrac32\\-\tfrac32\\-\tfrac32 \end{bmatrix}e^{2t}.
\end{equation*}
\end{solution}
\end{example}








\begin{example} \label{eg:decis2t} 
Use matrix analysis to find (by hand) a \idx{general solution} to the system of differential equations
\(\D ty=z\) and \(\D tz=-4y\)\,.
\begin{solution} 
Let vector~\(\yv=(y,z)\), and then form the differential equations into the matrix-vector system
\begin{equation*}
\D t\yv=\begin{bmatrix} \D ty\\\D tz \end{bmatrix}
=\begin{bmatrix} z\\-4y \end{bmatrix}
=\begin{bmatrix} 0&1\\-4&0 \end{bmatrix}
\begin{bmatrix} y\\z \end{bmatrix}
=\underbrace{\begin{bmatrix} 0&1\\-4&0 \end{bmatrix}}_A\yv\,.
\end{equation*}
To use \cref{thm:ddtsol} we need eigenvalues and eigenvectors of the matrix~\(A\).
Here the characteristic polynomial of~\(A\) is
\begin{equation*}
\det(A-\lambda I)=\det\begin{bmatrix} -\lambda&1\\-4&-\lambda \end{bmatrix}
=\lambda^2+4\,.
\end{equation*}
This determinant is only zero for \(\lambda^2=-4\), that is, \(\lambda=\pm2\i\) (where \(\i=\sqrt{-1}\))---a pair of complex conjugate eigenvalues.
\begin{itemize}
\item For eigenvalue \(\lambda=+2\i\) the corresponding eigenvectors~\(\pv\) satisfy
\begin{equation*}
(A-\lambda I)\pv=\begin{bmatrix} -2\i&1\\-4&-2\i \end{bmatrix}\pv
=\ov.
\end{equation*}
The second row of this matrix is \(-2\i\)~times the first row so we just need to satisfy the first row equation \(\begin{bmatrix} -2\i&1 \end{bmatrix}\pv=0\)\,.
This equation is \(-2\i p_1+p_2=0\), that is, \(p_2=2\i p_1\).
Hence all eigenvectors are of the form \(\pv=(1,2\i)p_1\)\,.
Choose any one, say \(\pv=(1,2\i)\).

\item Similarly, for eigenvalue \(\lambda=-2\i\) the corresponding eigenvectors~\(\pv\) satisfy
\begin{equation*}
(A-\lambda I)\pv=\begin{bmatrix} 2\i&1\\-4&2\i \end{bmatrix}\pv
=\ov.
\end{equation*}
The second row of this matrix is \(2\i\)~times the first row so we just need to satisfy the first row equation \(\begin{bmatrix} 2\i&1 \end{bmatrix}\pv=0\)\,.
This equation is \(2\i p_1+p_2=0\), that is, \(p_2=-2\i p_1\).
Hence all eigenvectors are of the form \(\pv=(1,-2\i)p_1\)\,.
Choose any one, say \(\pv=(1,-2\i)\).
\end{itemize}
With these two distinct eigenvalues, corresponding eigenvectors are linearly independent, and so \cref{thm:ddtsol} gives a general solution of the differential equations as
\begin{equation*}
\begin{bmatrix} y\\z \end{bmatrix}
=c_1\begin{bmatrix} 1\\2\i \end{bmatrix}e^{\i2t}
+c_2\begin{bmatrix} 1\\-2\i \end{bmatrix}e^{-\i2t}.
\end{equation*}
That is, \(y(t)=c_1e^{\i2t}+c_2e^{-\i2t}\) and \(z(t)=2\i c_1e^{\i2t}-2\i c_2e^{-\i2t}\) for every constants~\(c_1\) and~\(c_2\).
These formulas answer the exercise.
The next examples show that because of the complex exponentials, this solution describes oscillations in time~\(t\).
\end{solution}
\end{example}




\begin{example} \label{eg:decis2tb} 
Further consider \cref{eg:decis2t}.
Suppose we additionally know that \(y(0)=3\) and \(z(0)=0\)\,.  
Find the \idx{particular solution} that satisfies these two \idx{initial condition}s.
\begin{solution} 
Use the derived general solution that \(y(t)=c_1e^{\i2t}+c_2e^{-\i2t}\) and \(z(t)=2\i c_1e^{\i2t}-2\i c_2e^{-\i2t}\).
We find the constants~\(c_1\) and~\(c_2\) so that these satisfy the  conditions \(y(0)=3\) and \(z(0)=0\)\,.
Substitute \(t=0\) into the general solution to require, using \(e^0=1\)\,,
\begin{eqnarray*}
&&y(0)=3=c_1e^{\i2\cdot0}+c_2e^{-\i2\cdot0}=c_1+c_2\,,
\\&&z(0)=0=2\i c_1e^{\i2\cdot0}-2\i c_2e^{-\i2\cdot0}=2\i c_1-2\i c_2\,,
\end{eqnarray*}
The second of these equations requires \(2\i c_1=2\i c_2\)\,, that is, \(c_1=c_2\)\,.
The first, \(c_1+c_2=3\)\,, then requires that \(2c_1=3\)\,, that is, \(c_1=3/2\) and so \(c_2=3/2\)\,.
Hence the particular solution is
\begin{equation*}
y(t)=\frac32e^{\i2t}+\frac32e^{-\i2t} 
\quad\text{and}\quad
z(t)=3\i e^{\i2t}-3\i e^{-\i2t}.
\end{equation*}

\begin{wrapfigure}r{0pt}
\begin{tikzpicture}
\begin{axis}[footnotesize,font=\footnotesize
,axis lines=middle,xlabel={$t$},no marks,thick
,domain=0:10,smooth,legend pos=south east,legend columns=2]
\addplot+[]{3*cos(2*deg(x))};
\addlegendentry{$y(t)$};
\addplot+[]{-6*sin(2*deg(x))};
\addlegendentry{$z(t)$};
\end{axis}
\end{tikzpicture}
\end{wrapfigure}
But recall \idx{Euler's formula} that \(e^{\i\theta}=\cos\theta+\i\sin\theta\) for every~\(\theta\).
Invoking Euler's formula the above particular solution simplifies:
\begin{eqnarray*}
y(t)&=&\frac32[\cos2t+\i\sin2t]+\frac32[\cos(-2t)+\i\sin(-2t)]
\\&=&\frac32\cos2t+\frac32\i\sin2t+\frac32\cos2t-\frac32\i\sin2t
\\&=&3\cos2t\,,
\\z(t)&=&3\i[\cos2t+\i\sin2t]-3\i[\cos(-2t)+\i\sin(-2t)]
\\&=&3\i\cos2t-3\sin2t-3\i\cos2t-3\sin2t
\\&=&-6\sin2t\,.
\end{eqnarray*}
Because \(y(t)\) and~\(z(t)\) are just trigonometric functions of~\(t\), they oscillate in time~\(t\), as illustrated above-right.
\end{solution}
\end{example}



\begin{example} \label{eg:realcis}
In a real application the \idx{complex number}s of the \idx{general solution} to \cref{eg:decis2t} are usually inconvenient.  
Instead we typically express the solution solely in terms of real quantities as just done in the previous \cref{eg:decis2tb}.
Use \idx{Euler's formula}, that \(e^{\i\theta}=\cos\theta+\i\sin\theta\) for every~\(\theta\), to rewrite the general solution of \cref{eg:decis2t} in terms of \text{real functions.}
\begin{solution}
Here use Euler's formula in the expression for
\begin{eqnarray*}
y(t)&=&c_1e^{\i2t}+c_2e^{-\i2t}
\\&=&c_1[\cos 2t+\i\sin2t]+c_2[\cos(-2t)+\i\sin(-2t)]
\\&=&c_1\cos2t+\i c_1\sin2t +c_2\cos2t-\i c_2\sin2t
\\&=&(c_1+c_2)\cos2t +(\i c_1-\i c_2)\sin2t
\\&=&C_1\cos2t+C_2\sin2t
\end{eqnarray*}
for constants \(C_1=c_1+c_2\) and \(C_2=\i(c_1-c_2)\).
Let's view the arbitrariness in~\(c_1\) and~\(c_2\) as having being `transferred' to~\(C_1\) and~\(C_2\), then \(y(t)=C_1\cos2t+C_2\sin2t\) is a general solution to the differential equations, and is expressed purely in real factors.
In such a real form we explicitly see the oscillations in time~\(t\) through the trigonometric functions \(\cos 2t\) and~\(\sin2t\)\,.

The function \(z(t)=2\i c_1e^{\i2t}-2\i c_2e^{-\i2t}\) has a corresponding form found by replacing~\(c_1\) and~\(c_2\) in terms of the same~\(C_1\) and~\(C_2\).
From above, the constants \(C_1=c_1+c_2\) and \(C_2=\i(c_1-c_2)\).
Adding the first to \(\pm \i\)~times the second determines \(C_1-\i C_2=2c_1\) and \(C_1+\i C_2=2c_2\)\,, so that \(c_1=(C_1-\i C_2)/2\) and \(c_2=(C_1+\i C_2)/2\)\,.  
Then the expression for
\begin{eqnarray*}
z(t)&=&2\i c_1e^{\i2t}-2\i c_2e^{-\i2t}
\\&=&2\i\frac{C_1-\i C_2}2[\cos2t+\i\sin2t]
-2\i\frac{C_1+\i C_2}2[\cos2t-\i\sin2t]
\\&=&(\i C_1+C_2)[\cos2t+\i\sin2t]
+(-\i C_1+C_2)[\cos2t-\i\sin2t]
\\&=&\i C_1\cos2t-C_1\sin2t+C_2\cos2t+\i C_2\sin2t
\\&&{}
-\i C_1\cos2t-C_1\sin2t+C_2\cos2t-\i C_2\sin2t
\\&=&-2C_1\sin2t+2C_2\cos2t\,.
\end{eqnarray*}
That is, the corresponding general solution \(z(t)=-2C_1\sin2t+2C_2\cos2t\) is now also expressed in real factors.
\end{solution}
\end{example}




\begin{reduce}
\begingroup
\def\temp#1{\begin{tikzpicture}
\def\ugh{\ifnum0=#1 axis x line=none
\else axis x line=middle,xlabel={$x$},xtick={0}\fi}
\begin{axis}[footnotesize,font=\footnotesize
, axis y line=none,smooth,legend pos=south west
,domain=0:5.5,samples=40,axis equal,xmax=2,\ugh]
\addplot+[thick,no marks] ({2*x-cos(x*360)-12},{sin(x*360)});
\addlegendentry{spring $k$};
\addplot+[only marks,mark=*,mark size=1ex] coordinates {(0,0)};
\addlegendentry{mass $m$};
\end{axis}
\end{tikzpicture}}
\begin{wrapfigure}r{0pt} \temp0 \end{wrapfigure}
\begin{example}[oscillating applications] \label{eg:sprmas}
A huge variety of vibrating systems are analogous to the basic \idx{oscillation}s of a mass on a spring, 
illustrated schematically to the right.
The mass generally oscillates to and fro.  
Describe such a system mathematically with two differential equations, and solve the differential equations to confirm \text{it oscillates.}

\begin{wrapfigure}r{0pt} \temp1 \end{wrapfigure}
\begin{solution} 
At every time~\(t\), let the position of the mass relative to its rest position be denoted by~\(x(t)\): that is, we put an \(x\)-axis on the picture with \(x=0\) where the mass would stay at rest, as illustrated.
At every time~\(t\) let the mass be moving with velocity~\(v(t)\) (positive to the right, negative to the left).
Then we know one differential equation, that \(\D tx=v\).

\idx{Newton's law}, that \(\text{mass}\times\text{acceleration}=\text{force}\), provides another differential equation.
Here the mass is denoted by~\(m\), and the acceleration is~\(\D tv\).
The force on the mass comes from the spring: typically the force by the spring is proportional to the stretching of the spring, namely to~\(x\).
Different springs give different strength forces so let's denote the constant of proportionality by~\(k\)---a constant that differs depending upon the spring.
Then the force is~\(-kx\) as springs try to pull\slash push the mass back towards \(x=0\)\,.
Consequently Newton's law gives us the differential equation, \(m\D tv=-kx\)\,.  
Divide by mass~\(m\) and the differential equation is \(\D tv=-\frac kmx\)\,.

Write these two differential equations together as a matrix-vector system: 
\begin{equation*}
\begin{bmatrix} \D tx\\\D tv \end{bmatrix}
=\begin{bmatrix} v\\-\frac kmx \end{bmatrix},
\quad\text{that is,}\quad
\D t{}\begin{bmatrix} x\\v \end{bmatrix}
=\begin{bmatrix} 0&1\\-\frac km&0 \end{bmatrix}\begin{bmatrix} x\\v \end{bmatrix}.
\end{equation*}
\cref{thm:ddtsol} asserts a general solution comes from the eigenvalues and eigenvectors of the matrix.
\begin{itemize}
\item Here the characteristic polynomial of the matrix is
\begin{equation*}
\det\begin{bmatrix} -\lambda&1\\-\frac km&-\lambda \end{bmatrix}
=\lambda^2+\frac km\,.
\end{equation*}
This polynomial is zero only when the eigenvalues \(\lambda=\pm\sqrt{-k/m}=\pm \i\sqrt{k/m}\)\,: these are a complex conjugate pair of pure imaginary eigenvalues.

\item  The corresponding eigenvectors~\pv\ satisfy
\begin{equation*}
\begin{bmatrix} \mp \i\sqrt{k/m}&1\\- k/m&\mp \i\sqrt{k/m} \end{bmatrix}
\pv=\ov.
\end{equation*}
The two rows of this equation are satisfied by the corresponding eigenvectors \(\pv\propto (1,\pm \i\sqrt{k/m})\).
\end{itemize}
Then a general solution to the system of differential equations is
\begin{equation*}
\begin{bmatrix} x\\v \end{bmatrix}
=c_1\begin{bmatrix} 1\\\i\sqrt{k/m} \end{bmatrix}e^{\i\sqrt{k/m}t}
+c_2\begin{bmatrix} 1\\-\i\sqrt{k/m} \end{bmatrix}e^{-\i\sqrt{k/m}t}.
\end{equation*}
This formula shows that the mass on the spring generally oscillates as the complex exponentials are oscillatory.

However, in real applications we usually prefer a real algebraic expression.
Just as in \cref{eg:realcis}, we make the above formula real by changing from (complex) arbitrary constants~\(c_1\) and~\(c_2\) to new (real) arbitrary constants~\(C_1\) and~\(C_2\) where \(c_1=(C_1-\i C_2)/2\) and \(c_2=(C_1+\i C_2)/2\).
Substitute these relations into the above general solution, and using Euler's formula, gives that the position
\begin{eqnarray*}
x(t)&=&c_1e^{\i\sqrt{k/m}t}+c_2e^{-\i\sqrt{k/m}t}
\\&=&\frac{C_1-\i C_2}2\big[\cos(\sqrt{k/m}t)+\i\sin(\sqrt{k/m}t)\big]
+\frac{C_1+\i C_2}2\big[\cos(\sqrt{k/m}t)-\i\sin(\sqrt{k/m}t)\big]
\\&=&
\frac{C_1}2\cos(\sqrt{k/m}t)+\i\frac{C_1}2\sin(\sqrt{k/m}t)
-\i\frac{C_2}2\cos(\sqrt{k/m}t)+\frac{C_2}2\sin(\sqrt{k/m}t)
\\&&{}
+\frac{C_1}2\cos(\sqrt{k/m}t)-\i\frac{C_1}2\sin(\sqrt{k/m}t)
+\i\frac{C_2}2\cos(\sqrt{k/m}t)+\frac{C_2}2\sin(\sqrt{k/m}t)
\\&=&C_1\cos(\sqrt{k/m}t)+C_2\sin(\sqrt{k/m}t).
\end{eqnarray*}
For every value of the arbitrary constants~\(C_1\) and~\(C_2\), this formula describes the position~\(x(t)\) of the mass as purely real oscillations in time.
Similarly for the velocity~\(v(t)\) (\cref{ex:sprmas}).
\end{solution}
\end{example}
\endgroup
\end{reduce}




\index{differential equation|)}
\index{system|)}



\begin{comment}
Probable section includes Cayley--Hamilton theorem: recall \cref{sec:mpmev}, compute powers of matrices, then C-H.
Other applications include Markov chains.

Possibly develop a little more theory on similarity of matrices and coordinate transform of matrices.  Then Jordan form and solving non-diagonalizable systems of differential equations.

Possibly something on {Emergent quasi-stationary dynamics of metastable states} but perhaps done in 7.1.2--3. (Can one have too many buzz-words?)
\end{comment}















\sectionExercises


\begin{exercise} \label{ex:whichp} 
Which of the following matrices diagonalize the matrix \(Z=\begin{bmat} 7&12\\-2&-3 \end{bmat}\)?  
Show your working.
%\begin{verbatim}
%n=2
%for i=1:9999,P=0+round(randn(n)*2);if abs(det(P))==1,break,end,end
%P=P,D=diag(round(randn(1,n)*2)),A=round(P*D*inv(P))
%\end{verbatim}
\begin{Parts}
\item \(P_{\alph{enumii}}=\begin{bmatrix} -2&3\\1&-1 \end{bmatrix}\)
\answer{yes}

\item \(P_{\alph{enumii}}=\begin{bmatrix} 3&-2\\-1&1 \end{bmatrix}\)
\answer{yes}

\item \(P_{\alph{enumii}}=\begin{bmatrix} 1&-1\\-2&3 \end{bmatrix}\)
\answer{no}

\item \(P_{\alph{enumii}}=\begin{bmatrix} 1&3\\1&2 \end{bmatrix}\)
\answer{no}

\item \(P_{\alph{enumii}}=\begin{bmatrix} 4&3\\-2&-1 \end{bmatrix}\)
\answer{yes}

\item \(P_{\alph{enumii}}=\begin{bmatrix} -2&3\\2&-2 \end{bmatrix}\)
\answer{no}

\item \(P_{\alph{enumii}}=\begin{bmatrix} -1&1\\3&-2 \end{bmatrix}\)
\answer{no}

\item \(P_{\alph{enumii}}=\begin{bmatrix} 3&1\\2&1 \end{bmatrix}\)
\answer{no}

\end{Parts}
\end{exercise}



\begin{exercise} \label{ex:whichps} 
Redo \cref{ex:whichp} by finding which matrices~\(P_a,\ldots,P_h\) diagonalize each of the following matrices.
\begin{Parts}
\item \(\eAii=\begin{bmatrix} 5&12\\-2&-5 \end{bmatrix}\)
\item \(\eAii=\begin{bmatrix} -3&-12\\2&7 \end{bmatrix}\)
\item \(\eAii=\begin{bmatrix} -1&-1\\6&4 \end{bmatrix}\)
\item \(\eAii=\begin{bmatrix} 4&6\\-1&-1 \end{bmatrix}\)
\end{Parts}
\end{exercise}




\begin{exercise}  
Following \cref{eg:diagonaliseb}, prove that for every scalar~\(k\) the matrix \(\begin{bmat} k&1\\0&k \end{bmat}\) is not diagonalizable.
\end{exercise}




\begin{exercise}  
In each of the following cases, you are given three \idx{linearly independent} \idx{eigenvector}s and corresponding \idx{eigenvalue}s for some \(3\times3\) matrix~\(A\).
Write down three different matrices~\(P\) that diagonalize the matrix~\(A\), and for each write down the corresponding \idx{diagonal matrix} \(D=P^{-1}AP\).
\begin{enumerate}
\item \(\lambda_1=-1\), \(\pv_1=(3,2,-1)\);
\(\lambda_2=1\), \(\pv_2=(-4,-2,2)\);
\(\lambda_3=3\), \(\pv_3=(-1,0,2)\).

\item \(\lambda_1=-1\), \(\pv_1=( 2,1,2)\);
\(\lambda_2=-1\), \(\pv_2=(0,3,1)\);
\(\lambda_3=2\), \(\pv_3=(4,-2,-2)\).

\item \(\lambda_1=1\), \(\pv_1=(3,-7,2)\);
\(\lambda_2=-2\), \(\pv_2=(-4,-5,1)\);
\(\lambda_3=4\), \(\pv_3=(1,2,3)\).

\begin{reduce}
\item \(\lambda_1=3\), \(\pv_1=(2,-3,0)\);
\(\lambda_2=1\), \(\pv_2=(-1,2,-6)\);
\(\lambda_3=-4\), \(\pv_3=(-2,-1,-3)\).
\end{reduce}

\end{enumerate}
\end{exercise}





\begin{exercise}  
From the given information, are each of the matrices \idx{diagonalizable}?  
Give reasons.
\begin{enumerate}
\item The only eigenvalues of a \(2\times2\) matrix are \(2.2\) and~\(0.1\).
\answer{yes}

\item The only eigenvalues of a \(4\times4\) matrix are \(2.2\), \(1.9\), \(-1.8\) and~\(-1\).
\answer{yes}

\item The only eigenvalue of a \(2\times2\) matrix is~\(0.7\).
\answer{unknown}

\item The only eigenvalues of a \(6\times6\) matrix are \(-1.6\), \(0.3\), \(0.1\) and~\(-2.3\).
\answer{unknown}

\begin{reduce}
\item The only eigenvalues of a \(5\times5\) matrix are \(-1.7\), \(1.4\), \(1.3\), \(2.4\), \(0.5\) and~\(-2.3\).
\answer{error}

\item The only eigenvalues of a \(6\times6\) matrix are \(1.2\), \(-0.9\), \(-0.8\), \(2.2\), \(0.2\) and~\(-0.2\).
\answer{yes}

\item The \script\ function \verb|eig(A)| returns the result
\begin{verbatim}
ans =
   2.6816
  -0.1445
   0.0798
   0.3844
\end{verbatim}
\answer{yes}
\end{reduce}

\item The \script\ function \verb|eig(A)| returns the result
\begin{verbatim}
ans =
   3.0821 + 0.0000i
  -2.7996 + 0.0000i
  -0.7429 + 1.6123i
  -0.7429 - 1.6123i
\end{verbatim}
\answer{yes}

\item The \script\ function \verb|eig(A)| returns the result
\begin{verbatim}
ans =
  -1.0000
   1.0000
   2.0000
  -1.0000
\end{verbatim}
\answer{unknown}

\end{enumerate}
\end{exercise}




\begin{exercise}  
For each of the following \(3\times3\) matrices, show with hand algebra that each matrix has one \idx{eigenvalue} of \idx{multiplicity} three, and then determine the \idx{dimension} of the corresponding \idx{eigenspace}.
Also compute the eigenvalue and eigenvectors with \script: comment on any limitations in the computed `eigenvalues' and `eigenvectors'.
%\begin{verbatim}
%d=eye(3)*round(randn)+[0 round(rand) 0;0 0 round(rand); 0 0 0], for i=1:99999,p=round(randn(3)*3); if abs(det(p))==1, a=round(p*d*inv(p))+0, break, end, end, [v,d]=eig(a)
%\end{verbatim}
\begin{Parts}
\item \(\eAii=\begin{bmatrix} 2 & 0 & 0
\\0 & 2 & 0
\\0 & 0 & 2 \end{bmatrix}\)
\answer{\(\lambda=2\), three; all good}

\item \(\eAii=\begin{bmatrix} 2 & 1 & -2
\\-8 & -4 & 8
\\-2 & -1 & 2 \end{bmatrix}\)
\answer{\(\lambda=0\), two; errors~\(10^{-8}\) and two  eigenvectors effectively the same.}

\item \(\eAii=\begin{bmatrix} -2 & 1 & 0
\\0 & -2 & 0
\\0 & 1 & -2 \end{bmatrix}\)
\answer{\(\lambda=-2\), two; not lin.~indep.}

\item \(\eAii=\begin{bmatrix} 7 & 1 & 22
\\-6 & -1 & -20
\\-1 & 0 & -3 \end{bmatrix}\)
\answer{\(\lambda=1\), one; errors~\(10^{-5}\), all three eigenvectors effectively the same}

\begin{reduce}
\item \(\eAii=\begin{bmatrix} 3 & 5 & -1
\\-4 & -6 & 1
\\-5 & -6 & 0 \end{bmatrix}\)
\answer{\(\lambda=-1\), one; errors~\(10^{-6}\), all three eigenvectors \(\pm\)same}

\item \(\eAii=\begin{bmatrix} -6 & -5 & 3
\\-4 & -7 & 3
\\-12 & -15 & 7 \end{bmatrix}\)
\answer{\(\lambda=0\), two; errors~\(10^{-8}\) and two  eigenvectors effectively the same.}
\end{reduce}

\item \(\eAii=\begin{bmatrix} -3 & 2 & 7
\\2 & -3 & -9
\\-1 & 1 & 3 \end{bmatrix}\)
\answer{\(\lambda=-1\), one;  errors~\(10^{-6}\), all three eigenvectors \(\pm\)same}

\item \(\eAii=\begin{bmatrix} -1 & 0 & 0
\\0 & -1 & 0
\\0 & 0 & -1 \end{bmatrix}\)
\answer{\(\lambda=-1\), three; all good}


\end{Parts}
\end{exercise}







\begin{exercise} \label{ex:dimme} 
For a given \(n\times n\) square matrix~\(A\), suppose \(\lambda_1\) is an \idx{eigenvalue} with \(n\)~corresponding \idx{linearly independent} \idx{eigenvector}s \hlist\pv n\,.  
Adapt parts of the proof of \cref{thm:dimee} to prove that the \idx{characteristic polynomial} of matrix~\(A\) is \(\det(A-\lambda I)=(\lambda_1-\lambda)^n\). 
Then deduce that \(\lambda_1\)~is the only eigenvalue of~\(A\), and is of \idx{multiplicity}~\(n\).
\end{exercise}




\begin{exercise}  
For each of the following matrices, use \script\ to find the \idx{eigenvalue}s, their \idx{multiplicity}, and the \idx{dimension} of the \idx{eigenspace}s.  
Give reasons (remember computational error).
%\begin{verbatim}
%d=0+round(randn(1,2)*2);
%d=diag([d(1) d(1) d(1) d(2) d(2)]);
%d([2 6 8 12 20 24])=0+round(rand(6,1)*0.75).*[-1;1;-1;1;-1;1]
%for i=1:99999,p=round(randn(5)*3); if min(abs(abs(det(p))-[1 2 5]))<1e-7, a=p*d*inv(p), break, end, end,[V,D]=eig(a)
%\end{verbatim}
\begin{enumerate}
\item \(\eAii=\begin{bmatrix} -30.5 & 25.5 & 22 & -5 & -48.5
\\-111 & 88 & 76 & -20 & -173
\\87.5 & -70.5 & -61 & 16 & 136.5
\\51 & -43 & -36 & 9 & 81
\\-4.5 & 2.5 & 2 & -1 & -6.5 \end{bmatrix}\)
\setbox\ajrqrbox\hbox{\qrcode{% eig
A=[-30.5 25.5 22 -5 -48.5
 -111 88 76 -20 -173
 87.5 -70.5 -61 16 136.5
 51 -43 -36 9 81
 -4.5 2.5 2 -1 -6.5]
}}\marginajrbox%
\answer{\(\lambda=1\) twice, \(\dim\EE_{1}=1\); \(\lambda=-1\) thrice, \(\dim\EE_{-1}=2\)}

\begin{reduce}
\item \(\eAii=\begin{bmatrix} 2929 & 1140 & -1359 & 2352 & 406
\\-2441 & -950 & 1132 & -1960 & -338
\\-1070 & -416 & 495 & -858 & -148
\\-2929 & -1140 & 1359 & -2352 & -406
\\-895 & -347 & 412 & -715 & -124 \end{bmatrix}\)
\setbox\ajrqrbox\hbox{\qrcode{% eig
A=[2929 1140 -1359 2352 406
 -2441 -950 1132 -1960 -338
 -1070 -416 495 -858 -148
 -2929 -1140 1359 -2352 -406
 -895 -347 412 -715 -124]
}}\marginajrbox%
\answer{\(\lambda=-1\) twice, \(\dim\EE_{-1}=1\); \(\lambda=0\) thrice, \(\dim\EE_{0}=2\)}

\item \(\eAii=\begin{bmatrix} -168 & 115 & 305 & -120 & 70
\\-4710 & 3202 & 8510 & -3360 & 1990
\\450 & -305 & -813 & 320 & -190
\\-4545 & 3090 & 8205 & -3243 & 1920
\\-2415 & 1640 & 4355 & -1720 & 1017 \end{bmatrix}\)
\setbox\ajrqrbox\hbox{\qrcode{% eig
A=[-168 115 305 -120 70
 -4710 3202 8510 -3360 1990
 450 -305 -813 320 -190
 -4545 3090 8205 -3243 1920
 -2415 1640 4355 -1720 1017]
}}\marginajrbox%
\answer{\(\lambda=2\) twice, \(\dim\EE_{2}=2\); \(\lambda=-3\) thrice, \(\dim\EE_{-3}=3\)}
\end{reduce}

\item \(\eAii=\begin{bmatrix} -8 & 744 & -564 & 270 & 321
\\1 & -43 & 33 & -15 & -19
\\1 & 29 & -21 & 12 & 12
\\5 & -431 & 327 & -156 & -186
\\-5 & 535 & -405 & 195 & 231 \end{bmatrix}\)
\setbox\ajrqrbox\hbox{\qrcode{% eig
A=[-8 744 -564 270 321
 1 -43 33 -15 -19
 1 29 -21 12 12
 5 -431 327 -156 -186
 -5 535 -405 195 231]
}}\marginajrbox%
\answer{\(\lambda=0\) twice, \(\dim\EE_{0}=2\); \(\lambda=1\) thrice, \(\dim\EE_{1}=2\)}

\item \(\eAii=\begin{bmatrix}15.6 & 31.4 & -11.6 & -14.6 & -10.6
\\-16.6 & -32.4 & 11.6 & 14.6 & 10.6
\\33.6 & 56.4 & -19.6 & -27.6 & -15.6
\\38.6 & 75.4 & -28.6 & -35.6 & -26.6
\\-122 & -226 & 82 & 106 & 73 \end{bmatrix}\)
\setbox\ajrqrbox\hbox{\qrcode{% eig
A=[15.6 31.4 -11.6 -14.6 -10.6
 -16.6 -32.4 11.6 14.6 10.6
 33.6 56.4 -19.6 -27.6 -15.6
 38.6 75.4 -28.6 -35.6 -26.6
 -122 -226 82 106 73]
}}\marginajrbox%
\answer{\(\lambda=2\) twice, \(\dim\EE_{2}=1\); \(\lambda=-1\) thrice, \(\dim\EE_{-1}=3\)}

\item \(\eAii=\begin{bmatrix}-208 & 420 & -518 & 82 & 264
\\655 & -1336 & 1642 & -260 & -838
\\229 & -467 & 574 & -91 & -294
\\-171 & 348 & -428 & 66 & 218
\\-703 & 1431 & -1760 & 279 & 896 \end{bmatrix}\)
\setbox\ajrqrbox\hbox{\qrcode{% eig
A=[-208 420 -518 82 264
 655 -1336 1642 -260 -838
 229 -467 574 -91 -294
 -171 348 -428 66 218
 -703 1431 -1760 279 896]
}}\marginajrbox%
\answer{\(\lambda=-1\pm \i\) once each, \(\dim\EE_{-1\pm \i}=1\); \(\lambda=-2\) thrice, \(\dim\EE_{-2}=2\)}

\end{enumerate}
\end{exercise}






\begin{exercise}  
For each of the following systems of differential equations, find \idx{eigenvalue}s and \idx{eigenvector}s to derive a \idx{general solution} of the system of differential equations.  Show your working.
%\begin{verbatim}
%n=3,d=diag(0+round(randn(1,n)*(5-n)));for i=1:99999,p=round(randn(n)*3); if min(abs(abs(det(p))-[1 2 5 10]))<1e-7, a=round(p*d*inv(p)*100)/100, break, end, end,[V,D]=eig(a),V*diag(1./max(abs(V)))
%\end{verbatim}

\begin{Parts}
\item \(\D tx=x-1.5y\), \(\D ty=4x-4y\)
\answer{\((x,y)=c_1(3,4)e^{-t}+c_2(1,2)e^{-2t}\)}

\item \(\D tx=x\), \(\D ty=-12x+5y\)
\answer{\((x,y)=c_1(0,1)e^{5t}+c_2(1,3)e^{t}\)}

\item \(\D tx=7x-3y\)
\answer{Not possible: only one equation for two unknowns.}

\begin{reduce}
\item \(\D tu=2.8u-3.6v\), \(\D tv=-0.6u+2.2v\)
\answer{\((u,v)=c_1(3,-1)e^{4t}+c_2(2,1)e^{t}\)}

\item \(\D tp=14p+16q\), \(\D tq=-8p-10q\)
\answer{\((p,q)=c_1(2,-1)e^{6t}+c_2(-1,1)e^{-2t}\)}

\item \(\D tx=6.5x-0.6y-5.7z\), \(\D ty=-3x+4.4y+7.8z\)%, \(\D tz=1.5x-4.2y-7.9z\)
\answer{Not possible: only two equations for three unknowns.}

\item \(\D tx=-31x+26y-24z\), \(\D ty=-48x+39y-36z\), \(\D tz=-14x+10y-9z\)
\answer{\((x,y,z)=c_1(3,3,-1)e^{3t}+c_2(1,2,1)e^{-3t}+c_3(1,3,2)e^{-t}\)}
\end{reduce}

\item \(\D tx=0.2x+1.2z\), \(\D ty=-x\), \(\D tz=1.8x+0.8z\)
\answer{\((x,y,z)=c_1(0,1,0) +c_2(1,1,-1)e^{-t} +c_3(-2,1,3)e^{2t}\)}

\item \(\D tu=4.5u+7.5v+7.5w\), \(\D tv=3u+4v+5w\), \(\D tw=-7.5u-11.5v-12.5w\)
\answer{\((u,v,w)=c_1(-1,-1,2)e^{-3t} +c_2(-5,0,3) +c_3(0,1,-1)e^{-t}\)}

\item \(\D tp=-13p+30q+6r\), \(\D tq=-32p+69q+14r\), \(\D tr=125p-265q-54r\)
\answer{\((p,q,r)=c_1(2,2,-5)e^{2t}+c_2(3,1,2)e^{t}+c_3(0,-1,5)e^{-t}\)}

\end{Parts}
\end{exercise}





\begin{exercise}  
In each of the following, a general solution to a differential equation is given.  
Find the \idx{particular solution} that satisfies the specified \idx{initial condition}s.
Show your working.
\begin{enumerate}
\item \((x,y)=c_1(0,1)e^{-t}+c_2(1,3)e^{2t}\) where \(x(0)=2\) and \(y(0)=1\)
\answer{\((x,y)=-5(0,1)e^{-t}+2(1,3)e^{2t}\)}

\item \((x,y)=c_1(0,1)e^{-2t}+c_2(1,3)e^{t}\) where \(x(0)=0\) and \(y(0)=2\)
\answer{\((x,y)=(0,2)e^{-2t}\)}

\begin{reduce}
\item \(x=3c_1e^t+c_2e^{-t}\), \(y=5c_1e^t+2c_2e^{-t}\) where \(x(0)=0\) and \(y(0)=2\)
\answer{\(x=-6e^t+6e^{-t}\), \(y=-10e^t+12e^{-t}\)}

\item \(x=3c_1e^t+c_2e^{2t}\), \(y=5c_1e^t+2c_2e^{2t}\) where \(x(0)=1\) and \(y(0)=3\)
\answer{\(x=-3e^t+4e^{2t}\), \(y=-5e^t+8e^{2t}\)}
\end{reduce}

\item \((x,y,z)=c_1(0,0,1) +c_2(1,-1,1)e^{2t} +c_3(-2,3,1)e^{-2t}\) where \(x(0)=3\), \(y(0)=-4\) and \(z(0)=0\)
\answer{\((x,y,z)=(1,-1,1)e^{2t} +(2,-3,-1)e^{-2t}\)}

\item \((x,y,z)=c_1(0,0,1)e^{-3t} +c_2(1,-1,1) +c_3(-2,3,1)e^{-t}\) where \(x(0)=3\), \(y(0)=-4\) and \(z(0)=2\)
\answer{\((x,y,z)=(0,0,2)e^{-3t} +(1,-1,1) +(2,-3,-1)e^{-t}\)}

\end{enumerate}
\end{exercise}







\begin{exercise}  
For each of the following systems of differential equations, use \script\ to find \idx{eigenvalue}s and \idx{eigenvector}s and hence derive a \idx{general solution} of the system of differential equations.  
Record your working.
%\begin{verbatim}
%n=4,d=diag(0+round(randn(1,n)*10)/10); if rand<0.5,d(1:2,1:2)=[d(1) -d(2,2);d(2,2) d(1)],end,
%for i=1:99999,p=round(randn(n)*3); if min(abs(abs(det(p))-[1]))<1e-7, a=round(p*d*inv(p)*1e4)/1e4, break, end, end,
%[V,D]=eig(a),vd=num2str(V',2),dd=num2str(D',2)
%\end{verbatim}

\begin{enumerate}
\item \(\begin{array}{l} 
\D t{x_1}= 15.8x_1 +17.1x_2 -119.7x_3 +153.9x_4\,,\\
\D t{x_2}= 1.4x_1 +0.1x_2 +12.9x_3 -17.1x_4\,,\\
\D t{x_3}= 6.2x_1 +6.2x_2 -43x_3 +57.6x_4\,,\\
\D t{x_4}= 3.4x_1 +3.4x_2 -25x_3 +34x_4\,.
\end{array}\)
\setbox\ajrqrbox\hbox{\qrcode{% DE
A=[15.8 17.1 -119.7 153.9
 1.4 0.1 12.9 -17.1
 6.2 6.2 -43.0 57.6
 3.4 3.4 -25.0 34.0 ]
}}\marginajrbox%
\answer{\twodp\ \(\xv
=c_1(0.71,-0.71,0,0)e^{-1.3t} 
+c_2(-0.63,-0.63,-0.42,-0.21)e^{4.4t} 
+c_3(0,0.64,0.64,0.43)e^{1.6t} 
+c_4(-0.82,0.41,0.41)e^{2.2t}\)}

\item \(\begin{array}{l} 
\D t{x_1}= -12.2x_1 +53.7x_2 +50.1x_3 -22.8x_4\,,\\
\D t{x_2}= 0.6x_1 -2.3x_2 -2.7x_3 +1.2x_4\,,\\
\D t{x_3}= -20.4x_1 +93.8x_2 +90.2x_3 -40.8x_4\,,\\
\D t{x_4}= -38.2x_1 +177.9x_2 +170.7x_3 -77.2x_4\,.
\end{array}\)
\setbox\ajrqrbox\hbox{\qrcode{% DE
A=[ -12.2 53.7 50.1 -22.8
 0.6 -2.3 -2.7 1.2
 -20.4 93.8 90.2 -40.8
 -38.2 177.9 170.7 -77.2 ]
}}\marginajrbox%
\answer{\twodp\ \(\xv
=c_1(0.63,0.63,-0.42,0.21)e^{0.5t} 
+c_2(-0.23,-0.23,-0.23,-0.92)e^{0.4t} 
+c_3(0.64,0,0.43,0.64)e^{-1.6t} 
+c_4(-0.89,0,0,0.45)e^{0.8t}\)}

\item \(\begin{array}{l} 
\D t{x_1}= x_1 +29.4x_2 -3.2x_3 -12.9x_4\,,\\
\D t{x_2}= 1.4x_1 -38.4x_2 +5.6x_3 +18.2x_4\,,\\
\D t{x_3}= 2.3x_1 -80.3x_2 +12.3x_3 +36.7x_4\,,\\
\D t{x_4}= 2.4x_1 -65x_2 +9.2x_3 +31.1x_4\,.
\end{array}\)
\setbox\ajrqrbox\hbox{\qrcode{% DE
A=[ 1.0 29.4 -3.2 -12.9
 1.4 -38.4 5.6 18.2
 2.3 -80.3 12.3 36.7
 2.4 -65.0 9.2 31.1 ]
}}\marginajrbox%
\answer{\twodp\ \(\xv
=c_1(-0.63,0.32,0.32,0.63)e^{0.8t} 
+c_2(0.8,0.27,0,0.53)e^{2.2t} 
+c_3(0.76,0,-0.49-0.4\i,0.091+0.12\i)e^{(1.5-0.4\i)t} 
+c_4(0.76,0,-0.49+0.4\i,0.091-0.12\i)e^{(1.5+0.4\i)t}\)}

\begin{reduce}
\item \(\begin{array}{l}
\D t{x_1}= -50.2x_1 -39.5x_2 -20.2x_3 +68.9x_4\,,\\
\D t{x_2}= 62.8x_1 +50.2x_2 +28.4x_3 -85.2x_4\,,\\
\D t{x_3}= -17.3x_1 -13.7x_2 -8.9x_3 +22.7x_4\,,\\
\D t{x_4}= -6.4x_1 -4.6x_2 -1.4x_3 +9x_4\,.
\end{array}\)
\setbox\ajrqrbox\hbox{\qrcode{% DE
A=[-50.2 -39.5 -20.2 68.9
 62.8 50.2 28.4 -85.2
 -17.3 -13.7 -8.9 22.7
 -6.4 -4.6 -1.4 9 ]
}}\marginajrbox%
\answer{\twodp\ \(\xv
=c_1(-0.78,0.38+0.3\i,0.095+0.076\i,-0.32+0.21\i)e^{(0.1-1.1\i)t} 
+c_2(-0.78,0.38-0.3\i,0.095-0.076\i,-0.32-0.21\i)e^{(0.1+1.1\i)t} 
+c_3(0,-0.87,0.22,-0.44)e^{0.5t} 
+c_4(-0.5,-0.5,0.5,-0.5)e^{-0.6t}\)}
\end{reduce}

\end{enumerate}
\end{exercise}





\begin{reduce}
\begin{exercise} \label{ex:sprmas} 
Recall the general complex solution that \cref{eg:sprmas} derives for the \idx{oscillation}s of a mass on a spring.
Show that substituting \(c_1=(C_1-\i C_2)/2\) and \(c_2=(C_1+\i C_2)/2\) for real~\(C_1\) and~\(C_2\) results in the velocity~\(v(t)\) being expressed algebraically in purely \text{real terms.}
\end{exercise}
\end{reduce}






\begin{exercise} 
In a few sentences, answer\slash discuss each of the following.
\begin{enumerate}
\item What is the relation between diagonalization of a matrix and using non-standard basis vectors for a space?

\item  Why is diagonalization useful in population modelling? and systems of differential equations?

\item Suppose you generate a matrix at random: why would you expect the matrix to be diagonalizable?

\item How does it occur that the dimension of an eigenspace may be strictly less than the multiplicity of the corresponding eigenvalue?

\item Why is the possibility of complex valued eigenvalues and eigenvectors important in an otherwise real problem?

\end{enumerate}
\end{exercise}

\begin{comment}
%{ED498555.pdf}
why, what caused X?
how did X occur?
what-if? what-if-not?
how does X compare with Y?
what is the evidence for X?
why is X important?

Add exercises on perhaps simple SIR models, and perhaps mechanical models.
\end{comment}



\index{diagonalization|)}


%%!TEX root = ../larxxia.tex

\section{Summary of determinants}
\label{sec:sumd}


\begin{itemize}
\def\index#1{}% turn off indexing

\itemhi The `graphical formulas'~\eqref{eq:dets23b} for \(2\times2\) and \(3\times3\) determinants are useful for both theory and many practical small problems.


\subsubsection*{Geometry underlies determinants}

\item The term \index{nD-cube@$n$D-cube|textbf}\textbf{$n$D-cube}  generalizes a square and cube to \(n\)~dimensions (\(\RR^n)\).
The term  \index{nD-volume@$n$D-volume|textbf}\textbf{$n$D-volume} generalizes the notion of area and volume to \(n\)~dimensions.
When the dimension of the space is unspecified, then we may say \bfidx{hyper-cube} and \bfidx{hyper-volume}, respectively.

\itemhi Let \(A\) be an \(n\times n\) \idx{square matrix}, and let \(C\)~be the unit \index{nD-cube@$n$D-cube}$n$D-cube in~\(\RR^n\).
Transform the \(n\)D-cube~\(C\) by \(\xv\mapsto A\xv\) to its image~\(C'\) in~\(\RR^n\). 
Define the \bfidx{determinant} of~\(A\), denoted either~\(\det A\) (or sometimes~\(|A|\)) such that (\cref{def:detarea}):  \begin{itemize}
\item the \idx{magnitude}~\(|\det A|\) is the \index{nD-volume@$n$D-volume}$n$D-volume of~\(C'\); and 
\item the sign of~\(\det A\) to be negative iff the transformation reflects the orientation of the $n$D-cube.
\end{itemize}

\itemhi (\cref{thm:basicdet})
\begin{itemize}
\item   For every \(n\times n\) \idx{diagonal matrix}~\(D\),
the \idx{determinant} of~\(D\) is the product of the \idx{diagonal entries}: \(\det D=d_{11}d_{22}\cdots d_{nn}\)\,.
\item  For every \idx{orthogonal matrix}~\(Q\),  \(\det Q=\pm1\) (only one alternative, not both), and \(\det Q=\det(\tr Q)\).
\item   For every \(n\times n\) matrix~\(A\) and every scalar~\(k\), \(\det(kA)=k^n\det A\)\,.
\end{itemize}

\item Consider any \idx{bounded} smooth \index{nD-volume@$n$D-volume}$n$D-volume~\(C\) in~\(\RR^n\) and its image~\(C'\) after multiplication by \(n\times n\) matrix~\(A\).
Then (\cref{thm:detanyC})
\begin{equation*}
\det A=\pm\frac{\text{$n$D-volume of }C'}
{\text{$n$D-volume of }C}
\end{equation*}
with the negative sign when matrix~\(A\) changes the orientation.

\itemhi For every two \index{square matrix}\(n\times n\) matrices \(A\) and~\(B\), \(\det(AB)=\det(A)\det(B)\) (\cref{thm:detprod}).
Further, for \(n\times n\) matrices \hlist A\ell, \begin{equation*}
\det(A_1A_2\cdots A_\ell)=\det(A_1)\det(A_2)\cdots\det(A_\ell).
\end{equation*}


\itemme For every \idx{square matrix}~\(A\), \(\det(\tr A)=\det A\) (\cref{thm:dettr}).

\item For every \(n\times n\) \idx{square matrix}~\(A\), the \idx{magnitude} of its determinant \(|\det A|=\sigma_1\sigma_2\cdots\sigma_n\)\,, the product of all its \idx{singular value}s (\cref{thm:detsvd}).

\itemme A \idx{square matrix}~\(A\) is \idx{invertible} iff \(\det A\neq 0\) (\cref{thm:detinv}).
If a matrix~\(A\) is invertible, then \(\det(A^{-1})=1/(\det A)\).







\subsubsection*{Laplace expansion theorem for determinants}

\itemhi For every \(n\times n\) matrix~\(A\), the determinant has the following row and column properties (\cref{thm:ppdet}).
\begin{itemize}
\item
If \(A\) has a \idx{zero row} or \index{zero column}column, then \(\det A=0\)\,.
\item
If \(A\) has two \idx{identical rows} or two \idx{identical columns}, then  \(\det A=0\)\,.
\item
Let \(B\) be obtained by \idx{interchanging two rows} or \index{interchanging two columns}two columns of~\(A\), then \(\det B=-\det A\)\,.
\item
Let \(B\) be obtained by multiplying any one row or column of~\(A\) by a scalar~\(k\), then \(\det B=k\det A\)\,.
\end{itemize}

\item For every \(n\times n\) matrix~\(A\),
define the \((i,j)\)th~\bfidx{minor} \(A_{ij}\)~to be the \((n-1)\times(n-1)\) square matrix obtained from~\(A\) by omitting the \(i\)th~row and \(j\)th~column.  
If, except for the entry~\(a_{ij}\), the \(i\)th~row (or \(j\)th~column) of~\(A\) is all zero, then (\cref{thm:rpdet:vii})
\begin{equation*}
\det A=(-1)^{i+j}a_{ij}\det A_{ij}\,.
\end{equation*}

\item A \bfidx{triangular matrix} is a \idx{square matrix} where all entries are zero either to the lower-left of the diagonal or to the upper-right (\cref{def:trim}): an \idx{upper triangular} matrix has zeros in the lower-left; and a \idx{lower triangular} matrix has zeros in the upper-right.

\itemme For every \(n\times n\) \idx{triangular matrix}~\(A\),
the \idx{determinant} of~\(A\) is the product of the \idx{diagonal entries}, \(\det A=a_{11}a_{22}\cdots a_{nn}\) (\cref{thm:rpdet:vi}).

\item Let \(A\), \(B\) and~\(C\) be \(n\times n\) matrices.
If matrices~\(A\), \(B\) and~\(C\) are identical except for their \(i\)th~column, and that the \(i\)th~column of~\(A\) is the sum of the \(i\)th~columns of \(B\) and~\(C\), then \(\det A=\det B+\det C\) (\cref{thm:rpdet:v}).
Further, the same sum property holds when ``column'' is replaced by ``row'' throughout.

\itemme For every \(n\times n\) matrix \(A=\begin{bmatrix} a_{ij} \end{bmatrix}\) (\(n\geq2\)), and in terms of the \((i,j)\)th \idx{minor}~\(A_{ij}\), the \idx{determinant} of~\(A\) can be computed via expansion in any row~\(i\) or any column~\(j\) as, respectively, (\cref{thm:letdet})
\begin{eqnarray*}
\det A
&=&(-1)^{i+1}a_{i1}\det A_{i1}
+(-1)^{i+2}a_{i2}\det A_{i2}
\nonumber\\&&{}
+\cdots+(-1)^{i+n}a_{in}\det A_{in}
\nonumber\\&=&(-1)^{j+1}a_{1j}\det A_{1j}
+(-1)^{j+2}a_{2j}\det A_{2j}
\nonumber\\&&{}
+\cdots+(-1)^{j+n}a_{nj}\det A_{nj}\,.
\end{eqnarray*}

\itemme The determinant of every \(n\times n\) matrix expands to the sum of \index{factorial}\(n!\)~terms, where each term is~\(\pm1\) times a product of \(n\)~factors such that each factor comes from different rows and columns of the matrix (\cref{thm:detnfac}).



\end{itemize}



\makeanswers


\endinput
